\documentclass[10pt]{article}
\usepackage[T1]{fontenc}
\usepackage[utf8]{inputenc}
\usepackage[a4paper,margin=18mm]{geometry}
\usepackage{amsmath,amssymb,mathtools,array,booktabs,longtable,xcolor,listings,needspace,microtype,enumitem}
\usepackage[hidelinks]{hyperref}
\setlength{\parindent}{0pt}\setlength{\parskip}{4pt}\setlist{nosep,leftmargin=*}\sloppy
\definecolor{BaselineBlue}{HTML}{EAF2F8}\definecolor{DraftGray}{HTML}{F4F4F4}\definecolor{RewriteGreen}{HTML}{EDF7ED}
\newcommand{\Rocq}{Rocq}\newcommand{\coqin}[1]{\texttt{#1}}\newcommand{\PG}{\mathrm{PGL}(2,7)}\newcommand{\worddist}{\mu^{*L}}
\renewcommand{\cite}[1]{\textnormal{[\texttt{\detokenize{#1}}]}}\renewcommand{\label}[1]{}
\AtBeginDocument{\renewcommand{\ref}[1]{\texttt{\detokenize{#1}}}}
\lstset{basicstyle=\ttfamily\scriptsize,columns=fullflexible,keepspaces=true,breaklines=true,breakatwhitespace=false,showstringspaces=false,frame=single,backgroundcolor=\color{DraftGray},xleftmargin=2mm,xrightmargin=2mm,aboveskip=4pt,belowskip=5pt}
\newcommand{\Field}[2]{\textbf{#1}\ #2\par}\newcommand{\Tag}[1]{\fcolorbox{black!25}{black!4}{\strut\scriptsize #1}}
\title{Shinagawa 2021 Paragraph Baseline and WADT 2026 Rewrite Suggestions}\author{Comparative writing analysis}\date{22 August 2026}
\begin{document}\maketitle
\begin{abstract}
This file is a baseline reference and rewrite table. It records a purpose and narration structure for every included Shinagawa 2021 unit. It then groups every included WADT 2026 unit under a qualitatively similar baseline paragraph or under the unmatched category. Each WADT record contains its full source unit with footnotes omitted, a qualitative relationship, complexity and jargon tags, and a table-only rewrite suggestion. The suggestions do not modify the manuscript. No numeric similarity score is used, and no Shinagawa source prose is reproduced.
\end{abstract}
\section{Coverage and reading guide}
The baseline contains \textbf{299} included Shinagawa units. Of these, \textbf{193} paragraphs can receive a primary WADT match. The WADT census contains \textbf{225} included rewrite units. Every WADT unit has one qualitative relationship and one rewrite suggestion.
\begin{itemize}
\item \textbf{same purpose}: the units perform the same main writing job.
\item \textbf{analogous purpose}: the writing job transfers with a change of technical subject.
\item \textbf{partial analogue}: only a named structural move transfers.
\item \textbf{no baseline analogue}: forcing a match would distort the WADT unit.
\end{itemize}
Relationship totals: same purpose 6, analogous purpose 75, partial analogue 76, and no baseline analogue 68.

Complexity changes from low/medium/high = 90/94/41 to 157/63/5. Jargon changes from low/medium/high = 34/113/78 to 117/106/2.
\section{Shinagawa 2021 baseline catalog}
The catalog paraphrases each unit's writing purpose and narration structure. It does not reproduce the published prose. The match count shows how many WADT units are aggregated at each eligible paragraph.
\scriptsize\setlength{\tabcolsep}{3pt}\renewcommand{\arraystretch}{1.15}
\begin{longtable}{@{}p{0.105\textwidth}p{0.09\textwidth}p{0.35\textwidth}p{0.35\textwidth}p{0.055\textwidth}@{}}
\toprule\textbf{ID}&\textbf{Type}&\textbf{Purpose}&\textbf{Narration structure}&\textbf{WADT}\\\midrule\endfirsthead
\toprule\textbf{ID}&\textbf{Type}&\textbf{Purpose}&\textbf{Narration structure}&\textbf{WADT}\\\midrule\endhead
\texttt{S-ABS-P001}&abstract paragraph / contribution&Place the work in its research area and state the paper's technical and framework contributions.&field context $\rightarrow$ prior work context $\rightarrow$ evaluation criterion $\rightarrow$ contribution $\rightarrow$ result scope $\rightarrow$ framework contribution&1\\
\texttt{S-01-P001}&prose paragraph / background&Define the field and explain why its physical medium helps readers and learners.&field definition $\rightarrow$ implementation contrast $\rightarrow$ accessibility rationale $\rightarrow$ application context&2\\
\texttt{S-01-P002}&prose paragraph / problem&Show that the existing design handles a narrow function class and motivate a new design for efficient missing cases.&application examples $\rightarrow$ multi value need $\rightarrow$ prior design $\rightarrow$ efficiency question $\rightarrow$ limitation $\rightarrow$ open problem $\rightarrow$ solution preview&3\\
\texttt{S-01-P003}&prose paragraph / contribution&Present the new card construction and the broader protocol model as the paper's two main additions.&artifact contribution $\rightarrow$ protocol results $\rightarrow$ comparison summary $\rightarrow$ model contribution $\rightarrow$ reuse scope $\rightarrow$ future boundary&10\\
\texttt{S-02-P001}&prose paragraph / roadmap&State why the model needs both a card description and operations, then preview the section's explanatory order.&section goal $\rightarrow$ model requirements $\rightarrow$ expository choice $\rightarrow$ forward reference&5\\
\texttt{S-02-CAP001}&caption / comparison&Name the table's comparison target, abbreviations, and cost dimensions.&comparison target $\rightarrow$ abbreviation key $\rightarrow$ metric scope&--\\
\texttt{S-02-TC002}&table cell / comparison&Identify the implementation family used in each comparison row.&comparison dimension&--\\
\texttt{S-02-TC003}&table cell / comparison&Identify the resource count used as the first cost measure.&cost dimension&--\\
\texttt{S-02-TC004}&table cell / comparison&Identify the operation count used as the second cost measure.&cost dimension&--\\
\texttt{S-02-TC005}&table cell / comparison&Open the first task category in the protocol comparison.&task category&--\\
\texttt{S-02-TC006}&table cell / comparison&Identify the reference row for the first task category.&reference row&--\\
\texttt{S-02-TC014}&table cell / comparison&Open the second task category in the protocol comparison.&task category&--\\
\texttt{S-02-TC015}&table cell / comparison&Identify the reference row for the second task category.&reference row&--\\
\texttt{S-02-TC023}&table cell / comparison&Open the third task category in the protocol comparison.&task category&--\\
\texttt{S-02-TC024}&table cell / comparison&Identify the reference row for the third task category.&reference row&--\\
\texttt{S-02-TC032}&table cell / comparison&Open the fourth task category in the protocol comparison.&task category&--\\
\texttt{S-02-TC033}&table cell / comparison&Identify the reference row for the fourth task category.&reference row&--\\
\texttt{S-02-TC041}&table cell / comparison&Open the fifth task category in the protocol comparison.&task category&--\\
\texttt{S-02-TC042}&table cell / comparison&Identify the reference row for the fifth task category.&reference row&--\\
\texttt{S-02-P002}&prose paragraph / roadmap&Explain which ingredients the model must define and why a familiar case is used before the general case.&model requirements $\rightarrow$ expository example choice $\rightarrow$ later extension $\rightarrow$ appendix scope&5\\
\texttt{S-02-P003}&prose paragraph / problem&Review the earlier deck representation and show why it cannot describe the wider card class needed here.&prior definition $\rightarrow$ worked illustration $\rightarrow$ assumption $\rightarrow$ captured cases $\rightarrow$ expressiveness limit $\rightarrow$ definition handoff&0\\
\texttt{S-02-STM001}&statement / definition&Define the deck object as a tuple whose fields are explained immediately afterward.&object naming $\rightarrow$ tuple specification&--\\
\texttt{S-02-P004}&prose paragraph / definition&Assign a mathematical role to each field of the deck tuple.&carrier role $\rightarrow$ transformation role $\rightarrow$ symbol role $\rightarrow$ visibility role $\rightarrow$ multiset role&6\\
\texttt{S-02-P005}&prose paragraph / assumption&State the identity requirement and name the reduced tuple used as a card specification.&identity requirement $\rightarrow$ terminology assignment&1\\
\texttt{S-02-STM002}&statement / example&Introduce a familiar finite deck and announce its component-by-component specification.&example selection $\rightarrow$ application context $\rightarrow$ component list handoff&--\\
\texttt{S-02-LI001}&list item / example&Instantiate one component of a worked formal example.&example component assignment&--\\
\texttt{S-02-LI002}&list item / example&Instantiate the symbol component in the running formal example.&component identification $\rightarrow$ formal assignment&--\\
\texttt{S-02-LI003}&list item / example&Instantiate the transformation component in the running formal example.&component identification $\rightarrow$ formal assignment&--\\
\texttt{S-02-LI004}&list item / example&Instantiate the visibility component in the running formal example.&component identification $\rightarrow$ formal assignment&--\\
\texttt{S-02-LI005}&list item / example&Instantiate the deck-multiset component in the running formal example.&component identification $\rightarrow$ formal assignment&--\\
\texttt{S-02-P006}&prose paragraph / definition&Explain how orientation, transformation, and visibility determine the meaning of the running card instance.&orientation semantics $\rightarrow$ transformation effect $\rightarrow$ visibility semantics $\rightarrow$ instance naming&1\\
\texttt{S-02-STM003}&statement / definition&Define a sequence by applying allowed transformations to the cards of a deck.&input object $\rightarrow$ sequence construction&--\\
\texttt{S-02-P007}&prose paragraph / definition&State the transformation and multiset conditions, then name the set of all valid sequences.&transformation constraint $\rightarrow$ multiset constraint $\rightarrow$ collection notation&0\\
\texttt{S-02-STM004}&statement / example&Instantiate the sequence definition for the running deck.&example reference $\rightarrow$ sequence instance&--\\
\texttt{S-02-P008}&prose paragraph / example&Show which transformations generate the stated sequence instance.&construction rationale $\rightarrow$ transformation expansion&1\\
\texttt{S-02-P009}&prose paragraph / interpretation&State the observable form represented by the worked sequence.&representation interpretation&1\\
\texttt{S-02-P010}&prose paragraph / transition&Introduce the next derived object and hand its content to the formal statement.&concept naming $\rightarrow$ definition handoff&2\\
\texttt{S-02-STM005}&statement / definition&Define the observable sequence by applying the visibility map componentwise.&input sequence $\rightarrow$ componentwise observation $\rightarrow$ derived object&--\\
\texttt{S-02-P011}&prose paragraph / definition&Name the collection of all observable sequences associated with a deck.&collection naming $\rightarrow$ set specification&0\\
\texttt{S-02-STM006}&statement / example&Apply the visibility definition to the earlier sequence and record two shorthand forms.&example reference $\rightarrow$ derived observation $\rightarrow$ notation variants&--\\
\texttt{S-02-P012}&prose paragraph / definition&Fix a deck and sequence, then divide operations into state-changing and revealing classes.&object setup $\rightarrow$ two class partition&0\\
\texttt{S-02-LI006}&list item / definition&Define the operation class that changes the current sequence.&class naming $\rightarrow$ state change role&--\\
\texttt{S-02-P013}&prose paragraph / definition&Partition state-changing operations according to whether their output is fixed or random.&deterministic case $\rightarrow$ probabilistic case&1\\
\texttt{S-02-LI007}&list item / definition&Define the operation class that reveals information without changing the sequence.&class naming $\rightarrow$ revelation role $\rightarrow$ state preservation&--\\
\texttt{S-02-P014}&prose paragraph / method&Fix the common inputs and introduce the standard operation family in the order used below.&operation scope $\rightarrow$ common object setup $\rightarrow$ operation inventory $\rightarrow$ first operation lead in&0\\
\texttt{S-02-P015}&prose paragraph / definition&Explain the positional effect of a permutation and collect all such operations into one set.&operation effect $\rightarrow$ position interpretation $\rightarrow$ collection definition&1\\
\texttt{S-02-P016}&prose paragraph / definition&Specify the inputs and output shape for a transformation applied at selected positions.&position set setup $\rightarrow$ operation signature $\rightarrow$ output handoff&0\\
\texttt{S-02-P017}&prose paragraph / definition&Give the selected-position update rule, preserve all other positions, and name the operation set.&selected position rule $\rightarrow$ unchanged position rule $\rightarrow$ collection definition&0\\
\texttt{S-02-P018}&prose paragraph / comparison&Distinguish two operation classes by whether the observable sequence changes.&potential confusion $\rightarrow$ classification decision $\rightarrow$ distinguishing criterion&12\\
\texttt{S-02-P019}&prose paragraph / definition&Define a randomised sequence transformation through a permutation family and a distribution.&parameter tuple $\rightarrow$ sampling rule $\rightarrow$ sequence update $\rightarrow$ collection handoff&4\\
\texttt{S-02-P020}&prose paragraph / definition&Collect all shuffle specifications satisfying the stated family and distribution conditions.&collection definition&0\\
\texttt{S-02-P021}&prose paragraph / definition&Model one operation as a state transition paired with revealed information.&object setup $\rightarrow$ operation application $\rightarrow$ state transition $\rightarrow$ revelation output&1\\
\texttt{S-02-P022}&prose paragraph / definition&Specify what observers know before and after one operation and combine those observations into one view.&conversion case $\rightarrow$ opening case $\rightarrow$ pre state observation $\rightarrow$ post state observation $\rightarrow$ view assembly $\rightarrow$ example forward reference&0\\
\texttt{S-02-P023}&prose paragraph / definition&Extend the one-step observation model to an ordered list of operations and intermediate states.&operation list setup $\rightarrow$ state trace&2\\
\texttt{S-02-P024}&prose paragraph / definition&Pair each state observation with its revealed output, set the conversion convention, and name the set of views.&revealed output setup $\rightarrow$ paired trace $\rightarrow$ initial convention $\rightarrow$ conversion convention $\rightarrow$ view naming $\rightarrow$ collection definition&1\\
\texttt{S-02-STM007}&statement / example&Choose a deck and operation list that instantiate the view definition.&example objects $\rightarrow$ operation list&--\\
\texttt{S-02-P025}&prose paragraph / example&Apply the operation list to an initial sequence and display the successive states.&initial state $\rightarrow$ operation application $\rightarrow$ state trace&0\\
\texttt{S-02-P026}&prose paragraph / example&Derive the view associated with the preceding worked state trace.&view derivation $\rightarrow$ observer trace&1\\
\texttt{S-02-P027}&prose paragraph / scope&State when unrevealing outputs may be omitted from the displayed view.&notation condition $\rightarrow$ abbreviated form&3\\
\texttt{S-02-P028}&prose paragraph / scope&Introduce a more compact notation for the same worked view.&notation variant&0\\
\texttt{S-02-P029}&prose paragraph / transition&Name the protocol object and hand its content to the formal definition.&concept naming $\rightarrow$ definition handoff&1\\
\texttt{S-02-STM008}&statement / definition&Define a protocol as a tuple whose components are explained in the following list.&object naming $\rightarrow$ tuple specification $\rightarrow$ component list handoff&--\\
\texttt{S-02-LI008}&list item / definition&Mark the first component in the protocol tuple definition.&component marker&--\\
\texttt{S-02-LI009}&list item / definition&Mark the second component in the protocol tuple definition.&component marker&--\\
\texttt{S-02-LI010}&list item / definition&Mark the third component in the protocol tuple definition.&component marker&--\\
\texttt{S-02-LI011}&list item / definition&Mark the fourth component in the protocol tuple definition.&component marker&--\\
\texttt{S-02-LI012}&list item / definition&Mark the fifth component in the protocol tuple definition.&component marker&--\\
\texttt{S-02-P030}&prose paragraph / definition&Assign roles to the input count, domain, deck, operation set, and action function.&input count role $\rightarrow$ input domain role $\rightarrow$ deck role $\rightarrow$ operation set role $\rightarrow$ action function role&12\\
\texttt{S-02-P031}&prose paragraph / method&Fix a protocol and initial state, then introduce its execution as an ordered procedure.&protocol setup $\rightarrow$ initial state setup $\rightarrow$ procedure handoff&2\\
\texttt{S-02-LI013}&list item / method&Set the execution's initial sequence before the iterative rule begins.&initial state assignment&--\\
\texttt{S-02-P032}&prose paragraph / method&Initialise the current sequence and accumulated view variables.&current state initialisation $\rightarrow$ view initialisation $\rightarrow$ variable role explanation&0\\
\texttt{S-02-LI014}&list item / method&Compute the next action, update the sequence, and append the new observation until the stop value appears.&action evaluation $\rightarrow$ operation application $\rightarrow$ state update $\rightarrow$ view update $\rightarrow$ iteration condition&--\\
\texttt{S-02-LI015}&list item / method&Terminate the execution when the action function returns the stop value.&termination condition&--\\
\texttt{S-02-STM009}&statement / example&Introduce a modified published protocol as a complete instance of the formal model.&example selection $\rightarrow$ provenance context $\rightarrow$ tuple handoff&--\\
\texttt{S-02-P033}&prose paragraph / example&Specify the deck and operation family, then introduce the action function for the worked protocol.&deck assignment $\rightarrow$ operation set assignment $\rightarrow$ action function handoff&0\\
\texttt{S-02-LI016}&list item / example&Mark the first case in the worked action-function specification.&action case marker&--\\
\texttt{S-02-LI017}&list item / example&Mark the second case in the worked action-function specification.&action case marker&--\\
\texttt{S-02-LI018}&list item / example&Mark the third case in the worked action-function specification.&action case marker&--\\
\texttt{S-02-LI019}&list item / example&Mark the fourth case in the worked action-function specification.&action case marker&--\\
\texttt{S-02-LI020}&list item / example&Mark the fifth case in the worked action-function specification.&action case marker&--\\
\texttt{S-02-LI021}&list item / example&Mark the sixth case in the worked action-function specification.&action case marker&--\\
\texttt{S-02-P034}&prose paragraph / example&Enumerate the action-function cases and the view histories that select them.&action case inventory $\rightarrow$ stop case $\rightarrow$ view state setup $\rightarrow$ view history inventory&0\\
\texttt{S-02-LI022}&list item / example&Mark the first view-history case in the worked protocol.&view case marker&--\\
\texttt{S-02-LI023}&list item / example&Mark the second view-history case in the worked protocol.&view case marker&--\\
\texttt{S-02-LI024}&list item / example&Mark the third view-history case in the worked protocol.&view case marker&--\\
\texttt{S-02-LI025}&list item / example&Mark the fourth view-history case in the worked protocol.&view case marker&--\\
\texttt{S-02-LI026}&list item / example&Mark the fifth view-history case in the worked protocol.&view case marker&--\\
\texttt{S-02-P035}&prose paragraph / example&Choose an encoded input sequence and announce the step-by-step execution.&execution goal $\rightarrow$ initial sequence $\rightarrow$ step list handoff&1\\
\texttt{S-02-P036}&prose paragraph / definition&Define the two-card encoding used by the worked inputs and hand off to the procedure.&encoding domain $\rightarrow$ case based encoding $\rightarrow$ procedure handoff&0\\
\texttt{S-02-LI027}&list item / example&Apply the first deterministic rearrangement to the current sequence.&step identification $\rightarrow$ deterministic update&--\\
\texttt{S-02-LI028}&list item / example&Apply the randomised rearrangement and display its possible outcomes.&step identification $\rightarrow$ randomised update $\rightarrow$ outcome cases&--\\
\texttt{S-02-P037}&prose paragraph / definition&Assign a conventional name to the randomised operation used in the example.&terminology assignment&0\\
\texttt{S-02-LI029}&list item / example&Apply the inverse deterministic rearrangement to restore the intended positions.&step identification $\rightarrow$ inverse update&--\\
\texttt{S-02-LI030}&list item / example&Reveal the first encoded input and display the two observable cases.&step identification $\rightarrow$ revelation action $\rightarrow$ observable cases&--\\
\texttt{S-02-P038}&prose paragraph / method&Use the observed case to choose between immediate termination and continuation.&case test $\rightarrow$ termination branch $\rightarrow$ continuation branch&0\\
\texttt{S-02-LI031}&list item / example&Apply the final rearrangement in the continuing branch.&step identification $\rightarrow$ final update&--\\
\texttt{S-02-P039}&prose paragraph / result&State that both branches terminate and display the resulting output form.&branch reconciliation $\rightarrow$ termination result $\rightarrow$ final state&0\\
\texttt{S-02-P040}&prose paragraph / interpretation&Identify the logical property encoded by the final state and use it to classify the protocol.&output property $\rightarrow$ protocol classification&4\\
\texttt{S-02-P041}&prose paragraph / definition&Introduce an abstract input-output interface needed to state later protocol properties.&motivation $\rightarrow$ concept scope $\rightarrow$ formal handoff&0\\
\texttt{S-02-P042}&prose paragraph / definition&Give a compact symbolic representation of the running input-output contract.&object identification $\rightarrow$ display handoff&3\\
\texttt{S-02-P043}&prose paragraph / definition&Explain how unspecified positions can be suppressed when only part of an interface matters.&object setup $\rightarrow$ formal definition $\rightarrow$ scope boundary&0\\
\texttt{S-02-P044}&prose paragraph / definition&Define partial sequence descriptions and collect the descriptions compatible with a full sequence.&object setup $\rightarrow$ formal definition $\rightarrow$ scope boundary&1\\
\texttt{S-02-P045}&prose paragraph / definition&State when a complete sequence agrees with a partial sequence description.&premises $\rightarrow$ defining condition $\rightarrow$ scope closure&0\\
\texttt{S-02-STM010}&statement / example&Instantiate the partial-sequence construction on a short finite sequence.&example context $\rightarrow$ object instantiation $\rightarrow$ detail handoff&--\\
\texttt{S-02-P046}&prose paragraph / example&Show how one concrete full sequence satisfies a partial description.&instance setup $\rightarrow$ worked application $\rightarrow$ example outcome&1\\
\texttt{S-02-P047}&prose paragraph / definition&Lift full and partial sequences to families indexed by protocol inputs.&object setup $\rightarrow$ formal definition $\rightarrow$ scope boundary&0\\
\texttt{S-02-STM011}&statement / example&Introduce a concrete input-dependent sequence from the running protocol.&example context $\rightarrow$ object instantiation $\rightarrow$ detail handoff&--\\
\texttt{S-02-P048}&prose paragraph / example&Specify the running input family by cases.&instance setup $\rightarrow$ worked application $\rightarrow$ example outcome&1\\
\texttt{S-02-P049}&prose paragraph / example&Introduce the corresponding partial output family for the running protocol.&example context $\rightarrow$ object instantiation $\rightarrow$ detail handoff&0\\
\texttt{S-02-P050}&prose paragraph / definition&Signal the formal definition of the input-output interface.&motivation $\rightarrow$ concept scope $\rightarrow$ formal handoff&0\\
\texttt{S-02-STM012}&statement / definition&Define a functionality as paired input and partial-output families.&premises $\rightarrow$ defining condition $\rightarrow$ scope closure&--\\
\texttt{S-02-P051}&prose paragraph / definition&Clarify the domains and codomains of the input and output components.&object setup $\rightarrow$ formal definition $\rightarrow$ scope boundary&3\\
\texttt{S-02-P052}&prose paragraph / definition&Introduce the semantic criterion used to judge protocol outputs.&motivation $\rightarrow$ concept scope $\rightarrow$ formal handoff&6\\
\texttt{S-02-STM013}&statement / definition&Require every execution to terminate in a sequence compatible with the specified output.&premises $\rightarrow$ defining condition $\rightarrow$ scope closure&--\\
\texttt{S-02-P053}&prose paragraph / definition&Introduce a specialised correctness criterion for protocols whose values remain encoded.&motivation $\rightarrow$ concept scope $\rightarrow$ formal handoff&0\\
\texttt{S-02-STM014}&statement / definition&State the setup and the conditions for computing through commitments.&premises $\rightarrow$ defining condition $\rightarrow$ scope closure&--\\
\texttt{S-02-LI032}&list item / definition&Require the protocol to satisfy its underlying functionality.&criterion item&--\\
\texttt{S-02-LI033}&list item / definition&Require the input arrangement to consist of encoded arguments and fixed auxiliary material.&criterion item&--\\
\texttt{S-02-LI034}&list item / definition&Require the output arrangement to retain an encoding of the computed value.&criterion item&--\\
\texttt{S-02-P054}&prose paragraph / definition&Name the distribution induced by a protocol execution and identify its source of randomness.&object setup $\rightarrow$ formal definition $\rightarrow$ scope boundary&3\\
\texttt{S-02-P055}&prose paragraph / definition&Introduce the indistinguishability condition used for protocol security.&motivation $\rightarrow$ concept scope $\rightarrow$ formal handoff&3\\
\texttt{S-02-STM015}&statement / definition&Require the view distribution to be independent of the chosen input.&premises $\rightarrow$ defining condition $\rightarrow$ scope closure&--\\
\texttt{S-02-STM016}&statement / example&Set up a worked security argument for the running protocol.&example context $\rightarrow$ object instantiation $\rightarrow$ detail handoff&--\\
\texttt{S-02-P056}&prose paragraph / proof&Compute the possible views and their probabilities for an arbitrary input.&proof setup $\rightarrow$ calculation $\rightarrow$ intermediate result&3\\
\texttt{S-02-P057}&prose paragraph / proof&Use input-independent probabilities to discharge the security criterion.&condition verification $\rightarrow$ conclusion&5\\
\texttt{S-02-P058}&prose paragraph / definition&Present a subroutine as an atomic replacement for a complete protocol execution.&object setup $\rightarrow$ formal definition $\rightarrow$ scope boundary&1\\
\texttt{S-02-P059}&prose paragraph / definition&Introduce the formal syntax for invoking a protocol as one operation.&object identification $\rightarrow$ display handoff&0\\
\texttt{S-02-P060}&prose paragraph / definition&Constrain the selected positions and collect all valid invocations of one protocol.&object setup $\rightarrow$ formal definition $\rightarrow$ scope boundary&0\\
\texttt{S-02-P061}&prose paragraph / definition&Extend the operation set from one subordinate protocol to a finite family.&object setup $\rightarrow$ formal definition $\rightarrow$ scope boundary&0\\
\texttt{S-02-P062}&prose paragraph / definition&Introduce the compatibility condition needed for safe protocol composition.&motivation $\rightarrow$ concept scope $\rightarrow$ formal handoff&0\\
\texttt{S-02-STM017}&statement / definition&State the objects and conditions governing valid subroutine calls.&premises $\rightarrow$ defining condition $\rightarrow$ scope closure&--\\
\texttt{S-02-LI035}&list item / definition&Require every needed subroutine invocation to belong to the outer operation set.&criterion item&--\\
\texttt{S-02-LI036}&list item / definition&Require each call site to contain a valid subordinate input arrangement.&criterion item&--\\
\texttt{S-02-P063}&prose paragraph / definition&Complete the call-site condition and allow the subordinate input to vary between invocations.&object setup $\rightarrow$ formal definition $\rightarrow$ scope boundary&0\\
\texttt{S-02-STM018}&statement / example&Introduce a two-input protocol that will serve as a reusable subroutine.&example context $\rightarrow$ object instantiation $\rightarrow$ detail handoff&--\\
\texttt{S-02-P064}&prose paragraph / example&Specify the interface realised by the reusable subordinate protocol.&object identification $\rightarrow$ display handoff&0\\
\texttt{S-02-P065}&prose paragraph / example&Adapt the smaller protocol and use it to construct a larger protocol instance.&construction goal $\rightarrow$ mechanism $\rightarrow$ realisation&1\\
\texttt{S-02-P066}&prose paragraph / example&Specify the interface realised by the larger composed protocol.&object identification $\rightarrow$ display handoff&0\\
\texttt{S-02-P067}&prose paragraph / method&Move from the composed protocol specification to its ordered calls.&object recall $\rightarrow$ procedure handoff&0\\
\texttt{S-02-LI037}&list item / method&Apply the subordinate protocol to the first selected block.&step identifier $\rightarrow$ operation application $\rightarrow$ state update&--\\
\texttt{S-02-LI038}&list item / method&Apply the subordinate protocol to the second selected block.&step identifier $\rightarrow$ operation application $\rightarrow$ state update&--\\
\texttt{S-02-P068}&prose paragraph / proof&Verify the two defining composition conditions for the worked protocol.&condition verification $\rightarrow$ conclusion&1\\
\texttt{S-02-STM019}&statement / result&State that valid abstract subroutine calls can be expanded without losing correctness or security.&hypotheses $\rightarrow$ compatibility condition $\rightarrow$ result&--\\
\texttt{S-02-PRF001}&proof paragraph / proof&Replace each abstract call by its implementation and compare outputs and views.&construction $\rightarrow$ correctness transfer $\rightarrow$ security transfer $\rightarrow$ conclusion&--\\
\texttt{S-03-P001}&prose paragraph / definition&Introduce the parameterised card type through a list of its required capabilities.&motivation $\rightarrow$ concept scope $\rightarrow$ formal handoff&1\\
\texttt{S-03-LI001}&list item / definition&Mark the first component of the five-part capability list.&list item marker&--\\
\texttt{S-03-LI002}&list item / definition&Mark the second component of the five-part capability list.&list item marker&--\\
\texttt{S-03-LI003}&list item / definition&Mark the third component of the five-part capability list.&list item marker&--\\
\texttt{S-03-LI004}&list item / definition&Mark the fourth component of the five-part capability list.&list item marker&--\\
\texttt{S-03-LI005}&list item / definition&Mark the fifth component of the five-part capability list.&list item marker&--\\
\texttt{S-03-P002}&prose paragraph / definition&Specify the value domain, permitted transformations, and restricted observations.&object setup $\rightarrow$ formal definition $\rightarrow$ scope boundary&1\\
\texttt{S-03-P003}&prose paragraph / interpretation&Relate the abstract capability list to a geometric card shape and a concrete instance.&structural consequence $\rightarrow$ example handoff&2\\
\texttt{S-03-P004}&prose paragraph / method&Explain how visible geometry and hidden markings encode the card state.&construction goal $\rightarrow$ mechanism $\rightarrow$ realisation&1\\
\texttt{S-03-P005}&prose paragraph / method&Explain how rotation and reflection realise the allowed state changes.&construction goal $\rightarrow$ mechanism $\rightarrow$ realisation&0\\
\texttt{S-03-P006}&prose paragraph / example&Map the symmetry axes of a concrete card to transformation parameters.&instance setup $\rightarrow$ worked application $\rightarrow$ example outcome&0\\
\texttt{S-03-P007}&prose paragraph / example&Illustrate one parameterised reflection on the concrete card.&instance setup $\rightarrow$ worked application $\rightarrow$ example outcome&0\\
\texttt{S-03-P008}&prose paragraph / background&Provide brief material context for the hidden-marking mechanism.&term explanation $\rightarrow$ use context $\rightarrow$ boundary&0\\
\texttt{S-03-P009}&prose paragraph / method&Generalise the axis mapping and explain how the hidden marking supports one permitted observation.&construction goal $\rightarrow$ mechanism $\rightarrow$ realisation&0\\
\texttt{S-03-P010}&prose paragraph / example&Interpret one observation outcome and introduce the analogous residue observation.&instance setup $\rightarrow$ worked application $\rightarrow$ example outcome&0\\
\texttt{S-03-P011}&prose paragraph / definition&Interpret the second observation, name the two observable components, and begin the formal card specification.&object setup $\rightarrow$ formal definition $\rightarrow$ scope boundary&0\\
\texttt{S-03-P012}&prose paragraph / definition&Define the parameterised rotation action on one card.&operation label $\rightarrow$ parameter scope $\rightarrow$ formal form&0\\
\texttt{S-03-P013}&prose paragraph / definition&Define the parameterised reflection action on one card.&operation label $\rightarrow$ parameter scope $\rightarrow$ formal form&0\\
\texttt{S-03-P014}&prose paragraph / definition&Collect the identity, rotations, and reflections into one transformation set.&family signal $\rightarrow$ set definition&0\\
\texttt{S-03-P015}&prose paragraph / definition&Specify the public symbol alphabet of the card type.&family signal $\rightarrow$ set definition&0\\
\texttt{S-03-P016}&prose paragraph / definition&Define what public symbol is exposed by every card state.&premises $\rightarrow$ defining condition $\rightarrow$ scope closure&0\\
\texttt{S-03-P017}&prose paragraph / definition&Package the card states, transformations, symbols, and vision map into one specification.&premises $\rightarrow$ defining condition $\rightarrow$ scope closure&0\\
\texttt{S-03-P018}&prose paragraph / definition&Identify the single-card representation used to commit to a value.&premises $\rightarrow$ defining condition $\rightarrow$ scope closure&0\\
\texttt{S-03-P019}&prose paragraph / definition&Inventory the operations and align the permutation case with the earlier model.&operation inventory $\rightarrow$ prior definition link $\rightarrow$ definition&2\\
\texttt{S-03-P020}&prose paragraph / definition&Introduce the parameters and notation for a positional rotation.&operation label $\rightarrow$ parameter scope $\rightarrow$ formal form&1\\
\texttt{S-03-P021}&prose paragraph / definition&Describe the positional action and illustrate it on a mixed sequence.&parameter scope $\rightarrow$ state transformation $\rightarrow$ boundary or example&0\\
\texttt{S-03-P022}&prose paragraph / definition&Collect every valid positional rotation into one operation family.&family signal $\rightarrow$ set definition&0\\
\texttt{S-03-P023}&prose paragraph / definition&Introduce the selected positions for a shared random rotation.&operation label $\rightarrow$ parameter scope $\rightarrow$ formal form&0\\
\texttt{S-03-P024}&prose paragraph / definition&Specify the shared random parameter and its action on selected cards.&parameter scope $\rightarrow$ state transformation $\rightarrow$ boundary or example&0\\
\texttt{S-03-P025}&prose paragraph / example&Show one shared random rotation and preserve the unselected position.&instance setup $\rightarrow$ worked application $\rightarrow$ example outcome&0\\
\texttt{S-03-P026}&prose paragraph / definition&Collect every valid selected-position rotation shuffle.&family signal $\rightarrow$ set definition&0\\
\texttt{S-03-P027}&prose paragraph / definition&Introduce the axis and selected-position parameters for reflection.&operation label $\rightarrow$ parameter scope $\rightarrow$ formal form&0\\
\texttt{S-03-P028}&prose paragraph / definition&Specify the parameter domains and the induced sequence transformation.&parameter scope $\rightarrow$ state transformation $\rightarrow$ boundary or example&0\\
\texttt{S-03-P029}&prose paragraph / definition&Give the componentwise action, illustrate it, and collect the operation family.&parameter scope $\rightarrow$ state transformation $\rightarrow$ boundary or example&0\\
\texttt{S-03-P030}&prose paragraph / definition&Introduce the axes and disjoint position blocks for a coupled random reflection.&operation label $\rightarrow$ parameter scope $\rightarrow$ formal form&0\\
\texttt{S-03-P031}&prose paragraph / definition&Specify the shared random choice across blocks and illustrate its outcomes.&parameter scope $\rightarrow$ state transformation $\rightarrow$ boundary or example&0\\
\texttt{S-03-P032}&prose paragraph / method&Show the physical device used to realise the coupled random reflection.&construction goal $\rightarrow$ mechanism $\rightarrow$ realisation&0\\
\texttt{S-03-P033}&prose paragraph / definition&Collect every valid choice of axes and disjoint blocks.&family signal $\rightarrow$ set definition&0\\
\texttt{S-03-P034}&prose paragraph / definition&Introduce the selected positions for a shared binary rotation.&operation label $\rightarrow$ parameter scope $\rightarrow$ formal form&0\\
\texttt{S-03-P035}&prose paragraph / definition&Specify the position set and the induced sequence mapping.&parameter scope $\rightarrow$ state transformation $\rightarrow$ boundary or example&0\\
\texttt{S-03-P036}&prose paragraph / definition&Describe the shared random bit and illustrate the two possible outputs.&parameter scope $\rightarrow$ state transformation $\rightarrow$ boundary or example&0\\
\texttt{S-03-P037}&prose paragraph / method&Show the physical device used to realise the shared binary rotation.&construction goal $\rightarrow$ mechanism $\rightarrow$ realisation&0\\
\texttt{S-03-P038}&prose paragraph / definition&Collect every valid selected-position instance of the operation.&family signal $\rightarrow$ set definition&0\\
\texttt{S-03-P039}&prose paragraph / definition&Introduce the position parameter for revealing one sign component.&operation label $\rightarrow$ parameter scope $\rightarrow$ formal form&0\\
\texttt{S-03-P040}&prose paragraph / definition&Specify the revealed component, illustrate an output, and collect all positions.&parameter scope $\rightarrow$ state transformation $\rightarrow$ boundary or example&0\\
\texttt{S-03-P041}&prose paragraph / definition&Introduce the position parameter for revealing one residue component.&operation label $\rightarrow$ parameter scope $\rightarrow$ formal form&0\\
\texttt{S-03-P042}&prose paragraph / definition&Specify the revealed component, illustrate an output, and collect all positions.&parameter scope $\rightarrow$ state transformation $\rightarrow$ boundary or example&0\\
\texttt{S-03-P043}&prose paragraph / definition&Introduce the position parameter for revealing the complete card value.&operation label $\rightarrow$ parameter scope $\rightarrow$ formal form&0\\
\texttt{S-03-P044}&prose paragraph / definition&Specify complete revelation and relate it to successive partial openings.&parameter scope $\rightarrow$ state transformation $\rightarrow$ boundary or example&0\\
\texttt{S-03-P045}&prose paragraph / definition&Bind one notation to the union of the operation families used later.&notation signal $\rightarrow$ aggregate definition $\rightarrow$ forward use&0\\
\texttt{S-04-P001}&prose paragraph / definition&Introduce the formal target relation before the display that specifies it.&definition lead in $\rightarrow$ display handoff&1\\
\texttt{S-04-P002}&prose paragraph / definition&Complete the functionality domain and introduce the initialization protocol object.&domain qualification $\rightarrow$ protocol lead in $\rightarrow$ display handoff&0\\
\texttt{S-04-P003}&prose paragraph / transition&Open the ordered execution of the initialization protocol.&procedure announcement $\rightarrow$ list handoff&2\\
\texttt{S-04-LI001}&list item / method&Randomize the input card by a rotation shuffle.&ordered step $\rightarrow$ operation application $\rightarrow$ display handoff&--\\
\texttt{S-04-LI002}&list item / method&Reveal the shuffled card value and classify it as public information.&ordered step $\rightarrow$ operation application $\rightarrow$ output binding $\rightarrow$ information status&--\\
\texttt{S-04-LI003}&list item / method&Cancel the revealed offset by rotating the card through its additive inverse.&ordered step $\rightarrow$ operation application $\rightarrow$ display handoff&--\\
\texttt{S-04-P004}&prose paragraph / proof&Close correctness and begin security with the concrete transcript for an arbitrary input.&termination $\rightarrow$ correctness assessment $\rightarrow$ security setup $\rightarrow$ view instantiation&3\\
\texttt{S-04-P005}&prose paragraph / proof&Replace the input-shifted revealed value with an equivalent fresh uniform value.&relation instantiation $\rightarrow$ distributional equivalence $\rightarrow$ simulator definition&1\\
\texttt{S-04-P006}&prose paragraph / proof&Establish that the simulated transcript ignores the input and prepare the equality of views.&randomness qualification $\rightarrow$ independence claim $\rightarrow$ equality handoff&0\\
\texttt{S-04-P007}&prose paragraph / result&Conclude secure realization and assess card and random-operation costs.&security conclusion $\rightarrow$ efficiency assessment $\rightarrow$ resource breakdown&11\\
\texttt{S-04-P008}&prose paragraph / definition&Introduce the formal input-output relation for addition.&definition lead in $\rightarrow$ display handoff&0\\
\texttt{S-04-P009}&prose paragraph / definition&Specify the operands' domain and introduce the addition protocol tuple.&domain qualification $\rightarrow$ protocol lead in $\rightarrow$ display handoff&0\\
\texttt{S-04-P010}&prose paragraph / transition&Announce the ordered operations of the addition protocol.&procedure announcement $\rightarrow$ list handoff&0\\
\texttt{S-04-LI004}&list item / method&Reflect the left operand across the zero axis.&ordered step $\rightarrow$ operation application $\rightarrow$ display handoff&--\\
\texttt{S-04-LI005}&list item / method&Apply one shared random rotation to both cards.&ordered step $\rightarrow$ operation application $\rightarrow$ display handoff&--\\
\texttt{S-04-LI006}&list item / method&Open the shuffled left card and mark its value as revealed.&ordered step $\rightarrow$ operation application $\rightarrow$ output binding $\rightarrow$ information status&--\\
\texttt{S-04-LI007}&list item / method&Rotate both cards by the inverse of the opened left value.&ordered step $\rightarrow$ operation application $\rightarrow$ display handoff&--\\
\texttt{S-04-P011}&prose paragraph / proof&Derive the sum encoding and introduce the observed addition transcript.&premise instantiation $\rightarrow$ algebraic simplification $\rightarrow$ correctness close $\rightarrow$ security setup $\rightarrow$ view instantiation&2\\
\texttt{S-04-P012}&prose paragraph / proof&Identify the opened left value as uniform despite its dependence on the input.&relation instantiation $\rightarrow$ uniformity argument $\rightarrow$ continuation&0\\
\texttt{S-04-P013}&prose paragraph / proof&Complete the uniformity argument and define the input-free addition transcript.&continuation $\rightarrow$ equivalence claim $\rightarrow$ simulator definition&0\\
\texttt{S-04-P014}&prose paragraph / proof&State that the simulated addition transcript is input independent and lead to view equality.&randomness qualification $\rightarrow$ independence claim $\rightarrow$ equality handoff&2\\
\texttt{S-04-P015}&prose paragraph / result&Conclude that the addition protocol securely implements its functionality.&security conclusion&0\\
\texttt{S-04-P016}&prose paragraph / result&Report the two-card lower bound and the single randomized operation.&efficiency assessment $\rightarrow$ lower bound reason $\rightarrow$ resource breakdown&1\\
\texttt{S-04-P017}&prose paragraph / definition&Introduce the relation that maps a dihedral value to its residue modulo the half-size domain.&definition lead in $\rightarrow$ display handoff&0\\
\texttt{S-04-P018}&prose paragraph / definition&Bound the normalization input and introduce the protocol tuple.&domain qualification $\rightarrow$ protocol lead in $\rightarrow$ display handoff&0\\
\texttt{S-04-P019}&prose paragraph / transition&Open the three-step sign-normalization procedure.&procedure announcement $\rightarrow$ list handoff&0\\
\texttt{S-04-LI008}&list item / method&Randomly exchange the two sign cosets with a two-sided rotation shuffle.&ordered step $\rightarrow$ operation application $\rightarrow$ display handoff&--\\
\texttt{S-04-P020}&prose paragraph / definition&Define the shuffled exponent using a uniform mask bit.&state characterization $\rightarrow$ randomness scope&0\\
\texttt{S-04-LI009}&list item / method&Open the randomized sign and classify it as revealed information.&ordered step $\rightarrow$ operation application $\rightarrow$ output binding $\rightarrow$ information status&--\\
\texttt{S-04-LI010}&list item / method&Rotate by the revealed sign times the half modulus.&ordered step $\rightarrow$ operation application $\rightarrow$ display handoff&--\\
\texttt{S-04-P021}&prose paragraph / proof&Show that the mask and opened sign cancel, then instantiate the security transcript.&premise setup $\rightarrow$ state relations $\rightarrow$ algebraic simplification $\rightarrow$ correctness close $\rightarrow$ security setup $\rightarrow$ view instantiation&0\\
\texttt{S-04-P022}&prose paragraph / proof&Replace the masked revealed sign by a fresh uniform bit in an equivalent transcript.&relation instantiation $\rightarrow$ distributional equivalence $\rightarrow$ simulator definition&0\\
\texttt{S-04-P023}&prose paragraph / proof&Qualify the simulator's random bit and lead to pairwise equality of views.&randomness qualification $\rightarrow$ independence claim $\rightarrow$ equality handoff&0\\
\texttt{S-04-P024}&prose paragraph / result&Conclude security and count the minimal card use and two-sided shuffle cost.&security conclusion $\rightarrow$ efficiency assessment $\rightarrow$ resource breakdown&1\\
\texttt{S-04-P025}&prose paragraph / definition&Introduce the relation that converts a sign predicate into a card value.&definition lead in $\rightarrow$ display handoff&0\\
\texttt{S-04-P026}&prose paragraph / definition&Specify the input domain and introduce a protocol that includes initialization as a subroutine.&domain qualification $\rightarrow$ protocol lead in $\rightarrow$ display handoff&0\\
\texttt{S-04-P027}&prose paragraph / transition&Open the eight-step sign-to-value procedure.&procedure announcement $\rightarrow$ list handoff&0\\
\texttt{S-04-LI011}&list item / method&Mask the input sign with a two-sided rotation shuffle.&ordered step $\rightarrow$ operation application $\rightarrow$ display handoff&--\\
\texttt{S-04-P028}&prose paragraph / definition&Bind the first uniform bit introduced by the two-sided shuffle.&variable binding $\rightarrow$ randomness scope&0\\
\texttt{S-04-LI012}&list item / method&Open the left card's masked sign and relate it to the target predicate and first mask.&ordered step $\rightarrow$ operation application $\rightarrow$ output binding $\rightarrow$ state relation $\rightarrow$ information status&--\\
\texttt{S-04-LI013}&list item / method&Rotate the right card according to the first revealed sign.&ordered step $\rightarrow$ operation application $\rightarrow$ display handoff&--\\
\texttt{S-04-LI014}&list item / method&Apply initialization to the left card.&ordered step $\rightarrow$ subroutine application $\rightarrow$ display handoff&--\\
\texttt{S-04-LI015}&list item / method&Apply a random paired reflection to correlate the two card states.&ordered step $\rightarrow$ operation application $\rightarrow$ display handoff&--\\
\texttt{S-04-P029}&prose paragraph / definition&Bind the uniform bit introduced by the flipping shuffle.&variable binding $\rightarrow$ randomness scope&0\\
\texttt{S-04-LI016}&list item / method&Open the right card's sign, characterize the xor mask, and terminate the zero branch.&ordered step $\rightarrow$ operation application $\rightarrow$ output binding $\rightarrow$ state relation $\rightarrow$ branch decision&--\\
\texttt{S-04-LI017}&list item / method&Rotate the right card by the half modulus when the branch bit is one.&ordered step $\rightarrow$ condition scope $\rightarrow$ operation application $\rightarrow$ display handoff&--\\
\texttt{S-04-LI018}&list item / method&Reflect the left card when the branch bit is one.&ordered step $\rightarrow$ condition scope $\rightarrow$ operation application $\rightarrow$ display handoff&--\\
\texttt{S-04-P030}&prose paragraph / proof&Compute the left-card output when the branch bit is zero.&procedure close $\rightarrow$ case split $\rightarrow$ state evaluation $\rightarrow$ display handoff&0\\
\texttt{S-04-P031}&prose paragraph / proof&Compute the corrected left-card output when the branch bit is one.&case split $\rightarrow$ execution trace $\rightarrow$ state evaluation $\rightarrow$ display handoff&0\\
\texttt{S-04-P032}&prose paragraph / proof&Conclude branchwise correctness and introduce the input-indexed security transcript.&correctness close $\rightarrow$ security setup $\rightarrow$ view instantiation $\rightarrow$ branch structure&0\\
\texttt{S-04-P033}&prose paragraph / proof&Express both revealed signs through independent masks and replace the conditional transcript by fresh random bits.&state relations $\rightarrow$ conditional trace scope $\rightarrow$ distributional equivalence $\rightarrow$ simulator definition&0\\
\texttt{S-04-P034}&prose paragraph / proof&Qualify the simulator bits, retain their conditional suffix rule, and derive universal equality of views.&randomness scope $\rightarrow$ conditional trace scope $\rightarrow$ independence claim $\rightarrow$ equality handoff&0\\
\texttt{S-04-P035}&prose paragraph / result&Conclude security, count cards and subroutine use, note subroutine elimination, and total randomized operations.&security conclusion $\rightarrow$ efficiency assessment $\rightarrow$ composition consequence $\rightarrow$ resource breakdown&0\\
\texttt{S-04-P036}&prose paragraph / definition&Introduce the relation that outputs whether two residues overflow the modulus.&definition lead in $\rightarrow$ display handoff&0\\
\texttt{S-04-P037}&prose paragraph / definition&Specify the operand domain and introduce the carry protocol.&domain qualification $\rightarrow$ protocol lead in&0\\
\texttt{S-04-P038}&prose paragraph / definition&Define carry as a two-card protocol composed from addition and sign-to-value subroutines.&formal definition $\rightarrow$ composition structure&0\\
\texttt{S-04-P039}&prose paragraph / transition&Open the two-stage carry procedure.&procedure announcement $\rightarrow$ list handoff&0\\
\texttt{S-04-LI019}&list item / method&Apply addition to combine the two operands.&ordered step $\rightarrow$ subroutine application $\rightarrow$ display handoff&--\\
\texttt{S-04-LI020}&list item / method&Apply sign-to-value to the encoded sum.&ordered step $\rightarrow$ subroutine application $\rightarrow$ display handoff&--\\
\texttt{S-04-P040}&prose paragraph / proof&Close immediate correctness and write the fixed three-state view for an arbitrary input.&correctness close $\rightarrow$ security setup $\rightarrow$ view instantiation&0\\
\texttt{S-04-P041}&prose paragraph / proof&Use the transcript's fixed shape to establish equal views for all inputs.&constancy reason $\rightarrow$ independence claim $\rightarrow$ equality handoff&0\\
\texttt{S-04-P042}&prose paragraph / result&Conclude security, count two subroutine calls, note inlining, and total four randomized operations.&security conclusion $\rightarrow$ efficiency assessment $\rightarrow$ composition consequence $\rightarrow$ resource breakdown&0\\
\texttt{S-04-P043}&prose paragraph / definition&Introduce the relation that encodes whether an input equals zero.&definition lead in $\rightarrow$ display handoff&0\\
\texttt{S-04-P044}&prose paragraph / definition&Specify the zero-test domain and introduce its protocol tuple.&domain qualification $\rightarrow$ protocol lead in $\rightarrow$ display handoff&0\\
\texttt{S-04-P045}&prose paragraph / transition&Open the three-step zero-test procedure.&procedure announcement $\rightarrow$ list handoff&0\\
\texttt{S-04-LI021}&list item / method&Reflect the first card across the modulus axis.&ordered step $\rightarrow$ operation application $\rightarrow$ display handoff&--\\
\texttt{S-04-LI022}&list item / transition&Mark the second position in the zero-test operation list.&ordered step $\rightarrow$ continuation&--\\
\texttt{S-04-P046}&prose paragraph / method&Begin the second step by invoking sign-to-value on the first card.&step continuation $\rightarrow$ subroutine application $\rightarrow$ continuation&0\\
\texttt{S-04-P047}&prose paragraph / transition&Complete the second-step instruction by identifying the first card as the subroutine target.&continuation $\rightarrow$ target binding $\rightarrow$ display handoff&0\\
\texttt{S-04-LI023}&list item / method&Reflect the predicate card across the unit axis.&ordered step $\rightarrow$ operation application $\rightarrow$ display handoff&--\\
\texttt{S-04-P048}&prose paragraph / proof&Prove the transformed sign detects zero and introduce the fixed security transcript.&procedure close $\rightarrow$ equivalence argument $\rightarrow$ correctness close $\rightarrow$ security setup $\rightarrow$ view instantiation&0\\
\texttt{S-04-P049}&prose paragraph / proof&Use the constant transcript to derive equal views for every pair of inputs.&constancy reason $\rightarrow$ independence claim $\rightarrow$ equality handoff&0\\
\texttt{S-04-P050}&prose paragraph / result&Conclude security, count cards and one subroutine, note inlining, and total three randomized operations.&security conclusion $\rightarrow$ efficiency assessment $\rightarrow$ composition consequence $\rightarrow$ resource breakdown&0\\
\texttt{S-04-P051}&prose paragraph / definition&Introduce the relation that encodes equality of two inputs.&definition lead in $\rightarrow$ display handoff&0\\
\texttt{S-04-P052}&prose paragraph / definition&Specify the two-input domain and introduce the equality protocol.&domain qualification $\rightarrow$ protocol lead in&0\\
\texttt{S-04-P053}&prose paragraph / definition&Define equality by composing subtraction, sign normalization, and zero testing.&formal definition $\rightarrow$ composition structure&0\\
\texttt{S-04-P054}&prose paragraph / transition&Open the three-stage equality procedure.&procedure announcement $\rightarrow$ list handoff&0\\
\texttt{S-04-LI024}&list item / method&Apply subtraction to form the difference of the two inputs.&ordered step $\rightarrow$ subroutine application $\rightarrow$ display handoff&--\\
\texttt{S-04-LI025}&list item / method&Normalize the difference into the modulus-sized domain.&ordered step $\rightarrow$ subroutine application $\rightarrow$ display handoff&--\\
\texttt{S-04-LI026}&list item / method&Run the zero-test subroutine on the normalized difference.&ordered step $\rightarrow$ subroutine application $\rightarrow$ display handoff&--\\
\texttt{S-04-P055}&prose paragraph / proof&Show the normalized difference is zero exactly for equal inputs, then instantiate the fixed transcript.&intermediate state $\rightarrow$ composition check $\rightarrow$ equivalence argument $\rightarrow$ correctness close $\rightarrow$ security setup $\rightarrow$ view instantiation&0\\
\texttt{S-04-P056}&prose paragraph / proof&Use the constant four-state transcript to establish equal views for all inputs.&constancy reason $\rightarrow$ independence claim $\rightarrow$ equality handoff&0\\
\texttt{S-04-P057}&prose paragraph / result&Conclude security, count three subroutine calls, note inlining, and total five randomized operations.&security conclusion $\rightarrow$ efficiency assessment $\rightarrow$ composition consequence $\rightarrow$ resource breakdown&0\\
\texttt{S-04-P058}&prose paragraph / definition&Introduce the relation that encodes the ordering predicate between two inputs.&definition lead in $\rightarrow$ display handoff&0\\
\texttt{S-04-P059}&prose paragraph / definition&Specify the operand domain and introduce the greater-than protocol.&domain qualification $\rightarrow$ protocol lead in&0\\
\texttt{S-04-P060}&prose paragraph / definition&Define greater-than by composing subtraction with sign-to-value conversion.&formal definition $\rightarrow$ composition structure&0\\
\texttt{S-04-P061}&prose paragraph / transition&Open the three-step greater-than procedure.&procedure announcement $\rightarrow$ list handoff&0\\
\texttt{S-04-LI027}&list item / method&Apply subtraction to encode the second input minus the first.&ordered step $\rightarrow$ subroutine application $\rightarrow$ display handoff&--\\
\texttt{S-04-LI028}&list item / method&Convert the difference sign into the complement of the desired ordering bit.&ordered step $\rightarrow$ subroutine application $\rightarrow$ display handoff&--\\
\texttt{S-04-LI029}&list item / method&Reflect the comparison card across the unit axis.&ordered step $\rightarrow$ operation application $\rightarrow$ display handoff&--\\
\texttt{S-04-P062}&prose paragraph / transition&Mark completion of the greater-than execution.&procedure close&0\\
\texttt{S-04-P063}&prose paragraph / proof&Close immediate correctness and instantiate the fixed comparison transcript.&correctness close $\rightarrow$ security setup $\rightarrow$ view instantiation&0\\
\texttt{S-04-P064}&prose paragraph / proof&Use the constant comparison transcript to establish view equality for all inputs.&constancy reason $\rightarrow$ independence claim $\rightarrow$ equality handoff&0\\
\texttt{S-04-P065}&prose paragraph / result&Conclude security, count two subroutine calls, note inlining, and total four randomized operations.&security conclusion $\rightarrow$ efficiency assessment $\rightarrow$ composition consequence $\rightarrow$ resource breakdown&0\\
\texttt{S-05-P001}&prose paragraph / conclusion&Summarize the new card model and efficient predicate protocols, then propose extending invisible-ink designs to other objects.&contribution recap $\rightarrow$ interpretive claim $\rightarrow$ future work $\rightarrow$ illustrative extension&4\\
\texttt{S-06-P001}&prose paragraph / roadmap&Open the appendix, fix cyclic-card terminology, and announce the card-specification development.&scope statement $\rightarrow$ terminology choice $\rightarrow$ roadmap&0\\
\texttt{S-06-P002}&prose paragraph / definition&Define face-up and face-down notation, then introduce the finite cyclic card set.&notation definition $\rightarrow$ set definition lead in $\rightarrow$ display handoff&0\\
\texttt{S-06-P003}&prose paragraph / definition&Introduce rotation and turning as transformation classes, then begin the degree-indexed rotation definition.&category introduction $\rightarrow$ parameter binding $\rightarrow$ definition lead in $\rightarrow$ display handoff&0\\
\texttt{S-06-P004}&prose paragraph / definition&Complete the rotation rule by distinguishing face-up and face-down cards.&case split $\rightarrow$ branch definition $\rightarrow$ branch definition&0\\
\texttt{S-06-P005}&prose paragraph / definition&Explain the opposite orientation of face-down rotation and define the turning operation by cases.&meaning explanation $\rightarrow$ exception explanation $\rightarrow$ definition lead in $\rightarrow$ case split&0\\
\texttt{S-06-P006}&prose paragraph / definition&Introduce the set containing identity, rotations, and turning.&definition lead in $\rightarrow$ display handoff&0\\
\texttt{S-06-P007}&prose paragraph / definition&Introduce the visible labels and hidden marker for cyclic cards.&definition lead in $\rightarrow$ display handoff&0\\
\texttt{S-06-P008}&prose paragraph / definition&Define how face-up labels are observed and how all other card states appear hidden.&definition lead in $\rightarrow$ case split $\rightarrow$ output assignment&0\\
\texttt{S-06-P009}&prose paragraph / definition&Assemble the card set, transformation set, symbol set, and vision function into one specification.&definition lead in $\rightarrow$ tuple assembly $\rightarrow$ display handoff&4\\
\texttt{S-06-P010}&prose paragraph / comparison&Define cyclic-card operations by separating those inherited from binary cards from those inherited from dihedral cards.&comparison setup $\rightarrow$ first correspondence $\rightarrow$ second correspondence&1\\
\bottomrule\end{longtable}\normalsize\clearpage
\section{Aggregated WADT units and rewrite suggestions}
Each block starts from one eligible baseline paragraph. The baseline description is paraphrased. The original WADT unit appears as nonexecuting source text so that the full unit remains inspectable and cannot alter this report. Footnote commands and footnote content are omitted. The proposed rewrite is rendered as LaTeX so mathematical symbols remain readable.
\clearpage\subsection{Baseline S-ABS-P001: paper-level contribution summary}
\colorbox{BaselineBlue}{\parbox{0.97\textwidth}{\Field{Purpose.}{Place the work in its research area and state the paper's technical and framework contributions.}\Field{Primary purpose.}{contribution}\Field{Narration structure.}{field\_context $\rightarrow$ prior\_work\_context $\rightarrow$ evaluation\_criterion $\rightarrow$ contribution $\rightarrow$ result\_scope $\rightarrow$ framework\_contribution}\Field{Reusable move.}{Use brief field context to prepare a layered inventory of the paper's contributions.}}}
\needspace{10\baselineskip}\subsubsection*{W-ABS-P001: same purpose}
\Field{Location.}{Section ABS, abstract\_paragraph, lines 44--67.}
\Field{Draft purpose.}{Summarize the framework, main instance, exact and approximate guarantees, scope limits, and reuse.}
\Field{Qualitative closeness.}{same purpose. Both abstracts place the work, identify the central construction, summarize the main results, and close with broader reuse.}
\Field{Limits of the analogy.}{The WADT abstract carries exact theorem conditions, two shuffle laws, numerical bounds, exclusions, and trust information that the baseline cannot supply.}
\Field{Baseline narration used.}{full.}
\Field{Complexity and jargon.}{Complexity: \Tag{high} $\rightarrow$ \Tag{high}. Jargon: \Tag{high} $\rightarrow$ \Tag{medium}.}
\textbf{Original WADT unit. Footnotes omitted.}
\begin{lstlisting}
Card-based protocols use physical cards to compute with
information-theoretic security. I define a \Rocq{} framework that treats a
finite group, its permutation action, and its shuffle distribution as first-class
parameters. The framework connects group actions to executed traces and
word-shuffle mixing. The main instance uses the action of
$\PG$ on eight projective points.
Under the uniform shuffle distribution, the instance has exact correctness,
the recovery ramp $(t,r,n)=(3,7,8)$, and perfect view privacy for
coalitions of at most three players, which interpreter soundness transports
to the executed traces. These privacy results hold for a uniform Boolean
secret, a fixed encoded representative, and passive adversaries, with the
verifier learning the secret and with active security, security after the
reveal, and composition outside the model. Because no single physical move samples a uniform group element, I
analyze a finite implementation that draws a uniform word of length 200
over five generator letters. A checked fiber-counting certificate proves that
the resulting shuffle distribution is within $L_1$ distance $2^{-40}$ of the
uniform group distribution. This word-shuffle theorem transfers to single-card
marginals and to a product with any secret prior. Transport theorems therefore bound approximate coalition-view and
executed-trace privacy by $2^{-39}$ under the word distribution.
A five-card family and two spectral siblings show which algebraic,
operational, privacy, and mixing arguments the framework reuses, the
spectral siblings importing one analytic gap premise as their only
assumption beyond the classical principles used throughout.
\end{lstlisting}
\colorbox{RewriteGreen}{\parbox{0.97\textwidth}{\textbf{Rewrite suggestion.}\par\begingroup\small
Card-based protocols use physical cards to compute with information-theoretic security. I define a \Rocq{} framework with a finite group, its permutation action, and its shuffle distribution as explicit parameters. The framework connects group actions with executed traces and word-shuffle mixing. Its main instance uses the action of $\PG$ on eight projective points. Under the uniform shuffle distribution, this instance has exact correctness and the recovery ramp $(t,r,n)=(3,7,8)$, where $t$ is the privacy threshold, $r$ is the recovery threshold, and $n$ is the number of card positions. It also has perfect view privacy for coalitions of at most three players. Interpreter soundness transfers this privacy to executed traces. These privacy results assume a uniform Boolean secret, a fixed encoded representative, and passive adversaries. The verifier learns the secret. Active security, security after the reveal, and composition are outside the model. No single physical move samples a uniform group element. I therefore analyze a finite implementation that draws a uniform word of length 200 over five generator letters. A checked fiber-counting certificate counts the words that evaluate to each group element. It proves that the resulting word distribution is within $L_1$ distance $2^{-40}$ of the uniform group distribution. The theorem transfers to single-card marginals and to a product with any secret prior. Transport theorems then bound approximate coalition-view and executed-trace privacy by $2^{-39}$ under the word distribution. A five-card family and two related instances with spectral mixing proofs show which algebraic, operational, privacy, and mixing arguments the framework reuses. The two spectral instances import an analytic gap premise, an external assumption used only for their mixing proofs, as their only assumption beyond the classical principles used throughout.
\par\endgroup}}
\Field{Why this rewrite.}{The rewrite follows the full abstract sequence from field context through the framework, exact results, finite implementation, bounds, limits, and reuse. Shorter sentences explain the two thresholds, the card-position count, the counting certificate, and the imported premise without changing any condition or number.}
\Field{Terminology actions.}{Keep: information-theoretic security, uniform group distribution, word distribution, view privacy and trace privacy, Rocq | Explain: recovery ramp, fiber-counting certificate, analytic gap premise | Replace: first-class parameters -> explicit parameters, spectral siblings -> two related instances with spectral mixing proofs}
\par\smallskip
\clearpage\subsection{Baseline S-01-P001: field definition and accessibility context}
\colorbox{BaselineBlue}{\parbox{0.97\textwidth}{\Field{Purpose.}{Define the field and explain why its physical medium helps readers and learners.}\Field{Primary purpose.}{background}\Field{Narration structure.}{field\_definition $\rightarrow$ implementation\_contrast $\rightarrow$ accessibility\_rationale $\rightarrow$ application\_context}\Field{Reusable move.}{Establish the terms and context needed by the next technical move.}}}
\needspace{10\baselineskip}\subsubsection*{W-01-P001: analogous purpose}
\Field{Location.}{Section 01, prose\_paragraph, lines 71--78.}
\Field{Draft purpose.}{Introduce card-based cryptography and give the minimum historical context for the paper.}
\Field{Qualitative closeness.}{analogous purpose. Both paragraphs define card-based cryptography and establish context for the next technical problem.}
\Field{Limits of the analogy.}{The baseline contrasts implementations, motivates the medium through accessibility, and closes with applications. The WADT paragraph instead explains shuffles, reveals, and observation security before giving a short protocol history.}
\Field{Baseline narration used.}{adapted.}
\Field{Complexity and jargon.}{Complexity: \Tag{low} $\rightarrow$ \Tag{low}. Jargon: \Tag{medium} $\rightarrow$ \Tag{medium}.}
\textbf{Original WADT unit. Footnotes omitted.}
\begin{lstlisting}
Card-based cryptography uses face-down physical cards to compute functions
without computational assumptions. A shuffle hides a permutation, while a
reveal exposes selected card faces. Information-theoretic security requires
the resulting observation distribution to hide the protected secret. Den Boer's
five-card protocol began this line of work~\cite{denBoer1989}. Later protocols
reduced card counts and broadened the class of supported
computations~\cite{MizukiSone2009,MizukiSone2012,KochWalzerHartel2015}.
\end{lstlisting}
\colorbox{RewriteGreen}{\parbox{0.97\textwidth}{\textbf{Rewrite suggestion.}\par\begingroup\small
Card-based cryptography uses face-down physical cards to compute functions without computational assumptions. A shuffle hides a permutation. A reveal exposes selected card faces. Information-theoretic security requires the resulting observation distribution to hide the protected secret. Den Boer's five-card protocol computes Boolean conjunction and began this line of work~\cite{denBoer1989}. Later protocols used fewer cards and supported more classes of computations~\cite{MizukiSone2009,MizukiSone2012,KochWalzerHartel2015}.
\par\endgroup}}
\Field{Why this rewrite.}{The rewrite adapts the baseline's field-definition opening and keeps the WADT paragraph's operational explanation and short history. It separates the shuffle and reveal actions and identifies the five-card protocol's Boolean computation.}
\Field{Terminology actions.}{Keep: information-theoretic security, shuffle, observation distribution | Explain: five-card protocol}
\par\smallskip
\needspace{10\baselineskip}\subsubsection*{W-09-P001: partial analogue}
\Field{Location.}{Section 09, prose\_paragraph, lines 1796--1802.}
\Field{Draft purpose.}{Place the five-card trick historically and summarize later extensions.}
\Field{Qualitative closeness.}{partial analogue. The baseline contributes only the opening move that establishes field context before a narrower technical comparison.}
\Field{Limits of the analogy.}{The baseline defines card-based cryptography and motivates its physical accessibility. This paragraph gives a historical origin, one passive-security guarantee, and later protocol extensions.}
\Field{Baseline narration used.}{partial.}
\Field{Complexity and jargon.}{Complexity: \Tag{medium} $\rightarrow$ \Tag{low}. Jargon: \Tag{medium} $\rightarrow$ \Tag{low}.}
\textbf{Original WADT unit. Footnotes omitted.}
\begin{lstlisting}
Den Boer introduced the first card-based cryptographic protocol with the
five-card trick~\cite{denBoer1989,KochFUN2021}. The protocol uses five cards
to compute the AND of two secret bits with perfect security against passive
participants. Although the security is already perfect against passive participants,
later protocols reduced the number of cards and supported other
functions~\cite{MizukiSone2009,MizukiSone2012,KochWalzerHartel2015}.
\end{lstlisting}
\colorbox{RewriteGreen}{\parbox{0.97\textwidth}{\textbf{Rewrite suggestion.}\par\begingroup\small
Den Boer introduced the first card-based cryptographic protocol with the five-card trick~\cite{denBoer1989,KochFUN2021}. It uses five cards to compute the AND of two secret bits with perfect security against passive participants. Later protocols reduced the number of cards and supported other functions~\cite{MizukiSone2009,MizukiSone2012,KochWalzerHartel2015}.
\par\endgroup}}
\Field{Why this rewrite.}{Uses only the baseline's field-context opening. It preserves the historical claim, five-card AND result, passive-security scope, and later extensions.}
\Field{Terminology actions.}{Keep: perfect security, passive participants, AND | Explain: card-based cryptographic protocol}
\par\smallskip
\clearpage\subsection{Baseline S-01-P002: multi-valued computation limitation}
\colorbox{BaselineBlue}{\parbox{0.97\textwidth}{\Field{Purpose.}{Show that the existing design handles a narrow function class and motivate a new design for efficient missing cases.}\Field{Primary purpose.}{problem}\Field{Narration structure.}{application\_examples $\rightarrow$ multi\_value\_need $\rightarrow$ prior\_design $\rightarrow$ efficiency\_question $\rightarrow$ limitation $\rightarrow$ open\_problem $\rightarrow$ solution\_preview}\Field{Reusable move.}{Move from an existing capability to the limitation that motivates a new construction.}}}
\needspace{10\baselineskip}\subsubsection*{W-01-P002: analogous purpose}
\Field{Location.}{Section 01, prose\_paragraph, lines 79--91.}
\Field{Draft purpose.}{Contrast bounded protocol search with generic proof-assistant theorems and state what the broader method adds.}
\Field{Qualitative closeness.}{analogous purpose. Both paragraphs expose a limit of an earlier approach and use that limit to motivate a broader construction.}
\Field{Limits of the analogy.}{The baseline concerns function coverage and protocol efficiency. The WADT paragraph concerns the quantifier scope of bounded verification and adds formal privacy and mixing outputs.}
\Field{Baseline narration used.}{adapted.}
\Field{Complexity and jargon.}{Complexity: \Tag{medium} $\rightarrow$ \Tag{medium}. Jargon: \Tag{high} $\rightarrow$ \Tag{medium}.}
\textbf{Original WADT unit. Footnotes omitted.}
\begin{lstlisting}
Mechanized analysis must connect the physical operations to their probability
distributions. Koch, Schrempp, and Kirsten use software bounded model checking to
search finite protocol spaces and check lower bounds~\cite{Koch2019}.
Their method automates checks over a bounded state space, and a bound
checked for one finite protocol space says nothing about every
$t$-transitive group action. The development in this article instead uses a proof assistant whose generic
theorems quantify over the group, the action, and the shuffle distribution,
and it states reusable dealer distributions, participant views, and
executed traces. It combines a group-parametric model with perfect trace
privacy and a checked word-shuffle mixing bound.
\end{lstlisting}
\colorbox{RewriteGreen}{\parbox{0.97\textwidth}{\textbf{Rewrite suggestion.}\par\begingroup\small
Mechanized analysis must connect physical operations with their probability distributions. Koch, Schrempp, and Kirsten use bounded model checking to search finite protocol spaces and check lower bounds~\cite{Koch2019}. This method checks a bounded state space, meaning one finite collection of protocol states. The target class consists of $t$-transitive group actions. A bound for one such space does not establish a result for all such actions. The Rocq formalization instead uses a proof assistant whose generic theorems quantify over the group, action, and shuffle distribution. It defines reusable dealer distributions for sampling encoded initial decks, participant views, and executed traces. The resulting group-parametric model proves perfect trace privacy. It also checks a word-shuffle mixing bound that connects a finite generator process with the ideal shuffle law.
\par\endgroup}}
\Field{Why this rewrite.}{The rewrite adapts the baseline's capability-to-limitation sequence. It states the finite quantifier limit before the Rocq response and explains the dealer distribution and mixing bound in direct language.}
\Field{Terminology actions.}{Keep: bounded model checking, proof assistant, group-parametric model, executed trace | Explain: bounded state space, dealer distribution, word-shuffle mixing bound | Replace: the development -> the Rocq formalization, says nothing about every action -> does not establish a result for all such actions}
\par\smallskip
\needspace{10\baselineskip}\subsubsection*{W-09-P002: analogous purpose}
\Field{Location.}{Section 09, prose\_paragraph, lines 1803--1809.}
\Field{Draft purpose.}{Contrast bounded model checking with the reusable Rocq development.}
\Field{Qualitative closeness.}{analogous purpose. Both units expose the restricted reach of an earlier approach and motivate a broader formal construction.}
\Field{Limits of the analogy.}{The baseline targets a limited function class and protocol efficiency. This unit compares verification scope and reusable theorem infrastructure.}
\Field{Baseline narration used.}{adapted.}
\Field{Complexity and jargon.}{Complexity: \Tag{medium} $\rightarrow$ \Tag{low}. Jargon: \Tag{high} $\rightarrow$ \Tag{medium}.}
\textbf{Original WADT unit. Footnotes omitted.}
\begin{lstlisting}
Koch, Schrempp, and Kirsten encode bounded protocol spaces for software
bounded model checking~\cite{Koch2019}. Their method searches for protocols
and checks finite lower bounds. The present development, however, proves reusable group-action,
execution, privacy, and mixing lemmas in \Rocq{}. The Kim
instance also connects a biased shuffle model to a concrete word-shuffle
bound~\cite{KimCetinkaya2025}.
\end{lstlisting}
\colorbox{RewriteGreen}{\parbox{0.97\textwidth}{\textbf{Rewrite suggestion.}\par\begingroup\small
Koch, Schrempp, and Kirsten use software bounded model checking over encoded finite protocol spaces~\cite{Koch2019}. Their method searches for protocols and checks finite lower bounds. The \Rocq{} formalization instead proves reusable lemmas for group actions, execution, privacy, and mixing. The Kim instance also connects a biased shuffle model with a concrete word-shuffle bound~\cite{KimCetinkaya2025}.
\par\endgroup}}
\Field{Why this rewrite.}{Adapts the baseline's prior-capability-to-response sequence. It preserves the finite scope, four reusable theorem families, and the Kim instance.}
\Field{Terminology actions.}{Keep: bounded model checking, word-shuffle bound | Explain: bounded protocol space, reusable group-action lemmas | Replace: present development -> Rocq formalization}
\par\smallskip
\needspace{10\baselineskip}\subsubsection*{W-09-P006: partial analogue}
\Field{Location.}{Section 09, prose\_paragraph, lines 1846--1855.}
\Field{Draft purpose.}{Position exact fiber counting against spectral and coupling approaches to mixing.}
\Field{Qualitative closeness.}{partial analogue. The baseline contributes only the comparison move that states prior capabilities before explaining the paper's different response.}
\Field{Limits of the analogy.}{The baseline motivates a new protocol from a function-class limitation. This related-work paragraph surveys mixing methods and justifies exact fiber counting without claiming those methods are deficient.}
\Field{Baseline narration used.}{partial.}
\Field{Complexity and jargon.}{Complexity: \Tag{medium} $\rightarrow$ \Tag{low}. Jargon: \Tag{high} $\rightarrow$ \Tag{medium}.}
\textbf{Original WADT unit. Footnotes omitted.}
\begin{lstlisting}
Machine-checked probability libraries contain results that can support
mixing arguments. Karayel formalizes expander graphs, Cheeger's inequality,
and random-walk tail bounds in Isabelle/HOL~\cite{Karayel2023}. Barthe,
Gr\'egoire, Hsu, and Strub develop a relational program logic for couplings
and apply it to rapid-mixing examples~\cite{BartheEtAl2017}. Standard
mathematical treatments relate spectral information to convergence of
finite Markov chains~\cite{LevinPeresWilmer2009}. The PGL proof, however, uses another route. It counts the exact fibers of one finite word length and checks the
resulting integer inequality. This choice gives a concrete bound without a
general spectral development.
\end{lstlisting}
\colorbox{RewriteGreen}{\parbox{0.97\textwidth}{\textbf{Rewrite suggestion.}\par\begingroup\small
Machine-checked probability libraries provide results that can support mixing arguments. Karayel formalizes expander graphs, Cheeger's inequality, and random-walk tail bounds in Isabelle/HOL~\cite{Karayel2023}. Barthe, Gr\'egoire, Hsu, and Strub develop a relational program logic for couplings and apply it to rapid-mixing examples~\cite{BartheEtAl2017}. Standard mathematical treatments connect spectral information with convergence of finite Markov chains~\cite{LevinPeresWilmer2009}. The $\PG$ proof uses exact fiber counting instead. It counts the fibers at one finite word length and checks the resulting integer inequality. This method gives a concrete bound without requiring a general spectral development.
\par\endgroup}}
\Field{Why this rewrite.}{Uses only the baseline's prior-capability-to-response move. It preserves all three method families and states the exact-counting choice without criticizing the alternatives.}
\Field{Terminology actions.}{Keep: fiber counting, integer inequality | Explain: Cheeger inequality, coupling | Replace: uses another route -> uses exact fiber counting}
\par\smallskip
\clearpage\subsection{Baseline S-01-P003: two-part contribution account}
\colorbox{BaselineBlue}{\parbox{0.97\textwidth}{\Field{Purpose.}{Present the new card construction and the broader protocol model as the paper's two main additions.}\Field{Primary purpose.}{contribution}\Field{Narration structure.}{artifact\_contribution $\rightarrow$ protocol\_results $\rightarrow$ comparison\_summary $\rightarrow$ model\_contribution $\rightarrow$ reuse\_scope $\rightarrow$ future\_boundary}\Field{Reusable move.}{Use brief context to prepare a layered inventory of concrete additions.}}}
\needspace{10\baselineskip}\subsubsection*{W-01-P004: analogous purpose}
\Field{Location.}{Section 01, prose\_paragraph, lines 104--106.}
\Field{Draft purpose.}{Announce the central framework contribution and its two principal checked results.}
\Field{Qualitative closeness.}{analogous purpose. Both units state the paper's contribution at a high level before enumerating concrete results.}
\Field{Limits of the analogy.}{The baseline gives two contributions in a developed paragraph. The WADT unit is a short lead-in to two theorem statements.}
\Field{Baseline narration used.}{adapted.}
\Field{Complexity and jargon.}{Complexity: \Tag{low} $\rightarrow$ \Tag{low}. Jargon: \Tag{medium} $\rightarrow$ \Tag{low}.}
\textbf{Original WADT unit. Footnotes omitted.}
\begin{lstlisting}
Overall, this work contributes a group-parametric framework whose main
instance carries two machine-checked security results.
\end{lstlisting}
\colorbox{RewriteGreen}{\parbox{0.97\textwidth}{\textbf{Rewrite suggestion.}\par\begingroup\small
I contribute a group-parametric framework. Its primary $\PG$ instance has two machine-checked security results.
\par\endgroup}}
\Field{Why this rewrite.}{The rewrite adapts the baseline's contribution announcement and makes the authorial action explicit. It explains the main instance by naming it and separates the framework contribution from the result count.}
\Field{Terminology actions.}{Keep: group-parametric framework, machine-checked | Explain: main instance}
\par\smallskip
\needspace{10\baselineskip}\subsubsection*{W-01-LI001: partial analogue}
\Field{Location.}{Section 01, list\_item, lines 139--142.}
\Field{Draft purpose.}{Identify the common formal interface contributed by the paper.}
\Field{Qualitative closeness.}{partial analogue. Both units name a framework-level contribution before its concrete results.}
\Field{Limits of the analogy.}{The baseline combines two major contributions in one paragraph. This list item isolates the common Rocq interface and inventories its components.}
\Field{Baseline narration used.}{partial.}
\Field{Complexity and jargon.}{Complexity: \Tag{medium} $\rightarrow$ \Tag{medium}. Jargon: \Tag{high} $\rightarrow$ \Tag{medium}.}
\textbf{Original WADT unit. Footnotes omitted.}
\begin{lstlisting}
I give a common \Rocq{} interface for a finite group, its permutation
  action, the shuffle distribution, reconstruction, participant views, and executed
  traces.
\end{lstlisting}
\colorbox{RewriteGreen}{\parbox{0.97\textwidth}{\textbf{Rewrite suggestion.}\par\begingroup\small
I give a common \Rocq{} interface for a finite group, its permutation action, the shuffle distribution, reconstruction, participant views of the cards held by a chosen group of players, and executed traces.
\par\endgroup}}
\Field{Why this rewrite.}{The rewrite uses only the partial analogue's contribution-first move. It retains the six interface components and explains participant views without changing the list.}
\Field{Terminology actions.}{Keep: Rocq interface, shuffle distribution, executed trace | Explain: participant view}
\par\smallskip
\needspace{10\baselineskip}\subsubsection*{W-01-LI002: partial analogue}
\Field{Location.}{Section 01, list\_item, lines 142--148.}
\Field{Draft purpose.}{State the formal verification contribution for the five-card family.}
\Field{Qualitative closeness.}{partial analogue. Both units state a concrete contribution as evidence for the broader model.}
\Field{Limits of the analogy.}{The baseline gives a paper-level contribution account. This item combines two protocols, one bias parameter, and three verified security properties.}
\Field{Baseline narration used.}{partial.}
\Field{Complexity and jargon.}{Complexity: \Tag{high} $\rightarrow$ \Tag{medium}. Jargon: \Tag{high} $\rightarrow$ \Tag{medium}.}
\textbf{Original WADT unit. Footnotes omitted.}
\begin{lstlisting}
I verify the den Boer~\cite{denBoer1989} and the Kim and
  \c{C}etinkaya~\cite{KimCetinkaya2025} five-card protocols in \Rocq{} as
  one profile at bias $\varepsilon$, in the sense of
  Section~\ref{sec:framework}: executed-run correctness of the conjunction
  output, executed-trace secrecy for a single corrupted player, and
  input-hiding given the announced output.
\end{lstlisting}
\colorbox{RewriteGreen}{\parbox{0.97\textwidth}{\textbf{Rewrite suggestion.}\par\begingroup\small
I verify the den Boer~\cite{denBoer1989} and Kim and \c{C}etinkaya~\cite{KimCetinkaya2025} five-card protocols in \Rocq{} as one profile at bias $\varepsilon$, in the sense of Section~\ref{sec:framework}. A profile packages one instance's data and proofs. I prove executed-run correctness for the conjunction output, executed-trace privacy against one corrupted player, and input hiding given the announced output.
\par\endgroup}}
\Field{Why this rewrite.}{The rewrite uses the partial analogue's construction-then-properties move. It explains the profile and adversary scope, replaces secrecy with the canonical privacy term, and gives each verified property a clear role.}
\Field{Terminology actions.}{Keep: bias parameter, executed-run correctness, input hiding | Explain: profile, single corrupted player | Replace: executed-trace secrecy -> executed-trace privacy}
\par\smallskip
\needspace{10\baselineskip}\subsubsection*{W-01-LI003: partial analogue}
\Field{Location.}{Section 01, list\_item, lines 148--152.}
\Field{Draft purpose.}{State the exact leakage and finite-cut results for the five-card family.}
\Field{Qualitative closeness.}{partial analogue. Both units present a concrete result as part of the paper's contribution.}
\Field{Limits of the analogy.}{The baseline gives a broad contribution summary. This item reports exact leakage for all reveal patterns, a zero-bias recovery, and one checked numerical bound.}
\Field{Baseline narration used.}{partial.}
\Field{Complexity and jargon.}{Complexity: \Tag{medium} $\rightarrow$ \Tag{medium}. Jargon: \Tag{medium} $\rightarrow$ \Tag{medium}.}
\textbf{Original WADT unit. Footnotes omitted.}
\begin{lstlisting}
I prove exact mutual-information leakage values for every reveal
  pattern of the five-card family, recover den Boer's perfect dealing
  bound at bias zero, and check in the kernel a seven-cut single-card
  bound of $2^{-40}$ at bias $1/100$.
\end{lstlisting}
\colorbox{RewriteGreen}{\parbox{0.97\textwidth}{\textbf{Rewrite suggestion.}\par\begingroup\small
I prove the exact mutual-information leakage for every reveal pattern in the five-card family. At bias zero, I recover den Boer's perfect dealing bound. At bias $1/100$, I kernel-check that after seven cuts every single-card endpoint distribution is within $L_1$ distance $2^{-40}$ of uniform.
\par\endgroup}}
\Field{Why this rewrite.}{The rewrite uses the partial analogue's general-to-special sequence. It separates the exhaustive leakage result and the zero-bias case, then explains the checked single-card bound using the source-grounded endpoint statement from the local context.}
\Field{Terminology actions.}{Keep: mutual-information leakage, bias, kernel-checked | Explain: single-card bound}
\par\smallskip
\needspace{10\baselineskip}\subsubsection*{W-01-LI004: partial analogue}
\Field{Location.}{Section 01, list\_item, lines 152--155.}
\Field{Draft purpose.}{State the exact ideal-shuffle results for the main instance.}
\Field{Qualitative closeness.}{partial analogue. Both units identify a central construction and summarize the results it enables.}
\Field{Limits of the analogy.}{The baseline covers a new card model and protocol class. This item is limited to four ideal-shuffle guarantees for the main instance.}
\Field{Baseline narration used.}{partial.}
\Field{Complexity and jargon.}{Complexity: \Tag{medium} $\rightarrow$ \Tag{low}. Jargon: \Tag{high} $\rightarrow$ \Tag{medium}.}
\textbf{Original WADT unit. Footnotes omitted.}
\begin{lstlisting}
I prove correctness, the recovery ramp, coalition-view privacy, and
  executed-trace secrecy for the $\PG$ instance under the uniform shuffle
  distribution.
\end{lstlisting}
\colorbox{RewriteGreen}{\parbox{0.97\textwidth}{\textbf{Rewrite suggestion.}\par\begingroup\small
For the $\PG$ instance under the uniform shuffle distribution, I prove correctness and the recovery ramp. The ramp records the privacy threshold, the recovery threshold, and the number of card positions. I also prove coalition-view privacy and executed-trace privacy.
\par\endgroup}}
\Field{Why this rewrite.}{The rewrite uses only the partial analogue's instance-first move. It explains the two thresholds and the card-position count in the recovery ramp, uses the canonical privacy term, and separates the structural and privacy results.}
\Field{Terminology actions.}{Keep: uniform shuffle distribution, coalition-view privacy | Explain: recovery ramp | Replace: executed-trace secrecy -> executed-trace privacy}
\par\smallskip
\needspace{10\baselineskip}\subsubsection*{W-01-LI005: partial analogue}
\Field{Location.}{Section 01, list\_item, lines 155--167.}
\Field{Draft purpose.}{State the finite mixing bound and its two transfer consequences.}
\Field{Qualitative closeness.}{partial analogue. Both units state a paper contribution and identify the broader value of the result.}
\Field{Limits of the analogy.}{The baseline gives high-level contributions. This item states a precise distributional theorem and adds two transfer results.}
\Field{Baseline narration used.}{partial.}
\Field{Complexity and jargon.}{Complexity: \Tag{medium} $\rightarrow$ \Tag{low}. Jargon: \Tag{high} $\rightarrow$ \Tag{medium}.}
\textbf{Original WADT unit. Footnotes omitted.}
\begin{lstlisting}
I prove an $L_1$ bound between a 200-letter generator-word distribution
  and the uniform group distribution. The proof
  also transfers the bound to card
  endpoints and secret-prior product distributions.
\end{lstlisting}
\colorbox{RewriteGreen}{\parbox{0.97\textwidth}{\textbf{Rewrite suggestion.}\par\begingroup\small
I prove an $L_1$ bound between a 200-letter word distribution and the uniform group distribution. I also transfer this bound to the distribution of the card at each endpoint and to product distributions with any secret prior.
\par\endgroup}}
\Field{Why this rewrite.}{The rewrite uses the partial analogue's bound-then-transfer move. It excludes the footnote and explains card endpoints directly from the body claim while preserving both transported objects.}
\Field{Terminology actions.}{Keep: L1 distance, word distribution, secret prior | Explain: card endpoint}
\par\smallskip
\needspace{10\baselineskip}\subsubsection*{W-04-P010: analogous purpose}
\Field{Location.}{Section 04, prose\_paragraph, lines 808--822.}
\Field{Draft purpose.}{Synthesize the guarantees shared by the two five-card protocols and point to the instance table.}
\Field{Qualitative closeness.}{analogous purpose. Both units present a compact multi-part contribution account and then direct the reader to the broader construction landscape.}
\Field{Limits of the analogy.}{The baseline introduces a new card construction and model. This unit consolidates one family's reused proofs and security results.}
\Field{Baseline narration used.}{adapted.}
\Field{Complexity and jargon.}{Complexity: \Tag{high} $\rightarrow$ \Tag{medium}. Jargon: \Tag{high} $\rightarrow$ \Tag{medium}.}
\textbf{Original WADT unit. Footnotes omitted.}
\begin{lstlisting}
Because the two members share one executed program, correctness
transfers verbatim. The Kim program is the den Boer program, and its recovery
theorem is the den Boer proof reused. The committed inputs of the
two players realize the function-evaluation flow of
Section~\ref{sec:framework} with den Boer's input encoding. Den Boer's
uniform cut gives the exact endpoint distribution after one shuffle, and a
single corrupted player's executed trace leaves the secret's conditional
entropy equal to its plain
entropy.
Table~\ref{tab:instances} in Section~\ref{sec:instances} records the full
landscape.
\end{lstlisting}
\colorbox{RewriteGreen}{\parbox{0.97\textwidth}{\textbf{Rewrite suggestion.}\par\begingroup\small
The two family members share one executed program, so correctness carries over without change. The Kim program is the den Boer program, and its recovery theorem reuses the den Boer proof. The committed inputs of the two players realize the function-evaluation flow of Section~\ref{sec:framework} with den Boer's input encoding. Den Boer's uniform cut gives the exact endpoint distribution after one shuffle. For one corrupted player, the executed trace leaves the secret's conditional entropy equal to its plain entropy. Table~\ref{tab:instances} in Section~\ref{sec:instances} gives the full comparison.
\par\endgroup}}
\Field{Why this rewrite.}{Adapts the baseline's layered contribution sequence to proof reuse, framework realization, endpoint behavior, trace secrecy, and the comparison-table handoff.}
\Field{Terminology actions.}{Keep: executed trace, conditional entropy, uniform cut | Explain: correctness transfers, plain entropy | Replace: transfers verbatim -> carries over without change, full landscape -> full comparison}
\par\smallskip
\needspace{10\baselineskip}\subsubsection*{W-05-P001: analogous purpose}
\Field{Location.}{Section 05, prose\_paragraph, lines 831--845.}
\Field{Draft purpose.}{Open the main instance and summarize its exact and finite-shuffle guarantees.}
\Field{Qualitative closeness.}{analogous purpose. Both units present a main construction and a layered account of its technical contributions before the details.}
\Field{Limits of the analogy.}{The baseline introduction presents two paper-wide contributions. The WADT unit opens one instance section and includes exact thresholds, a sharp counterboundary, and a later approximate theorem.}
\Field{Baseline narration used.}{adapted.}
\Field{Complexity and jargon.}{Complexity: \Tag{high} $\rightarrow$ \Tag{medium}. Jargon: \Tag{high} $\rightarrow$ \Tag{medium}.}
\textbf{Original WADT unit. Footnotes omitted.}
\begin{lstlisting}
This section constructs the paper's main eight-player instance, a model of
the profile specification: it supplies the record fields of
Section~\ref{sec:framework} and discharges their obligations. Each player
receives one card, which acts as a physical share of the secret. Under the
fixed representative dealer and an independent uniform $\PG$ shuffle, every
passive coalition of at most three players has perfect view and
executed-trace privacy. Every reveal of seven positions determines the
secret class, while every reveal of at most six positions remains compatible
with both classes. A four-position view can leak information, so the privacy
cutoff at three is sharp. The resulting recovery ramp is
$(t,r,n)=(3,7,8)$. Decoding the full eight-card arrangement is correct for
every allowed group element. Theorem~A collects these exact guarantees, and
Theorem~B later replaces the uniform group element with a uniformly sampled
200-letter generator word and bounds the resulting privacy error.
\end{lstlisting}
\colorbox{RewriteGreen}{\parbox{0.97\textwidth}{\textbf{Rewrite suggestion.}\par\begingroup\small
I construct the paper's main eight-player instance in this section. It is an instance of the profile specification, which packages the data and proof obligations for one protocol instance. The construction supplies the record fields of Section~\ref{sec:framework} and proves their obligations. Each player receives one card as a physical share of the secret. Under the fixed representative dealer and an independent uniform $\PG$ shuffle, every passive coalition of at most three players has perfect view and executed-trace privacy. Every reveal of seven positions determines the secret class. Every reveal of at most six positions remains compatible with both classes. A view of four positions can leak information, so the privacy cutoff at three is sharp. These thresholds give the recovery ramp $(t,r,n)=(3,7,8)$. Decoding the full eight-card arrangement is correct for every allowed group element. Theorem~A collects these exact guarantees. Theorem~B later replaces the uniform group element with a uniformly sampled 200-letter generator word and bounds the resulting privacy error. To build the instance that supports these results, I first specify its eight card positions.
\par\endgroup}}
\Field{Why this rewrite.}{The rewrite retains the layered guarantee inventory and every exact condition. It connects the privacy and recovery facts through the recovery ramp, then explains how Theorems A and B divide the exact and finite-word results. The final sentence motivates the projective-line construction in the next unit.}
\Field{Terminology actions.}{Keep: fixed representative dealer, recovery ramp, executed-trace privacy | Explain: profile specification -> packages the data and proof obligations for one protocol instance, sharp privacy cutoff -> privacy through three positions is sharp because four positions can leak | Replace: model of the profile specification -> instance of the profile specification}
\par\smallskip
\needspace{10\baselineskip}\subsubsection*{W-08-P001: partial analogue}
\Field{Location.}{Section 08, prose\_paragraph, lines 1719--1732.}
\Field{Draft purpose.}{Explain which framework arguments transfer across five instances and how to read the coverage table.}
\Field{Qualitative closeness.}{partial analogue. The baseline supplies the move that separates a broader model contribution from concrete results and then states the model's reuse scope.}
\Field{Limits of the analogy.}{The baseline is an introduction-level contribution inventory. The WADT unit interprets a five-instance coverage table, distinguishes shared theorems from instance-specific endpoint evidence, and limits the meaning of its privacy and trace columns.}
\Field{Baseline narration used.}{partial.}
\Field{Complexity and jargon.}{Complexity: \Tag{high} $\rightarrow$ \Tag{medium}. Jargon: \Tag{high} $\rightarrow$ \Tag{medium}.}
\textbf{Original WADT unit. Footnotes omitted.}
\begin{lstlisting}
The framework's value shows in which arguments transfer across instances
unchanged and which need instance-specific finite or spectral evidence.
Beyond the instances of Sections~\ref{sec:fivecard} and~\ref{sec:pgl}, the
$S_5$ and $S_5\times S_5$ instances exercise the spectral part of the
framework, and Table~\ref{tab:instances} records the proved coverage
across all five. For example, the trace-lifting theorem of
Section~\ref{sec:framework} is consumed verbatim by all five instances,
while each instance supplies its own endpoint bound. Perfect privacy in
the table refers to the dealer distribution defined by each
instance.
The trace column records no second security property: it records the
interpreter-soundness transport of the same view privacy to the executed
trace of the stated corrupted players.
\end{lstlisting}
\colorbox{RewriteGreen}{\parbox{0.97\textwidth}{\textbf{Rewrite suggestion.}\par\begingroup\small
The framework shows which arguments transfer unchanged across instances and which require finite or spectral evidence for a specific instance. In addition to the instances in Sections~\ref{sec:fivecard} and~\ref{sec:pgl}, the $S_5$ and $S_5\times S_5$ instances use the spectral part of the framework. Table~\ref{tab:instances} records the formally proved result families for all five instances. The trace-lifting theorem of Section~\ref{sec:framework} is used unchanged by all five instances. Each instance supplies its own endpoint bound. Perfect privacy in the table refers to the dealer distribution defined by that instance. The trace column does not record a second security property. It records interpreter soundness carrying the same view privacy to the executed traces of the stated corrupted players.
\par\endgroup}}
\Field{Why this rewrite.}{Uses the baseline contribution-and-reuse move only to separate the shared theorem from instance-specific evidence, then follows the WADT table-interpretation contract.}
\Field{Terminology actions.}{Keep: trace-lifting theorem, endpoint bound, dealer distribution | Explain: trace transport -> interpreter soundness carries static view privacy to executed traces, proved coverage -> the formally proved result families in the table | Replace: is consumed verbatim -> is used unchanged}
\par\smallskip
\needspace{10\baselineskip}\subsubsection*{W-09-P007: partial analogue}
\Field{Location.}{Section 09, prose\_paragraph, lines 1856--1866.}
\Field{Draft purpose.}{Position the proof architecture among cryptographic proof assistants and identify the reused operational foundation.}
\Field{Qualitative closeness.}{partial analogue. The baseline supplies the contribution move of separating an inherited foundation from the new extension built on it.}
\Field{Limits of the analogy.}{The baseline presents two paper contributions. This paragraph primarily positions tools and prior work.}
\Field{Baseline narration used.}{partial.}
\Field{Complexity and jargon.}{Complexity: \Tag{high} $\rightarrow$ \Tag{medium}. Jargon: \Tag{high} $\rightarrow$ \Tag{medium}.}
\textbf{Original WADT unit. Footnotes omitted.}
\begin{lstlisting}
Proof-assistant frameworks for cryptography include CertiCrypt, CryptHOL, and
SSProve~\cite{CertiCrypt,CryptHOL,SSProve}. They support program semantics,
probabilistic reasoning, or modular cryptographic proofs at different
levels. The present work, however, states information-theoretic bounds rather
than game-based reductions, and it therefore uses the MathComp and
infotheo libraries for finite algebra, probability, independence, and
entropy~\cite{MathComp,InfoTheo}.
The earlier FORTE development supplies the interpreter and trace semantics
used here~\cite{WengEtAl2025}. The current paper connects that operational
foundation to group-parametric card shuffles and coalition traces.
\end{lstlisting}
\colorbox{RewriteGreen}{\parbox{0.97\textwidth}{\textbf{Rewrite suggestion.}\par\begingroup\small
Proof-assistant frameworks for cryptography include CertiCrypt, CryptHOL, and SSProve~\cite{CertiCrypt,CryptHOL,SSProve}. They support program semantics, probabilistic reasoning, or modular cryptographic proofs at different levels. This paper proves information-theoretic bounds rather than game-based reductions. It therefore uses MathComp and infotheo for finite algebra, probability, independence, and entropy~\cite{MathComp,InfoTheo}. The earlier FORTE development provides the interpreter and trace semantics used here~\cite{WengEtAl2025}. This paper extends those semantics with group-parametric card shuffles and coalition traces.
\par\endgroup}}
\Field{Why this rewrite.}{Uses only the baseline's inherited-foundation-to-extension move. It preserves the framework inventory, proof-style contrast, library roles, reused semantics, and new extension.}
\Field{Terminology actions.}{Keep: information-theoretic bounds, trace semantics, coalition traces | Explain: game-based reductions | Replace: operational foundation -> interpreter and trace semantics}
\par\smallskip
\clearpage\subsection{Baseline S-02-P001: unified-model section roadmap}
\colorbox{BaselineBlue}{\parbox{0.97\textwidth}{\Field{Purpose.}{State why the model needs both a card description and operations, then preview the section's explanatory order.}\Field{Primary purpose.}{roadmap}\Field{Narration structure.}{section\_goal $\rightarrow$ model\_requirements $\rightarrow$ expository\_choice $\rightarrow$ forward\_reference}\Field{Reusable move.}{Preview the section in the same order as its formal development.}}}
\needspace{10\baselineskip}\subsubsection*{W-01-P006: analogous purpose}
\Field{Location.}{Section 01, prose\_paragraph, lines 169--181.}
\Field{Draft purpose.}{Guide the reader through the paper's model, framework, instances, main theorems, trust account, and conclusion.}
\Field{Qualitative closeness.}{analogous purpose. Both paragraphs justify an exposition order and preview the concepts or results developed in sequence.}
\Field{Limits of the analogy.}{The baseline gives a short model-section roadmap. The WADT paragraph maps the entire paper, includes theorem destinations, and explains why the mixing section follows the ideal results.}
\Field{Baseline narration used.}{adapted.}
\Field{Complexity and jargon.}{Complexity: \Tag{medium} $\rightarrow$ \Tag{medium}. Jargon: \Tag{medium} $\rightarrow$ \Tag{medium}.}
\textbf{Original WADT unit. Footnotes omitted.}
\begin{lstlisting}
Section~\ref{sec:model} fixes the protocol flow, the two distributions, and
the security scope. Section~\ref{sec:framework} presents the framework components used by
the proofs. Section~\ref{sec:fivecard} instantiates them on the five-card
family, where the deck is small enough for direct case analysis.
Sections~\ref{sec:pgl} and~\ref{sec:exact} construct the $\PG$ instance,
where deck enumeration no longer suffices and three-transitivity carries
the privacy proofs, ending in Theorem A. Section~\ref{sec:mixing} proves
the word-shuffle results, ending in Theorem B, because a physical dealer
repeats elementary moves rather than sampling one uniform group element. Section~\ref{sec:instances} compares the sibling instances and
states the trust base. Section~\ref{sec:related} places the results among
related work. Section~\ref{sec:conclusion} states the remaining extensions
and other future directions.
\end{lstlisting}
\colorbox{RewriteGreen}{\parbox{0.97\textwidth}{\textbf{Rewrite suggestion.}\par\begingroup\small
Section~\ref{sec:model} fixes the protocol flow, the two distributions, and the security scope. Section~\ref{sec:framework} presents the framework components used in the proofs. Section~\ref{sec:fivecard} instantiates these components for the five-card family, whose small deck permits direct case analysis. Sections~\ref{sec:pgl} and~\ref{sec:exact} construct the $\PG$ instance. Deck enumeration no longer suffices, so three-transitivity carries the privacy proofs that end in Theorem A. Section~\ref{sec:mixing} proves the word-shuffle results that end in Theorem B. This section is needed because a physical dealer repeats elementary moves instead of sampling one uniform group element. Section~\ref{sec:instances} compares the additional instances and states the trust base. Section~\ref{sec:related} places the results among related work. Section~\ref{sec:conclusion} states the remaining extensions and other future directions.
\par\endgroup}}
\Field{Why this rewrite.}{The rewrite adapts the baseline's ordered roadmap. It preserves every section reference and theorem destination while separating the two causal explanations and clarifying that the sibling instances are additional profiles.}
\Field{Terminology actions.}{Keep: three-transitivity, Theorem A, Theorem B | Explain: sibling instances}
\par\smallskip
\needspace{10\baselineskip}\subsubsection*{W-03-P008: analogous purpose}
\Field{Location.}{Section 03, prose\_paragraph, lines 537--541.}
\Field{Draft purpose.}{Transition from component packaging to generic security arguments.}
\Field{Qualitative closeness.}{analogous purpose. Both paragraphs close one explanatory block and preview the order and purpose of the next block.}
\Field{Limits of the analogy.}{The baseline introduces a full model section. The WADT paragraph moves from one instance to representation-independent generic theorems.}
\Field{Baseline narration used.}{adapted.}
\Field{Complexity and jargon.}{Complexity: \Tag{low} $\rightarrow$ \Tag{low}. Jargon: \Tag{medium} $\rightarrow$ \Tag{low}.}
\textbf{Original WADT unit. Footnotes omitted.}
\begin{lstlisting}
Section~\ref{sec:fivecard} shows how one instance fills these components.
Once the data are packaged, the reusable security arguments depend on
mathematical hypotheses rather than the record representation. The next
subsection states these arguments directly.
\end{lstlisting}
\colorbox{RewriteGreen}{\parbox{0.97\textwidth}{\textbf{Rewrite suggestion.}\par\begingroup\small
Section~\ref{sec:fivecard} shows how one instance supplies these components. Once packaged, the reusable security arguments depend on mathematical hypotheses, not on the record representation. The next subsection states these arguments directly.
\par\endgroup}}
\Field{Why this rewrite.}{The rewrite adapts the baseline's ordered roadmap. It moves from a concrete instance to the representation-independent theorem principle and then to the next subsection.}
\Field{Terminology actions.}{Keep: mathematical hypotheses, record representation}
\par\smallskip
\needspace{10\baselineskip}\subsubsection*{W-06-P001: analogous purpose}
\Field{Location.}{Section 06, prose\_paragraph, lines 1279--1283.}
\Field{Draft purpose.}{Orient the reader to the section's three result layers and their relation to earlier construction facts.}
\Field{Qualitative closeness.}{analogous purpose. Both paragraphs explain why a section must combine several layers and preview the order in which it develops them.}
\Field{Limits of the analogy.}{The baseline opens the unified model. The WADT paragraph opens a results section that lifts construction facts to execution, recovery, and privacy before Theorem A.}
\Field{Baseline narration used.}{adapted.}
\Field{Complexity and jargon.}{Complexity: \Tag{low} $\rightarrow$ \Tag{low}. Jargon: \Tag{medium} $\rightarrow$ \Tag{low}.}
\textbf{Original WADT unit. Footnotes omitted.}
\begin{lstlisting}
The construction section established deck-level correctness and the
group-action premise for privacy. This section lifts correctness to the
executed protocol, determines the recovery ramp, and proves view and trace
privacy under a uniform shuffle. These results form Theorem A.
\end{lstlisting}
\colorbox{RewriteGreen}{\parbox{0.97\textwidth}{\textbf{Rewrite suggestion.}\par\begingroup\small
The construction section proved deck-level correctness and the group-action condition needed for privacy. This section lifts correctness from decks to full protocol executions. It also determines the recovery ramp and proves view and trace privacy under a uniform shuffle. Together, these results form Theorem A.
\par\endgroup}}
\Field{Why this rewrite.}{The rewrite adapts the baseline's ordered section preview. It explains what lifting correctness means and separates the execution, recovery, privacy, and synthesis moves.}
\Field{Terminology actions.}{Keep: group-action premise, uniform shuffle, Theorem A | Explain: lifts correctness}
\par\smallskip
\needspace{10\baselineskip}\subsubsection*{W-07-P001: analogous purpose}
\Field{Location.}{Section 07, prose\_paragraph, lines 1515--1524.}
\Field{Draft purpose.}{State the finite-word section's goal, scope limit, and route to security.}
\Field{Qualitative closeness.}{analogous purpose. Both units explain why several formal ingredients are needed and preview their order in the section.}
\Field{Limits of the analogy.}{The baseline organizes a unified protocol model. This unit organizes one proof chain and includes an explicit physical nonclaim.}
\Field{Baseline narration used.}{adapted.}
\Field{Complexity and jargon.}{Complexity: \Tag{high} $\rightarrow$ \Tag{medium}. Jargon: \Tag{high} $\rightarrow$ \Tag{medium}.}
\textbf{Original WADT unit. Footnotes omitted.}
\begin{lstlisting}
Section~\ref{sec:model} defines the generator word. This section proves a
machine-checked mixing certificate for the finitely generated shuffle: the
evaluated distribution of a uniform word is close to the uniform
distribution on the group. The letters are group generators rather than
cuts of the eight-card row, and I do not claim a physical realization of an
unobserved letter application. The certificate is the object of interest,
since it bounds the distance of the executable shuffle from the ideal one.
The proof transfers this bound to the security observables, and Theorem B
gives the final bound.
\end{lstlisting}
\colorbox{RewriteGreen}{\parbox{0.97\textwidth}{\textbf{Rewrite suggestion.}\par\begingroup\small
Section~\ref{sec:model} defines generator words. This section proves a machine-checked mixing certificate for the finitely generated shuffle. It shows that evaluating a uniform word gives a distribution close to the uniform distribution on the group. The letters are group generators, not cuts of the eight-card row. I do not claim a physical realization of an unobserved letter application. The main certificate bounds the distance between the executable and ideal shuffles. The proof then transfers this bound to security observables. Theorem B states the final bound.
\par\endgroup}}
\Field{Why this rewrite.}{Adapts the baseline's section roadmap to the proof chain. It preserves the mixing claim, physical nonclaim, certificate role, transfer, and theorem handoff.}
\Field{Terminology actions.}{Keep: mixing certificate, uniform group distribution, security observables | Explain: unobserved letter application | Replace: the object of interest -> the main certificate}
\par\smallskip
\needspace{10\baselineskip}\subsubsection*{W-07-P011: partial analogue}
\Field{Location.}{Section 07, prose\_paragraph, lines 1657--1661.}
\Field{Draft purpose.}{Organize the PGL proof into layers and introduce the formal source index.}
\Field{Qualitative closeness.}{partial analogue. The baseline supplies the move of previewing an ordered group of ingredients before detailed navigation.}
\Field{Limits of the analogy.}{The baseline previews a model section. This unit introduces a traceability table rather than mathematical exposition.}
\Field{Baseline narration used.}{partial.}
\Field{Complexity and jargon.}{Complexity: \Tag{medium} $\rightarrow$ \Tag{low}. Jargon: \Tag{high} $\rightarrow$ \Tag{medium}.}
\textbf{Original WADT unit. Footnotes omitted.}
\begin{lstlisting}
The PGL argument has four proof layers: construction, recovery, privacy,
and finite-word transfer. Table~\ref{tab:source-index} maps the
mathematical claims used by Theorems A and B to the \Rocq{} results that
establish them.
\end{lstlisting}
\colorbox{RewriteGreen}{\parbox{0.97\textwidth}{\textbf{Rewrite suggestion.}\par\begingroup\small
The $\PG$ argument has four proof layers: construction, recovery, privacy, and finite-word transfer. Table~\ref{tab:source-index} maps the mathematical claims used by Theorems A and B to the \Rocq{} results that establish them.
\par\endgroup}}
\Field{Why this rewrite.}{Uses only the baseline's ordered-preview move to introduce the four proof layers and their traceability table.}
\Field{Terminology actions.}{Keep: finite-word transfer, Rocq results | Explain: proof layers}
\par\smallskip
\clearpage\subsection{Baseline S-02-P002: model exposition plan}
\colorbox{BaselineBlue}{\parbox{0.97\textwidth}{\Field{Purpose.}{Explain which ingredients the model must define and why a familiar case is used before the general case.}\Field{Primary purpose.}{roadmap}\Field{Narration structure.}{model\_requirements $\rightarrow$ expository\_example\_choice $\rightarrow$ later\_extension $\rightarrow$ appendix\_scope}\Field{Reusable move.}{Preview the section in the same order as its formal development.}}}
\needspace{10\baselineskip}\subsubsection*{W-03-P001: analogous purpose}
\Field{Location.}{Section 03, prose\_paragraph, lines 333--349.}
\Field{Draft purpose.}{Explain why the framework gathers five instance choices into one profile record and how generic theorems consume them.}
\Field{Qualitative closeness.}{analogous purpose. Both paragraphs identify the ingredients needed by a formal model and explain why they are gathered before later theorems consume them.}
\Field{Limits of the analogy.}{The baseline prepares a card-and-operation model. The WADT paragraph describes a proof-oriented record, its theorem interface, one record-level privacy package, and the bridge table.}
\Field{Baseline narration used.}{adapted.}
\Field{Complexity and jargon.}{Complexity: \Tag{high} $\rightarrow$ \Tag{medium}. Jargon: \Tag{high} $\rightarrow$ \Tag{medium}.}
\textbf{Original WADT unit. Footnotes omitted.}
\begin{lstlisting}
Every instance in this paper supplies the same kinds of data: a group
with its action and generators, a secret type, the layout of a run, a
shuffle-security bound, and a decoder. The framework therefore gathers
these five choices into one record, the record type
\coqin{MonodromyProfile}, in order to reuse one executable protocol
across instances.
The generic theorems of this section take their hypotheses directly, and
the privacy theorem is additionally packaged once at the record level, so a
profile that satisfies the transitivity and distinct-deck obligations
inherits coalition privacy as one instantiation. The
record carries the data of Equation~\ref{eq:model-data} into the
executable protocol. Here $R$ is the real closed field of
Equation~\ref{eq:model-data}, not a protocol datum. The finite algebraic
structures use Mathematical Components~\cite{MathComp}, and
Table~\ref{tab:bridge} maps each model datum to its record role and to the
values the $\PG$ instance supplies.
\end{lstlisting}
\colorbox{RewriteGreen}{\parbox{0.97\textwidth}{\textbf{Rewrite suggestion.}\par\begingroup\small
Every instance supplies five kinds of data: a group with its action and generators, a secret type, a run layout, a shuffle-security bound, and a decoder. I gather these choices in one profile record so that one executable protocol can be reused across instances. The generic theorems take their hypotheses directly. I also package the privacy theorem once at the record level. As a result, a profile with a $t$-transitive action and distinct encoded decks, meaning that no card value is repeated, inherits privacy for coalitions of at most $t$ positions. The profile carries the data of Equation~\ref{eq:model-data} into the executable protocol. The scalar $R$ in Equation~\ref{eq:model-data} is the real closed field, not a protocol datum. The finite algebraic structures use Mathematical Components~\cite{MathComp}. Together, these data and obligations give the record an algebraic-specification interpretation. The next paragraph develops that interpretation, and Table~\ref{tab:bridge} then maps each model datum to its record role and to the values supplied by the $\PG$ instance.
\par\endgroup}}
\Field{Why this rewrite.}{The rewrite keeps the five profile components, their reuse purpose, the direct mathematical hypotheses, the separately packaged privacy theorem, and the transitivity consequence. It connects the scalar-field boundary, library citation, and table reference to the role of the profile instead of listing them as qualifications. The final two sentences lead first to the algebraic-specification interpretation and then to the bridge table.}
\Field{Terminology actions.}{Keep: profile record, real closed field, transitivity | Explain: distinct-deck obligation | Replace: record type MonodromyProfile -> profile record}
\par\smallskip
\needspace{10\baselineskip}\subsubsection*{W-03-P003: partial analogue}
\Field{Location.}{Section 03, prose\_paragraph, lines 387--391.}
\Field{Draft purpose.}{Connect the bridge table to the architecture figure.}
\Field{Qualitative closeness.}{partial analogue. Both units guide the reader from a component inventory to the dependency structure that follows.}
\Field{Limits of the analogy.}{The baseline previews a section's exposition. The WADT paragraph points backward to a table and forward to a figure instead of developing the model itself.}
\Field{Baseline narration used.}{partial.}
\Field{Complexity and jargon.}{Complexity: \Tag{low} $\rightarrow$ \Tag{low}. Jargon: \Tag{medium} $\rightarrow$ \Tag{low}.}
\textbf{Original WADT unit. Footnotes omitted.}
\begin{lstlisting}
The table identifies the role of each component.
Figure~\ref{fig:framework-architecture} shows how the components and
supporting records depend on one another, and how the profile determines
the executable protocol.
\end{lstlisting}
\colorbox{RewriteGreen}{\parbox{0.97\textwidth}{\textbf{Rewrite suggestion.}\par\begingroup\small
The table gives each component's role. Figure~\ref{fig:framework-architecture} shows the dependencies among the components and auxiliary records. It also shows how the profile determines the executable protocol.
\par\endgroup}}
\Field{Why this rewrite.}{The rewrite uses the partial analogue's inventory-to-interaction move. It replaces the vague supporting-record phrase with auxiliary records and separates dependency from derivation.}
\Field{Terminology actions.}{Keep: executable protocol | Explain: supporting records}
\par\smallskip
\needspace{10\baselineskip}\subsubsection*{W-03-CAP003: partial analogue}
\Field{Location.}{Section 03, caption, lines 514--518.}
\Field{Draft purpose.}{Explain the witness table and locate the evidence represented by its rows.}
\Field{Qualitative closeness.}{partial analogue. Both units orient the reader by naming what a structured comparison contains and where its details are developed.}
\Field{Limits of the analogy.}{The baseline gives a prose plan for one model section. The WADT caption summarizes optional witness slots and points to three later evidence locations.}
\Field{Baseline narration used.}{partial.}
\Field{Complexity and jargon.}{Complexity: \Tag{medium} $\rightarrow$ \Tag{medium}. Jargon: \Tag{high} $\rightarrow$ \Tag{medium}.}
\textbf{Original WADT unit. Footnotes omitted.}
\begin{lstlisting}
Realized combinations of the security witness's two optional
  slots. Section~\ref{sec:mixing} treats the word-shuffle counterpart of
  the $\PG$ row, Section~\ref{sec:fivecard} proves the den Boer and Kim
  rows, and Table~\ref{tab:instances} records the per-instance evidence.
\end{lstlisting}
\colorbox{RewriteGreen}{\parbox{0.97\textwidth}{\textbf{Rewrite suggestion.}\par\begingroup\small
Realized combinations of the security witness's two optional evidence fields. Section~\ref{sec:mixing} gives the word-shuffle counterpart of the $\PG$ row. Section~\ref{sec:fivecard} proves the den Boer and Kim rows. Table~\ref{tab:instances} records the evidence for each instance.
\par\endgroup}}
\Field{Why this rewrite.}{The rewrite uses the partial analogue's compact-roadmap move. It explains the optional slots as evidence fields and assigns each referenced source one role.}
\Field{Terminology actions.}{Keep: security witness, word-shuffle counterpart | Explain: optional slots}
\par\smallskip
\needspace{10\baselineskip}\subsubsection*{W-04-P011: analogous purpose}
\Field{Location.}{Section 04, prose\_paragraph, lines 823--828.}
\Field{Draft purpose.}{Close the five-card section and motivate the larger instance.}
\Field{Qualitative closeness.}{analogous purpose. Both units use a familiar manageable case to prepare a more general construction that needs additional machinery.}
\Field{Limits of the analogy.}{The baseline previews a general model after a familiar case. This unit contrasts proof methods and security strength across concrete instances.}
\Field{Baseline narration used.}{adapted.}
\Field{Complexity and jargon.}{Complexity: \Tag{medium} $\rightarrow$ \Tag{low}. Jargon: \Tag{high} $\rightarrow$ \Tag{medium}.}
\textbf{Original WADT unit. Footnotes omitted.}
\begin{lstlisting}
The five-card family keeps the group cyclic and the deck small enough for
kernel enumeration. The next sections instantiate the same records where
enumeration fails. The $\PG$ construction adds coalition privacy beyond
one card via three-transitivity, an orbit-class secret, and word
certificates proven without computing in the group.
\end{lstlisting}
\colorbox{RewriteGreen}{\parbox{0.97\textwidth}{\textbf{Rewrite suggestion.}\par\begingroup\small
The five-card family uses a cyclic group and a deck small enough for kernel enumeration. The next sections instantiate the same records when enumeration is not feasible. The $\PG$ construction strengthens coalition privacy beyond one card by using three-transitivity and an orbit-class secret. Its word certificates avoid computation inside the group.
\par\endgroup}}
\Field{Why this rewrite.}{Adapts the baseline's familiar-case-to-extension roadmap and makes the change in proof method and privacy strength explicit.}
\Field{Terminology actions.}{Keep: three-transitivity, orbit-class secret, word certificate | Explain: kernel enumeration}
\par\smallskip
\needspace{10\baselineskip}\subsubsection*{W-05-P006: analogous purpose}
\Field{Location.}{Section 05, prose\_paragraph, lines 880--884.}
\Field{Draft purpose.}{Give the construction order for the orbit encoding.}
\Field{Qualitative closeness.}{analogous purpose. Both units explain which ingredients a formal model needs and order their introduction before the details.}
\Field{Limits of the analogy.}{The baseline moves from a familiar deck case to a general model. The WADT unit stays within one instance and orders an encoding proof.}
\Field{Baseline narration used.}{adapted.}
\Field{Complexity and jargon.}{Complexity: \Tag{medium} $\rightarrow$ \Tag{low}. Jargon: \Tag{medium} $\rightarrow$ \Tag{low}.}
\textbf{Original WADT unit. Footnotes omitted.}
\begin{lstlisting}
To make the encoding precise, we first define the reachable patterns of four
positions. We then give an explicit classifier for the two orbits, choose one
deck from each class, and prove that the classifier is unchanged by every
allowed shuffle.
\end{lstlisting}
\colorbox{RewriteGreen}{\parbox{0.97\textwidth}{\textbf{Rewrite suggestion.}\par\begingroup\small
I make the encoding precise in four steps. First, I define the four-position heart patterns obtained from valid decks under allowed shuffles. Next, I give an explicit classifier for the two orbits. I then choose one deck from each class. Finally, I prove that every allowed shuffle preserves the classifier.
\par\endgroup}}
\Field{Why this rewrite.}{Adapts the baseline's ordered section preview to the four dependencies of the WADT encoding proof.}
\Field{Terminology actions.}{Keep: orbit classifier | Explain: reachable patterns -> four-position heart patterns obtained from valid decks under allowed shuffles}
\par\smallskip
\clearpage\subsection{Baseline S-02-P004: deck-field explanations}
\colorbox{BaselineBlue}{\parbox{0.97\textwidth}{\Field{Purpose.}{Assign a mathematical role to each field of the deck tuple.}\Field{Primary purpose.}{definition}\Field{Narration structure.}{carrier\_role $\rightarrow$ transformation\_role $\rightarrow$ symbol\_role $\rightarrow$ visibility\_role $\rightarrow$ multiset\_role}\Field{Reusable move.}{Introduce one formal component and state its role in the surrounding construction.}}}
\needspace{10\baselineskip}\subsubsection*{W-02-P004: same purpose}
\Field{Location.}{Section 02, prose\_paragraph, lines 204--210.}
\Field{Draft purpose.}{Assign a mathematical role to every component of the protocol model and distinguish its ideal and finite shuffle laws.}
\Field{Qualitative closeness.}{same purpose. Both units unpack formal data field by field so that later constructions can use every component without ambiguity.}
\Field{Limits of the analogy.}{The baseline fields describe card data. The WADT fields combine algebra, an action, probability scalars, and an iterated distribution.}
\Field{Baseline narration used.}{full.}
\Field{Complexity and jargon.}{Complexity: \Tag{medium} $\rightarrow$ \Tag{low}. Jargon: \Tag{low} $\rightarrow$ \Tag{low}.}
\textbf{Original WADT unit. Footnotes omitted.}
\begin{lstlisting}
Here $G$ is a finite shuffle group. The representation
$\rho:G\to S_n$ acts on $n$ card positions. The distribution $\mu$ selects one
generator letter, and $R$ is the real closed field used for probabilities.
The distribution $U_G$ is uniform on $G$. I write $\mu^{*L}$ for the word distribution at length $L$: the
distribution of the group element obtained by evaluating $L$ independent
letters.
\end{lstlisting}
\colorbox{RewriteGreen}{\parbox{0.97\textwidth}{\textbf{Rewrite suggestion.}\par\begingroup\small
Here $G$ is a finite shuffle group. The representation $\rho:G\to S_n$ means that each group element permutes the $n$ card positions. The distribution $\mu$ selects one generator letter, which denotes one elementary random move. The real closed field $R$ supplies the probability values. The distribution $U_G$ is uniform on $G$. I write $\mu^{*L}$ for the word distribution at length $L$. It is the distribution of the group element obtained by evaluating $L$ independent letters.
\par\endgroup}}
\Field{Why this rewrite.}{Reuses the baseline's field-by-field sequence and gives each algebraic or probabilistic component one explicit role.}
\Field{Terminology actions.}{Keep: finite shuffle group, real closed field, word distribution | Explain: permutation representation -> each group element permutes the card positions, generator letter -> one elementary random move}
\par\smallskip
\needspace{10\baselineskip}\subsubsection*{W-03-TC007: partial analogue}
\Field{Location.}{Section 03, table\_cell, lines 375--376.}
\Field{Draft purpose.}{Identify the algebraic action and the two shuffle-law choices in the second model-data row.}
\Field{Qualitative closeness.}{partial analogue. S-02-P004 supplies the limited field-role move. W-03-TC007 applies it to one compact field value containing the group action and two alternative probability laws.}
\Field{Limits of the analogy.}{The baseline assigns one role to each deck field in prose. This cell compresses algebra and probability into notation and depends on the endpoint-bound cell that follows.}
\Field{Baseline narration used.}{partial.}
\Field{Complexity and jargon.}{Complexity: \Tag{low} $\rightarrow$ \Tag{low}. Jargon: \Tag{medium} $\rightarrow$ \Tag{medium}.}
\textbf{Original WADT unit. Footnotes omitted.}
\begin{lstlisting}
$(G,\rho)$ with $\mu$ or $U_G$
\end{lstlisting}
\colorbox{RewriteGreen}{\parbox{0.97\textwidth}{\textbf{Rewrite suggestion.}\par\begingroup\small
group action $(G,\rho)$ with the generator distribution $\mu$ or the uniform group distribution $U_G$
\par\endgroup}}
\Field{Why this rewrite.}{The rewrite uses the partial analogue's field-first move. It names both symbols and explains the finite and ideal shuffle laws in the compact cell.}
\Field{Terminology actions.}{Keep: group action, uniform group distribution, generator distribution | Explain: mu or U\_G}
\par\smallskip
\needspace{10\baselineskip}\subsubsection*{W-03-CAP001: analogous purpose}
\Field{Location.}{Section 03, caption, lines 381--384.}
\Field{Draft purpose.}{Map the bridge table's four mathematical roles to their exact Rocq record names.}
\Field{Qualitative closeness.}{analogous purpose. Both units map formal components to their mathematical roles.}
\Field{Limits of the analogy.}{The baseline explains fields in running prose. The WADT caption summarizes a whole table and adds four source-level code identifiers.}
\Field{Baseline narration used.}{adapted.}
\Field{Complexity and jargon.}{Complexity: \Tag{medium} $\rightarrow$ \Tag{medium}. Jargon: \Tag{high} $\rightarrow$ \Tag{medium}.}
\textbf{Original WADT unit. Footnotes omitted.}
\begin{lstlisting}
From model data to framework records. In the sources the four
  roles are \coqin{PGGInterface}, \coqin{SecurityWitness},
  \coqin{ReconPlug}, and \coqin{MonodromyProfile}.
\end{lstlisting}
\colorbox{RewriteGreen}{\parbox{0.97\textwidth}{\textbf{Rewrite suggestion.}\par\begingroup\small
Mapping from model data to framework records. The source names the run-layout interface \coqin{PGGInterface}, the shuffle-evidence record \coqin{SecurityWitness}, the action-compatible reconstruction record \coqin{ReconPlug}, and the aggregate profile \coqin{MonodromyProfile}.
\par\endgroup}}
\Field{Why this rewrite.}{The rewrite adapts the baseline's field-to-role sequence. It retains every source identifier and explains the mathematical role of each one within a concise caption.}
\Field{Terminology actions.}{Keep: PGGInterface, SecurityWitness, ReconPlug, MonodromyProfile | Explain: SecurityWitness, MonodromyProfile}
\par\smallskip
\needspace{10\baselineskip}\subsubsection*{W-05-P003: analogous purpose}
\Field{Location.}{Section 05, prose\_paragraph, lines 851--865.}
\Field{Draft purpose.}{Explain the position labels, card labels, deck convention, and group action, then introduce the group-order display.}
\Field{Qualitative closeness.}{analogous purpose. Both units assign mathematical and operational roles to the components of a formal model before later constructions use them.}
\Field{Limits of the analogy.}{The WADT unit also compares deck conventions, cites prior work, and ends with a handoff to an unattached order display. The baseline only explains tuple fields.}
\Field{Baseline narration used.}{adapted.}
\Field{Complexity and jargon.}{Complexity: \Tag{high} $\rightarrow$ \Tag{medium}. Jargon: \Tag{high} $\rightarrow$ \Tag{medium}.}
\textbf{Original WADT unit. Footnotes omitted.}
\begin{lstlisting}
The symbols $0,\ldots,6$ are the seven field elements, and $\infty$ is the
additional projective point. Each point labels one physical position and
therefore the player who receives the card in that position. Card values are
separate labels drawn from $C=\{0,1,\ldots,7\}$. All eight cards therefore
carry distinct symbols. This is the standard-deck setting of the card-based
literature, which caters for commonly available decks of pairwise
distinguishable cards, in contrast with the two-symbol encodings in which
most protocol sizes are measured~\cite{NiemiRenvall1999,Koch2019}. The group $G=\PG$ acts on
$X$ by fractional linear transformations. Each group element therefore
permutes the eight card positions and specifies an allowed deck
rearrangement. Its order is
\end{lstlisting}
\colorbox{RewriteGreen}{\parbox{0.97\textwidth}{\textbf{Rewrite suggestion.}\par\begingroup\small
The symbols $0,\ldots,6$ are the seven field elements. The symbol $\infty$ is the additional projective point. Each point labels one physical position and the player who receives its card. Card values use the separate labels $C=\{0,1,\ldots,7\}$. Thus all eight cards carry distinct symbols. This standard-deck setting uses pairwise distinguishable cards of the kind found in common decks. Most protocol sizes are instead measured for two-symbol encodings~\cite{NiemiRenvall1999,Koch2019}. The group $G=\PG$ acts on $X$ by fractional linear transformations. Each group element therefore permutes the eight card positions and defines an allowed deck rearrangement. The group has order:
\par\endgroup}}
\Field{Why this rewrite.}{Adapts the baseline's component-by-component explanation to separate position labels, card labels, deck convention, group action, and order.}
\Field{Terminology actions.}{Keep: fractional linear transformation, pairwise distinguishable cards | Explain: projective point -> the seven field elements together with the additional point at infinity, standard-deck setting -> uses pairwise distinguishable cards | Replace: caters for commonly available decks -> uses cards of the kind found in common decks}
\par\smallskip
\needspace{10\baselineskip}\subsubsection*{W-05-STM002: analogous purpose}
\Field{Location.}{Section 05, statement, lines 912--944.}
\Field{Draft purpose.}{Define a Boolean deck decoder by extracting heart positions and classifying their cross ratio.}
\Field{Qualitative closeness.}{analogous purpose. Both units make a formal object usable by assigning a precise mathematical role to every component. The WADT roles form a longer dependent chain that ends in one decoder.}
\Field{Limits of the analogy.}{The baseline explains fields of a compact card tuple. The WADT statement defines a set map, a finite-field invariant, a special case, a piecewise classifier, and a composition inside a formal environment.}
\Field{Baseline narration used.}{adapted.}
\Field{Complexity and jargon.}{Complexity: \Tag{high} $\rightarrow$ \Tag{medium}. Jargon: \Tag{high} $\rightarrow$ \Tag{medium}.}
\textbf{Original WADT unit. Footnotes omitted.}
\begin{lstlisting}
A valid deck is a bijection $D:X\to C$. Its four heart values are
$0,1,2,3$, and their positions form
\[
  H(D)=\{i\in X\mid D(i)\in\{0,1,2,3\}\}\in\Omega_4.
\]
For $A=\{a<b<c<d\}$ under the order
$0<1<\cdots<6<\infty$, define the cross ratio over $\mathbb{F}_7$ by
\[
  [a,b;c,d]=\frac{(a-c)(b-d)}{(a-d)(b-c)},
  \qquad
  [a,b;c,\infty]=\frac{a-c}{b-c}.
\]
The classifier $\kappa:\Omega_4\to\{0,1\}$ is
\[
  \kappa(A)=
  \begin{cases}
    1, & [a,b;c,d]\in\{3,5\},\\
    0, & [a,b;c,d]\in\{2,4,6\}.
  \end{cases}
\]
The value one is equianharmonic, and the value zero is harmonic. The deck
decoder is
\[
  \operatorname{dec}(D)=\kappa(H(D)),
  \qquad
  D\xmapsto{H}\Omega_4\xmapsto{\kappa}\{0,1\}.
\]
\end{lstlisting}
\colorbox{RewriteGreen}{\parbox{0.97\textwidth}{\textbf{Rewrite suggestion.}\par\begingroup\small
A valid deck is a bijection $D:X\to C$. Its four heart-labeled card values are $0,1,2,3$. Their positions form
\[
  H(D)=\{i\in X\mid D(i)\in\{0,1,2,3\}\}\in\Omega_4.
\]
Let $A=\{a<b<c<d\}$ under the order $0<1<\cdots<6<\infty$. Define the cross ratio in the finite field $\mathbb{F}_7$ by
\[
  [a,b;c,d]=\frac{(a-c)(b-d)}{(a-d)(b-c)},
  \qquad
  [a,b;c,\infty]=\frac{a-c}{b-c}.
\]
The second formula covers the projective point at infinity. Define the classifier $\kappa:\Omega_4\to\{0,1\}$ by
\[
  \kappa(A)=
  \begin{cases}
    1, & [a,b;c,d]\in\{3,5\},\\
    0, & [a,b;c,d]\in\{2,4,6\}.
  \end{cases}
\]
The classifier maps the equianharmonic class to one and the harmonic class to zero. The deck decoder is
\[
  \operatorname{dec}(D)=\kappa(H(D)),
  \qquad
  D\xmapsto{H}\Omega_4\xmapsto{\kappa}\{0,1\}.
\]
\par\endgroup}}
\Field{Why this rewrite.}{Adapts the baseline's component-by-component exposition to define deck validity, heart positions, the cross ratio, the two classes, and the decoder in order. The excluded footnote does not enter the suggestion.}
\Field{Terminology actions.}{Keep: bijection, cross ratio, harmonic and equianharmonic, decoder | Explain: point at infinity -> the projective case covered by the second cross-ratio formula, classifier -> maps the two cross-ratio classes to Boolean values, finite field of order seven -> the arithmetic domain \$\textbackslash{}mathbb\{F\}\_7\$ | Replace: heart values -> heart-labeled card values}
\par\smallskip
\needspace{10\baselineskip}\subsubsection*{W-06-P005: analogous purpose}
\Field{Location.}{Section 06, prose\_paragraph, lines 1316--1330.}
\Field{Draft purpose.}{Define the recovery ramp and distinguish its two thresholds and the total player count.}
\Field{Qualitative closeness.}{analogous purpose. Both paragraphs assign a distinct mathematical role to each component of a compact formal tuple.}
\Field{Limits of the analogy.}{The baseline explains the fields of a deck tuple. The WADT paragraph defines privacy, recovery, and player-count coordinates, then adds their logical distinction, secret-sharing context, and the implemented-decoder boundary.}
\Field{Baseline narration used.}{adapted.}
\Field{Complexity and jargon.}{Complexity: \Tag{high} $\rightarrow$ \Tag{medium}. Jargon: \Tag{high} $\rightarrow$ \Tag{medium}.}
\textbf{Original WADT unit. Footnotes omitted.}
\begin{lstlisting}
Three parameters summarize the instance, and the two thresholds live on
different axes. The privacy threshold $t$ is defined as the largest
coalition size with perfect privacy, a probabilistic condition: the
coalition's view is independent of the secret. The recovery threshold $r$
is defined as the least number of revealed positions that always
determines the secret class, a possibilistic condition: below it, both
classes remain consistent with the revealed cards. The number of card
positions is $n$. I call the triple $(t,r,n)$ the recovery ramp, following
the ramp schemes of the secret-sharing literature, in which the gap between
the two thresholds is the graded regime between privacy and
recovery~\cite{BlakleyMeadows1984,Yamamoto1986}. The implemented decoder
reads all eight endpoints, so $r=7$ is a structural threshold of the
arrangement combinatorics rather than the reading rule of a smaller
implemented decoder.
\end{lstlisting}
\colorbox{RewriteGreen}{\parbox{0.97\textwidth}{\textbf{Rewrite suggestion.}\par\begingroup\small
I summarize the instance by the triple $(t,r,n)$. Its first two parameters answer different questions. The privacy threshold $t$ asks how large a coalition can be while its view remains independent of the secret. It is the largest coalition size with perfect privacy, so it is a probabilistic condition. By contrast, the recovery threshold $r$ asks how many revealed positions always determine the secret class. It is the least such number, so it is a possibilistic condition. Below $r$, both secret classes remain consistent with the revealed cards. Finally, $n$ is the number of card positions. Following ramp schemes in the secret-sharing literature, I call $(t,r,n)$ the recovery ramp~\cite{BlakleyMeadows1984,Yamamoto1986}. The gap between $t$ and $r$ is the graded regime between privacy and recovery. Although $r=7$, the implemented decoder reads all eight endpoints. Thus, $r=7$ is a structural property of the arrangements. It is not the reading rule of a smaller implemented decoder.
\par\endgroup}}
\Field{Why this rewrite.}{The rewrite starts from the triple and presents $t$ and $r$ as answers to different questions. The phrase By contrast marks the probabilistic and possibilistic distinction, while Finally introduces $n$. It preserves the recovery-ramp citations, the graded regime, all eight endpoints, and the distinction between the structural threshold and the implemented decoder.}
\Field{Terminology actions.}{Keep: privacy threshold, recovery threshold, recovery ramp | Explain: possibilistic condition, graded regime | Replace: thresholds live on different axes -> Its first two parameters answer different questions}
\par\smallskip
\clearpage\subsection{Baseline S-02-P005: identity and specification convention}
\colorbox{BaselineBlue}{\parbox{0.97\textwidth}{\Field{Purpose.}{State the identity requirement and name the reduced tuple used as a card specification.}\Field{Primary purpose.}{assumption}\Field{Narration structure.}{identity\_requirement $\rightarrow$ terminology\_assignment}\Field{Reusable move.}{State the constraint before deriving definitions that rely on it.}}}
\needspace{10\baselineskip}\subsubsection*{W-05-P004: analogous purpose}
\Field{Location.}{Section 05, prose\_paragraph, lines 869--870.}
\Field{Draft purpose.}{State the implementation convention for the projective point at infinity.}
\Field{Qualitative closeness.}{analogous purpose. Both units state a compact representation convention needed by later formulas.}
\Field{Limits of the analogy.}{The baseline fixes an identity element and a reduced tuple. The WADT unit maps one projective point to a numeric index.}
\Field{Baseline narration used.}{adapted.}
\Field{Complexity and jargon.}{Complexity: \Tag{low} $\rightarrow$ \Tag{low}. Jargon: \Tag{low} $\rightarrow$ \Tag{low}.}
\textbf{Original WADT unit. Footnotes omitted.}
\begin{lstlisting}
The formal representation writes $\infty$ as position seven.
\end{lstlisting}
\colorbox{RewriteGreen}{\parbox{0.97\textwidth}{\textbf{Rewrite suggestion.}\par\begingroup\small
The formal representation assigns the point at infinity $\infty$ to position seven.
\par\endgroup}}
\Field{Why this rewrite.}{Adapts the baseline's convention-setting move to the single projective-point index used later.}
\Field{Terminology actions.}{Keep: point at infinity}
\par\smallskip
\clearpage\subsection{Baseline S-02-P006: binary-card semantics and naming}
\colorbox{BaselineBlue}{\parbox{0.97\textwidth}{\Field{Purpose.}{Explain how orientation, transformation, and visibility determine the meaning of the running card instance.}\Field{Primary purpose.}{definition}\Field{Narration structure.}{orientation\_semantics $\rightarrow$ transformation\_effect $\rightarrow$ visibility\_semantics $\rightarrow$ instance\_naming}\Field{Reusable move.}{Introduce one formal component and state its role in the surrounding construction.}}}
\needspace{10\baselineskip}\subsubsection*{W-05-P007: same purpose}
\Field{Location.}{Section 05, prose\_paragraph, lines 900--906.}
\Field{Draft purpose.}{Interpret four-position sets and their group orbits as deck information.}
\Field{Qualitative closeness.}{same purpose. Both units explain how a formal representation determines the mathematical meaning of a concrete card object and fix the resulting terminology.}
\Field{Limits of the analogy.}{The baseline interprets card orientation, transformation, and visibility. The WADT unit interprets unlabeled heart positions and their orbit.}
\Field{Baseline narration used.}{full.}
\Field{Complexity and jargon.}{Complexity: \Tag{medium} $\rightarrow$ \Tag{low}. Jargon: \Tag{medium} $\rightarrow$ \Tag{low}.}
\textbf{Original WADT unit. Footnotes omitted.}
\begin{lstlisting}
Each set $A\in\Omega_4$ records the four positions occupied by hearts without
recording which heart occupies each position. Throughout this subsection, an
orbit means an orbit in $\Omega_4$. Two heart-position patterns are
equivalent when an allowed group action carries one to the other. A valid
labelled deck therefore determines one representative of a heart-position
orbit.
\end{lstlisting}
\colorbox{RewriteGreen}{\parbox{0.97\textwidth}{\textbf{Rewrite suggestion.}\par\begingroup\small
Each set $A\in\Omega_4$ records the four positions occupied by hearts but forgets the individual labels of those four hearts. In this subsection, every orbit is an orbit in $\Omega_4$. Two heart-position patterns are equivalent when an allowed group action carries one to the other. Each valid labelled deck therefore selects one orbit representative, namely its concrete set of heart positions.
\par\endgroup}}
\Field{Why this rewrite.}{Reuses the baseline's complete interpretation sequence to state what the set retains, what equivalence means, and what a deck supplies.}
\Field{Terminology actions.}{Keep: orbit representative | Explain: without recording which heart -> forgets the individual labels of the four hearts}
\par\smallskip
\clearpage\subsection{Baseline S-02-P008: sequence-construction explanation}
\colorbox{BaselineBlue}{\parbox{0.97\textwidth}{\Field{Purpose.}{Show which transformations generate the stated sequence instance.}\Field{Primary purpose.}{example}\Field{Narration structure.}{construction\_rationale $\rightarrow$ transformation\_expansion}\Field{Reusable move.}{Use a bounded worked instance to make an abstract definition concrete.}}}
\needspace{10\baselineskip}\subsubsection*{W-05-STM008: analogous purpose}
\Field{Location.}{Section 05, statement, lines 1219--1233.}
\Field{Draft purpose.}{Work through the generating maps and one word that transports an ordered triple.}
\Field{Qualitative closeness.}{analogous purpose. Both units explain how named transformations generate a concrete sequence from an initial arrangement.}
\Field{Limits of the analogy.}{The baseline works with one short card sequence. The WADT statement covers three maps, whole-group reachability, and one triple-transport word.}
\Field{Baseline narration used.}{adapted.}
\Field{Complexity and jargon.}{Complexity: \Tag{high} $\rightarrow$ \Tag{medium}. Jargon: \Tag{high} $\rightarrow$ \Tag{medium}.}
\textbf{Original WADT unit. Footnotes omitted.}
\begin{lstlisting}
Use projective-point labels to display the identity assignment, with label
$z$ at position $z$.
The translation map turns the row $(0,1,2,3,4,5,6,\infty)$ into
$(1,2,3,4,5,6,0,\infty)$ because position $i$ shows the image $g(i)$.
The scaling map fixes $0$ and $\infty$ and moves the six remaining
points in one six-cycle. The inversion map is an involution. It
exchanges $0$ with $\infty$, $1$ with $6$, $2$ with $3$, and $4$ with
$5$. Words in the three maps reach every one of the $336$ group
elements.
Lemma~\ref{thm:three-transitive} below gives one such word for every
ordered triple of distinct points. For example, the five-letter word that
applies the inverse scaling, then the inverse translation, then the
inversion, then the scaling, and then the inverse translation evaluates to
the map $z\mapsto(z-1)/(3-z)$, which carries $(0,1,2)$ to $(2,0,1)$.
\end{lstlisting}
\colorbox{RewriteGreen}{\parbox{0.97\textwidth}{\textbf{Rewrite suggestion.}\par\begingroup\small
Start with the identity row, which places label $z$ at position $z$. Translation changes $(0,1,2,3,4,5,6,\infty)$ to $(1,2,3,4,5,6,0,\infty)$ because position $i$ displays $g(i)$. Scaling fixes $0$ and $\infty$. It cycles the other six points in one six-cycle. Inversion is an involution, so applying it twice gives the identity. It exchanges $0$ with $\infty$, $1$ with $6$, $2$ with $3$, and $4$ with $5$. Words in these three maps reach all $336$ group elements. Lemma~\ref{thm:three-transitive} gives one such word for every ordered triple of distinct points. For example, apply inverse scaling, inverse translation, inversion, scaling, and inverse translation in that order. This five-letter generator word evaluates to $z\mapsto(z-1)/(3-z)$ and sends $(0,1,2)$ to $(2,0,1)$.
\par\endgroup}}
\Field{Why this rewrite.}{Adapts the baseline's worked-instance pattern to the three maps, whole-group reachability, and one explicit five-letter transport word.}
\Field{Terminology actions.}{Keep: involution, generator word | Explain: six-cycle -> cycles the six nonzero finite points, five-letter word -> the five named maps applied in the stated order}
\par\smallskip
\clearpage\subsection{Baseline S-02-P009: visible outcome of sequence instance}
\colorbox{BaselineBlue}{\parbox{0.97\textwidth}{\Field{Purpose.}{State the observable form represented by the worked sequence.}\Field{Primary purpose.}{interpretation}\Field{Narration structure.}{representation\_interpretation}\Field{Reusable move.}{Explain what the preceding formal outcome establishes.}}}
\needspace{10\baselineskip}\subsubsection*{W-04-P008: analogous purpose}
\Field{Location.}{Section 04, prose\_paragraph, lines 779--783.}
\Field{Draft purpose.}{Interpret the selected three-card and two-card cases after the figure.}
\Field{Qualitative closeness.}{analogous purpose. Both units interpret the observable meaning of a worked example after it has been displayed.}
\Field{Limits of the analogy.}{The baseline identifies one visible sequence. This unit compares several information values and one invariant across reveal shapes.}
\Field{Baseline narration used.}{adapted.}
\Field{Complexity and jargon.}{Complexity: \Tag{medium} $\rightarrow$ \Tag{low}. Jargon: \Tag{medium} $\rightarrow$ \Tag{low}.}
\textbf{Original WADT unit. Footnotes omitted.}
\begin{lstlisting}
Every three-card reveal leaks the same
$\tfrac65-\tfrac{9}{20}\log 3$ bits. The figure draws the consecutive
case, and the gapped case has the same value. Only the two-card value
depends on the pattern's shape.
\end{lstlisting}
\colorbox{RewriteGreen}{\parbox{0.97\textwidth}{\textbf{Rewrite suggestion.}\par\begingroup\small
Every three-card reveal leaks $\tfrac65-\tfrac{9}{20}\log 3$ bits. The consecutive and gapped cases therefore have the same value. Only two-card leakage depends on the relative revealed positions.
\par\endgroup}}
\Field{Why this rewrite.}{Adapts the baseline's interpretation move and presents the exact value before the invariance and contrast.}
\Field{Terminology actions.}{Keep: three-card reveal | Explain: gapped case | Replace: pattern's shape -> relative revealed positions}
\par\smallskip
\clearpage\subsection{Baseline S-02-P010: visible-sequence lead-in}
\colorbox{BaselineBlue}{\parbox{0.97\textwidth}{\Field{Purpose.}{Introduce the next derived object and hand its content to the formal statement.}\Field{Primary purpose.}{transition}\Field{Narration structure.}{concept\_naming $\rightarrow$ definition\_handoff}\Field{Reusable move.}{Use a short lead-in to name the role of what follows.}}}
\needspace{10\baselineskip}\subsubsection*{W-03-P011: analogous purpose}
\Field{Location.}{Section 03, prose\_paragraph, lines 585--588.}
\Field{Draft purpose.}{Motivate the theorem that transports view privacy to executed-trace privacy.}
\Field{Qualitative closeness.}{analogous purpose. Both units are short lead-ins that name the role of the formal object or result that immediately follows.}
\Field{Limits of the analogy.}{The baseline simply names a visible-sequence object before its definition. The WADT paragraph first contrasts static views with executed traces and adds the two-way determinacy condition required by the next theorem.}
\Field{Baseline narration used.}{adapted.}
\Field{Complexity and jargon.}{Complexity: \Tag{low} $\rightarrow$ \Tag{low}. Jargon: \Tag{medium} $\rightarrow$ \Tag{low}.}
\textbf{Original WADT unit. Footnotes omitted.}
\begin{lstlisting}
Three-transitivity establishes privacy of a coalition's card view, but the
interpreter records an executed trace. The next theorem lifts view
independence to trace secrecy when the trace and view determine each other.
\end{lstlisting}
\colorbox{RewriteGreen}{\parbox{0.97\textwidth}{\textbf{Rewrite suggestion.}\par\begingroup\small
Three-transitivity proves privacy for a coalition's card view. The interpreter instead records an executed trace. The next theorem derives executed-trace privacy when the trace is determined by the view and the view can be recovered from the trace.
\par\endgroup}}
\Field{Why this rewrite.}{The rewrite adapts the baseline's short concept handoff. It names the view-to-trace gap, replaces secrecy with the canonical privacy term, and explains the two-way determination condition.}
\Field{Terminology actions.}{Keep: executed trace, view independence | Explain: determine each other | Replace: trace secrecy -> executed-trace privacy}
\par\smallskip
\needspace{10\baselineskip}\subsubsection*{W-05-P008: same purpose}
\Field{Location.}{Section 05, prose\_paragraph, lines 907--910.}
\Field{Draft purpose.}{Motivate the cross-ratio classifier used by the decoder.}
\Field{Qualitative closeness.}{same purpose. Both units introduce the next derived object by explaining the need it satisfies and then hand its content to a formal definition.}
\Field{Limits of the analogy.}{The baseline introduces visible sequences. The WADT unit introduces an algebraic classifier and decoder.}
\Field{Baseline narration used.}{full.}
\Field{Complexity and jargon.}{Complexity: \Tag{low} $\rightarrow$ \Tag{low}. Jargon: \Tag{medium} $\rightarrow$ \Tag{low}.}
\textbf{Original WADT unit. Footnotes omitted.}
\begin{lstlisting}
To decode the orbit label from a deck, we first extract the positions
occupied by hearts. We then need an explicit test that identifies the orbit
of the resulting four-position set. The cross ratio supplies this test.
\end{lstlisting}
\colorbox{RewriteGreen}{\parbox{0.97\textwidth}{\textbf{Rewrite suggestion.}\par\begingroup\small
To decode an orbit label from a deck, I first extract the heart positions. I then need a test whose value is constant within each orbit and differs between the two orbits. The cross ratio supplies this test.
\par\endgroup}}
\Field{Why this rewrite.}{Reuses the baseline's complete motivation and definition handoff while stating the classifier requirement directly.}
\Field{Terminology actions.}{Keep: cross ratio | Explain: explicit test -> constant within each orbit and different between the two orbits}
\par\smallskip
\clearpage\subsection{Baseline S-02-P013: deterministic and probabilistic conversion distinction}
\colorbox{BaselineBlue}{\parbox{0.97\textwidth}{\Field{Purpose.}{Partition state-changing operations according to whether their output is fixed or random.}\Field{Primary purpose.}{definition}\Field{Narration structure.}{deterministic\_case $\rightarrow$ probabilistic\_case}\Field{Reusable move.}{Introduce one formal component and state its role in the surrounding construction.}}}
\needspace{10\baselineskip}\subsubsection*{W-03-P006: analogous purpose}
\Field{Location.}{Section 03, prose\_paragraph, lines 487--497.}
\Field{Draft purpose.}{Define the security witness's evidence slots and explain how the five instances fill them.}
\Field{Qualitative closeness.}{analogous purpose. Both paragraphs classify one formal component according to alternative evidence or behavior and explain the resulting cases.}
\Field{Limits of the analogy.}{The baseline divides conversions into deterministic and probabilistic classes. The WADT paragraph divides witness evidence into exact and asymptotic slots, gives instance examples, and records a deliberate omission.}
\Field{Baseline narration used.}{adapted.}
\Field{Complexity and jargon.}{Complexity: \Tag{high} $\rightarrow$ \Tag{medium}. Jargon: \Tag{high} $\rightarrow$ \Tag{medium}.}
\textbf{Original WADT unit. Footnotes omitted.}
\begin{lstlisting}
The record that carries the shuffle distribution and its endpoint bound
is called the security witness. It can carry exact evidence, asymptotic
evidence, or both. Table~\ref{tab:witness-mechanism} shows the combinations used by the
five instances in this paper. For example, the den Boer row fills only
the exact slot, since one uniform cut already gives the exact endpoint
distribution. The $\PG$ row also fills only the exact slot. The
word-shuffle bound of Theorem~B is proven in Section~\ref{sec:mixing} but
is not registered in the record's asymptotic slot, so the table reflects
the record contents and Theorem~B's evidence lives outside the witness
record.
\end{lstlisting}
\colorbox{RewriteGreen}{\parbox{0.97\textwidth}{\textbf{Rewrite suggestion.}\par\begingroup\small
The security witness records the shuffle distribution and its endpoint bound. It may record exact evidence, asymptotic evidence, or both. Table~\ref{tab:witness-mechanism} shows the combinations used by the five instances. The den Boer row records only exact evidence because one uniform cut gives the exact endpoint distribution. The $\PG$ row also records only exact evidence. Section~\ref{sec:mixing} proves the word-shuffle bound of Theorem~B. This bound is not stored in the record's optional asymptotic-evidence field. The table therefore reports the record contents, while the evidence for Theorem~B remains outside the security witness.
\par\endgroup}}
\Field{Why this rewrite.}{The rewrite adapts the baseline's case classification. It defines the witness, explains its optional evidence field, and separates table contents from the external theorem evidence without changing the record boundary.}
\Field{Terminology actions.}{Keep: security witness, exact evidence, asymptotic evidence | Explain: optional slot | Replace: fills only the exact slot -> records only exact evidence}
\par\smallskip
\clearpage\subsection{Baseline S-02-P015: permutation semantics and collection}
\colorbox{BaselineBlue}{\parbox{0.97\textwidth}{\Field{Purpose.}{Explain the positional effect of a permutation and collect all such operations into one set.}\Field{Primary purpose.}{definition}\Field{Narration structure.}{operation\_effect $\rightarrow$ position\_interpretation $\rightarrow$ collection\_definition}\Field{Reusable move.}{Introduce one formal component and state its role in the surrounding construction.}}}
\needspace{10\baselineskip}\subsubsection*{W-05-P014: analogous purpose}
\Field{Location.}{Section 05, prose\_paragraph, lines 1062--1065.}
\Field{Draft purpose.}{Explain the deck-action notation and state what the next lemma tracks.}
\Field{Qualitative closeness.}{analogous purpose. Both units explain the positional effect of a permutation before later results use that operation.}
\Field{Limits of the analogy.}{The baseline finishes by collecting a family of operations. The WADT unit instead changes level from deck entries to heart positions and bridges to an invariance lemma.}
\Field{Baseline narration used.}{adapted.}
\Field{Complexity and jargon.}{Complexity: \Tag{medium} $\rightarrow$ \Tag{low}. Jargon: \Tag{medium} $\rightarrow$ \Tag{low}.}
\textbf{Original WADT unit. Footnotes omitted.}
\begin{lstlisting}
The symbol $\star$ denotes this rearrangement of decks. Position $i$ in
$g\star D$ holds the value from position $\rho(g)(i)$ in $D$. The following
lemma tracks how this convention moves the heart positions.
\end{lstlisting}
\colorbox{RewriteGreen}{\parbox{0.97\textwidth}{\textbf{Rewrite suggestion.}\par\begingroup\small
The symbol $\star$ denotes this deck rearrangement. In $g\star D$, position $i$ contains the value that $D$ places at $\rho(g)(i)$. The set of new heart positions therefore pulls the old heart set back through $\rho(g)$, which produces the inverse position action. The following lemma states this heart-set transport.
\par\endgroup}}
\Field{Why this rewrite.}{Adapts the baseline's operation-and-position sequence and explains the exact change of level proved next.}
\Field{Terminology actions.}{Keep: heart-set transport | Explain: position convention -> pulling the old heart set back through the forward deck rule produces the inverse position action}
\par\smallskip
\clearpage\subsection{Baseline S-02-P018: turning-versus-opening distinction}
\colorbox{BaselineBlue}{\parbox{0.97\textwidth}{\Field{Purpose.}{Distinguish two operation classes by whether the observable sequence changes.}\Field{Primary purpose.}{comparison}\Field{Narration structure.}{potential\_confusion $\rightarrow$ classification\_decision $\rightarrow$ distinguishing\_criterion}\Field{Reusable move.}{Resolve a likely classification ambiguity by naming the decisive property.}}}
\needspace{10\baselineskip}\subsubsection*{W-01-P003: analogous purpose}
\Field{Location.}{Section 01, prose\_paragraph, lines 92--103.}
\Field{Draft purpose.}{Distinguish the ideal uniform shuffle from the finite generator-word model and preserve the physical-realization boundary.}
\Field{Qualitative closeness.}{analogous purpose. Both units distinguish two formal classes so that later claims use each one under the correct conditions.}
\Field{Limits of the analogy.}{The baseline separates turning from opening by observability. The WADT paragraph separates ideal and generator-word shuffle laws, and it must retain the stronger boundary that no physical realization is claimed for the main-instance letters.}
\Field{Baseline narration used.}{adapted.}
\Field{Complexity and jargon.}{Complexity: \Tag{high} $\rightarrow$ \Tag{medium}. Jargon: \Tag{high} $\rightarrow$ \Tag{medium}.}
\textbf{Original WADT unit. Footnotes omitted.}
\begin{lstlisting}
Two shuffle distributions organize the analysis. The uniform distribution
on the shuffle group is the ideal object of the security proofs. No dealer
samples a uniform group element in one physical move. A dealer instead
repeats elementary moves drawn from a small generator alphabet, and the
repeated moves follow what I call a word distribution over the generators
of the group. In the five-card family the moves are genuine cuts, while
the $\PG$ letters are group generators whose physical realization I do not
claim. The group parameters pay for this connection. Transitivity of the
action yields coalition privacy, and generation by a small alphabet yields
a finite executable shuffle. In this paper I prove the security results in the uniform-shuffle model of Section~\ref{sec:model} and then
prove that the word distribution approximates it.
\end{lstlisting}
\colorbox{RewriteGreen}{\parbox{0.97\textwidth}{\textbf{Rewrite suggestion.}\par\begingroup\small
The analysis needs two shuffle distributions because the ideal security proof and the finite implementation use different laws. The uniform distribution on the shuffle group is the ideal law for the security proofs. No dealer samples a uniform group element in one physical move. Instead, the dealer repeatedly draws elementary moves from a small generator alphabet. These draws define the word distribution over group generators. The physical status of the generator letters depends on the instance. In the five-card family, the letters are physical cuts. The $\PG$ letters are group generators, and I do not claim a physical realization for them. The group parameters connect the two analyses. Transitivity of the action gives coalition privacy under the ideal law. By contrast, a small generating alphabet gives a finite executable shuffle whose word distribution can approximate the ideal law. I first prove security in the uniform-shuffle model of Section~\ref{sec:model}. I then prove that the word distribution approximates the uniform law.
\par\endgroup}}
\Field{Why this rewrite.}{The rewrite states why the analysis needs two shuffle laws. It connects transitivity to ideal-law privacy and finite generation to executable word shuffles. It also preserves the distinction between physical five-card cuts and the $\PG$ generators, for which I claim no physical realization. The closing two sentences prepare the separate ideal and finite results announced next.}
\Field{Terminology actions.}{Keep: uniform distribution, word distribution, transitivity | Explain: generator alphabet, physical realization | Replace: the group parameters pay for this connection -> the group parameters connect the two analyses}
\par\smallskip
\needspace{10\baselineskip}\subsubsection*{W-02-P011: partial analogue}
\Field{Location.}{Section 02, prose\_paragraph, lines 271--277.}
\Field{Draft purpose.}{Place the model beside secret sharing and state its security exclusions.}
\Field{Qualitative closeness.}{partial analogue. The baseline provides a precise contrast that prevents two nearby operation classes from being conflated.}
\Field{Limits of the analogy.}{The WADT unit uses an external analogy and then sets several security boundaries. Only the boundary-setting move transfers.}
\Field{Baseline narration used.}{partial.}
\Field{Complexity and jargon.}{Complexity: \Tag{medium} $\rightarrow$ \Tag{low}. Jargon: \Tag{medium} $\rightarrow$ \Tag{low}.}
\textbf{Original WADT unit. Footnotes omitted.}
\begin{lstlisting}
The protected object is a dealt secret, as in a secret-sharing
scheme~\cite{Shamir1979}. The dealer holds the secret, the players hold
card observations in place of shares, and privacy bounds what coalitions of
curious players learn. The protocol has no private player inputs. Active
deviation, security after the verifier's reveal, and composition across
executions are outside the model.
\end{lstlisting}
\colorbox{RewriteGreen}{\parbox{0.97\textwidth}{\textbf{Rewrite suggestion.}\par\begingroup\small
The protocol protects a dealt secret, as in a secret-sharing scheme~\cite{Shamir1979}. The dealer holds the secret. The players hold card observations that play the role of shares. Privacy limits what coalitions of curious players learn. The protocol gives the players no private inputs. The model does not cover active deviation. It also makes no security claim after the full deck is publicly revealed, and it gives no composition theorem for multiple executions.
\par\endgroup}}
\Field{Why this rewrite.}{Uses the baseline boundary-setting move after the secret-sharing comparison and states each excluded security property separately.}
\Field{Terminology actions.}{Keep: secret-sharing scheme, active deviation | Explain: security after the verifier's reveal -> no security claim after the full deck is publicly revealed, composition across executions -> no composition theorem for multiple executions}
\par\smallskip
\needspace{10\baselineskip}\subsubsection*{W-04-P006: partial analogue}
\Field{Location.}{Section 04, prose\_paragraph, lines 716--728.}
\Field{Draft purpose.}{Distinguish output leakage from input privacy and state the required conditioning.}
\Field{Qualitative closeness.}{partial analogue. The baseline supplies the useful contrast move of distinguishing two observation properties by the information they preserve.}
\Field{Limits of the analogy.}{The baseline compares operation classes. This unit distinguishes two security notions and justifies a conditioning variable.}
\Field{Baseline narration used.}{partial.}
\Field{Complexity and jargon.}{Complexity: \Tag{high} $\rightarrow$ \Tag{medium}. Jargon: \Tag{high} $\rightarrow$ \Tag{medium}.}
\textbf{Original WADT unit. Footnotes omitted.}
\begin{lstlisting}
The leakage numbers above concern the announced output. The five-card
property of the literature is input-hiding given that output: the observer
learns nothing about the input pair beyond the conjunction it announces.
The development proves this property for every reveal pattern. The input
pair and the view of any reveal set are conditionally independent given the
announced conjunction, and the corresponding conditional mutual information
is zero.
The conditioning on the announced output is essential. The view contains
the announcement, so without the conditioning the independence fails.
\end{lstlisting}
\colorbox{RewriteGreen}{\parbox{0.97\textwidth}{\textbf{Rewrite suggestion.}\par\begingroup\small
The leakage values above concern the announced output. The five-card property in the literature is input-hiding given that output. An observer learns nothing about the input pair beyond the conjunction that the protocol announces. The formalization proves this property for every reveal pattern. Given the announced conjunction, the input pair and the view of any reveal set are conditionally independent, and their conditional mutual information is zero. Conditioning on the announced output is essential. The view contains the announcement, so independence fails without this conditioning.
\par\endgroup}}
\Field{Why this rewrite.}{Uses the baseline's distinguishing-criterion move to separate output leakage from input-hiding. It preserves the conditional-independence claim, its equivalent information claim, and the essential nonclaim.}
\Field{Terminology actions.}{Keep: conditionally independent, conditional mutual information, announced conjunction | Explain: given the announced conjunction | Replace: The development -> The formalization}
\par\smallskip
\needspace{10\baselineskip}\subsubsection*{W-05-P005: partial analogue}
\Field{Location.}{Section 05, prose\_paragraph, lines 871--877.}
\Field{Draft purpose.}{Explain the orbit-class encoding and separate correctness from privacy.}
\Field{Qualitative closeness.}{partial analogue. The baseline supplies the move of distinguishing two nearby concepts by the property that changes under an operation.}
\Field{Limits of the analogy.}{The WADT unit primarily introduces an invariant encoding and then separates correctness from privacy. Only the distinction move transfers.}
\Field{Baseline narration used.}{partial.}
\Field{Complexity and jargon.}{Complexity: \Tag{medium} $\rightarrow$ \Tag{low}. Jargon: \Tag{high} $\rightarrow$ \Tag{medium}.}
\textbf{Original WADT unit. Footnotes omitted.}
\begin{lstlisting}
The $\PG$ action on $X$ induces an action on sets of four card positions.
I store the secret bit as the orbit class of the four positions occupied by
hearts. Every allowed shuffle moves this set within the same orbit, so its
orbit label remains unchanged. This invariance will give correctness under
shuffling. Privacy is a separate property of the partial observations made
by a coalition.
\end{lstlisting}
\colorbox{RewriteGreen}{\parbox{0.97\textwidth}{\textbf{Rewrite suggestion.}\par\begingroup\small
The $\PG$ action on $X$ also moves sets of four card positions by moving every position in the set. I encode the secret bit as the orbit class of the four positions occupied by hearts. Every allowed shuffle keeps this set in the same orbit, so the orbit label does not change. This invariance gives correctness under shuffling. Privacy requires a separate argument about the card values that a coalition sees at its named positions.
\par\endgroup}}
\Field{Why this rewrite.}{Uses the baseline distinction move only to separate correctness from privacy after stating the encoding and invariant.}
\Field{Terminology actions.}{Keep: orbit class, invariance | Explain: induced action on sets -> moves every position in the set, partial observation -> the card values that a coalition sees at its named positions | Replace: I store the secret bit -> I encode the secret bit}
\par\smallskip
\needspace{10\baselineskip}\subsubsection*{W-05-P017: partial analogue}
\Field{Location.}{Section 05, prose\_paragraph, lines 1113--1118.}
\Field{Draft purpose.}{Separate correctness from privacy and name the privacy premise.}
\Field{Qualitative closeness.}{partial analogue. The baseline supplies the useful move of distinguishing two related properties by what the operation changes or reveals.}
\Field{Limits of the analogy.}{The WADT paragraph primarily separates a proved guarantee from an unproved one and adds a distributional premise.}
\Field{Baseline narration used.}{partial.}
\Field{Complexity and jargon.}{Complexity: \Tag{medium} $\rightarrow$ \Tag{low}. Jargon: \Tag{medium} $\rightarrow$ \Tag{low}.}
\textbf{Original WADT unit. Footnotes omitted.}
\begin{lstlisting}
By Corollary~\ref{cor:orbit-shuffle-correctness}, every allowed shuffle
preserves the decoded secret of an encoded deck. This correctness result
does not imply privacy. When $g\leftarrow U_G$ is independent of the secret,
privacy is a property of the distribution of partial observations. The
privacy proof uses three-transitivity as its group-action premise.
\end{lstlisting}
\colorbox{RewriteGreen}{\parbox{0.97\textwidth}{\textbf{Rewrite suggestion.}\par\begingroup\small
Corollary~\ref{cor:orbit-shuffle-correctness} shows that every allowed shuffle preserves the decoded secret of an encoded deck. Correctness and privacy require separate arguments. Privacy concerns the distribution of partial observations when $g\leftarrow U_G$ is uniform and independent of the secret. A coalition's partial observation contains only the card values at its named positions. The privacy proof uses three-transitivity, which lets the action carry any ordered triple of distinct positions to any other. Before I give that proof, I use the next figure to return from the deck equality to one complete execution.
\par\endgroup}}
\Field{Why this rewrite.}{The rewrite preserves the correctness-to-privacy distinction and every premise. Its final sentence explains why the execution figure intervenes before the privacy proof.}
\Field{Terminology actions.}{Keep: uniform group distribution, partial observation | Explain: three-transitivity -> the action carries any ordered triple of distinct positions to any other}
\par\smallskip
\needspace{10\baselineskip}\subsubsection*{W-05-P023: partial analogue}
\Field{Location.}{Section 05, prose\_paragraph, lines 1248--1252.}
\Field{Draft purpose.}{Formulate privacy as the distribution of named-position observations.}
\Field{Qualitative closeness.}{partial analogue. The baseline supplies the move of distinguishing two nearby notions through their different observational behavior.}
\Field{Limits of the analogy.}{The WADT paragraph also defines a probability question and explains why order matters.}
\Field{Baseline narration used.}{partial.}
\Field{Complexity and jargon.}{Complexity: \Tag{medium} $\rightarrow$ \Tag{low}. Jargon: \Tag{medium} $\rightarrow$ \Tag{low}.}
\textbf{Original WADT unit. Footnotes omitted.}
\begin{lstlisting}
Global orbit invariance gives correctness. Privacy asks a different question:
what does a small named set of positions observe when $g\leftarrow U_G$ is
independent and uniform? Ordered tuples matter because a coalition observes
card values at named positions.
\end{lstlisting}
\colorbox{RewriteGreen}{\parbox{0.97\textwidth}{\textbf{Rewrite suggestion.}\par\begingroup\small
Global orbit invariance proves correctness. Privacy asks what card values a fixed coalition sees when $g\leftarrow U_G$ is independent and uniform. The coalition observes named positions, so an ordered tuple must preserve which player saw each value.
\par\endgroup}}
\Field{Why this rewrite.}{Uses the baseline contrast move to separate global correctness from the local observation distribution that defines privacy.}
\Field{Terminology actions.}{Keep: ordered tuple | Explain: named set of positions -> a fixed coalition whose tuple records which player saw each value}
\par\smallskip
\needspace{10\baselineskip}\subsubsection*{W-06-P006: analogous purpose}
\Field{Location.}{Section 06, prose\_paragraph, lines 1344--1347.}
\Field{Draft purpose.}{Contrast the privacy and recovery cutoffs and direct the reader to their visual separation.}
\Field{Qualitative closeness.}{analogous purpose. Both paragraphs separate two closely related mechanisms by the different property each changes.}
\Field{Limits of the analogy.}{The baseline contrasts transformation and observation. The WADT paragraph contrasts privacy and recovery cutoffs and points to a two-track figure.}
\Field{Baseline narration used.}{adapted.}
\Field{Complexity and jargon.}{Complexity: \Tag{low} $\rightarrow$ \Tag{low}. Jargon: \Tag{medium} $\rightarrow$ \Tag{low}.}
\textbf{Original WADT unit. Footnotes omitted.}
\begin{lstlisting}
The recovery result and the privacy result have different cutoffs. Perfect
coalition privacy holds through size three, while deterministic recovery
begins at size seven. Figure~\ref{fig:ramp} places them on separate tracks.
\end{lstlisting}
\colorbox{RewriteGreen}{\parbox{0.97\textwidth}{\textbf{Rewrite suggestion.}\par\begingroup\small
Recovery and privacy have different cutoffs. Perfect coalition privacy holds for coalitions of at most three positions. Deterministic recovery starts with seven revealed positions. Figure~\ref{fig:ramp} shows these cutoffs on separate tracks.
\par\endgroup}}
\Field{Why this rewrite.}{The rewrite adapts the baseline's explicit comparison move. It scopes privacy by coalition size, scopes recovery by revealed positions, and preserves the figure reference.}
\Field{Terminology actions.}{Keep: perfect privacy, deterministic recovery}
\par\smallskip
\needspace{10\baselineskip}\subsubsection*{W-06-P013: analogous purpose}
\Field{Location.}{Section 06, prose\_paragraph, lines 1510--1512.}
\Field{Draft purpose.}{Transition from ideal-shuffle security to finite word-shuffle approximation.}
\Field{Qualitative closeness.}{analogous purpose. Both units mark a boundary between two related mechanisms before the text develops the second one.}
\Field{Limits of the analogy.}{The baseline contrasts two card operations. The WADT transition contrasts an ideal uniform element with a finite generator word and announces an error bound.}
\Field{Baseline narration used.}{adapted.}
\Field{Complexity and jargon.}{Complexity: \Tag{low} $\rightarrow$ \Tag{low}. Jargon: \Tag{medium} $\rightarrow$ \Tag{low}.}
\textbf{Original WADT unit. Footnotes omitted.}
\begin{lstlisting}
Theorem A uses a uniform group element. The next section replaces this
ideal shuffle with a finite generator word and bounds the resulting error.
\end{lstlisting}
\colorbox{RewriteGreen}{\parbox{0.97\textwidth}{\textbf{Rewrite suggestion.}\par\begingroup\small
Theorem A assumes a uniform group element. The next section replaces this ideal shuffle with a finite generator word and quantifies the resulting error.
\par\endgroup}}
\Field{Why this rewrite.}{The rewrite adapts the baseline's explicit distinction move. It preserves the ideal assumption, finite replacement, forward transition, and error-bound promise.}
\Field{Terminology actions.}{Keep: uniform group element, generator word, error bound}
\par\smallskip
\needspace{10\baselineskip}\subsubsection*{W-07-P007: partial analogue}
\Field{Location.}{Section 07, prose\_paragraph, lines 1604--1609.}
\Field{Draft purpose.}{Prevent two overreadings of the endpoint and product transfers.}
\Field{Qualitative closeness.}{partial analogue. The baseline supplies the contrast move of distinguishing formal objects by what their observations contain.}
\Field{Limits of the analogy.}{The baseline compares operation classes. This unit corrects two possible claims about proof provenance and sample space.}
\Field{Baseline narration used.}{partial.}
\Field{Complexity and jargon.}{Complexity: \Tag{medium} $\rightarrow$ \Tag{low}. Jargon: \Tag{high} $\rightarrow$ \Tag{medium}.}
\textbf{Original WADT unit. Footnotes omitted.}
\begin{lstlisting}
Equations~\ref{eq:endpoint-transfer} and~\ref{eq:product-transfer}
transfer the one certificate of Equation~\ref{eq:pgl-mixing} to two
further sample spaces. Neither is a separate mixing certificate, and
Equation~\ref{eq:product-transfer} still compares a prior with a shuffle
distribution rather than a distribution over executed runs.
\end{lstlisting}
\colorbox{RewriteGreen}{\parbox{0.97\textwidth}{\textbf{Rewrite suggestion.}\par\begingroup\small
Equations~\ref{eq:endpoint-transfer} and~\ref{eq:product-transfer} apply the single certificate in Equation~\ref{eq:pgl-mixing} to two additional sample spaces. Neither equation is a separate mixing certificate. Equation~\ref{eq:product-transfer} compares a secret prior paired with a shuffle distribution. It does not yet describe executed runs.
\par\endgroup}}
\Field{Why this rewrite.}{Uses only the baseline's distinguishing-criterion move to separate proof provenance and sample spaces without changing either nonclaim.}
\Field{Terminology actions.}{Keep: mixing certificate, executed run | Explain: prior with a shuffle distribution}
\par\smallskip
\needspace{10\baselineskip}\subsubsection*{W-08-P002: partial analogue}
\Field{Location.}{Section 08, prose\_paragraph, lines 1733--1736.}
\Field{Draft purpose.}{Identify the PGL-specific proof coverage and its assumption boundary.}
\Field{Qualitative closeness.}{partial analogue. The baseline supplies the move of distinguishing two nearby mechanisms by the condition that one does not use.}
\Field{Limits of the analogy.}{The WADT unit is an instance-coverage summary, not an operation comparison.}
\Field{Baseline narration used.}{partial.}
\Field{Complexity and jargon.}{Complexity: \Tag{medium} $\rightarrow$ \Tag{low}. Jargon: \Tag{high} $\rightarrow$ \Tag{medium}.}
\textbf{Original WADT unit. Footnotes omitted.}
\begin{lstlisting}
The PGL instance adds coalition views, coalition traces, the recovery ramp,
and the finite fiber certificate from Section~\ref{sec:mixing}. Its mixing
result does not depend on the $S_5$ Rayleigh premise.
\end{lstlisting}
\colorbox{RewriteGreen}{\parbox{0.97\textwidth}{\textbf{Rewrite suggestion.}\par\begingroup\small
The PGL instance adds coalition views, coalition traces, the recovery ramp, and the finite fiber certificate from Section~\ref{sec:mixing}. A coalition trace collects the executed traces of the corrupted players. The finite fiber certificate gives exact evidence for finite-word mixing. This mixing result does not use the imported Rayleigh premise required by the $S_5$ spectral route.
\par\endgroup}}
\Field{Why this rewrite.}{Uses the baseline assumption-boundary move after listing the PGL-specific coverage and explaining the two dense terms.}
\Field{Terminology actions.}{Keep: recovery ramp, finite fiber certificate | Explain: Rayleigh premise -> the imported analytic assumption required by the S5 spectral route, coalition trace -> the collection of executed traces of the corrupted players}
\par\smallskip
\needspace{10\baselineskip}\subsubsection*{W-08-P004: partial analogue}
\Field{Location.}{Section 08, prose\_paragraph, lines 1786--1793.}
\Field{Draft purpose.}{Explain the two spectral instances and the product instance's nonzero mixing floor.}
\Field{Qualitative closeness.}{partial analogue. The baseline supplies the move of contrasting related constructions through the state change that one operation cannot perform.}
\Field{Limits of the analogy.}{The WADT paragraph also reports spectral methods, an inherited assumption, asymptotic behavior, and an exact limitation.}
\Field{Baseline narration used.}{partial.}
\Field{Complexity and jargon.}{Complexity: \Tag{high} $\rightarrow$ \Tag{medium}. Jargon: \Tag{high} $\rightarrow$ \Tag{medium}.}
\textbf{Original WADT unit. Footnotes omitted.}
\begin{lstlisting}
The $S_5$ instance uses additive sharing modulo five. Its endpoint mixing
proof applies a spectral inequality at word length 286. The development
imports the required Rayleigh bound as an axiom. The $S_5\times S_5$ instance
uses two disjoint piles. It inherits the same Rayleigh premise for decay
within either pile. Because a card never changes piles, its distance from the uniform
distribution on all ten positions cannot converge to zero. The formal global
bound approaches the exact $L_1$ floor of one.
\end{lstlisting}
\colorbox{RewriteGreen}{\parbox{0.97\textwidth}{\textbf{Rewrite suggestion.}\par\begingroup\small
The $S_5$ instance represents its secret by additive sharing modulo five. Its endpoint mixing proof applies a spectral inequality at word length 286. This inequality controls convergence at that length. The development imports the required Rayleigh bound as an axiom. The $S_5\times S_5$ instance uses two disjoint piles and inherits the same Rayleigh premise. The premise gives decay within each five-card pile. A card never moves between the piles. Its endpoint distribution therefore cannot converge to the uniform distribution on all ten positions, which assigns mass to positions in both piles. The formal global bound approaches the exact $L_1$ floor of one.
\par\endgroup}}
\Field{Why this rewrite.}{Uses the baseline invariant-and-consequence move to explain the product obstruction after stating the two spectral methods and their shared premise.}
\Field{Terminology actions.}{Keep: additive sharing, Rayleigh bound, L1 floor of one | Explain: spectral inequality -> controls convergence at word length 286, within-pile decay -> decay within each five-card pile, uniform distribution on all ten positions -> a target law that assigns mass to positions in both piles}
\par\smallskip
\needspace{10\baselineskip}\subsubsection*{W-09-P003: partial analogue}
\Field{Location.}{Section 09, prose\_paragraph, lines 1810--1817.}
\Field{Draft purpose.}{Compare Shinagawa's unified card model with the paper's group-parametric framework.}
\Field{Qualitative closeness.}{partial analogue. The baseline supplies the move of distinguishing two formal classes by the feature that changes.}
\Field{Limits of the analogy.}{The baseline compares operations within one model. This unit positions two papers and adds proof-assistant outputs.}
\Field{Baseline narration used.}{partial.}
\Field{Complexity and jargon.}{Complexity: \Tag{medium} $\rightarrow$ \Tag{low}. Jargon: \Tag{high} $\rightarrow$ \Tag{medium}.}
\textbf{Original WADT unit. Footnotes omitted.}
\begin{lstlisting}
Shinagawa gives a unified model in which a card protocol is specified by a
deck and a set of operations~\cite{Shinagawa2021}. The model covers binary
cards, regular polygon cards, and dihedral cards. Although Shinagawa's model varies the card type and the allowed physical
operations, the framework in this paper varies the finite group, its
permutation action, the shuffle distribution, and the reconstruction
map. It connects these parameters to executable
traces and machine-checked privacy and mixing results.
\end{lstlisting}
\colorbox{RewriteGreen}{\parbox{0.97\textwidth}{\textbf{Rewrite suggestion.}\par\begingroup\small
Shinagawa specifies a card protocol through a deck and a set of operations in a unified model~\cite{Shinagawa2021}. This model covers binary, regular polygon, and dihedral cards. Shinagawa varies the card type and allowed physical operations. This paper instead varies the finite group, permutation action, shuffle distribution, and reconstruction map. It connects these parameters with executable traces and machine-checked privacy and mixing results.
\par\endgroup}}
\Field{Why this rewrite.}{Uses only the baseline's distinguishing-criterion move. It separates the two parameterizations and retains the formal outputs.}
\Field{Terminology actions.}{Keep: permutation action, shuffle distribution, reconstruction map | Explain: unified model, executable trace}
\par\smallskip
\clearpage\subsection{Baseline S-02-P019: shuffle sampling semantics}
\colorbox{BaselineBlue}{\parbox{0.97\textwidth}{\Field{Purpose.}{Define a randomised sequence transformation through a permutation family and a distribution.}\Field{Primary purpose.}{definition}\Field{Narration structure.}{parameter\_tuple $\rightarrow$ sampling\_rule $\rightarrow$ sequence\_update $\rightarrow$ collection\_handoff}\Field{Reusable move.}{Introduce one formal component and state its role in the surrounding construction.}}}
\needspace{10\baselineskip}\subsubsection*{W-02-P009: analogous purpose}
\Field{Location.}{Section 02, prose\_paragraph, lines 254--261.}
\Field{Draft purpose.}{Distinguish the fixed representative dealer from the all-decks dealer.}
\Field{Qualitative closeness.}{analogous purpose. Both units define a randomized transformation by separating the sampled object, its distribution, and its effect.}
\Field{Limits of the analogy.}{The baseline defines one shuffle operation. The WADT unit compares two dealer distributions that share a secret prior and later shuffle.}
\Field{Baseline narration used.}{adapted.}
\Field{Complexity and jargon.}{Complexity: \Tag{medium} $\rightarrow$ \Tag{low}. Jargon: \Tag{medium} $\rightarrow$ \Tag{low}.}
\textbf{Original WADT unit. Footnotes omitted.}
\begin{lstlisting}
The primary execution distribution samples a uniform Boolean secret and
deals one fixed representative arrangement $D_s$ for each secret $s$. It
then applies an independent uniform shuffle. I call this the fixed
representative dealer, and it is the dealer of the headline theorems. An
all-decks dealer instead
samples a uniform valid arrangement in the selected secret class before the
shuffle. Both execution distributions use a uniform Boolean prior.
\end{lstlisting}
\colorbox{RewriteGreen}{\parbox{0.97\textwidth}{\textbf{Rewrite suggestion.}\par\begingroup\small
The primary execution distribution samples a Boolean secret uniformly. For each secret $s$, it deals one fixed valid deck $D_s$ that represents the class of $s$. It then applies an independent uniform shuffle. I call this distribution the fixed representative dealer, and the headline theorems use it. The all-decks dealer instead samples a valid arrangement uniformly from the selected secret class before applying the shuffle. Both dealer distributions use a uniform Boolean prior, so the two secret bits have equal probability.
\par\endgroup}}
\Field{Why this rewrite.}{Adapts the baseline's sampling-rule sequence to contrast the two dealer distributions and states their common prior last.}
\Field{Terminology actions.}{Keep: fixed representative dealer, all-decks dealer | Explain: representative arrangement -> one fixed valid deck that represents each secret class, uniform Boolean prior -> the two secret bits have equal probability}
\par\smallskip
\needspace{10\baselineskip}\subsubsection*{W-02-P012: analogous purpose}
\Field{Location.}{Section 02, prose\_paragraph, lines 278--285.}
\Field{Draft purpose.}{Define the uniform-shuffle and word-shuffle models and introduce their distance measure.}
\Field{Qualitative closeness.}{analogous purpose. Both units define a randomized sequence transformation through an explicit distribution and its operational effect.}
\Field{Limits of the analogy.}{The baseline introduces a general shuffle operation. The WADT unit contrasts an ideal group law with an iterated word law and adds a distance measure.}
\Field{Baseline narration used.}{adapted.}
\Field{Complexity and jargon.}{Complexity: \Tag{medium} $\rightarrow$ \Tag{low}. Jargon: \Tag{high} $\rightarrow$ \Tag{medium}.}
\textbf{Original WADT unit. Footnotes omitted.}
\begin{lstlisting}
The model described so far applies one uniform group element, and I call
it the uniform-shuffle model. A second model describes the dealer
performing the shuffle as a sequence of elementary moves drawn from the
generator alphabet. I call this the
word-shuffle model. It
replaces $U_G$ by $\mu^{*L}$. I measure two distributions $P$
and $Q$ by the $L_1$ distance
\end{lstlisting}
\colorbox{RewriteGreen}{\parbox{0.97\textwidth}{\textbf{Rewrite suggestion.}\par\begingroup\small
The uniform-shuffle model applies one uniform group element. The word-shuffle model instead lets the dealer apply a sequence of elementary moves from a generator alphabet, whose letters name the permitted elementary shuffles. It replaces $U_G$ with $\mu^{*L}$, the distribution obtained by evaluating $L$ independent letters. I compare distributions $P$ and $Q$ using the $L_1$ distance:
\par\endgroup}}
\Field{Why this rewrite.}{Adapts the baseline's random-transformation sequence to contrast the ideal and word laws before introducing their distance.}
\Field{Terminology actions.}{Keep: uniform-shuffle model, word-shuffle model, L1 distance | Explain: generator alphabet -> letters that name the permitted elementary shuffles, convolution power -> the distribution obtained by evaluating independent letters}
\par\smallskip
\needspace{10\baselineskip}\subsubsection*{W-05-P019: analogous purpose}
\Field{Location.}{Section 05, prose\_paragraph, lines 1165--1172.}
\Field{Draft purpose.}{Define generators, generator words, and their observable execution effect.}
\Field{Qualitative closeness.}{analogous purpose. Both units define a randomized sequence transformation from a family of permutations and explain how it acts on a deck.}
\Field{Limits of the analogy.}{The WADT unit adds group generation, word composition, and a non-observation condition for intermediate states.}
\Field{Baseline narration used.}{adapted.}
\Field{Complexity and jargon.}{Complexity: \Tag{high} $\rightarrow$ \Tag{low}. Jargon: \Tag{high} $\rightarrow$ \Tag{medium}.}
\textbf{Original WADT unit. Footnotes omitted.}
\begin{lstlisting}
The protocol uses the permutation action $\rho:G\to S_X$ on the eight card
positions. A generator is one of a fixed set of group elements whose finite
products reach every element of $G$. Under $\rho$, each generator specifies
a rearrangement of the eight cards. A generator word specifies a
sequence of these rearrangements. In the passive execution model, no party
records the intermediate arrangements, so the word enters the execution
only through its product $g\in G$ and the resulting deck $g\star D_s$.
\end{lstlisting}
\colorbox{RewriteGreen}{\parbox{0.97\textwidth}{\textbf{Rewrite suggestion.}\par\begingroup\small
The protocol uses the permutation action $\rho:G\to S_X$ on the eight card positions. A generator is one element of a fixed set whose finite products reach every element of $G$. Under $\rho$, each generator rearranges the eight cards once. A generator word lists a sequence of these elementary rearrangements. Evaluating the word gives its product $g\in G$. In the passive execution model, no party observes an intermediate arrangement. The word therefore affects execution only through $g$ and the final deck $g\star D_s$.
\par\endgroup}}
\Field{Why this rewrite.}{Adapts the baseline's parameter-to-effect sequence to define generators, words, their product, and the execution observation boundary.}
\Field{Terminology actions.}{Keep: generator, generator word | Explain: word product -> the group element obtained by evaluating the word, passive execution model -> no party observes an intermediate arrangement}
\par\smallskip
\needspace{10\baselineskip}\subsubsection*{W-07-STM002: partial analogue}
\Field{Location.}{Section 07, statement, lines 1570--1581.}
\Field{Draft purpose.}{State that endpoint pushforward preserves the mixing bound.}
\Field{Qualitative closeness.}{partial analogue. The baseline supplies the move of defining a distribution after a randomized transformation of a sequence.}
\Field{Limits of the analogy.}{The baseline defines shuffle sampling. This statement proves a bound under a deterministic observation map.}
\Field{Baseline narration used.}{partial.}
\Field{Complexity and jargon.}{Complexity: \Tag{medium} $\rightarrow$ \Tag{low}. Jargon: \Tag{high} $\rightarrow$ \Tag{medium}.}
\textbf{Original WADT unit. Footnotes omitted.}
\begin{lstlisting}
By the $L_1$ data-processing inequality, for every card position $i$,
\begin{equation}
  \left\lVert
    (g\mapsto \rho(g)i)_\#\mu^{*200}
    -(g\mapsto \rho(g)i)_\#U_G
  \right\rVert_1
  \leq 2^{-40}.
  \label{eq:endpoint-transfer}
\end{equation}
\end{lstlisting}
\colorbox{RewriteGreen}{\parbox{0.97\textwidth}{\textbf{Rewrite suggestion.}\par\begingroup\small
By the $L_1$ data-processing inequality, for every card position $i$,
\begin{equation}
  \left\lVert
    (g\mapsto \rho(g)i)_\#\mu^{*200}
    -(g\mapsto \rho(g)i)_\#U_G
  \right\rVert_1
  \leq 2^{-40}.
  \label{eq:endpoint-transfer}
\end{equation}
\par\endgroup}}
\Field{Why this rewrite.}{Uses only the baseline's induced-distribution move and preserves the universal position quantifier, both pushforwards, equation label, and bound.}
\Field{Terminology actions.}{Keep: data-processing inequality, pushforward | Explain: rho(g)i}
\par\smallskip
\clearpage\subsection{Baseline S-02-P021: single-operation observation setup}
\colorbox{BaselineBlue}{\parbox{0.97\textwidth}{\Field{Purpose.}{Model one operation as a state transition paired with revealed information.}\Field{Primary purpose.}{definition}\Field{Narration structure.}{object\_setup $\rightarrow$ operation\_application $\rightarrow$ state\_transition $\rightarrow$ revelation\_output}\Field{Reusable move.}{Introduce one formal component and state its role in the surrounding construction.}}}
\needspace{10\baselineskip}\subsubsection*{W-07-P005: partial analogue}
\Field{Location.}{Section 07, prose\_paragraph, lines 1564--1568.}
\Field{Draft purpose.}{Motivate the map from group distributions to one-card endpoints.}
\Field{Qualitative closeness.}{partial analogue. The baseline supplies the move of defining an observation by pairing system evolution with revealed information.}
\Field{Limits of the analogy.}{The baseline defines a one-operation view. This paragraph only motivates a pushforward from shuffles to one endpoint.}
\Field{Baseline narration used.}{partial.}
\Field{Complexity and jargon.}{Complexity: \Tag{low} $\rightarrow$ \Tag{low}. Jargon: \Tag{medium} $\rightarrow$ \Tag{low}.}
\textbf{Original WADT unit. Footnotes omitted.}
\begin{lstlisting}
A coalition observes card endpoints rather than group elements. The
certificate therefore moves from the group to the endpoints first. The
first transfer, which I call the endpoint transfer, maps each
permutation to the endpoint of one fixed card.
\end{lstlisting}
\colorbox{RewriteGreen}{\parbox{0.97\textwidth}{\textbf{Rewrite suggestion.}\par\begingroup\small
A coalition observes card endpoints, not group elements. The proof therefore transfers the bound from group elements to endpoints. The endpoint transfer maps each permutation to the endpoint of one fixed card.
\par\endgroup}}
\Field{Why this rewrite.}{Uses only the baseline's observation-definition move and states the observer, required transfer, and endpoint map.}
\Field{Terminology actions.}{Keep: coalition, endpoint | Explain: endpoint transfer | Replace: moves from the group to the endpoints -> transfers the bound from group elements to endpoints}
\par\smallskip
\clearpage\subsection{Baseline S-02-P023: multi-operation trace setup}
\colorbox{BaselineBlue}{\parbox{0.97\textwidth}{\Field{Purpose.}{Extend the one-step observation model to an ordered list of operations and intermediate states.}\Field{Primary purpose.}{definition}\Field{Narration structure.}{operation\_list\_setup $\rightarrow$ state\_trace}\Field{Reusable move.}{Introduce one formal component and state its role in the surrounding construction.}}}
\needspace{10\baselineskip}\subsubsection*{W-02-P006: analogous purpose}
\Field{Location.}{Section 02, prose\_paragraph, lines 221--227.}
\Field{Draft purpose.}{Define executed traces and state the verifier's designed knowledge.}
\Field{Qualitative closeness.}{analogous purpose. Both units extend an observation model across an ordered execution and name the resulting trace object.}
\Field{Limits of the analogy.}{The baseline records intermediate states and operations. The WADT unit records process inputs and separates coalition traces from the verifier's complete view.}
\Field{Baseline narration used.}{adapted.}
\Field{Complexity and jargon.}{Complexity: \Tag{medium} $\rightarrow$ \Tag{low}. Jargon: \Tag{medium} $\rightarrow$ \Tag{low}.}
\textbf{Original WADT unit. Footnotes omitted.}
\begin{lstlisting}
The interpreter executes the dealer, player, and verifier processes. It
records each value received by a process. The resulting record is called the
executed trace $T_i$, and $T_C=(T_i)_{i\in C}$ is called the coalition
trace. The
verifier receives all endpoints and applies the decoder. It therefore learns
$S$ by design.
\end{lstlisting}
\colorbox{RewriteGreen}{\parbox{0.97\textwidth}{\textbf{Rewrite suggestion.}\par\begingroup\small
The interpreter runs the dealer, player, and verifier processes. It records every value that a process receives. This record is the player's executed trace $T_i$. For a coalition $C$, the coalition trace $T_C=(T_i)_{i\in C}$ collects the traces of its players. The verifier receives all endpoints and applies the decoder, which recovers $S$. The verifier therefore learns $S$ by design.
\par\endgroup}}
\Field{Why this rewrite.}{Adapts the baseline's trace-building sequence and names each execution record before stating the verifier's designed knowledge.}
\Field{Terminology actions.}{Keep: executed trace, decoder | Explain: interpreter -> runs the dealer, player, and verifier processes, coalition trace -> collects the traces of the coalition's players}
\par\smallskip
\needspace{10\baselineskip}\subsubsection*{W-07-P006: analogous purpose}
\Field{Location.}{Section 07, prose\_paragraph, lines 1583--1588.}
\Field{Draft purpose.}{Extend the endpoint comparison to a joint secret-and-shuffle law.}
\Field{Qualitative closeness.}{analogous purpose. Both units extend an observation model from one component to a richer joint object needed by later security reasoning.}
\Field{Limits of the analogy.}{The baseline extends one-step views to traces. This unit adds an independent secret to two shuffle laws.}
\Field{Baseline narration used.}{adapted.}
\Field{Complexity and jargon.}{Complexity: \Tag{medium} $\rightarrow$ \Tag{low}. Jargon: \Tag{high} $\rightarrow$ \Tag{medium}.}
\textbf{Original WADT unit. Footnotes omitted.}
\begin{lstlisting}
The endpoint transfer controls one card position, and transitivity makes
its uniform-shuffle endpoint distribution uniform on the eight positions.
The privacy theorem also tracks the secret. The next lemma takes the product
of both shuffle distributions with the same independent Boolean prior and
preserves the $L_1$ bound.
\end{lstlisting}
\colorbox{RewriteGreen}{\parbox{0.97\textwidth}{\textbf{Rewrite suggestion.}\par\begingroup\small
The endpoint transfer controls one card position. Transitivity makes the endpoint distribution under the uniform shuffle uniform on the eight positions. The privacy theorem also includes the secret variable. The next lemma forms the product of each shuffle distribution with the same independent Boolean prior and preserves the $L_1$ bound.
\par\endgroup}}
\Field{Why this rewrite.}{Adapts the baseline's extension from one observation to a richer joint object. It preserves the transitivity consequence, secret variable, common prior, and bound.}
\Field{Terminology actions.}{Keep: transitivity, Boolean prior | Explain: product distribution | Replace: tracks the secret -> includes the secret variable}
\par\smallskip
\clearpage\subsection{Baseline S-02-P024: execution-view trace and collection}
\colorbox{BaselineBlue}{\parbox{0.97\textwidth}{\Field{Purpose.}{Pair each state observation with its revealed output, set the conversion convention, and name the set of views.}\Field{Primary purpose.}{definition}\Field{Narration structure.}{revealed\_output\_setup $\rightarrow$ paired\_trace $\rightarrow$ initial\_convention $\rightarrow$ conversion\_convention $\rightarrow$ view\_naming $\rightarrow$ collection\_definition}\Field{Reusable move.}{Introduce one formal component and state its role in the surrounding construction.}}}
\needspace{10\baselineskip}\subsubsection*{W-02-P008: partial analogue}
\Field{Location.}{Section 02, prose\_paragraph, lines 243--253.}
\Field{Draft purpose.}{Explain the process notation and connect interpreter execution to trace privacy.}
\Field{Qualitative closeness.}{partial analogue. The baseline supplies the move that assembles state observations and revealed outputs into one named execution view.}
\Field{Limits of the analogy.}{The WADT unit also explains code notation and invokes interpreter soundness to transport privacy. Those functions have no full baseline counterpart.}
\Field{Baseline narration used.}{partial.}
\Field{Complexity and jargon.}{Complexity: \Tag{high} $\rightarrow$ \Tag{medium}. Jargon: \Tag{high} $\rightarrow$ \Tag{medium}.}
\textbf{Original WADT unit. Footnotes omitted.}
\begin{lstlisting}
In the listing, \coqin{\#} marks a dealt hand, \coqin{\$} marks the
public shuffle index, and \coqin{\&} marks a card position. The player
receives the hand and the index, selects one entry, and reveals it to the
verifier. The verifier receives this entry as its observation. The interpreter
records both received values in the player's trace.
The trace theorem therefore concerns the execution of this process. Its
content is interpreter soundness: under the trace-lifting hypothesis the
recorded trace and the static view determine each other, so the operational
execution reveals nothing beyond the static view, and view privacy
transports to the trace rather than being proved a second time.
\end{lstlisting}
\colorbox{RewriteGreen}{\parbox{0.97\textwidth}{\textbf{Rewrite suggestion.}\par\begingroup\small
In the listing, \coqin{\#} marks a dealt hand, \coqin{\$} marks the public shuffle index, and \coqin{\&} marks a card position. The player receives the hand and shuffle index. The player selects one entry and reveals it to the verifier. The verifier receives this entry as its observation. The interpreter records the hand and shuffle index in the player's trace. The trace theorem concerns this process execution. Under the trace-lifting hypothesis, the recorded trace and the static view determine each other. This is the interpreter-soundness property used here. The process therefore reveals no information beyond the static view. Independence of the static view from the secret consequently gives independence of the executed trace from the secret, so the proof reuses view privacy.
\par\endgroup}}
\Field{Why this rewrite.}{Uses only the baseline move that assembles execution observations into a named trace, then states the notation, equivalence premise, and privacy transport in separate steps.}
\Field{Terminology actions.}{Keep: trace-lifting hypothesis, static view | Explain: interpreter soundness -> the recorded trace and static view determine each other, privacy transports -> static-view independence gives executed-trace independence}
\par\smallskip
\clearpage\subsection{Baseline S-02-P026: worked observer trace}
\colorbox{BaselineBlue}{\parbox{0.97\textwidth}{\Field{Purpose.}{Derive the view associated with the preceding worked state trace.}\Field{Primary purpose.}{example}\Field{Narration structure.}{view\_derivation $\rightarrow$ observer\_trace}\Field{Reusable move.}{Use a bounded worked instance to make an abstract definition concrete.}}}
\needspace{10\baselineskip}\subsubsection*{W-06-STM004: analogous purpose}
\Field{Location.}{Section 06, statement, lines 1424--1434.}
\Field{Draft purpose.}{Illustrate why three cards cannot distinguish the two secrets.}
\Field{Qualitative closeness.}{analogous purpose. Both units work through a concrete observer trace to make an abstract security claim visible.}
\Field{Limits of the analogy.}{The baseline derives one protocol view from prior states. The WADT example compares two secret encodings under related shuffles and uses three-transitivity to align the observed triples.}
\Field{Baseline narration used.}{adapted.}
\Field{Complexity and jargon.}{Complexity: \Tag{high} $\rightarrow$ \Tag{medium}. Jargon: \Tag{medium} $\rightarrow$ \Tag{medium}.}
\textbf{Original WADT unit. Footnotes omitted.}
\begin{lstlisting}
Let the dealt secret be $s=0$, so the dealt arrangement is $D_0$, and let
the shuffle be the translation $x\mapsto x+1$, which fixes the point
$\infty$ at position seven. Each position then shows the card that $D_0$
places one position later, and the coalition $\{0,1,2\}$ sees the card
values $(1,2,3)$. At $s=1$ the same shuffle would show $(1,2,4)$. By
Lemma~\ref{thm:three-transitive} some shuffle instead sends the
coalition's three positions to the positions $1$, $2$, and $4$ of $D_1$,
which hold the values $1$, $2$, and $3$. A coalition of three cards
therefore never separates the two secrets, and the proposition above makes
this uniformity exact.
\end{lstlisting}
\colorbox{RewriteGreen}{\parbox{0.97\textwidth}{\textbf{Rewrite suggestion.}\par\begingroup\small
Let the dealt secret be $s=0$, so the dealt arrangement is $D_0$. Let the shuffle be the translation $x\mapsto x+1$, which leaves the point $\infty$ fixed at position seven. Each position then shows the card that $D_0$ places one position later. The coalition $\{0,1,2\}$ sees the values $(1,2,3)$. For $s=1$, the same shuffle would show $(1,2,4)$. By Lemma~\ref{thm:three-transitive}, another shuffle sends the coalition's positions to positions $1$, $2$, and $4$ of $D_1$. These positions contain the values $1$, $2$, and $3$. Thus, three observed cards never distinguish the two secrets. The preceding proposition makes this equality of view distributions exact.
\par\endgroup}}
\Field{Why this rewrite.}{The rewrite adapts the baseline's worked-instance sequence. It separates each computation, explains the fixed point, and preserves both observed triples, the lemma reference, and the exact-uniformity conclusion.}
\Field{Terminology actions.}{Keep: three-transitivity, coalition view | Explain: translation fixes infinity}
\par\smallskip
\clearpage\subsection{Baseline S-02-P027: omitted-output convention}
\colorbox{BaselineBlue}{\parbox{0.97\textwidth}{\Field{Purpose.}{State when unrevealing outputs may be omitted from the displayed view.}\Field{Primary purpose.}{scope}\Field{Narration structure.}{notation\_condition $\rightarrow$ abbreviated\_form}\Field{Reusable move.}{State the convention or boundary exactly where later notation depends on it.}}}
\needspace{10\baselineskip}\subsubsection*{W-02-P013: partial analogue}
\Field{Location.}{Section 02, prose\_paragraph, lines 290--296.}
\Field{Draft purpose.}{State that only the product of a generator word is observed.}
\Field{Qualitative closeness.}{partial analogue. The baseline supplies the useful move of stating when unrevealed outputs are omitted from an observation.}
\Field{Limits of the analogy.}{The WADT unit imposes a structural non-observation assumption on an entire generator sequence and compares two models.}
\Field{Baseline narration used.}{partial.}
\Field{Complexity and jargon.}{Complexity: \Tag{high} $\rightarrow$ \Tag{low}. Jargon: \Tag{medium} $\rightarrow$ \Tag{low}.}
\textbf{Original WADT unit. Footnotes omitted.}
\begin{lstlisting}
The model observes the word shuffle only through its product. The dealer
performs the 200 physical moves and no party records an intermediate
arrangement. This assumption is structural in the formalization. The dealer process takes a list of alternative cuts and deals one hand
per cut. A word hence enters the execution only through its evaluated
group element. The uniform-shuffle model makes the same assumption, since its single
uniform shuffle is likewise unobserved.
\end{lstlisting}
\colorbox{RewriteGreen}{\parbox{0.97\textwidth}{\textbf{Rewrite suggestion.}\par\begingroup\small
The word-shuffle model exposes only the product of the word. The dealer performs 200 physical moves, but no party records an intermediate arrangement. This non-observation rule is a structural assumption of the formalization. The dealer process takes a list of alternative cuts and deals one hand for each cut. Thus a word affects execution only through its evaluated group element. The uniform-shuffle model uses the same observation rule because its single uniform shuffle is also unobserved.
\par\endgroup}}
\Field{Why this rewrite.}{Uses the baseline convention-setting move to name the unobserved intermediate arrangements, then follows the WADT execution contract.}
\Field{Terminology actions.}{Keep: evaluated group element, intermediate arrangement | Explain: list of alternative cuts -> a list from which the dealer process deals one hand for each cut, structural assumption -> the model's non-observation rule}
\par\smallskip
\needspace{10\baselineskip}\subsubsection*{W-06-P010: analogous purpose}
\Field{Location.}{Section 06, prose\_paragraph, lines 1459--1464.}
\Field{Draft purpose.}{Define the shuffle-free dealer and delimit the priors covered by its view and trace results.}
\Field{Qualitative closeness.}{analogous purpose. Both paragraphs state a scope convention that prevents the reader from attributing a broader result than the model supports.}
\Field{Limits of the analogy.}{The baseline permits omission of unrevealing outputs. The WADT paragraph distinguishes prior-generic view privacy from uniform-prior trace privacy and explicitly denies a broader trace claim.}
\Field{Baseline narration used.}{adapted.}
\Field{Complexity and jargon.}{Complexity: \Tag{medium} $\rightarrow$ \Tag{medium}. Jargon: \Tag{high} $\rightarrow$ \Tag{medium}.}
\textbf{Original WADT unit. Footnotes omitted.}
\begin{lstlisting}
A third dealer samples a uniform valid deck in the chosen class and applies
no later shuffle. I call the resulting distribution the shuffle-free deck
distribution, and view independence holds under it for
every Boolean secret prior. The available executed-trace theorem uses the
uniform Boolean prior. No prior-generic trace theorem is claimed.
\end{lstlisting}
\colorbox{RewriteGreen}{\parbox{0.97\textwidth}{\textbf{Rewrite suggestion.}\par\begingroup\small
A third dealer samples a valid deck uniformly from the chosen secret class and applies no later shuffle. I call this law the shuffle-free deck distribution. Under this law, view independence holds for every Boolean secret prior. The available executed-trace theorem assumes the uniform Boolean prior. I do not claim a prior-generic trace theorem.
\par\endgroup}}
\Field{Why this rewrite.}{The rewrite adapts the baseline's scope-convention move. It defines the shuffle-free law first, then separates the prior-generic view result from the uniform-prior trace theorem and explicit nonclaim.}
\Field{Terminology actions.}{Keep: any Boolean prior, uniform Boolean prior, no prior-generic trace theorem | Explain: shuffle-free deck distribution}
\par\smallskip
\needspace{10\baselineskip}\subsubsection*{W-07-P010: analogous purpose}
\Field{Location.}{Section 07, prose\_paragraph, lines 1652--1656.}
\Field{Draft purpose.}{Limit the word-shuffle privacy theorem to the fixed representative dealer.}
\Field{Qualitative closeness.}{analogous purpose. Both units state the precise condition under which a previously introduced representation or theorem applies.}
\Field{Limits of the analogy.}{The baseline limits when outputs may be omitted. This unit limits a privacy transfer by dealer sample space.}
\Field{Baseline narration used.}{adapted.}
\Field{Complexity and jargon.}{Complexity: \Tag{medium} $\rightarrow$ \Tag{low}. Jargon: \Tag{high} $\rightarrow$ \Tag{medium}.}
\textbf{Original WADT unit. Footnotes omitted.}
\begin{lstlisting}
The word-shuffle privacy results cover the fixed representative dealer only,
since the all-decks and shuffle-free dealers use different sample spaces
that the product-distribution transfer does not reach. The trust base for
these results is stated in Section~\ref{sec:instances}.
\end{lstlisting}
\colorbox{RewriteGreen}{\parbox{0.97\textwidth}{\textbf{Rewrite suggestion.}\par\begingroup\small
The word-shuffle privacy results cover only the fixed representative dealer. The all-decks and shuffle-free dealers use different sample spaces, so the product-distribution transfer does not apply to them. Section~\ref{sec:instances} states the trust base for these results.
\par\endgroup}}
\Field{Why this rewrite.}{Adapts the baseline's local boundary statement and gives the dealer restriction, its sample-space reason, and the trust-base reference.}
\Field{Terminology actions.}{Keep: fixed representative dealer, all-decks and shuffle-free dealers | Explain: different sample spaces | Replace: does not reach -> does not apply to}
\par\smallskip
\clearpage\subsection{Baseline S-02-P029: protocol-definition lead-in}
\colorbox{BaselineBlue}{\parbox{0.97\textwidth}{\Field{Purpose.}{Name the protocol object and hand its content to the formal definition.}\Field{Primary purpose.}{transition}\Field{Narration structure.}{concept\_naming $\rightarrow$ definition\_handoff}\Field{Reusable move.}{Use a short lead-in to name the role of what follows.}}}
\needspace{10\baselineskip}\subsubsection*{W-02-P003: same purpose}
\Field{Location.}{Section 02, prose\_paragraph, lines 199--200.}
\Field{Draft purpose.}{Introduce the formal data that describes the protocol.}
\Field{Qualitative closeness.}{same purpose. Both units are short formal lead-ins that name an object and hand its content to the next definition.}
\Field{Limits of the analogy.}{The baseline leads to a protocol tuple, while the WADT fragment leads to a broader framework data display.}
\Field{Baseline narration used.}{full.}
\Field{Complexity and jargon.}{Complexity: \Tag{low} $\rightarrow$ \Tag{low}. Jargon: \Tag{low} $\rightarrow$ \Tag{low}.}
\textbf{Original WADT unit. Footnotes omitted.}
\begin{lstlisting}
The framework describes a protocol by the data
\end{lstlisting}
\colorbox{RewriteGreen}{\parbox{0.97\textwidth}{\textbf{Rewrite suggestion.}\par\begingroup\small
The framework describes a protocol using the following data:
\par\endgroup}}
\Field{Why this rewrite.}{Reuses the baseline's complete concept-naming and display-handoff sequence in a shorter lead-in.}
\Field{Terminology actions.}{Keep: framework}
\par\smallskip
\clearpage\subsection{Baseline S-02-P030: protocol-component explanations}
\colorbox{BaselineBlue}{\parbox{0.97\textwidth}{\Field{Purpose.}{Assign roles to the input count, domain, deck, operation set, and action function.}\Field{Primary purpose.}{definition}\Field{Narration structure.}{input\_count\_role $\rightarrow$ input\_domain\_role $\rightarrow$ deck\_role $\rightarrow$ operation\_set\_role $\rightarrow$ action\_function\_role}\Field{Reusable move.}{Introduce one formal component and state its role in the surrounding construction.}}}
\needspace{10\baselineskip}\subsubsection*{W-02-P005: analogous purpose}
\Field{Location.}{Section 02, prose\_paragraph, lines 211--217.}
\Field{Draft purpose.}{Define the dealt arrangement, group action, player view, and coalition view.}
\Field{Qualitative closeness.}{analogous purpose. Both units assign operational roles to the components of a protocol model before execution is defined.}
\Field{Limits of the analogy.}{The baseline explains a protocol tuple. The WADT unit specializes those roles to dealing, shuffling, and coalition observations.}
\Field{Baseline narration used.}{adapted.}
\Field{Complexity and jargon.}{Complexity: \Tag{medium} $\rightarrow$ \Tag{low}. Jargon: \Tag{medium} $\rightarrow$ \Tag{low}.}
\textbf{Original WADT unit. Footnotes omitted.}
\begin{lstlisting}
A dealer samples a secret $S$ and a valid dealt arrangement $D$. The action
of a sampled shuffle $g$ produces the arrangement $\rho(g)D$. Position $i$
of the shuffled arrangement holds the card that $D$ places at position
$\rho(g)(i)$. Player $i$
observes a view $V_i$ derived from the card at position $i$. A coalition
$C\subseteq\{0,\ldots,n-1\}$ observes the joint view
\end{lstlisting}
\colorbox{RewriteGreen}{\parbox{0.97\textwidth}{\textbf{Rewrite suggestion.}\par\begingroup\small
A dealer samples a secret $S$ and a dealt arrangement $D$ that is valid for the instance. A sampled shuffle $g$ acts on the deck by its permutation action and produces $\rho(g)D$. Position $i$ of the shuffled arrangement holds the card that $D$ places at position $\rho(g)(i)$. Player $i$ observes the view $V_i$ derived from the card at position $i$. A coalition $C\subseteq\{0,\ldots,n-1\}$ observes the following joint view, which collects the observations at its named positions:
\par\endgroup}}
\Field{Why this rewrite.}{Adapts the baseline's role-by-role exposition to dealing, shuffling, individual observation, and coalition observation.}
\Field{Terminology actions.}{Keep: permutation action, coalition view | Explain: valid dealt arrangement -> a dealt arrangement that is valid for the instance, joint view -> the observations at the coalition's named positions}
\par\smallskip
\needspace{10\baselineskip}\subsubsection*{W-03-TC004: partial analogue}
\Field{Location.}{Section 03, table\_cell, lines 373--374.}
\Field{Draft purpose.}{Identify the model datum represented by the first row of the bridge table.}
\Field{Qualitative closeness.}{partial analogue. S-02-P030 supplies the limited move of naming a protocol component by its role. W-03-TC004 uses that move only to identify the run-layout datum.}
\Field{Limits of the analogy.}{The baseline explains five fields in connected prose. This cell is only a first-column label, and its meaning depends on the process-role cell that follows.}
\Field{Baseline narration used.}{partial.}
\Field{Complexity and jargon.}{Complexity: \Tag{low} $\rightarrow$ \Tag{low}. Jargon: \Tag{low} $\rightarrow$ \Tag{low}.}
\textbf{Original WADT unit. Footnotes omitted.}
\begin{lstlisting}
layout of a run
\end{lstlisting}
\colorbox{RewriteGreen}{\parbox{0.97\textwidth}{\textbf{Rewrite suggestion.}\par\begingroup\small
layout of one protocol run
\par\endgroup}}
\Field{Why this rewrite.}{The rewrite uses the partial analogue's component-naming move and makes the singular run scope explicit.}
\Field{Terminology actions.}{Keep: layout of a run}
\par\smallskip
\needspace{10\baselineskip}\subsubsection*{W-03-TC005: partial analogue}
\Field{Location.}{Section 03, table\_cell, lines 373--374.}
\Field{Draft purpose.}{State the framework role corresponding to the run-layout datum.}
\Field{Qualitative closeness.}{partial analogue. S-02-P030 supplies the limited component-to-role move. W-03-TC005 applies it to the three processes that realize one run layout.}
\Field{Limits of the analogy.}{The baseline develops a complete protocol tuple. This cell gives only the process-role side of one table row and relies on the preceding layout label and following instance value.}
\Field{Baseline narration used.}{partial.}
\Field{Complexity and jargon.}{Complexity: \Tag{low} $\rightarrow$ \Tag{low}. Jargon: \Tag{low} $\rightarrow$ \Tag{low}.}
\textbf{Original WADT unit. Footnotes omitted.}
\begin{lstlisting}
dealer, player, and verifier processes
\end{lstlisting}
\colorbox{RewriteGreen}{\parbox{0.97\textwidth}{\textbf{Rewrite suggestion.}\par\begingroup\small
dealer, player, and verifier processes
\par\endgroup}}
\Field{Why this rewrite.}{The rewrite uses the partial analogue's role-inventory move and retains the three operational roles as a concise table cell.}
\Field{Terminology actions.}{Keep: dealer, player, and verifier processes}
\par\smallskip
\needspace{10\baselineskip}\subsubsection*{W-03-TC006: partial analogue}
\Field{Location.}{Section 03, table\_cell, lines 373--375.}
\Field{Draft purpose.}{Give the main instance's concrete participant and card allocation for the run layout.}
\Field{Qualitative closeness.}{partial analogue. S-02-P030 supplies only the move of giving a formal component its concrete role. W-03-TC006 uses that move for the main instance's player count and card allocation.}
\Field{Limits of the analogy.}{The baseline explains generic tuple fields. This cell is an instance value, and its role is complete only when read with the preceding layout and process cells.}
\Field{Baseline narration used.}{partial.}
\Field{Complexity and jargon.}{Complexity: \Tag{low} $\rightarrow$ \Tag{low}. Jargon: \Tag{medium} $\rightarrow$ \Tag{low}.}
\textbf{Original WADT unit. Footnotes omitted.}
\begin{lstlisting}
eight
      players, one card each
\end{lstlisting}
\colorbox{RewriteGreen}{\parbox{0.97\textwidth}{\textbf{Rewrite suggestion.}\par\begingroup\small
eight players, each holding one card
\par\endgroup}}
\Field{Why this rewrite.}{The rewrite uses the partial analogue's exact instantiation move and states both the player count and per-player allocation directly.}
\Field{Terminology actions.}{Keep: eight players, one card each}
\par\smallskip
\needspace{10\baselineskip}\subsubsection*{W-03-TC008: partial analogue}
\Field{Location.}{Section 03, table\_cell, lines 375--377.}
\Field{Draft purpose.}{State the security-witness role associated with the group action and shuffle law.}
\Field{Qualitative closeness.}{partial analogue. S-02-P030 supplies only the move of mapping a component to its protocol role. W-03-TC008 uses that move to identify the endpoint-bound obligation of the security witness.}
\Field{Limits of the analogy.}{The baseline explains a full protocol tuple. This cell states one guarantee and relies on the preceding shuffle-law datum and following instance evidence.}
\Field{Baseline narration used.}{partial.}
\Field{Complexity and jargon.}{Complexity: \Tag{low} $\rightarrow$ \Tag{low}. Jargon: \Tag{medium} $\rightarrow$ \Tag{low}.}
\textbf{Original WADT unit. Footnotes omitted.}
\begin{lstlisting}
endpoint bound for the shuffle
      distribution
\end{lstlisting}
\colorbox{RewriteGreen}{\parbox{0.97\textwidth}{\textbf{Rewrite suggestion.}\par\begingroup\small
bound on each card position's endpoint distribution under the shuffle law
\par\endgroup}}
\Field{Why this rewrite.}{The rewrite uses the partial analogue's role-description move and explains that an endpoint is the resulting distribution at one card position.}
\Field{Terminology actions.}{Keep: endpoint bound, shuffle distribution | Explain: endpoint distribution}
\par\smallskip
\needspace{10\baselineskip}\subsubsection*{W-03-TC009: partial analogue}
\Field{Location.}{Section 03, table\_cell, lines 376--377.}
\Field{Draft purpose.}{Identify the main instance's algebraic and finite-word evidence for the endpoint bound.}
\Field{Qualitative closeness.}{partial analogue. S-02-P030 supplies only the move of assigning concrete content to a formal role. W-03-TC009 uses it to list the two distinct kinds of evidence supplied by the main instance.}
\Field{Limits of the analogy.}{The baseline gives generic component roles. This cell combines an action property and a distributional bound, and the surrounding row is needed to identify both as endpoint-security evidence.}
\Field{Baseline narration used.}{partial.}
\Field{Complexity and jargon.}{Complexity: \Tag{low} $\rightarrow$ \Tag{low}. Jargon: \Tag{medium} $\rightarrow$ \Tag{medium}.}
\textbf{Original WADT unit. Footnotes omitted.}
\begin{lstlisting}
three-transitive action, word bound
\end{lstlisting}
\colorbox{RewriteGreen}{\parbox{0.97\textwidth}{\textbf{Rewrite suggestion.}\par\begingroup\small
three-transitive action and a checked bound for the generator-word distribution
\par\endgroup}}
\Field{Why this rewrite.}{The rewrite uses the partial analogue's evidence-inventory move. It distinguishes the algebraic evidence from the finite-distribution bound and explains the word bound.}
\Field{Terminology actions.}{Keep: three-transitive action, word bound | Explain: word bound}
\par\smallskip
\needspace{10\baselineskip}\subsubsection*{W-03-TC011: partial analogue}
\Field{Location.}{Section 03, table\_cell, lines 377--378.}
\Field{Draft purpose.}{State how the reconstruction component must interact with the group action.}
\Field{Qualitative closeness.}{partial analogue. S-02-P030 supplies only the move of stating a formal component's role. W-03-TC011 applies it to action-invariant reconstruction.}
\Field{Limits of the analogy.}{The baseline explains a protocol field without a proof obligation. This cell names a compatibility condition and relies on the omitted decoder instance cell for its concrete realization.}
\Field{Baseline narration used.}{partial.}
\Field{Complexity and jargon.}{Complexity: \Tag{low} $\rightarrow$ \Tag{low}. Jargon: \Tag{medium} $\rightarrow$ \Tag{low}.}
\textbf{Original WADT unit. Footnotes omitted.}
\begin{lstlisting}
reconstruction against the action
\end{lstlisting}
\colorbox{RewriteGreen}{\parbox{0.97\textwidth}{\textbf{Rewrite suggestion.}\par\begingroup\small
reconstruction invariant under the group action
\par\endgroup}}
\Field{Why this rewrite.}{The rewrite uses the partial analogue's component-and-obligation move and states the decoder's action compatibility directly.}
\Field{Terminology actions.}{Keep: reconstruction, group action}
\par\smallskip
\needspace{10\baselineskip}\subsubsection*{W-03-DIA002: partial analogue}
\Field{Location.}{Section 03, diagram\_text, lines 426--428.}
\Field{Draft purpose.}{Label the supporting node that supplies exact or asymptotic shuffle-security evidence.}
\Field{Qualitative closeness.}{partial analogue. The baseline supplies only the move of naming a formal component by what it contributes. This node applies that move to the evidence feeding shuffle security.}
\Field{Limits of the analogy.}{The baseline uses full prose and defines a protocol field. This intentional fragment distinguishes two evidence modes, while its edge and color carry the relation to the shuffle-bound node.}
\Field{Baseline narration used.}{partial.}
\Field{Complexity and jargon.}{Complexity: \Tag{low} $\rightarrow$ \Tag{low}. Jargon: \Tag{medium} $\rightarrow$ \Tag{medium}.}
\textbf{Original WADT unit. Footnotes omitted.}
\begin{lstlisting}
exact and asymptotic\\[-2pt]
      security evidence
\end{lstlisting}
\colorbox{RewriteGreen}{\parbox{0.97\textwidth}{\textbf{Rewrite suggestion.}\par\begingroup\small
exact and asymptotic\\[-2pt] shuffle-bound evidence
\par\endgroup}}
\Field{Why this rewrite.}{The rewrite uses the partial analogue's classification move. It retains both evidence fields, which may occur together, and connects their guarantees to shuffle bounds in a concise diagram label.}
\Field{Terminology actions.}{Keep: exact security evidence, asymptotic security evidence | Explain: security evidence}
\par\smallskip
\needspace{10\baselineskip}\subsubsection*{W-03-DIA004: partial analogue}
\Field{Location.}{Section 03, diagram\_text, lines 431--433.}
\Field{Draft purpose.}{Label the component record that fixes the protocol's dealer, player, and verifier processes.}
\Field{Qualitative closeness.}{partial analogue. The baseline supplies only the component-role naming move. This node uses it for the executable processes bundled by the protocol layout.}
\Field{Limits of the analogy.}{The baseline explains a tuple field in sentences. This node is a fragment whose blue color and bundle arrow establish required-profile status.}
\Field{Baseline narration used.}{partial.}
\Field{Complexity and jargon.}{Complexity: \Tag{low} $\rightarrow$ \Tag{low}. Jargon: \Tag{medium} $\rightarrow$ \Tag{low}.}
\textbf{Original WADT unit. Footnotes omitted.}
\begin{lstlisting}
protocol layout\\[-2pt]
      \footnotesize dealer, players, verifier
\end{lstlisting}
\colorbox{RewriteGreen}{\parbox{0.97\textwidth}{\textbf{Rewrite suggestion.}\par\begingroup\small
protocol layout\\[-2pt] \footnotesize dealer, players, verifier
\par\endgroup}}
\Field{Why this rewrite.}{The rewrite uses the partial analogue's component-before-roles move and retains the diagram's concise label structure.}
\Field{Terminology actions.}{Keep: protocol layout, dealer, players, verifier}
\par\smallskip
\needspace{10\baselineskip}\subsubsection*{W-03-DIA005: partial analogue}
\Field{Location.}{Section 03, diagram\_text, lines 433--435.}
\Field{Draft purpose.}{Label the component record that bounds one-position endpoint distributions after shuffling.}
\Field{Qualitative closeness.}{partial analogue. The baseline supplies only the move of naming a formal component and its role. This node applies that move to the endpoint guarantee.}
\Field{Limits of the analogy.}{The baseline does not encode dependencies visually. Here the incoming evidence edge and outgoing profile edge carry claims not stated by the fragment alone.}
\Field{Baseline narration used.}{partial.}
\Field{Complexity and jargon.}{Complexity: \Tag{low} $\rightarrow$ \Tag{low}. Jargon: \Tag{medium} $\rightarrow$ \Tag{low}.}
\textbf{Original WADT unit. Footnotes omitted.}
\begin{lstlisting}
shuffle-security bound\\[-2pt]
      \footnotesize endpoint bound
\end{lstlisting}
\colorbox{RewriteGreen}{\parbox{0.97\textwidth}{\textbf{Rewrite suggestion.}\par\begingroup\small
shuffle-security bound\\[-2pt] \footnotesize bound on each endpoint distribution
\par\endgroup}}
\Field{Why this rewrite.}{The rewrite uses the partial analogue's component-and-guarantee move. It retains the shuffle-security bound and explains each endpoint as a post-shuffle distribution.}
\Field{Terminology actions.}{Keep: shuffle-security bound, endpoint bound | Explain: endpoint}
\par\smallskip
\needspace{10\baselineskip}\subsubsection*{W-03-DIA006: partial analogue}
\Field{Location.}{Section 03, diagram\_text, lines 435--437.}
\Field{Draft purpose.}{Label the reconstruction component that connects the group action to the decoder.}
\Field{Qualitative closeness.}{partial analogue. The baseline supplies only the move of assigning a role to one formal component. This node applies it to action-compatible decoding.}
\Field{Limits of the analogy.}{The baseline states a field role without a proof relation. This fragment compresses both reconstruction and invariance, with the arrow supplying the aggregate dependency.}
\Field{Baseline narration used.}{partial.}
\Field{Complexity and jargon.}{Complexity: \Tag{low} $\rightarrow$ \Tag{low}. Jargon: \Tag{medium} $\rightarrow$ \Tag{low}.}
\textbf{Original WADT unit. Footnotes omitted.}
\begin{lstlisting}
reconstruction\\[-2pt]
      \footnotesize group action to decoder
\end{lstlisting}
\colorbox{RewriteGreen}{\parbox{0.97\textwidth}{\textbf{Rewrite suggestion.}\par\begingroup\small
reconstruction\\[-2pt] \footnotesize decoder invariant under the group action
\par\endgroup}}
\Field{Why this rewrite.}{The rewrite uses the partial analogue's relation-naming move and states the decoder's action compatibility directly.}
\Field{Terminology actions.}{Keep: reconstruction, group action, decoder | Explain: group action to decoder}
\par\smallskip
\needspace{10\baselineskip}\subsubsection*{W-03-P005: analogous purpose}
\Field{Location.}{Section 03, prose\_paragraph, lines 481--486.}
\Field{Draft purpose.}{Explain how profile components determine the executable processes and decoder.}
\Field{Qualitative closeness.}{analogous purpose. Both paragraphs assign operational roles to components of a formal protocol object.}
\Field{Limits of the analogy.}{The baseline explains a protocol tuple field by field. The WADT paragraph traces how layout, action, distribution, and reconstruction determine an executable instance.}
\Field{Baseline narration used.}{adapted.}
\Field{Complexity and jargon.}{Complexity: \Tag{medium} $\rightarrow$ \Tag{low}. Jargon: \Tag{medium} $\rightarrow$ \Tag{low}.}
\textbf{Original WADT unit. Footnotes omitted.}
\begin{lstlisting}
The layout supplies the executable processes described in
Section~\ref{sec:model}. The group action and shuffle distribution fix the
dealing rule and endpoint bound. The reconstruction component chooses the
decoder, so profiles over the same group can use different reconstruction
schemes.
\end{lstlisting}
\colorbox{RewriteGreen}{\parbox{0.97\textwidth}{\textbf{Rewrite suggestion.}\par\begingroup\small
The run layout supplies the executable dealer, player, and verifier processes described in Section~\ref{sec:model}. The group action and shuffle distribution determine the dealing rule and endpoint bound. The reconstruction component selects the decoder. Profiles over the same group can therefore use different reconstruction schemes.
\par\endgroup}}
\Field{Why this rewrite.}{The rewrite adapts the baseline's field-by-field mapping. It explains layout through its processes and gives each component one direct consequence.}
\Field{Terminology actions.}{Keep: endpoint bound, reconstruction component | Explain: layout}
\par\smallskip
\clearpage\subsection{Baseline S-02-P031: protocol-execution setup}
\colorbox{BaselineBlue}{\parbox{0.97\textwidth}{\Field{Purpose.}{Fix a protocol and initial state, then introduce its execution as an ordered procedure.}\Field{Primary purpose.}{method}\Field{Narration structure.}{protocol\_setup $\rightarrow$ initial\_state\_setup $\rightarrow$ procedure\_handoff}\Field{Reusable move.}{Present a procedure as ordered state changes with an explicit continuation or stopping rule.}}}
\needspace{10\baselineskip}\subsubsection*{W-02-P001: analogous purpose}
\Field{Location.}{Section 02, prose\_paragraph, lines 184--191.}
\Field{Draft purpose.}{Explain one complete protocol run before the formal model.}
\Field{Qualitative closeness.}{analogous purpose. Both units explain an execution as an ordered procedure before later details use it.}
\Field{Limits of the analogy.}{The baseline fixes an initial state and formal execution variables. The WADT unit gives an actor-level overview and ends at an instance diagram.}
\Field{Baseline narration used.}{adapted.}
\Field{Complexity and jargon.}{Complexity: \Tag{medium} $\rightarrow$ \Tag{low}. Jargon: \Tag{low} $\rightarrow$ \Tag{low}.}
\textbf{Original WADT unit. Footnotes omitted.}
\begin{lstlisting}
A run of the protocol handles one secret and a deck of face-down cards. A
dealer encodes the secret as a valid arrangement of the deck. The dealer
then samples a shuffle from a finite group and rearranges the deck by its
permutation action. Each player receives the card at one position. The
players reveal their cards to a verifier, and the verifier decodes the
secret from the revealed arrangement. Each instance section shows its own
run as a flow diagram.
\end{lstlisting}
\colorbox{RewriteGreen}{\parbox{0.97\textwidth}{\textbf{Rewrite suggestion.}\par\begingroup\small
A run uses one secret and a deck of face-down cards. The dealer encodes the secret as a deck arrangement that is valid for the instance. The dealer then samples a finite-group shuffle and applies its permutation action to the deck. Each player receives the card at one position. The players reveal their cards to the verifier. The verifier decodes the secret from the revealed arrangement. Each instance section depicts this sequence in a flow diagram.
\par\endgroup}}
\Field{Why this rewrite.}{Adapts the baseline's ordered execution sequence while retaining the WADT actors, deck states, and diagram handoff.}
\Field{Terminology actions.}{Keep: permutation action, verifier | Explain: valid arrangement -> a deck arrangement that is valid for the instance}
\par\smallskip
\needspace{10\baselineskip}\subsubsection*{W-06-P002: analogous purpose}
\Field{Location.}{Section 06, prose\_paragraph, lines 1286--1292.}
\Field{Draft purpose.}{Explain the executed protocol and derive its decoder correctness.}
\Field{Qualitative closeness.}{analogous purpose. Both paragraphs fix the protocol execution and trace the ordered state changes needed for correctness.}
\Field{Limits of the analogy.}{The baseline introduces general execution. The WADT paragraph follows one fixed interpreter run and ends with an orbit-invariance correctness argument.}
\Field{Baseline narration used.}{adapted.}
\Field{Complexity and jargon.}{Complexity: \Tag{medium} $\rightarrow$ \Tag{medium}. Jargon: \Tag{medium} $\rightarrow$ \Tag{low}.}
\textbf{Original WADT unit. Footnotes omitted.}
\begin{lstlisting}
The interpreter deals $D_s$, applies an allowed shuffle, and sends one card
to each of eight players. The players return their cards to the verifier. The
verifier's endpoint sequence is the shuffled dealt arrangement. Orbit-class
invariance and the encoder equation then give correctness.
\end{lstlisting}
\colorbox{RewriteGreen}{\parbox{0.97\textwidth}{\textbf{Rewrite suggestion.}\par\begingroup\small
The interpreter deals $D_s$, applies an allowed shuffle, and sends one card to each of eight players. The players return their cards to the verifier. The verifier's endpoint sequence is the shuffled dealt arrangement. Orbit-class invariance and the encoder equation, which identifies the decoded orbit class with the encoded secret, then prove correctness.
\par\endgroup}}
\Field{Why this rewrite.}{The rewrite adapts the baseline's ordered execution procedure. It excludes the footnote, traces the cards through every role, and explains the encoder equation before the conclusion.}
\Field{Terminology actions.}{Keep: endpoint sequence, orbit-class invariance | Explain: encoder equation}
\par\smallskip
\clearpage\subsection{Baseline S-02-P035: worked execution lead-in}
\colorbox{BaselineBlue}{\parbox{0.97\textwidth}{\Field{Purpose.}{Choose an encoded input sequence and announce the step-by-step execution.}\Field{Primary purpose.}{example}\Field{Narration structure.}{execution\_goal $\rightarrow$ initial\_sequence $\rightarrow$ step\_list\_handoff}\Field{Reusable move.}{Use a bounded worked instance to make an abstract definition concrete.}}}
\needspace{10\baselineskip}\subsubsection*{W-04-P007: partial analogue}
\Field{Location.}{Section 04, prose\_paragraph, lines 729--732.}
\Field{Draft purpose.}{Tell the reader what the leakage figure is selected to show.}
\Field{Qualitative closeness.}{partial analogue. The baseline contributes only the move of selecting a bounded example to make an abstract result visible.}
\Field{Limits of the analogy.}{The baseline announces a stepwise protocol execution from one encoded input. This paragraph previews static leakage cases chosen to show shape dependence and the entropy maximum.}
\Field{Baseline narration used.}{partial.}
\Field{Complexity and jargon.}{Complexity: \Tag{low} $\rightarrow$ \Tag{low}. Jargon: \Tag{medium} $\rightarrow$ \Tag{low}.}
\textbf{Original WADT unit. Footnotes omitted.}
\begin{lstlisting}
Figure~\ref{fig:fivecard-leakage} selects cases that expose two features of
the leakage theorem. Two-card leakage depends on the reveal shape, while
four and five revealed cards reach the secret's entropy.
\end{lstlisting}
\colorbox{RewriteGreen}{\parbox{0.97\textwidth}{\textbf{Rewrite suggestion.}\par\begingroup\small
Figure~\ref{fig:fivecard-leakage} shows two features of the leakage theorem. For two revealed cards, leakage depends on their relative positions. Four and five revealed cards reach the secret's entropy.
\par\endgroup}}
\Field{Why this rewrite.}{Uses only the baseline's bounded-example handoff and states the two visual comparisons directly.}
\Field{Terminology actions.}{Keep: secret entropy | Explain: reveal shape | Replace: expose two features -> show two features}
\par\smallskip
\clearpage\subsection{Baseline S-02-P040: protocol-purpose interpretation}
\colorbox{BaselineBlue}{\parbox{0.97\textwidth}{\Field{Purpose.}{Identify the logical property encoded by the final state and use it to classify the protocol.}\Field{Primary purpose.}{interpretation}\Field{Narration structure.}{output\_property $\rightarrow$ protocol\_classification}\Field{Reusable move.}{Close a worked protocol by linking its final encoding to the function it realises.}}}
\needspace{10\baselineskip}\subsubsection*{W-04-DIA008: partial analogue}
\Field{Location.}{Section 04, diagram\_text, lines 663--664.}
\Field{Draft purpose.}{State the predicate that the verifier reads from the revealed deck.}
\Field{Qualitative closeness.}{partial analogue. The baseline supplies the move of naming the logical property read from a final card state.}
\Field{Limits of the analogy.}{The baseline interprets a worked protocol outcome in prose. This unit is an imperative diagram fragment and does not classify the full protocol.}
\Field{Baseline narration used.}{partial.}
\Field{Complexity and jargon.}{Complexity: \Tag{low} $\rightarrow$ \Tag{low}. Jargon: \Tag{low} $\rightarrow$ \Tag{low}.}
\textbf{Original WADT unit. Footnotes omitted.}
\begin{lstlisting}
read: are the three\\[-1pt] hearts consecutive
\end{lstlisting}
\colorbox{RewriteGreen}{\parbox{0.97\textwidth}{\textbf{Rewrite suggestion.}\par\begingroup\small
read whether the three\\[-1pt]hearts are consecutive
\par\endgroup}}
\Field{Why this rewrite.}{Uses only the baseline's output-property move and converts the prompt-like label into a direct readout action.}
\Field{Terminology actions.}{Keep: consecutive hearts}
\par\smallskip
\needspace{10\baselineskip}\subsubsection*{W-05-P010: analogous purpose}
\Field{Location.}{Section 05, prose\_paragraph, lines 976--982.}
\Field{Draft purpose.}{Interpret the two orbit classes as Boolean secret values and preview their guarantees.}
\Field{Qualitative closeness.}{analogous purpose. Both units interpret a formal outcome as the intended logical value and use that interpretation to classify the protocol's role.}
\Field{Limits of the analogy.}{The baseline closes one worked protocol. The WADT unit opens an encoding construction and previews two later proof obligations.}
\Field{Baseline narration used.}{adapted.}
\Field{Complexity and jargon.}{Complexity: \Tag{medium} $\rightarrow$ \Tag{low}. Jargon: \Tag{medium} $\rightarrow$ \Tag{low}.}
\textbf{Original WADT unit. Footnotes omitted.}
\begin{lstlisting}
The two orbit classes provide two stable labels and can therefore represent
the values of a Boolean secret. The dealer fixes one valid deck $D_s$ for
each label $s\in\{0,1\}$ and later distributes the shuffled cards, one to
each player. Orbit invariance will establish correctness. Privacy of the
players' partial observations is proved later from the shuffle distribution
and three-transitivity.
\end{lstlisting}
\colorbox{RewriteGreen}{\parbox{0.97\textwidth}{\textbf{Rewrite suggestion.}\par\begingroup\small
I use the two stable orbit labels to represent a Boolean secret. For each label $s\in\{0,1\}$, the dealer fixes one valid representative deck $D_s$. The dealer later shuffles this deck and gives one card to each player. Orbit invariance will prove correctness. Privacy will follow later from the shuffle distribution and three-transitivity, which makes each ordered observation of at most three positions uniform. The following definition gives the two representative decks and the encoder.
\par\endgroup}}
\Field{Why this rewrite.}{The rewrite preserves the Boolean interpretation, dealer action, and separate correctness and privacy premises. Its final sentence explains why the two representative decks and the encoder are defined next.}
\Field{Terminology actions.}{Keep: representative deck, orbit invariance | Explain: three-transitivity -> makes each ordered observation of at most three positions uniform}
\par\smallskip
\needspace{10\baselineskip}\subsubsection*{W-06-P003: analogous purpose}
\Field{Location.}{Section 06, prose\_paragraph, lines 1303--1309.}
\Field{Draft purpose.}{Explain the distributional consequence of universal correctness and distinguish it from approximate privacy.}
\Field{Qualitative closeness.}{analogous purpose. Both paragraphs interpret a proved protocol fact and use it to classify the resulting guarantee.}
\Field{Limits of the analogy.}{The baseline classifies a protocol by its logical output. The WADT paragraph interprets universal correctness as probability-one correctness under any shuffle law and separates privacy error.}
\Field{Baseline narration used.}{adapted.}
\Field{Complexity and jargon.}{Complexity: \Tag{medium} $\rightarrow$ \Tag{low}. Jargon: \Tag{medium} $\rightarrow$ \Tag{low}.}
\textbf{Original WADT unit. Footnotes omitted.}
\begin{lstlisting}
Correctness quantifies over every group element. It therefore holds with
probability one under the uniform distribution and under any word distribution, and only
privacy carries an approximation error. A corollary instantiates the
executed run at the evaluated 200-letter word.
\end{lstlisting}
\colorbox{RewriteGreen}{\parbox{0.97\textwidth}{\textbf{Rewrite suggestion.}\par\begingroup\small
Because correctness holds for every group element, it holds with probability one under both the uniform distribution and any word distribution. In particular, the executed run is correct for the evaluated 200-letter word. By contrast, only privacy incurs an approximation error. Having settled full-run correctness for both shuffle laws, I next ask how many revealed positions suffice to determine the secret.
\par\endgroup}}
\Field{Why this rewrite.}{The rewrite preserves universal correctness, both probability-one consequences, the evaluated 200-letter corollary, and the privacy-only approximation boundary. It then states the new recovery question, which connects full eight-endpoint decoding to the finer structural threshold in the next paragraph.}
\Field{Terminology actions.}{Keep: probability one, word distribution | Explain: only privacy carries an approximation error}
\par\smallskip
\needspace{10\baselineskip}\subsubsection*{W-07-P004: partial analogue}
\Field{Location.}{Section 07, prose\_paragraph, lines 1557--1563.}
\Field{Draft purpose.}{Interpret the displayed inequality as one mixing certificate and list its inputs.}
\Field{Qualitative closeness.}{partial analogue. The baseline contributes only the move of interpreting a formal result by naming what it establishes.}
\Field{Limits of the analogy.}{The baseline reads the final protocol state and classifies the computed function. This paragraph names a mixing certificate, inventories its proof inputs, and limits it to group-element distributions.}
\Field{Baseline narration used.}{partial.}
\Field{Complexity and jargon.}{Complexity: \Tag{medium} $\rightarrow$ \Tag{low}. Jargon: \Tag{high} $\rightarrow$ \Tag{medium}.}
\textbf{Original WADT unit. Footnotes omitted.}
\begin{lstlisting}
Equation~\ref{eq:pgl-mixing} is one checked certificate, which I call
the mixing certificate. Its inputs are the
five-letter generation theorem, the group order in
Equation~\ref{eq:pgl-order}, the $336$ fiber counts, and the final integer
inequality. The certificate concerns the distribution of shuffle group
elements.
\end{lstlisting}
\colorbox{RewriteGreen}{\parbox{0.97\textwidth}{\textbf{Rewrite suggestion.}\par\begingroup\small
Equation~\ref{eq:pgl-mixing} is a checked mixing certificate. Its inputs are the five-letter generation theorem, the group order in Equation~\ref{eq:pgl-order}, the $336$ fiber counts, and the final integer inequality. The certificate concerns only the distribution of shuffle group elements.
\par\endgroup}}
\Field{Why this rewrite.}{Uses only the baseline's interpretation move. It names the certificate, lists its four inputs, and preserves its sample-space limit.}
\Field{Terminology actions.}{Keep: mixing certificate, fiber counts, integer inequality | Explain: five-letter generation theorem}
\par\smallskip
\clearpage\subsection{Baseline S-02-P042: symbolic functionality representation}
\colorbox{BaselineBlue}{\parbox{0.97\textwidth}{\Field{Purpose.}{Give a compact symbolic representation of the running input-output contract.}\Field{Primary purpose.}{definition}\Field{Narration structure.}{object\_identification $\rightarrow$ display\_handoff}\Field{Reusable move.}{Use a short prose cue to state the role of the following formal object.}}}
\needspace{10\baselineskip}\subsubsection*{W-02-P014: partial analogue}
\Field{Location.}{Section 02, prose\_paragraph, lines 297--300.}
\Field{Draft purpose.}{Introduce the figure that separates the exact and approximate proof paths.}
\Field{Qualitative closeness.}{partial analogue. The baseline supplies the short lead-in move that states the role of a following formal object.}
\Field{Limits of the analogy.}{The baseline only introduces a symbolic functionality display. The WADT unit introduces a proof-path figure, distinguishes two distributional routes, and states the lower route's privacy endpoint.}
\Field{Baseline narration used.}{partial.}
\Field{Complexity and jargon.}{Complexity: \Tag{low} $\rightarrow$ \Tag{low}. Jargon: \Tag{medium} $\rightarrow$ \Tag{low}.}
\textbf{Original WADT unit. Footnotes omitted.}
\begin{lstlisting}
Figure~\ref{fig:models} separates the proved paths for the two distributions. The
lower path reaches approximate coalition-view and executed-trace privacy at
$2^{-39}$.
\end{lstlisting}
\colorbox{RewriteGreen}{\parbox{0.97\textwidth}{\textbf{Rewrite suggestion.}\par\begingroup\small
Figure~\ref{fig:models} separates the proved paths for the two distributions. The lower path starts from the finite word distribution and its mixing certificate. It yields approximate coalition-view and executed-trace privacy at $2^{-39}$.
\par\endgroup}}
\Field{Why this rewrite.}{Uses the baseline's visual handoff move and adds only the WADT figure's lower-path origin and endpoint.}
\Field{Terminology actions.}{Keep: coalition-view privacy, executed-trace privacy | Explain: lower path -> the path from the finite word distribution and mixing certificate}
\par\smallskip
\needspace{10\baselineskip}\subsubsection*{W-05-P011: analogous purpose}
\Field{Location.}{Section 05, prose\_paragraph, lines 996--998.}
\Field{Draft purpose.}{Direct the reader from the encoder definition to its visual representation.}
\Field{Qualitative closeness.}{analogous purpose. Both units use a short prose cue to state the role of the formal object that follows.}
\Field{Limits of the analogy.}{The baseline introduces a symbolic functionality display. The WADT sentence introduces a figure that compares the heart-position patterns of two chosen decks.}
\Field{Baseline narration used.}{adapted.}
\Field{Complexity and jargon.}{Complexity: \Tag{low} $\rightarrow$ \Tag{low}. Jargon: \Tag{low} $\rightarrow$ \Tag{low}.}
\textbf{Original WADT unit. Footnotes omitted.}
\begin{lstlisting}
Figure~\ref{fig:encoding} shows how the two chosen decks place the four
hearts in different position patterns.
\end{lstlisting}
\colorbox{RewriteGreen}{\parbox{0.97\textwidth}{\textbf{Rewrite suggestion.}\par\begingroup\small
Figure~\ref{fig:encoding} shows the different heart-position patterns in the two representative decks.
\par\endgroup}}
\Field{Why this rewrite.}{Adapts the baseline's short visual handoff to name the exact comparison shown by the encoding figure.}
\Field{Terminology actions.}{Keep: position patterns}
\par\smallskip
\needspace{10\baselineskip}\subsubsection*{W-05-P018: analogous purpose}
\Field{Location.}{Section 05, prose\_paragraph, lines 1119--1120.}
\Field{Draft purpose.}{Introduce the execution figure.}
\Field{Qualitative closeness.}{analogous purpose. Both units use a short prose cue to state the role of the formal object that follows.}
\Field{Limits of the analogy.}{The baseline introduces a symbolic functionality display. The WADT sentence introduces an execution figure that locates encoding and decoding in one run.}
\Field{Baseline narration used.}{adapted.}
\Field{Complexity and jargon.}{Complexity: \Tag{low} $\rightarrow$ \Tag{low}. Jargon: \Tag{low} $\rightarrow$ \Tag{low}.}
\textbf{Original WADT unit. Footnotes omitted.}
\begin{lstlisting}
Figure~\ref{fig:run} shows how the encoder and decoder enter one execution.
\end{lstlisting}
\colorbox{RewriteGreen}{\parbox{0.97\textwidth}{\textbf{Rewrite suggestion.}\par\begingroup\small
Figure~\ref{fig:run} shows where the encoder and decoder act during one execution.
\par\endgroup}}
\Field{Why this rewrite.}{Adapts the baseline's short visual handoff to the two operations located in the run diagram.}
\Field{Terminology actions.}{Keep: execution}
\par\smallskip
\clearpage\subsection{Baseline S-02-P044: dummy-sequence abstraction}
\colorbox{BaselineBlue}{\parbox{0.97\textwidth}{\Field{Purpose.}{Define partial sequence descriptions and collect the descriptions compatible with a full sequence.}\Field{Primary purpose.}{definition}\Field{Narration structure.}{object\_setup $\rightarrow$ formal\_definition $\rightarrow$ scope\_boundary}\Field{Reusable move.}{Introduce the object, state its rule, and close with the boundary needed for later use.}}}
\needspace{10\baselineskip}\subsubsection*{W-05-STM001: analogous purpose}
\Field{Location.}{Section 05, statement, lines 886--898.}
\Field{Draft purpose.}{Define four-position patterns, group orbits, and orbit equivalence.}
\Field{Qualitative closeness.}{analogous purpose. Both units define a collection of structured objects and a relation that says when another object is compatible with one of them.}
\Field{Limits of the analogy.}{The baseline concerns partial and full sequences. The WADT statement concerns subsets and group-action equivalence.}
\Field{Baseline narration used.}{adapted.}
\Field{Complexity and jargon.}{Complexity: \Tag{medium} $\rightarrow$ \Tag{low}. Jargon: \Tag{medium} $\rightarrow$ \Tag{low}.}
\textbf{Original WADT unit. Footnotes omitted.}
\begin{lstlisting}
Let
\[
  \Omega_4=\{A\subseteq X\mid |A|=4\}.
\]
For $A,B\in\Omega_4$, define
\[
  \operatorname{Orb}_G(A)=\{g\mathbin{\cdot}A\mid g\in G\},
  \qquad
  A\sim_G B
  \quad\Longleftrightarrow\quad
  B\in\operatorname{Orb}_G(A).
\]
\end{lstlisting}
\colorbox{RewriteGreen}{\parbox{0.97\textwidth}{\textbf{Rewrite suggestion.}\par\begingroup\small
Let $\Omega_4$ be the set of all four-position subsets of $X$:
\[
  \Omega_4=\{A\subseteq X\mid |A|=4\}.
\]
For $A,B\in\Omega_4$, define the orbit and orbit equivalence by
\[
  \operatorname{Orb}_G(A)=\{g\mathbin{\cdot}A\mid g\in G\},
  \qquad
  A\sim_G B
  \quad\Longleftrightarrow\quad
  B\in\operatorname{Orb}_G(A).
\]
\par\endgroup}}
\Field{Why this rewrite.}{Adapts the baseline's object-then-rule sequence and keeps the formal statement terse.}
\Field{Terminology actions.}{Keep: orbit, orbit equivalence | Explain: Omega four -> the set of all four-position subsets of X}
\par\smallskip
\clearpage\subsection{Baseline S-02-P046: matching-relation example}
\colorbox{BaselineBlue}{\parbox{0.97\textwidth}{\Field{Purpose.}{Show how one concrete full sequence satisfies a partial description.}\Field{Primary purpose.}{example}\Field{Narration structure.}{instance\_setup $\rightarrow$ worked\_application $\rightarrow$ example\_outcome}\Field{Reusable move.}{Use one bounded instance to make an abstract rule inspectable.}}}
\needspace{10\baselineskip}\subsubsection*{W-05-STM005: same purpose}
\Field{Location.}{Section 05, statement, lines 1028--1051.}
\Field{Draft purpose.}{Verify decoding on both chosen representatives.}
\Field{Qualitative closeness.}{same purpose. Both units use a concrete finite calculation to show that one full object satisfies a previously defined relation.}
\Field{Limits of the analogy.}{The baseline checks sequence matching. The WADT statement checks both encoder representatives and composes two maps.}
\Field{Baseline narration used.}{full.}
\Field{Complexity and jargon.}{Complexity: \Tag{medium} $\rightarrow$ \Tag{low}. Jargon: \Tag{medium} $\rightarrow$ \Tag{low}.}
\textbf{Original WADT unit. Footnotes omitted.}
\begin{lstlisting}
The two rows in Figure~\ref{fig:encoding} satisfy
\[
  H(D_0)=\{0,1,2,3\},
  \qquad
  \kappa(H(D_0))=0,
\]
and
\[
  H(D_1)=\{0,1,2,4\},
  \qquad
  \kappa(H(D_1))=1.
\]
Thus the full encoding path is
\[
  s\longmapsto D_s
   \longmapsto H(D_s)
   \xmapsto{\kappa}s,
  \qquad
  \operatorname{dec}(\operatorname{enc}(s))=s.
\]
\end{lstlisting}
\colorbox{RewriteGreen}{\parbox{0.97\textwidth}{\textbf{Rewrite suggestion.}\par\begingroup\small
The rows in Figure~\ref{fig:encoding} satisfy
\[
  H(D_0)=\{0,1,2,3\},
  \qquad
  \kappa(H(D_0))=0,
\]
and
\[
  H(D_1)=\{0,1,2,4\},
  \qquad
  \kappa(H(D_1))=1.
\]
Thus the complete encoding path applies the encoder, heart-set extraction, and classifier in order:
\[
  s\longmapsto D_s
   \longmapsto H(D_s)
   \xmapsto{\kappa}s,
  \qquad
  \operatorname{dec}(\operatorname{enc}(s))=s.
\]
\par\endgroup}}
\Field{Why this rewrite.}{Reuses the baseline's full worked-instance sequence for the two representatives and their composed round-trip equality. The excluded footnote does not enter the suggestion.}
\Field{Terminology actions.}{Keep: heart set, encoding path | Explain: map composition -> applies the encoder, heart-set extraction, and classifier in order}
\par\smallskip
\clearpage\subsection{Baseline S-02-P048: piecewise input-family instance}
\colorbox{BaselineBlue}{\parbox{0.97\textwidth}{\Field{Purpose.}{Specify the running input family by cases.}\Field{Primary purpose.}{example}\Field{Narration structure.}{instance\_setup $\rightarrow$ worked\_application $\rightarrow$ example\_outcome}\Field{Reusable move.}{Use one bounded instance to make an abstract rule inspectable.}}}
\needspace{10\baselineskip}\subsubsection*{W-05-STM004: analogous purpose}
\Field{Location.}{Section 05, statement, lines 984--994.}
\Field{Draft purpose.}{Define the encoder by choosing one concrete deck for each bit.}
\Field{Qualitative closeness.}{analogous purpose. Both units instantiate an input-dependent formal family by explicit cases.}
\Field{Limits of the analogy.}{The baseline specifies a variable sequence. The WADT statement selects two deck permutations as orbit representatives.}
\Field{Baseline narration used.}{adapted.}
\Field{Complexity and jargon.}{Complexity: \Tag{low} $\rightarrow$ \Tag{low}. Jargon: \Tag{low} $\rightarrow$ \Tag{low}.}
\textbf{Original WADT unit. Footnotes omitted.}
\begin{lstlisting}
The encoder chooses the valid decks
\[
  D_0=(0,1,2,3,4,5,6,7),
  \qquad
  D_1=(0,1,2,4,3,5,6,7),
\]
and sets $\operatorname{enc}(s)=D_s$.
\end{lstlisting}
\colorbox{RewriteGreen}{\parbox{0.97\textwidth}{\textbf{Rewrite suggestion.}\par\begingroup\small
The encoder chooses the two valid representative decks
\[
  D_0=(0,1,2,3,4,5,6,7),
  \qquad
  D_1=(0,1,2,4,3,5,6,7),
\]
and sets $\operatorname{enc}(s)=D_s$.
\par\endgroup}}
\Field{Why this rewrite.}{Adapts the baseline's bounded-instance sequence to the two explicit representatives and the encoder definition. The excluded footnote does not enter the suggestion.}
\Field{Terminology actions.}{Keep: encoder}
\par\smallskip
\clearpage\subsection{Baseline S-02-P051: functionality component types}
\colorbox{BaselineBlue}{\parbox{0.97\textwidth}{\Field{Purpose.}{Clarify the domains and codomains of the input and output components.}\Field{Primary purpose.}{definition}\Field{Narration structure.}{object\_setup $\rightarrow$ formal\_definition $\rightarrow$ scope\_boundary}\Field{Reusable move.}{Introduce the object, state its rule, and close with the boundary needed for later use.}}}
\needspace{10\baselineskip}\subsubsection*{W-03-TC022: partial analogue}
\Field{Location.}{Section 03, table\_cell, lines 506--508.}
\Field{Draft purpose.}{Describe the exact-slot mechanism used by the den Boer and main-group rows.}
\Field{Qualitative closeness.}{partial analogue. S-02-P051 supplies only the move of stating a formal rule with its scope. This cell applies it to exact equality under a zero-bias uniform-law condition.}
\Field{Limits of the analogy.}{The baseline explains input and output component domains. This cell reports evidence for two instance rows and relies on the slot and realized-by columns.}
\Field{Baseline narration used.}{partial.}
\Field{Complexity and jargon.}{Complexity: \Tag{low} $\rightarrow$ \Tag{low}. Jargon: \Tag{medium} $\rightarrow$ \Tag{low}.}
\textbf{Original WADT unit. Footnotes omitted.}
\begin{lstlisting}
exact equality at $\varepsilon=0$ under the uniform
      group distribution
\end{lstlisting}
\colorbox{RewriteGreen}{\parbox{0.97\textwidth}{\textbf{Rewrite suggestion.}\par\begingroup\small
exact equality of the endpoint and target distributions at $\varepsilon=0$ under the uniform group distribution
\par\endgroup}}
\Field{Why this rewrite.}{The rewrite uses the partial analogue's rule-and-scope move. It explains what is equal and preserves the zero-bias and uniform-law conditions.}
\Field{Terminology actions.}{Keep: exact equality, uniform group distribution, bias epsilon | Explain: exact equality}
\par\smallskip
\needspace{10\baselineskip}\subsubsection*{W-03-TC026: partial analogue}
\Field{Location.}{Section 03, table\_cell, lines 508--510.}
\Field{Draft purpose.}{Describe Kim's mechanism for filling both exact and asymptotic witness slots.}
\Field{Qualitative closeness.}{partial analogue. S-02-P051 supplies only the move of stating a component rule and its scope. This cell applies it to count-based decay as word length changes.}
\Field{Limits of the analogy.}{The baseline defines interface component types. This cell summarizes a quantitative mechanism and depends on the row's two present-slot markers and Kim label.}
\Field{Baseline narration used.}{partial.}
\Field{Complexity and jargon.}{Complexity: \Tag{low} $\rightarrow$ \Tag{low}. Jargon: \Tag{medium} $\rightarrow$ \Tag{low}.}
\textbf{Original WADT unit. Footnotes omitted.}
\begin{lstlisting}
exact count with geometric decay in the word
      length
\end{lstlisting}
\colorbox{RewriteGreen}{\parbox{0.97\textwidth}{\textbf{Rewrite suggestion.}\par\begingroup\small
exact finite count giving an error bound that decays geometrically with word length
\par\endgroup}}
\Field{Why this rewrite.}{The rewrite uses the partial analogue's certificate-then-rate move and explains geometric decay as a decreasing error bound.}
\Field{Terminology actions.}{Keep: exact count, geometric decay, word length | Explain: geometric decay}
\par\smallskip
\needspace{10\baselineskip}\subsubsection*{W-03-TC030: partial analogue}
\Field{Location.}{Section 03, table\_cell, lines 510--511.}
\Field{Draft purpose.}{Describe the asymptotic-only mechanism used by the two symmetric-group instances.}
\Field{Qualitative closeness.}{partial analogue. S-02-P051 supplies only the move of stating a rule with its boundary. This cell applies it to a spectral certificate whose validity depends on an imported premise.}
\Field{Limits of the analogy.}{The baseline sets formal component domains. This cell reports an assumption-bearing proof mechanism and relies on the adjacent carrier cell to identify its instances.}
\Field{Baseline narration used.}{partial.}
\Field{Complexity and jargon.}{Complexity: \Tag{low} $\rightarrow$ \Tag{low}. Jargon: \Tag{medium} $\rightarrow$ \Tag{medium}.}
\textbf{Original WADT unit. Footnotes omitted.}
\begin{lstlisting}
spectral certificate with an imported gap premise
\end{lstlisting}
\colorbox{RewriteGreen}{\parbox{0.97\textwidth}{\textbf{Rewrite suggestion.}\par\begingroup\small
spectral certificate based on an imported spectral-gap premise
\par\endgroup}}
\Field{Why this rewrite.}{The rewrite uses the partial analogue's method-and-assumption move and states that the certificate depends on an imported spectral-gap premise.}
\Field{Terminology actions.}{Keep: spectral certificate, gap premise | Explain: spectral gap}
\par\smallskip
\clearpage\subsection{Baseline S-02-P052: correctness definition lead-in}
\colorbox{BaselineBlue}{\parbox{0.97\textwidth}{\Field{Purpose.}{Introduce the semantic criterion used to judge protocol outputs.}\Field{Primary purpose.}{definition}\Field{Narration structure.}{motivation $\rightarrow$ concept\_scope $\rightarrow$ formal\_handoff}\Field{Reusable move.}{Explain why a concept is needed before introducing its formal specification.}}}
\needspace{10\baselineskip}\subsubsection*{W-03-LI001: partial analogue}
\Field{Location.}{Section 03, list\_item, lines 469--472.}
\Field{Draft purpose.}{State the profile's single-position shuffle-security obligation.}
\Field{Qualitative closeness.}{partial analogue. S-02-P052 supplies only the move of stating a formal criterion before later results use it. This item applies that move to a per-position distance bound.}
\Field{Limits of the analogy.}{The baseline introduces whole-protocol correctness. This item states a distributional profile obligation and includes two code locators.}
\Field{Baseline narration used.}{partial.}
\Field{Complexity and jargon.}{Complexity: \Tag{medium} $\rightarrow$ \Tag{medium}. Jargon: \Tag{high} $\rightarrow$ \Tag{medium}.}
\textbf{Original WADT unit. Footnotes omitted.}
\begin{lstlisting}
Every single card position lands close to uniform: \coqin{sw\_bound}
  bounds the distance of each position's endpoint distribution from the
  uniform distribution by \coqin{sw\_bound\_eps}.
\end{lstlisting}
\colorbox{RewriteGreen}{\parbox{0.97\textwidth}{\textbf{Rewrite suggestion.}\par\begingroup\small
For every card position, the endpoint-distance field bounds the distance between that position's post-shuffle card distribution and the uniform distribution by \coqin{sw\_bound\_eps}. The field is \coqin{sw\_bound}.
\par\endgroup}}
\Field{Why this rewrite.}{The rewrite uses only the partial analogue's obligation-first move. It explains endpoint distribution and maps the mathematical bound to both formal field names.}
\Field{Terminology actions.}{Keep: endpoint distribution, uniform distribution, sw\_bound | Explain: endpoint | Replace: sw\_bound bounds -> the endpoint-distance field bounds}
\par\smallskip
\needspace{10\baselineskip}\subsubsection*{W-03-LI002: partial analogue}
\Field{Location.}{Section 03, list\_item, lines 472--476.}
\Field{Draft purpose.}{State the profile's paired threshold reconstruction and hiding obligations.}
\Field{Qualitative closeness.}{partial analogue. S-02-P052 supplies only the criterion-stating move. This item applies it to two coordinated threshold conditions rather than one protocol-correctness definition.}
\Field{Limits of the analogy.}{The baseline introduces correctness alone. This item pairs a universal decoder property with a below-threshold ambiguity property and gives separate code locators.}
\Field{Baseline narration used.}{partial.}
\Field{Complexity and jargon.}{Complexity: \Tag{medium} $\rightarrow$ \Tag{medium}. Jargon: \Tag{high} $\rightarrow$ \Tag{medium}.}
\textbf{Original WADT unit. Footnotes omitted.}
\begin{lstlisting}
The threshold scheme recovers and hides: \coqin{ts\_correct} decodes
  every valid share tuple to its secret, and \coqin{ts\_private} gives
  every coalition below the threshold a share pattern that is equally
  consistent with either secret.
\end{lstlisting}
\colorbox{RewriteGreen}{\parbox{0.97\textwidth}{\textbf{Rewrite suggestion.}\par\begingroup\small
The threshold scheme provides reconstruction and threshold privacy. The field \coqin{ts\_correct} decodes every valid share tuple to its secret. The field \coqin{ts\_private} ensures that every coalition below the threshold sees a share pattern compatible with each secret in the same way.
\par\endgroup}}
\Field{Why this rewrite.}{The rewrite uses the partial analogue's recovery-then-hiding sequence. It replaces the compressed heading, explains equal consistency, and assigns one sentence to each formal field.}
\Field{Terminology actions.}{Keep: threshold scheme, ts\_correct, ts\_private | Explain: equally consistent | Replace: recovers and hides -> provides reconstruction and threshold privacy}
\par\smallskip
\needspace{10\baselineskip}\subsubsection*{W-03-LI003: partial analogue}
\Field{Location.}{Section 03, list\_item, lines 476--479.}
\Field{Draft purpose.}{State the profile's shuffle-invariant reconstruction obligation.}
\Field{Qualitative closeness.}{partial analogue. S-02-P052 supplies only the move of stating a formal criterion for later use. This item applies it to compatibility between reconstruction and the group action.}
\Field{Limits of the analogy.}{The baseline prepares a correctness definition. This item states a universal invariance obligation and includes one code-facing theorem field.}
\Field{Baseline narration used.}{partial.}
\Field{Complexity and jargon.}{Complexity: \Tag{medium} $\rightarrow$ \Tag{medium}. Jargon: \Tag{high} $\rightarrow$ \Tag{medium}.}
\textbf{Original WADT unit. Footnotes omitted.}
\begin{lstlisting}
Reconstruction is shuffle-invariant: for any allowed shuffle and any
  valid share tuple, \coqin{rp\_recon\_invariant} states that permuting the
  shares by the shuffle leaves the recovered secret unchanged.
\end{lstlisting}
\colorbox{RewriteGreen}{\parbox{0.97\textwidth}{\textbf{Rewrite suggestion.}\par\begingroup\small
Reconstruction is invariant under every allowed shuffle. For any valid share tuple, permuting the shares by the shuffle leaves the recovered secret unchanged. The field \coqin{rp\_recon\_invariant} states this property.
\par\endgroup}}
\Field{Why this rewrite.}{The rewrite uses the partial analogue's property-before-identifier move. It preserves both universal quantifiers and explains the share permutation before naming the field.}
\Field{Terminology actions.}{Keep: shuffle invariance, rp\_recon\_invariant | Explain: permuting the shares}
\par\smallskip
\needspace{10\baselineskip}\subsubsection*{W-05-P009: analogous purpose}
\Field{Location.}{Section 05, prose\_paragraph, lines 946--950.}
\Field{Draft purpose.}{State what the classifier theorem must prove.}
\Field{Qualitative closeness.}{analogous purpose. Both units state the semantic criterion that a construction must satisfy immediately before its formal verification.}
\Field{Limits of the analogy.}{The baseline introduces protocol correctness. The WADT unit introduces orbit classification and a quantitative count.}
\Field{Baseline narration used.}{adapted.}
\Field{Complexity and jargon.}{Complexity: \Tag{medium} $\rightarrow$ \Tag{low}. Jargon: \Tag{medium} $\rightarrow$ \Tag{low}.}
\textbf{Original WADT unit. Footnotes omitted.}
\begin{lstlisting}
For $\kappa$ to serve as an orbit label, its fibers must coincide with the
orbits of $G$. Thus $\kappa$ must be constant on each orbit and must assign
different values to distinct orbits. The next theorem verifies both
requirements and counts the two fibers.
\end{lstlisting}
\colorbox{RewriteGreen}{\parbox{0.97\textwidth}{\textbf{Rewrite suggestion.}\par\begingroup\small
For $\kappa$ to be an orbit label, each fiber of $\kappa$ must equal one orbit of $G$. A fiber is the set of patterns assigned one classifier value. Thus $\kappa$ must be constant on each orbit and must give different values to distinct orbits. The next theorem proves both requirements and counts the two fibers.
\par\endgroup}}
\Field{Why this rewrite.}{Adapts the baseline's requirement-to-formal-statement sequence and defines the counted object before the theorem handoff.}
\Field{Terminology actions.}{Keep: fiber, orbit label | Explain: classifier fiber -> the set of patterns assigned one classifier value}
\par\smallskip
\needspace{10\baselineskip}\subsubsection*{W-05-P012: partial analogue}
\Field{Location.}{Section 05, prose\_paragraph, lines 1022--1026.}
\Field{Draft purpose.}{Set the two correctness checks for the encoder and decoder.}
\Field{Qualitative closeness.}{partial analogue. The baseline supplies the move of stating a correctness criterion before a formal check.}
\Field{Limits of the analogy.}{The WADT unit divides correctness into representative recovery and shuffle invariance, while the baseline gives one general output criterion.}
\Field{Baseline narration used.}{partial.}
\Field{Complexity and jargon.}{Complexity: \Tag{medium} $\rightarrow$ \Tag{low}. Jargon: \Tag{low} $\rightarrow$ \Tag{low}.}
\textbf{Original WADT unit. Footnotes omitted.}
\begin{lstlisting}
The encoder must place each bit in the class that the decoder assigns to
that bit. The following calculation checks this round trip on both
representatives. Shuffle invariance will then extend it to every deck
reachable from either representative.
\end{lstlisting}
\colorbox{RewriteGreen}{\parbox{0.97\textwidth}{\textbf{Rewrite suggestion.}\par\begingroup\small
The decoder must assign each representative deck to the bit that its encoder received. The following calculation checks this encoder-decoder round trip for both representatives. Shuffle invariance will then extend the equality to every deck reachable from either representative.
\par\endgroup}}
\Field{Why this rewrite.}{Uses the baseline move that states the exact correctness obligation before the calculation and retains the separate invariance step.}
\Field{Terminology actions.}{Keep: round trip, shuffle invariance}
\par\smallskip
\needspace{10\baselineskip}\subsubsection*{W-07-P008: partial analogue}
\Field{Location.}{Section 07, prose\_paragraph, lines 1610--1614.}
\Field{Draft purpose.}{Summarize the proved distribution chain and hand it to the protocol theorem.}
\Field{Qualitative closeness.}{partial analogue. The baseline contributes only the explanatory handoff from established concepts to the formal statement that uses them.}
\Field{Limits of the analogy.}{The baseline motivates a correctness criterion before defining it. This paragraph summarizes three distributional levels and invokes interpreter transport before a theorem.}
\Field{Baseline narration used.}{partial.}
\Field{Complexity and jargon.}{Complexity: \Tag{medium} $\rightarrow$ \Tag{low}. Jargon: \Tag{high} $\rightarrow$ \Tag{medium}.}
\textbf{Original WADT unit. Footnotes omitted.}
\begin{lstlisting}
The proved chain reaches group-distribution proximity, every single-card
marginal, and the product with any secret prior. Transport through the
protocol interpreter turns the certificate into the second headline
theorem.
\end{lstlisting}
\colorbox{RewriteGreen}{\parbox{0.97\textwidth}{\textbf{Rewrite suggestion.}\par\begingroup\small
The proof chain establishes proximity of the group distributions, every single-card marginal, and the product with any secret prior. Transport through the protocol interpreter then turns the certificate into Theorem B.
\par\endgroup}}
\Field{Why this rewrite.}{Uses only the baseline's explanatory handoff and preserves all three distributional levels and the interpreter transport.}
\Field{Terminology actions.}{Keep: single-card marginal, protocol interpreter | Explain: transport | Replace: headline theorem -> Theorem B}
\par\smallskip
\clearpage\subsection{Baseline S-02-P054: view-distribution notation}
\colorbox{BaselineBlue}{\parbox{0.97\textwidth}{\Field{Purpose.}{Name the distribution induced by a protocol execution and identify its source of randomness.}\Field{Primary purpose.}{definition}\Field{Narration structure.}{object\_setup $\rightarrow$ formal\_definition $\rightarrow$ scope\_boundary}\Field{Reusable move.}{Introduce the object, state its rule, and close with the boundary needed for later use.}}}
\needspace{10\baselineskip}\subsubsection*{W-02-P002: partial analogue}
\Field{Location.}{Section 02, prose\_paragraph, lines 192--198.}
\Field{Draft purpose.}{Explain why shuffle randomness controls privacy and prepare a distributional account.}
\Field{Qualitative closeness.}{partial analogue. The baseline supplies the useful move of naming the execution distribution and its random source before security is stated.}
\Field{Limits of the analogy.}{The WADT unit also gives privacy intuition and compares individual with coalition views, so the baseline does not guide the whole paragraph.}
\Field{Baseline narration used.}{partial.}
\Field{Complexity and jargon.}{Complexity: \Tag{medium} $\rightarrow$ \Tag{low}. Jargon: \Tag{medium} $\rightarrow$ \Tag{low}.}
\textbf{Original WADT unit. Footnotes omitted.}
\begin{lstlisting}
Security rests on the shuffle alone, because the dealer's encoding of
the secret is fixed and the shuffle carries all the randomness. A player
sees one card, and a coalition sees the cards at its positions. The shuffle spreads every small
coalition view over many arrangements, and the privacy theorems make this
spreading exact. The rest of this section states the same flow as
probability distributions.
\end{lstlisting}
\colorbox{RewriteGreen}{\parbox{0.97\textwidth}{\textbf{Rewrite suggestion.}\par\begingroup\small
The fixed secret encoding leaves the shuffle as the only source of randomness. One player observes one card. A coalition view contains the cards at the coalition's positions. The shuffle distributes each small coalition view across many deck arrangements. The privacy theorems state exactly that the resulting coalition-view distribution is independent of the secret. The rest of this section expresses the same flow as probability distributions.
\par\endgroup}}
\Field{Why this rewrite.}{Uses only the baseline move that names the distribution and its random source, then follows the WADT contract for the privacy explanation and formal handoff.}
\Field{Terminology actions.}{Keep: coalition view | Explain: spreads every small coalition view -> distributes each small coalition view across many deck arrangements | Replace: this spreading exact -> the resulting coalition-view distribution is independent of the secret}
\par\smallskip
\needspace{10\baselineskip}\subsubsection*{W-06-P008: partial analogue}
\Field{Location.}{Section 06, prose\_paragraph, lines 1401--1406.}
\Field{Draft purpose.}{Specify the random experiment for the main privacy result and state the independence consequence.}
\Field{Qualitative closeness.}{partial analogue. The baseline supplies the random-experiment setup move by naming the execution law before security is stated.}
\Field{Limits of the analogy.}{The baseline does not provide a fixed representative dealer, a coalition bound, three-transitivity, or an independence conclusion.}
\Field{Baseline narration used.}{partial.}
\Field{Complexity and jargon.}{Complexity: \Tag{medium} $\rightarrow$ \Tag{medium}. Jargon: \Tag{high} $\rightarrow$ \Tag{medium}.}
\textbf{Original WADT unit. Footnotes omitted.}
\begin{lstlisting}
The main uniform-shuffle execution samples $S$ uniformly from $\{0,1\}$. It deals the
fixed representative $D_S$ and samples an independent shuffle $g\leftarrow
U_G$. Three-transitivity gives independence for every coalition $C$ with
$\lvert C\rvert\leq3$.
\end{lstlisting}
\colorbox{RewriteGreen}{\parbox{0.97\textwidth}{\textbf{Rewrite suggestion.}\par\begingroup\small
The main uniform-shuffle execution samples $S$ uniformly from $\{0,1\}$. It deals the fixed representative $D_S$, meaning the chosen encoded deck for secret $S$. It then samples an independent shuffle $g\leftarrow
U_G$. Three-transitivity means that any ordered triple of distinct positions can be moved to any other. This property makes the view independent of the secret for every coalition $C$ with $\lvert C\rvert\leq3$.
\par\endgroup}}
\Field{Why this rewrite.}{The rewrite uses the partial analogue's random-choices-before-security move. It excludes the footnote, explains the fixed representative and three-transitivity, and preserves the independence and coalition conditions.}
\Field{Terminology actions.}{Keep: fixed representative, three-transitivity, independence | Explain: fixed representative dealer, three-transitivity}
\par\smallskip
\needspace{10\baselineskip}\subsubsection*{W-06-P009: analogous purpose}
\Field{Location.}{Section 06, prose\_paragraph, lines 1436--1442.}
\Field{Draft purpose.}{Define the all-decks dealer distribution and explain what robustness question it tests.}
\Field{Qualitative closeness.}{analogous purpose. Both paragraphs define the random source behind a protocol-view distribution before the next security result uses it.}
\Field{Limits of the analogy.}{The baseline names general view-distribution notation and its randomness. The WADT paragraph motivates a robustness check and specifies a concrete secret, class-conditional deck, and independent group-shuffle experiment.}
\Field{Baseline narration used.}{adapted.}
\Field{Complexity and jargon.}{Complexity: \Tag{low} $\rightarrow$ \Tag{low}. Jargon: \Tag{medium} $\rightarrow$ \Tag{low}.}
\textbf{Original WADT unit. Footnotes omitted.}
\begin{lstlisting}
Two further dealers act as robustness checks on the fixed representative.
The all-decks dealer tests whether the fixed representative hides a special
choice. It samples $S$ uniformly. Conditional on $S=s$, it samples a uniform
valid arrangement of class $s$. It then applies an independent uniform
$\PG$ shuffle.
\end{lstlisting}
\colorbox{RewriteGreen}{\parbox{0.97\textwidth}{\textbf{Rewrite suggestion.}\par\begingroup\small
Two additional dealers test the robustness of the fixed representative choice. The all-decks dealer checks whether that choice is special. It samples $S$ uniformly. Given $S=s$, it samples a valid arrangement uniformly from class $s$. It then applies an independent uniform $\PG$ shuffle.
\par\endgroup}}
\Field{Why this rewrite.}{The rewrite adapts the baseline's distribution-definition sequence. It excludes the footnote, explains the robustness question, and states each random choice in order.}
\Field{Terminology actions.}{Keep: all-decks dealer, class-conditional | Explain: hides a special choice}
\par\smallskip
\clearpage\subsection{Baseline S-02-P055: security definition lead-in}
\colorbox{BaselineBlue}{\parbox{0.97\textwidth}{\Field{Purpose.}{Introduce the indistinguishability condition used for protocol security.}\Field{Primary purpose.}{definition}\Field{Narration structure.}{motivation $\rightarrow$ concept\_scope $\rightarrow$ formal\_handoff}\Field{Reusable move.}{Explain why a concept is needed before introducing its formal specification.}}}
\needspace{10\baselineskip}\subsubsection*{W-02-P010: analogous purpose}
\Field{Location.}{Section 02, prose\_paragraph, lines 262--266.}
\Field{Draft purpose.}{State the passive adversary model and the complete view-privacy condition, then introduce trace privacy.}
\Field{Qualitative closeness.}{analogous purpose. Both units introduce an independence criterion immediately before its formal statement.}
\Field{Limits of the analogy.}{The baseline introduces one general view criterion. The WADT unit states coalition-view independence in prose and only announces the separate trace condition, whose formula is outside the registered span.}
\Field{Baseline narration used.}{adapted.}
\Field{Complexity and jargon.}{Complexity: \Tag{medium} $\rightarrow$ \Tag{low}. Jargon: \Tag{medium} $\rightarrow$ \Tag{low}.}
\textbf{Original WADT unit. Footnotes omitted.}
\begin{lstlisting}
The security theorems concern passive, honest-but-curious adversaries. A
corrupted player follows the protocol and retains its view and executed
trace. Perfect coalition privacy states that $S$ and $V_C$ are independent.
Perfect trace privacy states
\end{lstlisting}
\colorbox{RewriteGreen}{\parbox{0.97\textwidth}{\textbf{Rewrite suggestion.}\par\begingroup\small
The security theorems assume passive, honest-but-curious adversaries. A corrupted player observes the protocol but does not deviate from it. The player retains its view and executed trace. Perfect coalition privacy means that $S$ and $V_C$ are independent. Perfect trace privacy is stated by the following equality:
\par\endgroup}}
\Field{Why this rewrite.}{Adapts the baseline's explanation-to-formula sequence and defines the adversary behavior before the two privacy conditions.}
\Field{Terminology actions.}{Keep: passive, honest-but-curious adversaries, perfect coalition privacy, trace privacy | Explain: corrupted player -> observes the protocol but does not deviate, trace privacy -> the equality in the following display}
\par\smallskip
\needspace{10\baselineskip}\subsubsection*{W-03-P009: analogous purpose}
\Field{Location.}{Section 03, prose\_paragraph, lines 544--547.}
\Field{Draft purpose.}{Motivate the generic privacy theorem and identify how an instance discharges its main hypothesis.}
\Field{Qualitative closeness.}{analogous purpose. Both paragraphs introduce a privacy result by naming the condition that will support it.}
\Field{Limits of the analogy.}{The baseline leads to a security definition. The WADT paragraph leads to a theorem and includes a concrete three-transitive instance.}
\Field{Baseline narration used.}{adapted.}
\Field{Complexity and jargon.}{Complexity: \Tag{low} $\rightarrow$ \Tag{low}. Jargon: \Tag{medium} $\rightarrow$ \Tag{low}.}
\textbf{Original WADT unit. Footnotes omitted.}
\begin{lstlisting}
The first generic theorem yields coalition privacy from transitivity. An
instance discharges its hypothesis by exhibiting a $t$-transitive action on
decks of distinct cards, as Section~\ref{sec:pgl} does for $t=3$.
\end{lstlisting}
\colorbox{RewriteGreen}{\parbox{0.97\textwidth}{\textbf{Rewrite suggestion.}\par\begingroup\small
The first generic theorem derives coalition privacy from transitivity. An instance must provide a $t$-transitive action on decks whose cards are all distinct. Section~\ref{sec:pgl} proves this requirement for $t=3$.
\par\endgroup}}
\Field{Why this rewrite.}{The rewrite adapts the baseline's motivation-to-formal-handoff sequence. It explains distinct cards and separates the generic mechanism from the main instance's proof.}
\Field{Terminology actions.}{Keep: t-transitive action, coalition privacy | Explain: distinct cards}
\par\smallskip
\needspace{10\baselineskip}\subsubsection*{W-04-P009: analogous purpose}
\Field{Location.}{Section 04, prose\_paragraph, lines 794--797.}
\Field{Draft purpose.}{Move from the fixed-reveal example to the shuffle-security proposition.}
\Field{Qualitative closeness.}{analogous purpose. Both units identify the security condition needed next and hand the reader to its formal statement.}
\Field{Limits of the analogy.}{The baseline introduces input-indistinguishability. This unit introduces a quantitative shuffle witness for one concrete dealer.}
\Field{Baseline narration used.}{adapted.}
\Field{Complexity and jargon.}{Complexity: \Tag{low} $\rightarrow$ \Tag{low}. Jargon: \Tag{medium} $\rightarrow$ \Tag{low}.}
\textbf{Original WADT unit. Footnotes omitted.}
\begin{lstlisting}
The example describes what a fixed reveal set exposes. The profile also
needs a shuffle-security witness for Kim's repeated biased cut. The next
proposition supplies that witness at the concrete parameters used here.
\end{lstlisting}
\colorbox{RewriteGreen}{\parbox{0.97\textwidth}{\textbf{Rewrite suggestion.}\par\begingroup\small
The example describes the information exposed by a fixed reveal set. The profile also needs a shuffle-security witness for Kim's repeated biased cut. The next proposition gives this witness at the parameters used here.
\par\endgroup}}
\Field{Why this rewrite.}{Adapts the baseline's motivation and formal handoff from the completed example to the remaining security obligation.}
\Field{Terminology actions.}{Keep: fixed reveal set, biased cut | Explain: shuffle-security witness}
\par\smallskip
\clearpage\subsection{Baseline S-02-P056: view-distribution calculation}
\colorbox{BaselineBlue}{\parbox{0.97\textwidth}{\Field{Purpose.}{Compute the possible views and their probabilities for an arbitrary input.}\Field{Primary purpose.}{proof}\Field{Narration structure.}{proof\_setup $\rightarrow$ calculation $\rightarrow$ intermediate\_result}\Field{Reusable move.}{Reduce a universal claim to an arbitrary case and expose the decisive calculation.}}}
\needspace{10\baselineskip}\subsubsection*{W-05-P024: analogous purpose}
\Field{Location.}{Section 05, prose\_paragraph, lines 1264--1273.}
\Field{Draft purpose.}{Explain the finite proof of three-transitivity and connect it to privacy.}
\Field{Qualitative closeness.}{analogous purpose. Both units expose a finite distributional calculation so the reader can see the evidence used by the next formal security conclusion.}
\Field{Limits of the analogy.}{The baseline enumerates a small set of views and probabilities. The WADT unit enumerates transport witnesses and passes the result through a generic privacy theorem.}
\Field{Baseline narration used.}{adapted.}
\Field{Complexity and jargon.}{Complexity: \Tag{high} $\rightarrow$ \Tag{medium}. Jargon: \Tag{high} $\rightarrow$ \Tag{medium}.}
\textbf{Original WADT unit. Footnotes omitted.}
\begin{lstlisting}
The formal proof enumerates generator words that send a fixed base triple to
each distinct target triple. The kernel checks the resulting finite table. For example, the base
triple is $(0,1,2)$, and the table stores for every other ordered triple
of distinct points one generator word that carries $(0,1,2)$ to it.
Lemma~\ref{thm:three-transitive} makes every ordered view of at most three
positions uniform under $U_G$. It supplies the group-action premise of
Theorem~\ref{thm:generic-privacy} at $t=3$. That theorem then gives
uniform-shuffle coalition privacy. The nonuniform word distribution is treated
separately in Section~\ref{sec:mixing}.
\end{lstlisting}
\colorbox{RewriteGreen}{\parbox{0.97\textwidth}{\textbf{Rewrite suggestion.}\par\begingroup\small
The formal proof enumerates generator words from one fixed base triple to every distinct target triple. The Rocq kernel checks the resulting finite table. The base triple is $(0,1,2)$. For every other ordered triple of distinct points, the table stores one generator word that sends $(0,1,2)$ to that triple. Lemma~\ref{thm:three-transitive} then makes every ordered view of at most three positions uniform under $U_G$. This lemma supplies the group-action premise of Theorem~\ref{thm:generic-privacy} at $t=3$. The generic theorem gives uniform-shuffle coalition privacy. Section~\ref{sec:mixing} treats the nonuniform word distribution separately because it requires a mixing argument.
\par\endgroup}}
\Field{Why this rewrite.}{Adapts the baseline's finite-case proof sequence to the checked witness table, view uniformity, generic privacy theorem, and separate word-law argument.}
\Field{Terminology actions.}{Keep: kernel checks, three-transitivity | Explain: ordered view uniform -> every ordered view of at most three positions is uniform, nonuniform word distribution -> requires a separate mixing argument}
\par\smallskip
\needspace{10\baselineskip}\subsubsection*{W-06-P007: analogous purpose}
\Field{Location.}{Section 06, prose\_paragraph, lines 1372--1398.}
\Field{Draft purpose.}{Classify the post-threshold reveal patterns and derive their exact leakage values.}
\Field{Qualitative closeness.}{analogous purpose. Both paragraphs turn a finite observation distribution into an inspectable probability or count calculation that supports a security conclusion.}
\Field{Limits of the analogy.}{The baseline enumerates a small protocol's views. The WADT paragraph uses group orbits, collision counts, a pen-and-paper formula, exact rational leakages, and two explicit formalization limits.}
\Field{Baseline narration used.}{adapted.}
\Field{Complexity and jargon.}{Complexity: \Tag{high} $\rightarrow$ \Tag{high}. Jargon: \Tag{high} $\rightarrow$ \Tag{high}.}
\textbf{Original WADT unit. Footnotes omitted.}
\begin{lstlisting}
The witness coalition
makes the privacy cutoff three sharp. The leakage for reveal sets of sizes
four through six reduces to four numbers, one per orbit of reveal sets
under $G$. The four-subsets fall into two orbits, of sizes $42$ and $28$,
separated by the classical cross-ratio invariant, and the five-subsets and
the six-subsets each form a single orbit. For a representative $C$ with at least three positions, the two
deals restrict to $C$ injectively, so, given either secret value, the
coalition's observation is uniform over the $336$ equally likely
possibilities, and the mutual information between the secret class
and the revealed cards equals $1-m/336$ bits, where $m$ counts the
observations shared by the two deals. This identity is a pen-and-paper step
justified by sharp three-transitivity, and it fails below three positions,
where the true leakage is zero. The kernel-checked collision counts of the
four representatives are $m=96$, $72$, $36$, and $12$, so the leakage
values are $5/7$, $11/14$, $25/28$, and $27/28$ bits, evaluated from the
counts outside the kernel by the committed script
\path{pgg-smc/instances/pgl27/pgl_leakage_targets.py}. The four-heart
witness coalition is the representative of the $42$-element orbit, and its
leakage is exactly $5/7$ bits. Extending each number from its
representative to the whole orbit uses invariance of leakage under
relabeling by a group element, a pen-and-paper step that the development
does not formalize.
\end{lstlisting}
\colorbox{RewriteGreen}{\parbox{0.97\textwidth}{\textbf{Rewrite suggestion.}\par\begingroup\small
The witness coalition shows that the privacy cutoff of three is sharp. Leakage for reveal sets of sizes four through six reduces to four values, one for each orbit under $G$. Four-position reveal sets form two orbits of sizes $42$ and $28$, separated by the classical cross-ratio invariant. Five-position reveal sets form one orbit, and six-position reveal sets form one orbit. Fix an orbit representative $C$ with at least three positions. For either secret value, restriction of the two deals to $C$ is injective. The coalition's observation is therefore uniform over $336$ equally likely possibilities. Its mutual information with the secret class is $1-m/336$ bits, where $m$ is the number of observations shared by the two deals. Sharp three-transitivity justifies this identity by a pen-and-paper argument. The identity does not apply below three positions, where the true leakage is zero. The kernel-checked collision counts for the four representatives are $m=96$, $72$, $36$, and $12$. The corresponding leakage values are $5/7$, $11/14$, $25/28$, and $27/28$ bits. The committed script \path{pgg-smc/instances/pgl27/pgl_leakage_targets.py} evaluates these values from the counts outside the kernel. The four-heart witness represents the orbit of size $42$ and has leakage exactly $5/7$ bits. A pen-and-paper relabeling argument extends each value from its representative to the full orbit because a group element preserves leakage. The formalization does not include this relabeling-invariance argument.
\par\endgroup}}
\Field{Why this rewrite.}{The rewrite adapts the baseline's arbitrary-case calculation sequence. It excludes the footnote, explains sharp three-transitivity and relabeling invariance, and preserves every orbit size, count, leakage value, verification boundary, and script path.}
\Field{Terminology actions.}{Keep: orbit, collision count, mutual information | Explain: sharp three-transitivity, relabeling invariance | Replace: the four-subsets fall into two orbits -> four-position reveal sets form two orbits}
\par\smallskip
\needspace{10\baselineskip}\subsubsection*{W-07-P003: analogous purpose}
\Field{Location.}{Section 07, prose\_paragraph, lines 1534--1546.}
\Field{Draft purpose.}{Explain how exact word counts produce the finite mixing certificate.}
\Field{Qualitative closeness.}{analogous purpose. Both units expose a finite distribution calculation before the formal conclusion that uses it.}
\Field{Limits of the analogy.}{The baseline enumerates a small view distribution. This unit summarizes a very large word space through exact fiber counts and an integer certificate.}
\Field{Baseline narration used.}{adapted.}
\Field{Complexity and jargon.}{Complexity: \Tag{high} $\rightarrow$ \Tag{medium}. Jargon: \Tag{medium} $\rightarrow$ \Tag{low}.}
\textbf{Original WADT unit. Footnotes omitted.}
\begin{lstlisting}
For a word length $L$, the shuffle distribution evaluates a uniform word in
$\{0,1,2,3,4\}^L$. For instance, at the Theorem B length $L=200$ the sample space contains
$5^{200}$ words.
The proof partitions these words by evaluated group element. I call the
set of words that evaluate to $g$ the fiber of $g$, and there are $336$
fibers. Roughly speaking, the certificate counts how often each of the
$336$ group elements occurs among the $5^{200}$ words and checks that no
element occurs much more or much less often than average. For example, at $L=1$ the word consisting of the translation letter
evaluates to $z\mapsto z+1$ and belongs to the fiber of that group element.
A computation over binary naturals records the size of every
fiber. The final arithmetic check sums the deviations of these fiber counts
from $5^{200}/336$.
\end{lstlisting}
\colorbox{RewriteGreen}{\parbox{0.97\textwidth}{\textbf{Rewrite suggestion.}\par\begingroup\small
For word length $L$, the shuffle distribution evaluates a uniform word in $\{0,1,2,3,4\}^L$. At the length $L=200$ used in Theorem B, the sample space contains $5^{200}$ words. The proof partitions these words by their evaluated group element. I call the words that evaluate to $g$ the fiber of $g$. The $336$ group elements therefore define $336$ fibers. Informally, the certificate counts how often each group element occurs among the $5^{200}$ words and checks its deviation from the average. For example, at $L=1$, the translation letter evaluates to $z\mapsto z+1$ and belongs to that group element's fiber. A computation using binary natural numbers records every fiber size. The final arithmetic check sums the deviations from the average size $5^{200}/336$.
\par\endgroup}}
\Field{Why this rewrite.}{Adapts the baseline's calculation sequence to the word space, fiber definition, one-letter example, exact count representation, and final deviation sum.}
\Field{Terminology actions.}{Keep: fiber, fiber count, word distribution | Explain: binary naturals, average fiber size | Replace: the Theorem B length -> the word length used in Theorem B, Roughly speaking -> Informally}
\par\smallskip
\clearpage\subsection{Baseline S-02-P057: security example conclusion}
\colorbox{BaselineBlue}{\parbox{0.97\textwidth}{\Field{Purpose.}{Use input-independent probabilities to discharge the security criterion.}\Field{Primary purpose.}{proof}\Field{Narration structure.}{condition\_verification $\rightarrow$ conclusion}\Field{Reusable move.}{Conclude by matching the established facts directly to the earlier definition.}}}
\needspace{10\baselineskip}\subsubsection*{W-03-STM001: partial analogue}
\Field{Location.}{Section 03, statement, lines 549--558.}
\Field{Draft purpose.}{State the generic transitivity-to-view-independence theorem.}
\Field{Qualitative closeness.}{partial analogue. The baseline supplies one reusable result move: conclude security after identifying why the observation law is independent of the protected input.}
\Field{Limits of the analogy.}{The baseline closes a worked example from calculated probabilities. The WADT unit is a generic theorem statement whose transitivity, distinct-card, prior, uniform-shuffle, and coalition hypotheses precede an independence conclusion.}
\Field{Baseline narration used.}{partial.}
\Field{Complexity and jargon.}{Complexity: \Tag{medium} $\rightarrow$ \Tag{low}. Jargon: \Tag{medium} $\rightarrow$ \Tag{medium}.}
\textbf{Original WADT unit. Footnotes omitted.}
\begin{lstlisting}
Let $G$ act $t$-transitively on the card positions. Let $S$ be a Boolean
secret with any prior, and suppose that every encoded deck has distinct
cards. If $g$ is uniform on $G$, then
\[
  S \mathrel{\perp\!\!\!\perp} V_C
\]
for every coalition $C$ with $|C|\leq t$.
\end{lstlisting}
\colorbox{RewriteGreen}{\parbox{0.97\textwidth}{\textbf{Rewrite suggestion.}\par\begingroup\small
Let $G$ act $t$-transitively on the card positions. Let $S$ be a Boolean secret with any prior. Suppose that every encoded deck assigns distinct card values to its positions. If $g$ is uniform on $G$, then
\[
  S \mathrel{\perp\!\!\!\perp} V_C
\]
for every coalition $C$ with $|C|\leq t$.
\par\endgroup}}
\Field{Why this rewrite.}{The rewrite uses the partial analogue's premises-before-conclusion move. It excludes the footnote, explains the distinct-card condition, and preserves the prior, shuffle, coalition quantifier, and displayed independence claim.}
\Field{Terminology actions.}{Keep: t-transitivity, independence, uniform group element | Explain: encoded deck has distinct cards}
\par\smallskip
\needspace{10\baselineskip}\subsubsection*{W-03-STM002: partial analogue}
\Field{Location.}{Section 03, statement, lines 590--599.}
\Field{Draft purpose.}{State the theorem that transports view independence to executed-trace privacy.}
\Field{Qualitative closeness.}{partial analogue. The baseline supplies the move of deriving a security conclusion from an input-independent observation.}
\Field{Limits of the analogy.}{The baseline works with explicit view probabilities. The WADT theorem states a general entropy equality from mutually determining trace and view variables.}
\Field{Baseline narration used.}{partial.}
\Field{Complexity and jargon.}{Complexity: \Tag{medium} $\rightarrow$ \Tag{medium}. Jargon: \Tag{high} $\rightarrow$ \Tag{medium}.}
\textbf{Original WADT unit. Footnotes omitted.}
\begin{lstlisting}
Let $S$, $V$, and $T$ be finite random variables. Suppose that
$T=\tau(V)$ and that a map $\nu$ satisfies $\nu(\tau(v))=v$ for every $v$.
If $V$ is independent of $S$, then
\begin{equation}
  H(S\mid T)=H(S).
  \label{eq:view-to-trace}
\end{equation}
\end{lstlisting}
\colorbox{RewriteGreen}{\parbox{0.97\textwidth}{\textbf{Rewrite suggestion.}\par\begingroup\small
Let $S$, $V$, and $T$ be finite random variables. Suppose that $T=\tau(V)$. Suppose also that a recovery map $\nu$ satisfies $\nu(\tau(v))=v$ for every $v$. If $V$ is independent of $S$, then
\begin{equation}
  H(S\mid T)=H(S).
  \label{eq:view-to-trace}
\end{equation}
\par\endgroup}}
\Field{Why this rewrite.}{The rewrite uses the partial analogue's premise-to-information-equality move. It excludes the footnote, names the role of both maps, and preserves every formula, quantifier, and label.}
\Field{Terminology actions.}{Keep: conditional entropy, independence, round-trip condition | Explain: nu and tau}
\par\smallskip
\needspace{10\baselineskip}\subsubsection*{W-03-STM003: partial analogue}
\Field{Location.}{Section 03, statement, lines 606--612.}
\Field{Draft purpose.}{State contraction of L1 distance under a deterministic map.}
\Field{Qualitative closeness.}{partial analogue. The baseline supplies the general proof move that applying a deterministic observation does not create an input-dependent distribution.}
\Field{Limits of the analogy.}{The baseline proves equality in one worked security example. The WADT statement gives a general L1 contraction inequality for pushforward distributions.}
\Field{Baseline narration used.}{partial.}
\Field{Complexity and jargon.}{Complexity: \Tag{low} $\rightarrow$ \Tag{low}. Jargon: \Tag{medium} $\rightarrow$ \Tag{medium}.}
\textbf{Original WADT unit. Footnotes omitted.}
\begin{lstlisting}
For finite distributions $P$ and $Q$ on $A$ and a map $f:A\to B$,
\[
  \lVert f_*P-f_*Q\rVert_1\leq\lVert P-Q\rVert_1.
\]
\end{lstlisting}
\colorbox{RewriteGreen}{\parbox{0.97\textwidth}{\textbf{Rewrite suggestion.}\par\begingroup\small
Let $P$ and $Q$ be finite distributions on $A$, and let $f:A\to B$. Their pushforward distributions satisfy
\[
  \lVert f_*P-f_*Q\rVert_1\leq\lVert P-Q\rVert_1.
\]
\par\endgroup}}
\Field{Why this rewrite.}{The rewrite uses the partial analogue's objects-before-inequality move. It excludes the footnote, names \$f\_*P\$ and \$f\_*Q\$ as pushforwards, and preserves the full contraction formula.}
\Field{Terminology actions.}{Keep: L1 distance, pushforward | Explain: f\_*P}
\par\smallskip
\needspace{10\baselineskip}\subsubsection*{W-04-STM002: partial analogue}
\Field{Location.}{Section 04, statement, lines 800--806.}
\Field{Draft purpose.}{State the checked finite biased-cut guarantee.}
\Field{Qualitative closeness.}{partial analogue. The baseline supplies the move of turning a distribution calculation into a formal security conclusion.}
\Field{Limits of the analogy.}{The baseline closes a small input-independence example. This statement gives a numerical endpoint bound for a biased shuffle.}
\Field{Baseline narration used.}{partial.}
\Field{Complexity and jargon.}{Complexity: \Tag{medium} $\rightarrow$ \Tag{low}. Jargon: \Tag{high} $\rightarrow$ \Tag{medium}.}
\textbf{Original WADT unit. Footnotes omitted.}
\begin{lstlisting}
Let the shuffle distribution be the biased cut $w_{1/100}$ repeated seven
times. Then every single-card endpoint distribution is within $L_1$
distance $2^{-40}$ of uniform, and the kernel checks the computation.
\end{lstlisting}
\colorbox{RewriteGreen}{\parbox{0.97\textwidth}{\textbf{Rewrite suggestion.}\par\begingroup\small
Let the shuffle distribution be the biased cut $w_{1/100}$ repeated seven times. Then every single-card endpoint distribution is within $L_1$ distance $2^{-40}$ of uniform. The kernel checks this computation.
\par\endgroup}}
\Field{Why this rewrite.}{Uses only the baseline's calculation-to-conclusion move and retains the bias, iteration count, universal endpoint claim, bound, and kernel check.}
\Field{Terminology actions.}{Keep: L1 distance, single-card endpoint distribution, kernel checks | Explain: biased cut repeated seven times}
\par\smallskip
\needspace{10\baselineskip}\subsubsection*{W-07-STM001: partial analogue}
\Field{Location.}{Section 07, statement, lines 1548--1555.}
\Field{Draft purpose.}{State the checked mixing bound for the symmetric word law.}
\Field{Qualitative closeness.}{partial analogue. The baseline supplies the move of stating the formal distributional conclusion after a finite calculation.}
\Field{Limits of the analogy.}{The baseline proves security from input-independent probabilities. This statement concerns proximity to uniformity, not protocol privacy.}
\Field{Baseline narration used.}{partial.}
\Field{Complexity and jargon.}{Complexity: \Tag{low} $\rightarrow$ \Tag{low}. Jargon: \Tag{medium} $\rightarrow$ \Tag{low}.}
\textbf{Original WADT unit. Footnotes omitted.}
\begin{lstlisting}
Let $\mu$ be uniform on the five-letter symmetric generator tuple. Then
\begin{equation}
  \left\lVert \mu^{*200}-U_G\right\rVert_1 \leq 2^{-40}.
  \label{eq:pgl-mixing}
\end{equation}
\end{lstlisting}
\colorbox{RewriteGreen}{\parbox{0.97\textwidth}{\textbf{Rewrite suggestion.}\par\begingroup\small
Let $\mu$ be uniform on the five-letter symmetric generator tuple. Then
\begin{equation}
  \left\lVert \mu^{*200}-U_G\right\rVert_1 \leq 2^{-40}.
  \label{eq:pgl-mixing}
\end{equation}
\par\endgroup}}
\Field{Why this rewrite.}{Uses only the baseline's calculation-to-conclusion move and preserves the distribution, convolution length, uniform law, label, and exact bound.}
\Field{Terminology actions.}{Keep: L1 distance, uniform group law | Explain: convolution power 200}
\par\smallskip
\clearpage\subsection{Baseline S-02-P058: subroutine intuition and diagram}
\colorbox{BaselineBlue}{\parbox{0.97\textwidth}{\Field{Purpose.}{Present a subroutine as an atomic replacement for a complete protocol execution.}\Field{Primary purpose.}{definition}\Field{Narration structure.}{object\_setup $\rightarrow$ formal\_definition $\rightarrow$ scope\_boundary}\Field{Reusable move.}{Introduce the object, state its rule, and close with the boundary needed for later use.}}}
\needspace{10\baselineskip}\subsubsection*{W-03-P012: partial analogue}
\Field{Location.}{Section 03, prose\_paragraph, lines 601--604.}
\Field{Draft purpose.}{Motivate transport of a group-distribution bound to protocol observations.}
\Field{Qualitative closeness.}{partial analogue. The baseline supplies the move of connecting an abstract object to the protocol-level object that consumes it.}
\Field{Limits of the analogy.}{The baseline presents a subroutine as an atomic protocol replacement. The WADT paragraph relates a group-level distance to endpoint, view, and trace observables.}
\Field{Baseline narration used.}{partial.}
\Field{Complexity and jargon.}{Complexity: \Tag{low} $\rightarrow$ \Tag{low}. Jargon: \Tag{medium} $\rightarrow$ \Tag{low}.}
\textbf{Original WADT unit. Footnotes omitted.}
\begin{lstlisting}
The word-shuffle certificate bounds distributions of group elements, while
protocol security concerns endpoints, coalition views, and traces. The next
theorem transfers such a bound through any deterministic observable.
\end{lstlisting}
\colorbox{RewriteGreen}{\parbox{0.97\textwidth}{\textbf{Rewrite suggestion.}\par\begingroup\small
The word-shuffle certificate bounds distributions of group elements. Protocol security instead concerns card positions after the shuffle, coalition views, and traces. The next theorem transfers the bound through any deterministic observable.
\par\endgroup}}
\Field{Why this rewrite.}{The rewrite uses the partial analogue's abstraction-gap move. It explains endpoints as post-shuffle card positions and names deterministic transport as the bridge.}
\Field{Terminology actions.}{Keep: deterministic observable, coalition view | Explain: endpoint}
\par\smallskip
\clearpage\subsection{Baseline S-02-P065: composed-protocol construction}
\colorbox{BaselineBlue}{\parbox{0.97\textwidth}{\Field{Purpose.}{Adapt the smaller protocol and use it to construct a larger protocol instance.}\Field{Primary purpose.}{example}\Field{Narration structure.}{construction\_goal $\rightarrow$ mechanism $\rightarrow$ realisation}\Field{Reusable move.}{Connect an abstract operation to the concrete mechanism that realises it.}}}
\needspace{10\baselineskip}\subsubsection*{W-03-P007: analogous purpose}
\Field{Location.}{Section 03, prose\_paragraph, lines 521--536.}
\Field{Draft purpose.}{Explain how input encoding extends the common flow to committed inputs and why correctness is preserved.}
\Field{Qualitative closeness.}{analogous purpose. Both paragraphs explain how an existing protocol component is adapted or invoked to construct a larger protocol while preserving its interface.}
\Field{Limits of the analogy.}{The baseline works through one composed protocol instance. The WADT paragraph describes a generic commit prologue, two correctness obligations, a reused program, and the empty-input reduction.}
\Field{Baseline narration used.}{adapted.}
\Field{Complexity and jargon.}{Complexity: \Tag{high} $\rightarrow$ \Tag{medium}. Jargon: \Tag{high} $\rightarrow$ \Tag{medium}.}
\textbf{Original WADT unit. Footnotes omitted.}
\begin{lstlisting}
With an \coqin{InputEncoding}, the flow gains a phase that I call the
commit prologue. It collects the players' inputs and assembles the dealt
deck from them, so the same flow evaluates a function of committed
inputs. The realized encoding is den Boer's. Its
orbit obligation places the assembled layouts of equal-output inputs in one
shuffle orbit. This property relates different inputs that produce the same
result. Correctness uses two other facts: the assembled layout is valid, and
reconstruction is invariant under every allowed shuffle. It follows that the
shuffled layout reconstructs the intended output. The Kim variant reuses the
den Boer program unchanged. With an empty input list
the prologue reduces by computation to the plain dealer, which is the
secret-sharing case, and the $\PG$ instance passes exactly this empty
list. Section~\ref{sec:fivecard} realizes the committed-input flow, and
the instance table of Section~\ref{sec:instances} records the full
comparison.
\end{lstlisting}
\colorbox{RewriteGreen}{\parbox{0.97\textwidth}{\textbf{Rewrite suggestion.}\par\begingroup\small
With an input-encoding component, the flow gains a commit prologue. This phase collects the players' inputs and assembles the dealt deck, so the flow computes a function of committed inputs. The realized encoding is den Boer's. Its orbit obligation places assembled layouts for equal-output inputs in one shuffle orbit. Thus, it relates different inputs with the same result. Correctness also requires a valid assembled layout and reconstruction invariant under every allowed shuffle. These facts imply that the shuffled layout reconstructs the intended output. The Kim variant reuses the den Boer program unchanged. With an empty input list, the prologue computes to the plain dealer used for secret sharing. The $\PG$ instance supplies exactly this empty list. Section~\ref{sec:fivecard} realizes the committed-input flow. The instance table in Section~\ref{sec:instances} gives the full comparison.
\par\endgroup}}
\Field{Why this rewrite.}{The rewrite adapts the baseline's mechanism-to-realization sequence. It explains the assembled layout, separates the three correctness obligations, and preserves Kim reuse, the empty-input reduction, and both cross-references.}
\Field{Terminology actions.}{Keep: commit prologue, orbit obligation, shuffle-invariant reconstruction | Explain: assembled layout | Replace: With an InputEncoding -> With an input-encoding component}
\par\smallskip
\clearpage\subsection{Baseline S-02-P068: composition example verification}
\colorbox{BaselineBlue}{\parbox{0.97\textwidth}{\Field{Purpose.}{Verify the two defining composition conditions for the worked protocol.}\Field{Primary purpose.}{proof}\Field{Narration structure.}{condition\_verification $\rightarrow$ conclusion}\Field{Reusable move.}{Conclude by matching the established facts directly to the earlier definition.}}}
\needspace{10\baselineskip}\subsubsection*{W-05-STM003: analogous purpose}
\Field{Location.}{Section 05, statement, lines 953--974.}
\Field{Draft purpose.}{Characterize orbit equivalence by classifier equality and count both classes.}
\Field{Qualitative closeness.}{analogous purpose. Both units verify the defining conditions of a constructed object and close with the resulting formal property.}
\Field{Limits of the analogy.}{The baseline verifies a composed protocol example. The WADT statement proves an orbit-class equivalence and adds exact cardinalities.}
\Field{Baseline narration used.}{adapted.}
\Field{Complexity and jargon.}{Complexity: \Tag{high} $\rightarrow$ \Tag{medium}. Jargon: \Tag{high} $\rightarrow$ \Tag{medium}.}
\textbf{Original WADT unit. Footnotes omitted.}
\begin{lstlisting}
For all $A,B\in\Omega_4$,
\[
  \kappa(A)=\kappa(B)
  \quad\Longleftrightarrow\quad
  \exists g\in G,\ B=g\mathbin{\cdot}A.
\]
Hence
\[
  \Omega_4/G\cong\{0,1\},
  \qquad
  |\kappa^{-1}(0)|=42,
  \qquad
  |\kappa^{-1}(1)|=28.
\]
The fiber $\kappa^{-1}(0)$ is harmonic, and
$\kappa^{-1}(1)$ is equianharmonic.
\end{lstlisting}
\colorbox{RewriteGreen}{\parbox{0.97\textwidth}{\textbf{Rewrite suggestion.}\par\begingroup\small
For all $A,B\in\Omega_4$,
\[
  \kappa(A)=\kappa(B)
  \quad\Longleftrightarrow\quad
  \exists g\in G,\ B=g\mathbin{\cdot}A.
\]
Thus the quotient by orbit equivalence has two classes:
\[
  \Omega_4/G\cong\{0,1\},
  \qquad
  |\kappa^{-1}(0)|=42,
  \qquad
  |\kappa^{-1}(1)|=28.
\]
The fiber $\kappa^{-1}(0)$ is harmonic. The fiber $\kappa^{-1}(1)$ is equianharmonic.
\par\endgroup}}
\Field{Why this rewrite.}{Adapts the baseline's condition-to-conclusion sequence to the orbit equivalence, quotient, exact fiber cardinalities, and class names. The excluded footnote does not enter the suggestion.}
\Field{Terminology actions.}{Keep: orbit quotient, fiber cardinality | Explain: quotient by the group -> the quotient by orbit equivalence}
\par\smallskip
\clearpage\subsection{Baseline S-03-P001: dihedral-card capability lead-in}
\colorbox{BaselineBlue}{\parbox{0.97\textwidth}{\Field{Purpose.}{Introduce the parameterised card type through a list of its required capabilities.}\Field{Primary purpose.}{definition}\Field{Narration structure.}{motivation $\rightarrow$ concept\_scope $\rightarrow$ formal\_handoff}\Field{Reusable move.}{Explain why a concept is needed before introducing its formal specification.}}}
\needspace{10\baselineskip}\subsubsection*{W-04-P002: analogous purpose}
\Field{Location.}{Section 04, prose\_paragraph, lines 673--681.}
\Field{Draft purpose.}{Introduce the first concrete profile and prepare its formal definition.}
\Field{Qualitative closeness.}{analogous purpose. Both units motivate a formal object, delimit its components, and hand the reader to its specification.}
\Field{Limits of the analogy.}{The baseline introduces a card type through capabilities. This unit introduces a proof-oriented protocol profile.}
\Field{Baseline narration used.}{adapted.}
\Field{Complexity and jargon.}{Complexity: \Tag{low} $\rightarrow$ \Tag{low}. Jargon: \Tag{medium} $\rightarrow$ \Tag{low}.}
\textbf{Original WADT unit. Footnotes omitted.}
\begin{lstlisting}
\needspace{5\baselineskip}
The five-card family is the first concrete example of how the framework
components are packed into one profile. The following definition combines
the cyclic action, the biased-cut security witness, and the three-hearts
decoder.
\end{lstlisting}
\colorbox{RewriteGreen}{\parbox{0.97\textwidth}{\textbf{Rewrite suggestion.}\par\begingroup\small
\needspace{5\baselineskip} The five-card family first shows how one profile combines the framework components. The following definition combines the cyclic action, the security witness for the biased cut, and the three-hearts decoder.
\par\endgroup}}
\Field{Why this rewrite.}{Adapts the baseline's motivation and formal handoff. It names the profile's three components directly and omits the excluded footnote.}
\Field{Terminology actions.}{Keep: cyclic action, decoder | Explain: profile, biased-cut security witness | Replace: packed into one profile -> combined in one profile}
\par\smallskip
\clearpage\subsection{Baseline S-03-P002: dihedral-card capabilities}
\colorbox{BaselineBlue}{\parbox{0.97\textwidth}{\Field{Purpose.}{Specify the value domain, permitted transformations, and restricted observations.}\Field{Primary purpose.}{definition}\Field{Narration structure.}{object\_setup $\rightarrow$ formal\_definition $\rightarrow$ scope\_boundary}\Field{Reusable move.}{Introduce the object, state its rule, and close with the boundary needed for later use.}}}
\needspace{10\baselineskip}\subsubsection*{W-05-P020: partial analogue}
\Field{Location.}{Section 05, prose\_paragraph, lines 1173--1187.}
\Field{Draft purpose.}{Present the three generating maps and connect group order to sharp three-transitivity.}
\Field{Qualitative closeness.}{partial analogue. The baseline supplies the move of inventorying the allowed transformations and assigning each an operational role.}
\Field{Limits of the analogy.}{The WADT unit also explains a figure, expands an executable alphabet, and derives a uniqueness consequence from cardinality.}
\Field{Baseline narration used.}{partial.}
\Field{Complexity and jargon.}{Complexity: \Tag{high} $\rightarrow$ \Tag{medium}. Jargon: \Tag{high} $\rightarrow$ \Tag{medium}.}
\textbf{Original WADT unit. Footnotes omitted.}
\begin{lstlisting}
For $G=\PG$, three fractional linear maps generate the group. They are the
translation $z\mapsto z+1$, the scaling $z\mapsto 3z$, and the inversion
$z\mapsto -1/z$.
Figure~\ref{fig:pgl-generators} shows each generator as a rearrangement of
the row of eight cards, one card per projective point. The five-letter
alphabet of the executable shuffle in Section~\ref{sec:mixing} contains
these three maps together with the inverses of translation and scaling.
The equality $|G|=336=8\cdot7\cdot6$ also matches the number of ordered
triples of distinct points. Lemma~\ref{thm:three-transitive} produces a
group element for each ordered destination, and the count equality makes
that element unique.
\end{lstlisting}
\colorbox{RewriteGreen}{\parbox{0.97\textwidth}{\textbf{Rewrite suggestion.}\par\begingroup\small
For $G=\PG$, three fractional linear transformations generate the group. Translation is $z\mapsto z+1$. Scaling is $z\mapsto 3z$. Inversion is $z\mapsto -1/z$. On the projective line, translation and scaling fix $\infty$, while inversion exchanges $0$ and $\infty$. Figure~\ref{fig:pgl-generators} shows each map as a rearrangement of eight cards, one at each projective point. The five-letter alphabet used by the executable shuffle in Section~\ref{sec:mixing} adds the inverses of translation and scaling to these three maps. The group order is $|G|=336=8\cdot7\cdot6$, which is also the number of ordered triples of distinct points. Lemma~\ref{thm:three-transitive} provides a group element for each ordered destination triple. The count then makes this element unique.
\par\endgroup}}
\Field{Why this rewrite.}{Uses the baseline listing move for the three transformations, then follows the WADT contract for the projective action, figure, executable alphabet, and uniqueness count. The excluded footnote does not enter the suggestion.}
\Field{Terminology actions.}{Keep: translation, scaling, and inversion, ordered triple | Explain: fractional linear transformation -> translation and scaling fix infinity while inversion exchanges zero and infinity, three-transitive -> a group element exists for each ordered destination triple}
\par\smallskip
\clearpage\subsection{Baseline S-03-P003: polygonal representation consequence}
\colorbox{BaselineBlue}{\parbox{0.97\textwidth}{\Field{Purpose.}{Relate the abstract capability list to a geometric card shape and a concrete instance.}\Field{Primary purpose.}{interpretation}\Field{Narration structure.}{structural\_consequence $\rightarrow$ example\_handoff}\Field{Reusable move.}{State the structural consequence of a definition before giving a concrete illustration.}}}
\needspace{10\baselineskip}\subsubsection*{W-04-P001: partial analogue}
\Field{Location.}{Section 04, prose\_paragraph, lines 616--623.}
\Field{Draft purpose.}{Introduce the five-card family and distinguish its two dealer variants.}
\Field{Qualitative closeness.}{partial analogue. The baseline contributes only the move from an abstract structural account to a concrete instance and figure.}
\Field{Limits of the analogy.}{The baseline derives a geometric representation from formal capabilities. This paragraph instead introduces a protocol family, compares two dealer laws, and previews a run.}
\Field{Baseline narration used.}{partial.}
\Field{Complexity and jargon.}{Complexity: \Tag{medium} $\rightarrow$ \Tag{low}. Jargon: \Tag{medium} $\rightarrow$ \Tag{low}.}
\textbf{Original WADT unit. Footnotes omitted.}
\begin{lstlisting}
The simplest instantiation of the framework covers two classic protocols
with one record. Five cards carry the deck, the cyclic group $C_5$ acts by
rotation, and the decoder reads whether the three hearts are consecutive.
Den Boer's five-card protocol~\cite{denBoer1989} and the Kim and
\c{C}etinkaya variant~\cite{KimCetinkaya2025} differ only in the dealing
distribution, so one profile at bias $\varepsilon$ specifies both.
Figure~\ref{fig:fivecard-run} shows a run.
\end{lstlisting}
\colorbox{RewriteGreen}{\parbox{0.97\textwidth}{\textbf{Rewrite suggestion.}\par\begingroup\small
The five-card family is the simplest concrete instance of the framework. One profile covers two classic protocols. The cyclic group $C_5$ rotates the five-card deck, and the decoder tests whether the three hearts are consecutive. Den Boer's protocol~\cite{denBoer1989} and the Kim and \c{C}etinkaya variant~\cite{KimCetinkaya2025} differ only in their dealing distributions. A bias parameter $\varepsilon$ therefore distinguishes them within the same profile. Figure~\ref{fig:fivecard-run} shows one run.
\par\endgroup}}
\Field{Why this rewrite.}{Uses the baseline's concrete-instance handoff, then retains the two protocols, shared profile, bias parameter, and figure reference in a direct sequence.}
\Field{Terminology actions.}{Keep: cyclic group, decoder | Explain: bias parameter | Replace: with one record -> in one profile}
\par\smallskip
\needspace{10\baselineskip}\subsubsection*{W-05-P021: partial analogue}
\Field{Location.}{Section 05, prose\_paragraph, lines 1214--1217.}
\Field{Draft purpose.}{Move from generator tables to a worked physical and algebraic example.}
\Field{Qualitative closeness.}{partial analogue. The baseline supplies the move of relating an abstract formal representation to a concrete physical interpretation.}
\Field{Limits of the analogy.}{The WADT unit is a transition to an example and also prepares a group-word transport argument.}
\Field{Baseline narration used.}{partial.}
\Field{Complexity and jargon.}{Complexity: \Tag{low} $\rightarrow$ \Tag{low}. Jargon: \Tag{medium} $\rightarrow$ \Tag{low}.}
\textbf{Original WADT unit. Footnotes omitted.}
\begin{lstlisting}
The figure gives the three permutation tables. The following example reads
them as physical row movements and connects words in these generators to
the triple transport used below.
\end{lstlisting}
\colorbox{RewriteGreen}{\parbox{0.97\textwidth}{\textbf{Rewrite suggestion.}\par\begingroup\small
The figure gives the three permutation tables. The following example first reads each table as a physical movement of the card row. It then shows how a generator word sends one ordered triple of source points to an ordered triple of target points.
\par\endgroup}}
\Field{Why this rewrite.}{Uses the baseline concrete-interpretation move before the WADT example connects the tables to triple transport.}
\Field{Terminology actions.}{Keep: permutation table | Explain: triple transport -> a generator word sends an ordered source triple to an ordered target triple}
\par\smallskip
\clearpage\subsection{Baseline S-03-P004: physical encoding design}
\colorbox{BaselineBlue}{\parbox{0.97\textwidth}{\Field{Purpose.}{Explain how visible geometry and hidden markings encode the card state.}\Field{Primary purpose.}{method}\Field{Narration structure.}{construction\_goal $\rightarrow$ mechanism $\rightarrow$ realisation}\Field{Reusable move.}{Connect an abstract operation to the concrete mechanism that realises it.}}}
\needspace{10\baselineskip}\subsubsection*{W-04-CAP001: partial analogue}
\Field{Location.}{Section 04, caption, lines 667--670.}
\Field{Draft purpose.}{Explain how one physical five-card run corresponds to the formal one-card-per-input encoding.}
\Field{Qualitative closeness.}{partial analogue. The baseline contributes the correspondence move that connects an abstract encoding with its physical realization.}
\Field{Limits of the analogy.}{The baseline explains one card mechanism in running prose. This caption contrasts a two-card physical commitment with a one-value formal encoding for a specific figure.}
\Field{Baseline narration used.}{partial.}
\Field{Complexity and jargon.}{Complexity: \Tag{medium} $\rightarrow$ \Tag{low}. Jargon: \Tag{medium} $\rightarrow$ \Tag{low}.}
\textbf{Original WADT unit. Footnotes omitted.}
\begin{lstlisting}
One five-card run with committed inputs. The physical protocol
  commits each input bit as two face-down cards, and the formalization
  encodes each committed bit as one card value.
\end{lstlisting}
\colorbox{RewriteGreen}{\parbox{0.97\textwidth}{\textbf{Rewrite suggestion.}\par\begingroup\small
One five-card run with committed inputs. The physical protocol represents each input bit with two face-down cards. The formalization represents each committed bit with one card value.
\par\endgroup}}
\Field{Why this rewrite.}{Uses the baseline's abstract-to-concrete correspondence and separates the physical and formal encodings into parallel sentences.}
\Field{Terminology actions.}{Keep: committed input, card value | Explain: formal encoding}
\par\smallskip
\clearpage\subsection{Baseline S-03-P019: operation inventory and permutation}
\colorbox{BaselineBlue}{\parbox{0.97\textwidth}{\Field{Purpose.}{Inventory the operations and align the permutation case with the earlier model.}\Field{Primary purpose.}{definition}\Field{Narration structure.}{operation\_inventory $\rightarrow$ prior\_definition\_link $\rightarrow$ definition}\Field{Reusable move.}{Open a technical catalogue with the complete inventory and reuse prior definitions where possible.}}}
\needspace{10\baselineskip}\subsubsection*{W-01-P005: analogous purpose}
\Field{Location.}{Section 01, prose\_paragraph, lines 137--138.}
\Field{Draft purpose.}{Prepare the reader for an enumerated contribution list.}
\Field{Qualitative closeness.}{analogous purpose. Both units introduce a finite inventory before the text names each member.}
\Field{Limits of the analogy.}{The baseline inventories operation families within a model. The WADT sentence introduces a list of paper contributions.}
\Field{Baseline narration used.}{adapted.}
\Field{Complexity and jargon.}{Complexity: \Tag{low} $\rightarrow$ \Tag{low}. Jargon: \Tag{low} $\rightarrow$ \Tag{low}.}
\textbf{Original WADT unit. Footnotes omitted.}
\begin{lstlisting}
There are five contributions in this work.
\end{lstlisting}
\colorbox{RewriteGreen}{\parbox{0.97\textwidth}{\textbf{Rewrite suggestion.}\par\begingroup\small
I make five contributions.
\par\endgroup}}
\Field{Why this rewrite.}{The rewrite adapts the baseline's inventory opening and uses first-person singular for the sole author's contribution.}
\Field{Terminology actions.}{Keep: contributions}
\par\smallskip
\needspace{10\baselineskip}\subsubsection*{W-07-P002: analogous purpose}
\Field{Location.}{Section 07, prose\_paragraph, lines 1525--1533.}
\Field{Draft purpose.}{Define the symmetric generator alphabet and justify its inverse closure.}
\Field{Qualitative closeness.}{analogous purpose. Both units inventory an operation family and explain how it relates to an already defined generating set.}
\Field{Limits of the analogy.}{The baseline aligns card operations with an earlier model. This unit adds inverse generators for a counting algorithm and proves the group is unchanged.}
\Field{Baseline narration used.}{adapted.}
\Field{Complexity and jargon.}{Complexity: \Tag{medium} $\rightarrow$ \Tag{low}. Jargon: \Tag{high} $\rightarrow$ \Tag{medium}.}
\textbf{Original WADT unit. Footnotes omitted.}
\begin{lstlisting}
The executable shuffle samples a word over five letters. Three letters are
the projective translation, scaling, and inversion generators. The alphabet
also contains the inverses of translation and scaling. Inversion is already
self-inverse. The alphabet includes the two inverse letters because the fiber
computation below steps backward through a letter's inverse and must stay
inside the alphabet. The symmetric tuple still generates the same group
as the original three generators, so the fiber count concerns the same
$336$ elements.
\end{lstlisting}
\colorbox{RewriteGreen}{\parbox{0.97\textwidth}{\textbf{Rewrite suggestion.}\par\begingroup\small
The executable shuffle samples a word over five letters. Three letters are the projective translation, scaling, and inversion generators. The other two are the inverses of translation and scaling. Inversion is self-inverse. The fiber computation steps backward through each letter's inverse, so every inverse must remain in the alphabet. This symmetric five-letter generator tuple generates the same group as the original three generators. The fiber count therefore still concerns the same $336$ group elements.
\par\endgroup}}
\Field{Why this rewrite.}{Adapts the baseline's operation inventory and prior-definition link. It lists the five letters, explains inverse closure, and preserves the generation and group-order claims.}
\Field{Terminology actions.}{Keep: translation, scaling, and inversion, inverse letters, same group | Explain: symmetric generator tuple}
\par\smallskip
\clearpage\subsection{Baseline S-03-P020: rotation-operation syntax}
\colorbox{BaselineBlue}{\parbox{0.97\textwidth}{\Field{Purpose.}{Introduce the parameters and notation for a positional rotation.}\Field{Primary purpose.}{definition}\Field{Narration structure.}{operation\_label $\rightarrow$ parameter\_scope $\rightarrow$ formal\_form}\Field{Reusable move.}{Present an operation name and parameter scope before explaining its action.}}}
\needspace{10\baselineskip}\subsubsection*{W-04-P003: analogous purpose}
\Field{Location.}{Section 04, prose\_paragraph, lines 690--692.}
\Field{Draft purpose.}{Introduce notation for the five-card dealing law.}
\Field{Qualitative closeness.}{analogous purpose. Both units introduce parameters and notation immediately before a formal specification.}
\Field{Limits of the analogy.}{The baseline defines a positional rotation. This unit attributes a probability distribution and prepares a displayed law.}
\Field{Baseline narration used.}{adapted.}
\Field{Complexity and jargon.}{Complexity: \Tag{low} $\rightarrow$ \Tag{low}. Jargon: \Tag{medium} $\rightarrow$ \Tag{low}.}
\textbf{Original WADT unit. Footnotes omitted.}
\begin{lstlisting}
The dealing distribution is the biased cut of Kim and \c{C}etinkaya,
and I write $w_\varepsilon$ for it:
\end{lstlisting}
\colorbox{RewriteGreen}{\parbox{0.97\textwidth}{\textbf{Rewrite suggestion.}\par\begingroup\small
Kim and \c{C}etinkaya's biased cut defines the dealing distribution. I denote this distribution by $w_\varepsilon$:
\par\endgroup}}
\Field{Why this rewrite.}{Adapts the baseline's parameter-to-notation sequence and states the attribution before introducing the symbol.}
\Field{Terminology actions.}{Keep: dealing distribution, w\_epsilon | Explain: biased cut}
\par\smallskip
\clearpage\subsection{Baseline S-04-P001: initialization functionality lead-in}
\colorbox{BaselineBlue}{\parbox{0.97\textwidth}{\Field{Purpose.}{Introduce the formal target relation before the display that specifies it.}\Field{Primary purpose.}{definition}\Field{Narration structure.}{definition\_lead\_in $\rightarrow$ display\_handoff}\Field{Reusable move.}{Use one prose lead-in to identify the role of an immediately following formal display.}}}
\needspace{10\baselineskip}\subsubsection*{W-05-P002: analogous purpose}
\Field{Location.}{Section 05, prose\_paragraph, lines 846--848.}
\Field{Draft purpose.}{Introduce the set that indexes the eight card positions.}
\Field{Qualitative closeness.}{analogous purpose. Both units state a formal target object immediately before the display that specifies it.}
\Field{Limits of the analogy.}{The baseline introduces an input-output relation. The WADT unit introduces a geometric index set.}
\Field{Baseline narration used.}{adapted.}
\Field{Complexity and jargon.}{Complexity: \Tag{low} $\rightarrow$ \Tag{low}. Jargon: \Tag{medium} $\rightarrow$ \Tag{low}.}
\textbf{Original WADT unit. Footnotes omitted.}
\begin{lstlisting}
The construction indexes the eight card positions by the projective line
over $\mathbb{F}_7$:
\end{lstlisting}
\colorbox{RewriteGreen}{\parbox{0.97\textwidth}{\textbf{Rewrite suggestion.}\par\begingroup\small
I index the eight card positions by the projective line over $\mathbb{F}_7$, which consists of seven field elements and one point at infinity:
\par\endgroup}}
\Field{Why this rewrite.}{Adapts the baseline's formal-display handoff to the geometric index set and explains why it has eight points.}
\Field{Terminology actions.}{Keep: projective line | Explain: projective line over the seven-element field -> seven field elements and one point at infinity}
\par\smallskip
\clearpage\subsection{Baseline S-04-P003: initialization procedure handoff}
\colorbox{BaselineBlue}{\parbox{0.97\textwidth}{\Field{Purpose.}{Open the ordered execution of the initialization protocol.}\Field{Primary purpose.}{transition}\Field{Narration structure.}{procedure\_announcement $\rightarrow$ list\_handoff}\Field{Reusable move.}{Use a short handoff after a formal tuple so the operational account begins cleanly.}}}
\needspace{10\baselineskip}\subsubsection*{W-02-P007: analogous purpose}
\Field{Location.}{Section 02, prose\_paragraph, lines 228--233.}
\Field{Draft purpose.}{Hand the abstract execution model to a concrete player process.}
\Field{Qualitative closeness.}{analogous purpose. Both units move from a formal object to the ordered procedure that realizes it.}
\Field{Limits of the analogy.}{The baseline opens a complete protocol procedure. The WADT unit introduces only the player-process listing, while its protected footnote remains excluded evidence.}
\Field{Baseline narration used.}{adapted.}
\Field{Complexity and jargon.}{Complexity: \Tag{low} $\rightarrow$ \Tag{low}. Jargon: \Tag{medium} $\rightarrow$ \Tag{low}.}
\textbf{Original WADT unit. Footnotes omitted.}
\begin{lstlisting}
The following player process makes this execution model
concrete.
\end{lstlisting}
\colorbox{RewriteGreen}{\parbox{0.97\textwidth}{\textbf{Rewrite suggestion.}\par\begingroup\small
The following player process makes the execution model concrete by showing the values that the player receives and reveals:
\par\endgroup}}
\Field{Why this rewrite.}{Adapts the baseline's short procedure handoff to the specific player-process listing. The excluded footnote does not enter the suggestion.}
\Field{Terminology actions.}{Keep: player process | Explain: execution model -> the values that the player receives and reveals}
\par\smallskip
\needspace{10\baselineskip}\subsubsection*{W-04-DIA006: partial analogue}
\Field{Location.}{Section 04, diagram\_text, lines 654--655.}
\Field{Draft purpose.}{Name the dealer actions in their execution order.}
\Field{Qualitative closeness.}{partial analogue. The baseline supplies the useful move of presenting protocol actions in execution order.}
\Field{Limits of the analogy.}{The baseline opens a prose procedure. This unit is a two-line action node inside a diagram.}
\Field{Baseline narration used.}{partial.}
\Field{Complexity and jargon.}{Complexity: \Tag{low} $\rightarrow$ \Tag{low}. Jargon: \Tag{medium} $\rightarrow$ \Tag{low}.}
\textbf{Original WADT unit. Footnotes omitted.}
\begin{lstlisting}
assemble the five-card deck\\[-1pt] draw the cut
\end{lstlisting}
\colorbox{RewriteGreen}{\parbox{0.97\textwidth}{\textbf{Rewrite suggestion.}\par\begingroup\small
assemble the five-card deck\\[-1pt]draw the cut
\par\endgroup}}
\Field{Why this rewrite.}{Uses only the baseline's ordered-action move and preserves the two diagram actions.}
\Field{Terminology actions.}{Keep: five-card deck | Explain: cut}
\par\smallskip
\clearpage\subsection{Baseline S-04-P004: initialization correctness and observed-view setup}
\colorbox{BaselineBlue}{\parbox{0.97\textwidth}{\Field{Purpose.}{Close correctness and begin security with the concrete transcript for an arbitrary input.}\Field{Primary purpose.}{proof}\Field{Narration structure.}{termination $\rightarrow$ correctness\_assessment $\rightarrow$ security\_setup $\rightarrow$ view\_instantiation}\Field{Reusable move.}{Join a brief correctness close to the exact transcript that launches security reasoning.}}}
\needspace{10\baselineskip}\subsubsection*{W-05-P013: analogous purpose}
\Field{Location.}{Section 05, prose\_paragraph, lines 1053--1059.}
\Field{Draft purpose.}{Extend representative correctness to shuffled executions and prepare the deck-action definition.}
\Field{Qualitative closeness.}{analogous purpose. Both units close an initial correctness fact and prepare the execution-level object needed for the next proof.}
\Field{Limits of the analogy.}{The baseline also begins a security transcript. The WADT unit instead ends with a handoff to an unattached deck-action display.}
\Field{Baseline narration used.}{adapted.}
\Field{Complexity and jargon.}{Complexity: \Tag{medium} $\rightarrow$ \Tag{low}. Jargon: \Tag{high} $\rightarrow$ \Tag{low}.}
\textbf{Original WADT unit. Footnotes omitted.}
\begin{lstlisting}
The example establishes
$\operatorname{dec}(\operatorname{enc}(s))=s$ before shuffling. Correctness
of a full run also requires the decoder to return the same value after every
allowed shuffle. Let $\rho:G\to S_X$ denote the induced permutation of card
positions. Since a deck $D:X\to C$ maps each position to its card value,
define
\end{lstlisting}
\colorbox{RewriteGreen}{\parbox{0.97\textwidth}{\textbf{Rewrite suggestion.}\par\begingroup\small
The example proves $\operatorname{dec}(\operatorname{enc}(s))=s$ before shuffling. Correctness of a full run also requires the decoder to return the same value after every allowed shuffle. Let $\rho:G\to S_X$ be the permutation of card positions induced by the group action. A deck $D:X\to C$ maps each position to its card value. Define the shuffled deck by:
\par\endgroup}}
\Field{Why this rewrite.}{Adapts the baseline's correctness-to-execution transition and leads directly to the deck-action display.}
\Field{Terminology actions.}{Keep: decoder, induced permutation | Explain: induced permutation -> the permutation of card positions induced by the group action}
\par\smallskip
\needspace{10\baselineskip}\subsubsection*{W-05-P016: analogous purpose}
\Field{Location.}{Section 05, prose\_paragraph, lines 1089--1096.}
\Field{Draft purpose.}{Connect deck invariance to one complete protocol run.}
\Field{Qualitative closeness.}{analogous purpose. Both units connect a concrete execution with a correctness result and then prepare a distinct later property.}
\Field{Limits of the analogy.}{The baseline begins a security transcript immediately. The WADT unit points to interpreter lifting and defers recovery and privacy.}
\Field{Baseline narration used.}{adapted.}
\Field{Complexity and jargon.}{Complexity: \Tag{high} $\rightarrow$ \Tag{medium}. Jargon: \Tag{high} $\rightarrow$ \Tag{low}.}
\textbf{Original WADT unit. Footnotes omitted.}
\begin{lstlisting}
In one run, the dealer encodes $s$ as $D_s$, applies an allowed shuffle, and
distributes one card to each player. When all eight players reveal their
cards, the verifier sees $g\star D_s$. The following corollary gives the
deck-level correctness required by this execution. Proposition~\ref{thm:pgl-correctness}
later lifts this equality to the executable interpreter. The recovery
threshold and coalition privacy are separate properties proved in
Section~\ref{sec:exact}.
\end{lstlisting}
\colorbox{RewriteGreen}{\parbox{0.97\textwidth}{\textbf{Rewrite suggestion.}\par\begingroup\small
In one run, the dealer encodes $s$ as $D_s$, applies an allowed shuffle, and gives one card to each player. When all eight players reveal their cards, the verifier sees $g\star D_s$. Proposition~\ref{thm:pgl-correctness} later lifts deck-level correctness to the executable interpreter. The recovery threshold asks how many revealed positions determine the secret. Coalition privacy asks what smaller named groups learn. Section~\ref{sec:exact} proves these two separate properties. Those later arguments depend first on deck-level correctness. The following corollary states the equality required by this execution.
\par\endgroup}}
\Field{Why this rewrite.}{The rewrite retains the complete execution, the later interpreter lift, and the separate recovery and privacy questions. It orders the roadmap by dependency and ends with the deck-level equality that the adjacent corollary proves.}
\Field{Terminology actions.}{Keep: deck-level correctness, interpreter lift | Explain: recovery threshold -> how many revealed positions determine the secret, coalition privacy -> what smaller named groups learn}
\par\smallskip
\needspace{10\baselineskip}\subsubsection*{W-05-STM007: analogous purpose}
\Field{Location.}{Section 05, statement, lines 1099--1111.}
\Field{Draft purpose.}{State correctness of decoding every shuffled encoded deck.}
\Field{Qualitative closeness.}{analogous purpose. Both units close a correctness argument for an arbitrary execution state before the text turns to privacy.}
\Field{Limits of the analogy.}{The baseline presents the conclusion in prose and begins transcript analysis. The WADT unit is a terse universal equality.}
\Field{Baseline narration used.}{adapted.}
\Field{Complexity and jargon.}{Complexity: \Tag{low} $\rightarrow$ \Tag{low}. Jargon: \Tag{low} $\rightarrow$ \Tag{low}.}
\textbf{Original WADT unit. Footnotes omitted.}
\begin{lstlisting}
For every $s\in\{0,1\}$ and $g\in G$,
\[
  \operatorname{dec}(g\star\operatorname{enc}(s))
  =\operatorname{dec}(\operatorname{enc}(s))
  =s.
\]
\end{lstlisting}
\colorbox{RewriteGreen}{\parbox{0.97\textwidth}{\textbf{Rewrite suggestion.}\par\begingroup\small
For every $s\in\{0,1\}$ and $g\in G$,
\[
  \operatorname{dec}(g\star\operatorname{enc}(s))
  =\operatorname{dec}(\operatorname{enc}(s))
  =s.
\]
\par\endgroup}}
\Field{Why this rewrite.}{Adapts the baseline's correctness close to the exact universal shuffle-correctness equality. The excluded footnote does not enter the suggestion.}
\Field{Terminology actions.}{Keep: shuffle correctness}
\par\smallskip
\clearpage\subsection{Baseline S-04-P005: initialization shifted-view reparameterization}
\colorbox{BaselineBlue}{\parbox{0.97\textwidth}{\Field{Purpose.}{Replace the input-shifted revealed value with an equivalent fresh uniform value.}\Field{Primary purpose.}{proof}\Field{Narration structure.}{relation\_instantiation $\rightarrow$ distributional\_equivalence $\rightarrow$ simulator\_definition}\Field{Reusable move.}{Isolate the random shift and restate the transcript using a fresh uniform variable.}}}
\needspace{10\baselineskip}\subsubsection*{W-05-P015: analogous purpose}
\Field{Location.}{Section 05, prose\_paragraph, lines 1085--1088.}
\Field{Draft purpose.}{Explain why the shuffled heart set stays in the same orbit.}
\Field{Qualitative closeness.}{analogous purpose. Both units reparameterize an input-dependent transformed object by an equivalent group or random parameter and then state the invariant consequence.}
\Field{Limits of the analogy.}{The baseline uses a fresh uniform value for a security transcript. The WADT unit uses the inverse group element for orbit membership.}
\Field{Baseline narration used.}{adapted.}
\Field{Complexity and jargon.}{Complexity: \Tag{medium} $\rightarrow$ \Tag{low}. Jargon: \Tag{medium} $\rightarrow$ \Tag{low}.}
\textbf{Original WADT unit. Footnotes omitted.}
\begin{lstlisting}
Because $\rho(g)^{-1}=\rho(g^{-1})$, the heart set of the shuffled deck
remains in the same $G$-orbit as $H(D)$. The shuffle may change the
representative in $\Omega_4$, but it cannot change the orbit label.
\end{lstlisting}
\colorbox{RewriteGreen}{\parbox{0.97\textwidth}{\textbf{Rewrite suggestion.}\par\begingroup\small
The equality $\rho(g)^{-1}=\rho(g^{-1})$ shows that the inverse permutation comes from the inverse group element. Thus the shuffled heart set remains in the same $G$-orbit as $H(D)$. The shuffle may change the orbit representative in $\Omega_4$, but the orbit label remains fixed.
\par\endgroup}}
\Field{Why this rewrite.}{Adapts the baseline's algebraic-relation sequence to identify the inverse group element and derive orbit-label invariance.}
\Field{Terminology actions.}{Keep: orbit representative, orbit label | Explain: inverse representation -> the inverse permutation comes from the inverse group element}
\par\smallskip
\clearpage\subsection{Baseline S-04-P007: initialization security and resource counts}
\colorbox{BaselineBlue}{\parbox{0.97\textwidth}{\Field{Purpose.}{Conclude secure realization and assess card and random-operation costs.}\Field{Primary purpose.}{result}\Field{Narration structure.}{security\_conclusion $\rightarrow$ efficiency\_assessment $\rightarrow$ resource\_breakdown}\Field{Reusable move.}{Pair the security conclusion with concrete resource counts before changing protocols.}}}
\needspace{10\baselineskip}\subsubsection*{W-01-STM001: analogous purpose}
\Field{Location.}{Section 01, statement, lines 108--114.}
\Field{Draft purpose.}{State the exact ideal-shuffle guarantees of the main instance.}
\Field{Qualitative closeness.}{analogous purpose. Both units consolidate correctness and security results into a concise result statement after the supporting development.}
\Field{Limits of the analogy.}{The baseline concludes one initialization protocol and reports resources. The WADT statement combines universal recovery, a recovery ramp, and two privacy notions under a larger assumption set.}
\Field{Baseline narration used.}{adapted.}
\Field{Complexity and jargon.}{Complexity: \Tag{high} $\rightarrow$ \Tag{medium}. Jargon: \Tag{high} $\rightarrow$ \Tag{medium}.}
\textbf{Original WADT unit. Footnotes omitted.}
\begin{lstlisting}
The $\PG$ instance recovers the secret from all eight endpoints for every
group element. Under the uniform shuffle distribution, with a uniform
Boolean secret and the fixed representative dealer, it has recovery
parameters $(t,r,n)=(3,7,8)$ and perfect view and executed-trace privacy
against every passive coalition of at most three
players.
\end{lstlisting}
\colorbox{RewriteGreen}{\parbox{0.97\textwidth}{\textbf{Rewrite suggestion.}\par\begingroup\small
For every group element, the $\PG$ instance recovers the secret from all eight endpoints. Under the uniform shuffle distribution, a uniform Boolean secret, and the fixed representative dealer that always encodes each secret with its chosen deck, the instance has recovery parameters $(t,r,n)=(3,7,8)$. It also has perfect view and executed-trace privacy against every passive coalition of at most three players.
\par\endgroup}}
\Field{Why this rewrite.}{The rewrite adapts the baseline's result-and-scope sequence. It gives universal recovery first, explains the fixed representative dealer, and then states the shared assumptions and privacy bound.}
\Field{Terminology actions.}{Keep: passive coalition, recovery parameters, executed-trace privacy | Explain: fixed representative dealer, uniform Boolean secret}
\par\smallskip
\needspace{10\baselineskip}\subsubsection*{W-01-STM002: analogous purpose}
\Field{Location.}{Section 01, statement, lines 123--130.}
\Field{Draft purpose.}{State the finite word-shuffle guarantees of the main instance.}
\Field{Qualitative closeness.}{analogous purpose. Both units close a technical development with a compact guarantee and its operational consequence.}
\Field{Limits of the analogy.}{The baseline reports exact security and resource use. The WADT statement reports an approximation bound, transported view and trace bounds, and wordwise correctness.}
\Field{Baseline narration used.}{adapted.}
\Field{Complexity and jargon.}{Complexity: \Tag{high} $\rightarrow$ \Tag{medium}. Jargon: \Tag{high} $\rightarrow$ \Tag{medium}.}
\textbf{Original WADT unit. Footnotes omitted.}
\begin{lstlisting}
The 200-letter word shuffle over the five generator letters is within
$L_1$ distance $2^{-40}$ of the uniform shuffle distribution. For
the fixed representative dealer, the coalition view distributions at the two
secrets are within $L_1$ distance $2^{-39}$ for every passive
coalition of at most three players, and the executed coalition traces
satisfy the same bound. Executed correctness holds for every
word.
\end{lstlisting}
\colorbox{RewriteGreen}{\parbox{0.97\textwidth}{\textbf{Rewrite suggestion.}\par\begingroup\small
The 200-letter word shuffle over five generator letters is within $L_1$ distance $2^{-40}$ of the uniform shuffle distribution. The fixed representative dealer uses one chosen deck for each secret. For every passive coalition of at most three players, the two secret-conditioned view distributions are within $L_1$ distance $2^{-39}$ of each other. The corresponding executed coalition-trace distributions satisfy the same bound. Executed correctness holds for every word.
\par\endgroup}}
\Field{Why this rewrite.}{The rewrite adapts the baseline's bounded result sequence. It explains the dealer, explicitly compares the two secret-conditioned distributions, and separates the shuffle, view, trace, and exact-correctness claims while preserving both bounds and every condition.}
\Field{Terminology actions.}{Keep: L1 distance, word shuffle, executed correctness | Explain: distance between the two secrets, fixed representative dealer}
\par\smallskip
\needspace{10\baselineskip}\subsubsection*{W-04-P004: partial analogue}
\Field{Location.}{Section 04, prose\_paragraph, lines 698--702.}
\Field{Draft purpose.}{Interpret the zero-bias boundary of the witness bound.}
\Field{Qualitative closeness.}{partial analogue. The baseline supplies the closing move from a proved guarantee to a concrete protocol consequence.}
\Field{Limits of the analogy.}{The baseline closes a security proof and reports resource costs. This unit interprets only one boundary parameter.}
\Field{Baseline narration used.}{partial.}
\Field{Complexity and jargon.}{Complexity: \Tag{low} $\rightarrow$ \Tag{low}. Jargon: \Tag{medium} $\rightarrow$ \Tag{low}.}
\textbf{Original WADT unit. Footnotes omitted.}
\begin{lstlisting}
At bias zero the witness bound collapses to zero for any positive word
length. The unbiased member is therefore den Boer's protocol in
precisely this sense.
\end{lstlisting}
\colorbox{RewriteGreen}{\parbox{0.97\textwidth}{\textbf{Rewrite suggestion.}\par\begingroup\small
At bias zero, the witness bound is zero for every positive word length. In this precise sense, the unbiased member is den Boer's protocol.
\par\endgroup}}
\Field{Why this rewrite.}{Uses the baseline's result-to-consequence move and states the exact condition that defines the comparison.}
\Field{Terminology actions.}{Keep: bias zero, positive word length | Explain: witness bound | Replace: in precisely this sense -> in the sense that its witness bound is zero for every positive word length}
\par\smallskip
\needspace{10\baselineskip}\subsubsection*{W-04-P005: partial analogue}
\Field{Location.}{Section 04, prose\_paragraph, lines 703--715.}
\Field{Draft purpose.}{Summarize the exact leakage theorem and explain the leakage ramp.}
\Field{Qualitative closeness.}{partial analogue. The baseline supplies the move of closing a formal result with an interpretation of its measured consequences.}
\Field{Limits of the analogy.}{The baseline reports security and operational cost. This unit reports an information curve, a closed form, and terminology for its maximum.}
\Field{Baseline narration used.}{partial.}
\Field{Complexity and jargon.}{Complexity: \Tag{medium} $\rightarrow$ \Tag{low}. Jargon: \Tag{high} $\rightarrow$ \Tag{medium}.}
\textbf{Original WADT unit. Footnotes omitted.}
\begin{lstlisting}
What I call the master theorem gives the exact leakage for every fixed reveal set. One revealed card carries no information about the conjunction, which
is the information-theoretic counterpart of the scheme's privacy
threshold, and the ramp climbs to the secret's own entropy
$2-\tfrac34\log 3\approx 0.811$, which I call the cap. The decimals are
evaluations of the proven closed forms.
\end{lstlisting}
\colorbox{RewriteGreen}{\parbox{0.97\textwidth}{\textbf{Rewrite suggestion.}\par\begingroup\small
The reveal-set leakage theorem gives the exact leakage for every fixed reveal set. One revealed card carries no information about the conjunction. This is the information-theoretic counterpart of the scheme's privacy threshold. The leakage curve reaches the secret's entropy, $2-\tfrac34\log 3\approx 0.811$. I call this maximum leakage the cap. The reported decimals evaluate the proven closed forms.
\par\endgroup}}
\Field{Why this rewrite.}{Uses the baseline's quantitative-result interpretation. It separates the theorem scope, zero-leakage endpoint, entropy maximum, terminology, and closed-form evidence while omitting the footnote.}
\Field{Terminology actions.}{Keep: entropy, fixed reveal set, closed form | Explain: the ramp, privacy threshold | Replace: master theorem -> reveal-set leakage theorem, the cap -> the maximum leakage}
\par\smallskip
\needspace{10\baselineskip}\subsubsection*{W-06-STM001: analogous purpose}
\Field{Location.}{Section 06, statement, lines 1294--1301.}
\Field{Draft purpose.}{State universal endpoint recovery for the main instance.}
\Field{Qualitative closeness.}{analogous purpose. Both units state the correctness conclusion after an execution argument.}
\Field{Limits of the analogy.}{The baseline concludes a concrete protocol and also reports resources. The WADT proposition is a universal recovery statement over secrets and group elements.}
\Field{Baseline narration used.}{adapted.}
\Field{Complexity and jargon.}{Complexity: \Tag{low} $\rightarrow$ \Tag{low}. Jargon: \Tag{medium} $\rightarrow$ \Tag{low}.}
\textbf{Original WADT unit. Footnotes omitted.}
\begin{lstlisting}
\par
For every $s\in\{0,1\}$ and $g\in\PG$, decoding all eight verifier endpoints
after the shuffle $g$ returns $s$.
\end{lstlisting}
\colorbox{RewriteGreen}{\parbox{0.97\textwidth}{\textbf{Rewrite suggestion.}\par\begingroup\small
For every $s\in\{0,1\}$ and $g\in\PG$, decoding all eight verifier endpoints after the shuffle $g$ returns $s$.
\par\endgroup}}
\Field{Why this rewrite.}{The rewrite adapts the baseline's concise result sequence. It excludes the footnote and preserves both universal quantifiers and the full-endpoint recovery claim.}
\Field{Terminology actions.}{Keep: all eight endpoints, group element}
\par\smallskip
\needspace{10\baselineskip}\subsubsection*{W-06-STM003: partial analogue}
\Field{Location.}{Section 06, statement, lines 1408--1421.}
\Field{Draft purpose.}{State perfect view privacy and executed-trace privacy for the fixed-dealer ideal experiment.}
\Field{Qualitative closeness.}{partial analogue. S-04-P007 supplies the limited result-closing move of stating security after the supporting conditions have been fixed.}
\Field{Limits of the analogy.}{The baseline concludes one initialization protocol and adds resource counts. This WADT proposition states two privacy equations for a fixed representative dealer and a bounded coalition, without an efficiency assessment.}
\Field{Baseline narration used.}{partial.}
\Field{Complexity and jargon.}{Complexity: \Tag{medium} $\rightarrow$ \Tag{medium}. Jargon: \Tag{high} $\rightarrow$ \Tag{medium}.}
\textbf{Original WADT unit. Footnotes omitted.}
\begin{lstlisting}
Under the uniform Boolean prior, the fixed representative dealer, and an
independent uniform $\PG$ shuffle,
\[
  S\mathrel{\perp}V_C
  \quad\text{and}\quad
  H(S\mid T_C)=H(S)
  \qquad\text{for every }\lvert C\rvert\leq3.
\]
\end{lstlisting}
\colorbox{RewriteGreen}{\parbox{0.97\textwidth}{\textbf{Rewrite suggestion.}\par\begingroup\small
Assume the uniform Boolean prior, the fixed representative dealer that uses one chosen deck for each secret, and an independent uniform $\PG$ shuffle. Then
\[
  S\mathrel{\perp}V_C
  \quad\text{and}\quad
  H(S\mid T_C)=H(S)
  \qquad\text{for every }\lvert C\rvert\leq3.
\]
\par\endgroup}}
\Field{Why this rewrite.}{The rewrite uses the partial analogue's experiment-before-conclusions move. It excludes the footnote, explains the dealer, and preserves both privacy formulas and the universal coalition bound.}
\Field{Terminology actions.}{Keep: view independence, conditional entropy, coalition bound | Explain: fixed representative dealer}
\par\smallskip
\needspace{10\baselineskip}\subsubsection*{W-06-STM005: partial analogue}
\Field{Location.}{Section 06, statement, lines 1444--1457.}
\Field{Draft purpose.}{State perfect view privacy and executed-trace privacy for the all-decks dealer.}
\Field{Qualitative closeness.}{partial analogue. S-04-P007 supplies the limited move of closing an analysis with its established security result.}
\Field{Limits of the analogy.}{The baseline concerns one initialization protocol and resource counts. This proposition states view privacy and executed-trace privacy for the all-decks dealer under a uniform shuffle and coalition bound.}
\Field{Baseline narration used.}{partial.}
\Field{Complexity and jargon.}{Complexity: \Tag{medium} $\rightarrow$ \Tag{medium}. Jargon: \Tag{high} $\rightarrow$ \Tag{medium}.}
\textbf{Original WADT unit. Footnotes omitted.}
\begin{lstlisting}
Under the uniform Boolean prior, the uniform class-conditional deck dealer,
and an independent uniform $\PG$ shuffle,
\[
  S\mathrel{\perp}V_C
  \quad\text{and}\quad
  H(S\mid T_C)=H(S)
  \qquad\text{for every }\lvert C\rvert\leq3.
\]
\end{lstlisting}
\colorbox{RewriteGreen}{\parbox{0.97\textwidth}{\textbf{Rewrite suggestion.}\par\begingroup\small
Assume the uniform Boolean prior, a dealer that samples a deck uniformly from the chosen secret class, and an independent uniform $\PG$ shuffle. Then
\[
  S\mathrel{\perp}V_C
  \quad\text{and}\quad
  H(S\mid T_C)=H(S)
  \qquad\text{for every }\lvert C\rvert\leq3.
\]
\par\endgroup}}
\Field{Why this rewrite.}{The rewrite uses the partial analogue's experiment-before-conclusions move. It excludes the footnote, explains class-conditional sampling, and preserves both privacy formulas and the coalition scope.}
\Field{Terminology actions.}{Keep: uniform class-conditional deck, view independence, executed-trace privacy | Explain: class-conditional}
\par\smallskip
\needspace{10\baselineskip}\subsubsection*{W-06-STM006: partial analogue}
\Field{Location.}{Section 06, statement, lines 1466--1474.}
\Field{Draft purpose.}{State the exact privacy guarantees of the shuffle-free dealer.}
\Field{Qualitative closeness.}{partial analogue. S-04-P007 supplies the limited move of presenting the security conclusion after its conditions.}
\Field{Limits of the analogy.}{The baseline gives one security conclusion with efficiency data. This proposition separates a prior-generic view theorem from a uniform-prior executed-trace privacy theorem for a shuffle-free class-conditional deck.}
\Field{Baseline narration used.}{partial.}
\Field{Complexity and jargon.}{Complexity: \Tag{medium} $\rightarrow$ \Tag{medium}. Jargon: \Tag{high} $\rightarrow$ \Tag{medium}.}
\textbf{Original WADT unit. Footnotes omitted.}
\begin{lstlisting}
For any Boolean prior, the uniform class-conditional deck gives
$S\mathrel{\perp}V_C$ whenever $\lvert C\rvert\leq3$. Under the uniform
Boolean prior, it also gives $H(S\mid T_C)=H(S)$.
\end{lstlisting}
\colorbox{RewriteGreen}{\parbox{0.97\textwidth}{\textbf{Rewrite suggestion.}\par\begingroup\small
For any Boolean prior, a deck sampled uniformly from the chosen secret class gives
$S\mathrel{\perp}V_C$ whenever $\lvert C\rvert\leq3$. Thus, the coalition view is independent of the secret. Under the uniform Boolean prior, the same dealer also gives $H(S\mid T_C)=H(S)$, so the executed trace preserves the secret's entropy.
\par\endgroup}}
\Field{Why this rewrite.}{The rewrite uses the partial analogue's general-result-before-restricted-result move. It excludes the footnote, explains class-conditional sampling, view independence, and conditional entropy, and preserves both prior scopes.}
\Field{Terminology actions.}{Keep: any Boolean prior, uniform Boolean prior, class-conditional deck, view independence, conditional entropy | Explain: class-conditional deck, view independence, conditional entropy}
\par\smallskip
\needspace{10\baselineskip}\subsubsection*{W-06-STM007: analogous purpose}
\Field{Location.}{Section 06, statement, lines 1483--1496.}
\Field{Draft purpose.}{State Theorem A's complete correctness, recovery, and privacy guarantees.}
\Field{Qualitative closeness.}{analogous purpose. Both units close a development with a concise correctness-and-security conclusion.}
\Field{Limits of the analogy.}{The baseline concerns one small initialization protocol and resources. The WADT theorem combines universal recovery, a recovery ramp, and two ideal-shuffle privacy equations.}
\Field{Baseline narration used.}{adapted.}
\Field{Complexity and jargon.}{Complexity: \Tag{high} $\rightarrow$ \Tag{medium}. Jargon: \Tag{high} $\rightarrow$ \Tag{medium}.}
\textbf{Original WADT unit. Footnotes omitted.}
\begin{lstlisting}
The $\PG$ instance recovers the secret from all eight endpoints for every
group element. Under the uniform Boolean prior, the fixed
representative dealer, and an independent uniform $\PG$ shuffle, the
instance has recovery parameters
\[
  (t,r,n)=(3,7,8),
\]
and for every passive coalition $C$ with $\lvert C\rvert\leq3$,
\[
  S\mathrel{\perp}V_C
  \quad\text{and}\quad
  H(S\mid T_C)=H(S).
\]
\end{lstlisting}
\colorbox{RewriteGreen}{\parbox{0.97\textwidth}{\textbf{Rewrite suggestion.}\par\begingroup\small
For every group element, the $\PG$ instance recovers the secret from all eight endpoints. Under the uniform Boolean prior, the fixed representative dealer that uses one chosen deck for each secret, and an independent uniform $\PG$ shuffle, the instance has recovery parameters
\[
  (t,r,n)=(3,7,8),
\]
and for every passive coalition $C$ with $\lvert C\rvert\leq3$,
\[
  S\mathrel{\perp}V_C
  \quad\text{and}\quad
  H(S\mid T_C)=H(S).
\]
\par\endgroup}}
\Field{Why this rewrite.}{The rewrite adapts the baseline's result-and-scope sequence. It excludes the footnote mark, explains the dealer and endpoint scope, and preserves universal recovery, the recovery tuple, both privacy formulas, and the passive-coalition bound.}
\Field{Terminology actions.}{Keep: recovery parameters, passive coalition, conditional entropy | Explain: fixed representative dealer, all eight endpoints}
\par\smallskip
\needspace{10\baselineskip}\subsubsection*{W-07-STM004: analogous purpose}
\Field{Location.}{Section 07, statement, lines 1616--1634.}
\Field{Draft purpose.}{State Theorem B's finite-word mixing, privacy, trace, and correctness guarantees.}
\Field{Qualitative closeness.}{analogous purpose. Both units consolidate a formal security conclusion and its conditions after the preceding proof development.}
\Field{Limits of the analogy.}{The baseline closes one protocol proof and reports resources. This theorem combines mixing, approximate privacy, trace privacy, and correctness.}
\Field{Baseline narration used.}{adapted.}
\Field{Complexity and jargon.}{Complexity: \Tag{high} $\rightarrow$ \Tag{high}. Jargon: \Tag{high} $\rightarrow$ \Tag{high}.}
\textbf{Original WADT unit. Footnotes omitted.}
\begin{lstlisting}
Let $\mu$ be uniform on the five-letter symmetric generator
tuple. Then
\[
  \left\lVert \mu^{*200}-U_G\right\rVert_1 \leq 2^{-40}.
\]
For the fixed representative dealer, every passive coalition $C$ with
$\lvert C\rvert\leq3$, and all secrets $s,s'\in\{0,1\}$, where the coalition
view at dealt secret $s$ and shuffle $g$ is denoted $V_C(s,g)$,
\begin{equation}
  \left\lVert
    (g\mapsto V_C(s,g))_\#\mu^{*200}
    -(g\mapsto V_C(s',g))_\#\mu^{*200}
  \right\rVert_1
  \leq 2^{-39},
  \label{eq:word-view-privacy}
\end{equation}
and the executed coalition trace satisfies the same bound. Decoding all
eight endpoints returns the dealt secret for every word.
\end{lstlisting}
\colorbox{RewriteGreen}{\parbox{0.97\textwidth}{\textbf{Rewrite suggestion.}\par\begingroup\small
Let $\mu$ be uniform on the five-letter symmetric generator tuple. Then
\[
  \left\lVert \mu^{*200}-U_G\right\rVert_1 \leq 2^{-40}.
\]
For the fixed representative dealer, every passive coalition $C$ with $\lvert C\rvert\leq3$, and all secrets $s,s'\in\{0,1\}$, let $V_C(s,g)$ denote the coalition view at dealt secret $s$ and shuffle $g$. Then
\begin{equation}
  \left\lVert
    (g\mapsto V_C(s,g))_\#\mu^{*200}
    -(g\mapsto V_C(s',g))_\#\mu^{*200}
  \right\rVert_1
  \leq 2^{-39}.
  \label{eq:word-view-privacy}
\end{equation}
The executed coalition trace satisfies the same bound. Decoding all eight endpoints returns the dealt secret for every word.
\par\endgroup}}
\Field{Why this rewrite.}{Adapts the baseline's result consolidation. It keeps the mixing premise, dealer scope, coalition and secret quantifiers, view definition, privacy bounds, and wordwise correctness in terse mathematical form.}
\Field{Terminology actions.}{Keep: fixed representative dealer, coalition view, executed trace | Explain: for every word}
\par\smallskip
\needspace{10\baselineskip}\subsubsection*{W-10-P002: analogous purpose}
\Field{Location.}{Section 10, prose\_paragraph, lines 1881--1887.}
\Field{Draft purpose.}{Summarize Theorem B's finite-shuffle guarantees and interpret their cost.}
\Field{Qualitative closeness.}{analogous purpose. Both units conclude a formal security development and interpret the resulting quantitative cost.}
\Field{Limits of the analogy.}{The baseline reports protocol card and operation costs. This unit reports approximation and privacy distances.}
\Field{Baseline narration used.}{adapted.}
\Field{Complexity and jargon.}{Complexity: \Tag{medium} $\rightarrow$ \Tag{low}. Jargon: \Tag{high} $\rightarrow$ \Tag{medium}.}
\textbf{Original WADT unit. Footnotes omitted.}
\begin{lstlisting}
For a finite shuffle implementation, Theorem B gives an $L_1$
distance of at most $2^{-40}$ after 200 generator letters, the same bound
for every single-card endpoint and for a product with any Boolean secret
prior, and approximate coalition-view and executed-trace bounds at
$2^{-39}$. The word-shuffle model hence replaces the ideal uniform
shuffle at a quantified and machine-checked cost.
\end{lstlisting}
\colorbox{RewriteGreen}{\parbox{0.97\textwidth}{\textbf{Rewrite suggestion.}\par\begingroup\small
For a finite shuffle implementation, Theorem B gives an $L_1$ distance of at most $2^{-40}$ after 200 generator letters. The same bound holds for every single-card endpoint and for the product with any Boolean secret prior. The approximate coalition-view and executed-trace bounds are $2^{-39}$. The word-shuffle model therefore replaces the ideal uniform shuffle at a quantified, machine-checked cost.
\par\endgroup}}
\Field{Why this rewrite.}{Adapts the baseline's result-and-cost sequence and gives each numerical guarantee and its scope separately.}
\Field{Terminology actions.}{Keep: L1 distance, coalition-view and executed-trace bounds | Explain: product with any Boolean secret prior}
\par\smallskip
\clearpage\subsection{Baseline S-04-P011: addition algebra and transcript setup}
\colorbox{BaselineBlue}{\parbox{0.97\textwidth}{\Field{Purpose.}{Derive the sum encoding and introduce the observed addition transcript.}\Field{Primary purpose.}{proof}\Field{Narration structure.}{premise\_instantiation $\rightarrow$ algebraic\_simplification $\rightarrow$ correctness\_close $\rightarrow$ security\_setup $\rightarrow$ view\_instantiation}\Field{Reusable move.}{Make mask cancellation explicit, then reuse the same random variable to motivate security.}}}
\needspace{10\baselineskip}\subsubsection*{W-05-STM006: analogous purpose}
\Field{Location.}{Section 05, statement, lines 1067--1083.}
\Field{Draft purpose.}{State heart-set transport, deck validity, and decoder invariance under shuffling.}
\Field{Qualitative closeness.}{analogous purpose. Both units derive an algebraic transformation identity and use it to establish the property needed by the next execution claim.}
\Field{Limits of the analogy.}{The baseline derives an encoded sum and opens a transcript argument. The WADT lemma derives set transport and two invariants.}
\Field{Baseline narration used.}{adapted.}
\Field{Complexity and jargon.}{Complexity: \Tag{high} $\rightarrow$ \Tag{medium}. Jargon: \Tag{high} $\rightarrow$ \Tag{medium}.}
\textbf{Original WADT unit. Footnotes omitted.}
\begin{lstlisting}
For every $g\in G$ and valid deck $D$,
\[
  H(g\star D)=\rho(g)^{-1}\mathbin{\cdot}H(D).
\]
The shuffled deck also satisfies
\[
  g\star D\text{ is valid},
  \qquad
  \operatorname{dec}(g\star D)=\operatorname{dec}(D).
\]
\end{lstlisting}
\colorbox{RewriteGreen}{\parbox{0.97\textwidth}{\textbf{Rewrite suggestion.}\par\begingroup\small
For every $g\in G$ and valid deck $D$, the heart positions move by the inverse positional action:
\[
  H(g\star D)=\rho(g)^{-1}\mathbin{\cdot}H(D).
\]
The shuffled deck remains valid, and its decoded value does not change:
\[
  g\star D\text{ is valid},
  \qquad
  \operatorname{dec}(g\star D)=\operatorname{dec}(D).
\]
\par\endgroup}}
\Field{Why this rewrite.}{Adapts the baseline's relation-to-consequence sequence while keeping the lemma terse and mathematical. The excluded footnote does not enter the suggestion.}
\Field{Terminology actions.}{Keep: decoder invariance, valid deck | Explain: inverse action on the heart set -> the heart positions move by the inverse positional action}
\par\smallskip
\needspace{10\baselineskip}\subsubsection*{W-05-P022: partial analogue}
\Field{Location.}{Section 05, prose\_paragraph, lines 1237--1247.}
\Field{Draft purpose.}{State sharp three-transitivity and explain its four main consequences.}
\Field{Qualitative closeness.}{partial analogue. The baseline supplies the move of stating one algebraic fact before deriving the property needed by the next security argument.}
\Field{Limits of the analogy.}{The WADT paragraph draws several independent design and security consequences from sharp transitivity. The baseline does not support that full fan-out.}
\Field{Baseline narration used.}{partial.}
\Field{Complexity and jargon.}{Complexity: \Tag{high} $\rightarrow$ \Tag{medium}. Jargon: \Tag{high} $\rightarrow$ \Tag{medium}.}
\textbf{Original WADT unit. Footnotes omitted.}
\begin{lstlisting}
The organizing fact of this instance is that the action of $\PG$ on the
eight projective points is sharply three-transitive. Any ordered triple of
distinct points is carried to any other by exactly one group element,
because the order $336=8\cdot7\cdot6$ equals the number of ordered triples.
Under a uniform shuffle, three revealed positions therefore carry no
invariant at all, and four positions carry exactly one, the cross ratio.
This single fact organizes the instance: it makes the cross-ratio decoder
exist, caps perfect privacy at coalitions of size three, forces
four-position coalitions to leak, and is the reason the construction uses
$\PG$.
\end{lstlisting}
\colorbox{RewriteGreen}{\parbox{0.97\textwidth}{\textbf{Rewrite suggestion.}\par\begingroup\small
The action of $\PG$ on the eight projective points is sharply three-transitive. Thus exactly one group element sends any ordered triple of distinct points to any other. The uniqueness follows because the group order is $336=8\cdot7\cdot6$, which is the number of ordered triples. Under a uniform shuffle, every ordered view of three positions is uniform and carries no group invariant. A view of four positions carries one invariant, the cross ratio. This algebraic fact has four consequences. It makes the cross-ratio decoder possible. It gives perfect privacy to coalitions of at most three players. It forces coalitions of four players to leak information. It also explains the choice of $\PG$ for the construction.
\par\endgroup}}
\Field{Why this rewrite.}{Uses only the baseline move that states the controlling algebraic fact before deriving the WADT decoder, privacy, leakage, and design consequences.}
\Field{Terminology actions.}{Keep: sharply three-transitive, cross ratio | Explain: carries no invariant -> every ordered three-position view is uniform, privacy cutoff -> perfect privacy holds through three positions while four positions must leak}
\par\smallskip
\clearpage\subsection{Baseline S-04-P014: addition simulator independence and view equality}
\colorbox{BaselineBlue}{\parbox{0.97\textwidth}{\Field{Purpose.}{State that the simulated addition transcript is input independent and lead to view equality.}\Field{Primary purpose.}{proof}\Field{Narration structure.}{randomness\_qualification $\rightarrow$ independence\_claim $\rightarrow$ equality\_handoff}\Field{Reusable move.}{Follow a simulator definition with its independence property and formal consequence.}}}
\needspace{10\baselineskip}\subsubsection*{W-07-STM003: partial analogue}
\Field{Location.}{Section 07, statement, lines 1590--1602.}
\Field{Draft purpose.}{State that multiplying both shuffle laws by the same secret prior preserves the bound.}
\Field{Qualitative closeness.}{partial analogue. The baseline supplies the move of defining an input-free comparison distribution and then stating equality of observations.}
\Field{Limits of the analogy.}{The baseline proves exact equality of simulated views. This statement preserves an approximate bound under a product construction.}
\Field{Baseline narration used.}{partial.}
\Field{Complexity and jargon.}{Complexity: \Tag{medium} $\rightarrow$ \Tag{low}. Jargon: \Tag{high} $\rightarrow$ \Tag{medium}.}
\textbf{Original WADT unit. Footnotes omitted.}
\begin{lstlisting}
Let $\pi$ be any distribution on Boolean secrets. Independence is built
into the product distributions, and
\begin{equation}
  \left\lVert
    \pi\mathbin{\times}\mu^{*200}
    -\pi\mathbin{\times}U_G
  \right\rVert_1
  \leq 2^{-40}.
  \label{eq:product-transfer}
\end{equation}
\end{lstlisting}
\colorbox{RewriteGreen}{\parbox{0.97\textwidth}{\textbf{Rewrite suggestion.}\par\begingroup\small
Let $\pi$ be any distribution on Boolean secrets. The product distributions make the secret and shuffle independent, and
\begin{equation}
  \left\lVert
    \pi\mathbin{\times}\mu^{*200}
    -\pi\mathbin{\times}U_G
  \right\rVert_1
  \leq 2^{-40}.
  \label{eq:product-transfer}
\end{equation}
\par\endgroup}}
\Field{Why this rewrite.}{Uses only the baseline's independence-to-formal-consequence move and preserves the arbitrary prior, both product laws, label, and exact bound.}
\Field{Terminology actions.}{Keep: Boolean prior, independent product distributions | Explain: pi times mu}
\par\smallskip
\needspace{10\baselineskip}\subsubsection*{W-07-P009: analogous purpose}
\Field{Location.}{Section 07, prose\_paragraph, lines 1642--1651.}
\Field{Draft purpose.}{Explain the proof of the word-shuffle privacy bound and state a companion joint-distribution bound.}
\Field{Qualitative closeness.}{analogous purpose. Both units construct an input-independent intermediate distribution and use it to derive equality or proximity of views.}
\Field{Limits of the analogy.}{The baseline gives an exact simulator equality. This unit uses two approximate triangle legs and then states a companion joint-law bound.}
\Field{Baseline narration used.}{adapted.}
\Field{Complexity and jargon.}{Complexity: \Tag{high} $\rightarrow$ \Tag{medium}. Jargon: \Tag{high} $\rightarrow$ \Tag{medium}.}
\textbf{Original WADT unit. Footnotes omitted.}
\begin{lstlisting}
The proof is a triangle inequality through the uniform group distribution. The
coalition-view distribution under the uniform shuffle does not depend on the secret,
by three-transitivity. The data-processing inequality applies the
certificate in Equation~\ref{eq:pgl-mixing} on each side of the triangle. A
companion theorem bounds the joint view-and-secret distribution under the word
shuffle against the product of its uniform-shuffle marginals by $2^{-40}$ at
any Boolean prior.
\end{lstlisting}
\colorbox{RewriteGreen}{\parbox{0.97\textwidth}{\textbf{Rewrite suggestion.}\par\begingroup\small
The proof uses a triangle inequality through the uniform group distribution. By three-transitivity, the coalition-view distribution under the uniform shuffle does not depend on the secret. The data-processing inequality applies the certificate in Equation~\ref{eq:pgl-mixing} to both sides of the triangle. A companion theorem also gives a bound of $2^{-40}$ for any Boolean prior. It compares the joint view-and-secret distribution under the word shuffle with the product of its uniform-shuffle marginals.
\par\endgroup}}
\Field{Why this rewrite.}{Adapts the baseline's independence-to-consequence proof sequence to two approximate triangle legs and the companion joint-law bound.}
\Field{Terminology actions.}{Keep: triangle inequality, three-transitivity | Explain: product of its uniform-shuffle marginals}
\par\smallskip
\clearpage\subsection{Baseline S-04-P016: addition card optimality and shuffle count}
\colorbox{BaselineBlue}{\parbox{0.97\textwidth}{\Field{Purpose.}{Report the two-card lower bound and the single randomized operation.}\Field{Primary purpose.}{result}\Field{Narration structure.}{efficiency\_assessment $\rightarrow$ lower\_bound\_reason $\rightarrow$ resource\_breakdown}\Field{Reusable move.}{Support an optimality statement with the input-count lower bound before listing randomness cost.}}}
\needspace{10\baselineskip}\subsubsection*{W-04-STM001: partial analogue}
\Field{Location.}{Section 04, statement, lines 785--792.}
\Field{Draft purpose.}{State the quantitative two-card leakage result and its transition to shape independence.}
\Field{Qualitative closeness.}{partial analogue. The baseline contributes only the move of reporting exact quantities together with the cases that determine them.}
\Field{Limits of the analogy.}{The baseline assesses card minimality and randomness cost. This statement compares two leakage values, gives one closed form, and states when reveal geometry stops affecting leakage.}
\Field{Baseline narration used.}{partial.}
\Field{Complexity and jargon.}{Complexity: \Tag{high} $\rightarrow$ \Tag{medium}. Jargon: \Tag{high} $\rightarrow$ \Tag{medium}.}
\textbf{Original WADT unit. Footnotes omitted.}
\begin{lstlisting}
Let the reveal expose two of the five positions. Two adjacent positions
reveal $0.154$ bits about the conjunction, and two positions at cyclic
distance two reveal $0.119$ bits. Both decimals evaluate proven closed
forms, for instance the adjacent value is
$\tfrac{27}{10}-\tfrac14\log 5-\tfrac{7}{10}\log 7$. The shape stops
mattering at three revealed cards, because the closed form there depends
only on how many cards the reveal exposes.
\end{lstlisting}
\colorbox{RewriteGreen}{\parbox{0.97\textwidth}{\textbf{Rewrite suggestion.}\par\begingroup\small
Let the reveal expose two of the five positions. Two adjacent positions reveal $0.154$ bits about the conjunction. Two positions at cyclic distance two reveal $0.119$ bits. Both decimals evaluate proven closed forms. In particular, the adjacent value is $\tfrac{27}{10}-\tfrac14\log 5-\tfrac{7}{10}\log 7$. For three revealed cards, the relative positions no longer affect leakage because the closed form depends only on the number of revealed cards.
\par\endgroup}}
\Field{Why this rewrite.}{Uses only the baseline's exact-quantity move. The statement gives each geometric case and value separately, then states the three-card threshold.}
\Field{Terminology actions.}{Keep: cyclic distance two, closed form | Explain: shape independence at three cards | Replace: The shape stops mattering -> The relative positions no longer affect leakage}
\par\smallskip
\clearpage\subsection{Baseline S-04-P024: sign-normalization security and resource counts}
\colorbox{BaselineBlue}{\parbox{0.97\textwidth}{\Field{Purpose.}{Conclude security and count the minimal card use and two-sided shuffle cost.}\Field{Primary purpose.}{result}\Field{Narration structure.}{security\_conclusion $\rightarrow$ efficiency\_assessment $\rightarrow$ resource\_breakdown}\Field{Reusable move.}{End the protocol analysis by pairing confidentiality with its exact card and randomness costs.}}}
\needspace{10\baselineskip}\subsubsection*{W-06-P011: analogous purpose}
\Field{Location.}{Section 06, prose\_paragraph, lines 1476--1481.}
\Field{Draft purpose.}{Synthesize the ingredients of the first headline theorem.}
\Field{Qualitative closeness.}{analogous purpose. Both paragraphs collect preceding correctness and security results into a final theorem-level conclusion.}
\Field{Limits of the analogy.}{The baseline closes one protocol with security and resource counts. The WADT paragraph combines executed correctness, recovery thresholds, a leakage witness, and two privacy notions.}
\Field{Baseline narration used.}{adapted.}
\Field{Complexity and jargon.}{Complexity: \Tag{medium} $\rightarrow$ \Tag{medium}. Jargon: \Tag{high} $\rightarrow$ \Tag{medium}.}
\textbf{Original WADT unit. Footnotes omitted.}
\begin{lstlisting}
Executed correctness supplies full-run decoding. The recovery-ramp
proposition supplies the recovery threshold and the four-card leakage
witness. Fixed-dealer privacy supplies view and trace privacy through size
three, which makes the four-card witness sharp. Together these results give
the first headline theorem.
\end{lstlisting}
\colorbox{RewriteGreen}{\parbox{0.97\textwidth}{\textbf{Rewrite suggestion.}\par\begingroup\small
Executed correctness gives decoding for the full run. The recovery-ramp proposition gives the recovery threshold and the four-card leakage witness. Fixed-dealer privacy gives view and trace privacy through coalition size three. This makes the four-card witness sharp. Together, these results prove Theorem A.
\par\endgroup}}
\Field{Why this rewrite.}{The rewrite adapts the baseline's result-synthesis sequence. It assigns each conclusion to its source result, replaces the vague headline phrase with Theorem A, and preserves the sharpness inference.}
\Field{Terminology actions.}{Keep: executed correctness, four-card leakage witness, fixed-dealer privacy | Explain: first headline theorem | Replace: headline theorem -> Theorem A}
\par\smallskip
\clearpage\subsection{Baseline S-05-P001: dihedral-card contributions and future object designs}
\colorbox{BaselineBlue}{\parbox{0.97\textwidth}{\Field{Purpose.}{Summarize the new card model and efficient predicate protocols, then propose extending invisible-ink designs to other objects.}\Field{Primary purpose.}{conclusion}\Field{Narration structure.}{contribution\_recap $\rightarrow$ interpretive\_claim $\rightarrow$ future\_work $\rightarrow$ illustrative\_extension}\Field{Reusable move.}{Move from a concise contribution recap to one future direction grounded in the paper's design principle.}}}
\needspace{10\baselineskip}\subsubsection*{W-03-P002: partial analogue}
\Field{Location.}{Section 03, prose\_paragraph, lines 350--365.}
\Field{Draft purpose.}{Interpret the profile record as a loose specification, explain the implementation relation, justify its name, and delimit future instances.}
\Field{Qualitative closeness.}{partial analogue. Both units end by stating a grounded future direction and illustrating it with concrete candidate extensions.}
\Field{Limits of the analogy.}{Only the future direction and its examples align. The baseline moves there after a contribution recap. This paragraph mainly defines and interprets a profile specification, explains a refinement relation, and justifies terminology.}
\Field{Baseline narration used.}{partial.}
\Field{Complexity and jargon.}{Complexity: \Tag{high} $\rightarrow$ \Tag{medium}. Jargon: \Tag{high} $\rightarrow$ \Tag{medium}.}
\textbf{Original WADT unit. Footnotes omitted.}
\begin{lstlisting}
In the vocabulary of algebraic specification, the record type is a
parameterized loose specification: its fields are the signature, its stated
obligations are the axioms, and an instance that supplies the fields and
discharges the obligations is a model of the
specification~\cite{GoguenBurstall1992,SannellaTarlecki2012}. The five
instances of this paper are five such models, and the relation between the
uniform-shuffle model and its word-shuffle implementation is an approximate
refinement whose distance the $L_1$ bounds of Section~\ref{sec:mixing}
quantify. The record's name borrows monodromy for the way one run
transports the deal along a group element, and prose outside code contexts
says profile record. The specification admits further models beyond the
five of this paper: the affine group $\mathrm{AGL}(1,7)$, which carries
any ordered pair of distinct positions to any other by exactly one
element, and $\mathrm{PGL}(2,11)$ at the present three-transitive
parameters are the natural next instances.
\end{lstlisting}
\colorbox{RewriteGreen}{\parbox{0.97\textwidth}{\textbf{Rewrite suggestion.}\par\begingroup\small
In algebraic specification, the profile record is a parameterized loose specification. Its fields form the signature, which lists the available data and operations. Its stated obligations form the axioms that the data must satisfy. An instance that supplies the fields and proves the obligations is a model of the specification~\cite{GoguenBurstall1992,SannellaTarlecki2012}. This interpretation explains the five instances in this paper. Each instance is a model of the same specification. It also explains the relation between the ideal and finite shuffle laws. The word-shuffle implementation approximately refines the uniform-shuffle model, and the $L_1$ bounds of Section~\ref{sec:mixing} quantify the distance. The formal name invokes monodromy because one run transports the deal along a group element. Outside code contexts, I use the term profile record. Because the specification is loose, it also admits models beyond the five in this paper. Natural next instances are the affine group $\mathrm{AGL}(1,7)$, which maps any ordered pair of distinct positions to any other by exactly one element, and $\mathrm{PGL}(2,11)$ with the present three-transitive parameters. With this interpretation in place, the following table maps the model data to their record roles.
\par\endgroup}}
\Field{Why this rewrite.}{The rewrite makes the specification interpretation the reason for the current-model claim, the refinement claim, and the two prospective models. It preserves both citations, the quantified $L_1$ refinement, the monodromy explanation, the prose term, and the exact affine-group property. The prospective instances illustrate the looseness of the specification instead of serving as an isolated future-work list. The final sentence prepares the bridge table.}
\Field{Terminology actions.}{Keep: loose specification, approximate refinement, monodromy | Explain: signature and axioms, AGL(1,7), PGL(2,11) | Replace: the record's name borrows monodromy -> The formal name invokes monodromy because one run transports the deal along a group element}
\par\smallskip
\needspace{10\baselineskip}\subsubsection*{W-09-P005: partial analogue}
\Field{Location.}{Section 09, prose\_paragraph, lines 1825--1845.}
\Field{Draft purpose.}{Compare exact physical group shuffles with the paper's approximate word shuffle and identify open physical realization and vocabulary costs.}
\Field{Qualitative closeness.}{partial analogue. The baseline supplies only the future-design move that turns a physical limitation into a concrete research direction.}
\Field{Limits of the analogy.}{The baseline concludes one card design and proposes new objects. This paragraph is a multi-source comparison with quantitative and vocabulary tradeoffs.}
\Field{Baseline narration used.}{partial.}
\Field{Complexity and jargon.}{Complexity: \Tag{high} $\rightarrow$ \Tag{medium}. Jargon: \Tag{high} $\rightarrow$ \Tag{medium}.}
\textbf{Original WADT unit. Footnotes omitted.}
\begin{lstlisting}
Group-based shuffle models appear in work on closed shuffles and graph
automorphisms. Miyamoto and Shinagawa realize a uniformly chosen graph
automorphism by pile-scramble shuffles~\cite{MiyamotoShinagawa2022}.
Shinagawa and Miyamoto extend the construction to graphs and hypergraphs and
give applications to cyclic group shuffles~\cite{ShinagawaMiyamoto2023}.
These results describe physical protocols that implement structured
uniform shuffle groups exactly, and the word approach of this paper trades
against them. A pile-scramble realization samples its group uniformly in
one physical procedure, while the word shuffle repeats elementary moves
and pays an $L_1$ error of $2^{-40}$ that the kernel certifies. Whether
the five $\PG$ letters admit unobserved physical realizations is the
question their line answers for structured shuffles and this paper leaves
open. The eight-card instance also pays in deck vocabulary against
two-symbol encodings: it stores a Boolean secret in the arrangement of
eight pairwise distinct cards, where committed formats store one bit in
two two-symbol cards~\cite{MizukiSone2009}.
A single uniform closed shuffle suffices to
compute any Boolean circuit~\cite{ShinagawaNuida2021}. A later compiler maps
simple finite-runtime card protocols to private simultaneous-message
protocols~\cite{ShinagawaNuida2025}.
\end{lstlisting}
\colorbox{RewriteGreen}{\parbox{0.97\textwidth}{\textbf{Rewrite suggestion.}\par\begingroup\small
Work on closed shuffles and graph automorphisms provides a physical point of comparison for group-based shuffle models. Miyamoto and Shinagawa realize a uniformly chosen graph automorphism with pile-scramble shuffles~\cite{MiyamotoShinagawa2022}. Shinagawa and Miyamoto extend this construction to graphs and hypergraphs and apply it to cyclic group shuffles~\cite{ShinagawaMiyamoto2023}. These physical protocols implement structured uniform shuffle groups exactly. This paper's word approach offers a different tradeoff. A pile-scramble realization samples its group uniformly in one physical procedure, while a word shuffle repeats elementary moves and incurs a kernel-certified $L_1$ error of $2^{-40}$. That line of work provides unobserved physical realizations for structured shuffles. I leave open whether the five $\PG$ letters admit such realizations. The eight-card instance also has a larger deck vocabulary than two-symbol encodings. It stores a Boolean secret in the arrangement of eight pairwise distinct cards, whereas committed formats store one bit in two two-symbol cards~\cite{MizukiSone2009}. Broader results show what physical shuffle models can support. A single uniform closed shuffle can compute any Boolean circuit~\cite{ShinagawaNuida2021}. A later compiler maps simple finite-runtime card protocols to private simultaneous-message protocols~\cite{ShinagawaNuida2025}. Because this paper uses an approximate word shuffle, the next comparison concerns the formal tools used to certify its mixing bound.
\par\endgroup}}
\Field{Why this rewrite.}{The rewrite organizes the paragraph around two reader questions. It first compares exact physical sampling with the approximate word shuffle and states the exact unresolved question about physical realizations of the five $\PG$ letters. It then explains the representation cost and uses the circuit and compiler results to show why physical shuffle models matter beyond this instance. The closing sentence returns to the approximate shuffle and prepares the next paragraph's comparison of mixing proof methods. Every original claim, numerical bound, citation, scope limit, and first-person responsibility is preserved.}
\Field{Terminology actions.}{Keep: pile-scramble shuffle, L1 error, pairwise distinct cards | Explain: unobserved physical realization, closed shuffle | Replace: trades against them -> offers a different tradeoff}
\par\smallskip
\needspace{10\baselineskip}\subsubsection*{W-10-P001: analogous purpose}
\Field{Location.}{Section 10, prose\_paragraph, lines 1869--1880.}
\Field{Draft purpose.}{Recap the framework, main instance, security scope, and transfer conditions.}
\Field{Qualitative closeness.}{analogous purpose. Both units synthesize the central construction and principal results in the conclusion.}
\Field{Limits of the analogy.}{The baseline turns next to physical future work. This unit instead audits theorem scope and states conditional transfer.}
\Field{Baseline narration used.}{adapted.}
\Field{Complexity and jargon.}{Complexity: \Tag{high} $\rightarrow$ \Tag{medium}. Jargon: \Tag{high} $\rightarrow$ \Tag{medium}.}
\textbf{Original WADT unit. Footnotes omitted.}
\begin{lstlisting}
I presented a group-parametric \Rocq{} framework for card-based protocols
and instantiated it with the action of $\PG$ on eight positions. Theorem A
gives executed correctness, recovery parameters $(3,7,8)$, and perfect view
and trace privacy for coalitions of at most three players under the uniform
group distribution. These security results use a uniform Boolean prior, the
fixed representative dealer, passive adversaries, and the dealer
distributions stated in Section~\ref{sec:model}. The verifier learns the
secret. Theorem~\ref{thm:generic-privacy} consumes only the
three-transitive action, the fixed encoding into decks of distinct
cards, and the uniform group distribution, and it therefore transfers to
any instance that supplies these three inputs.
\end{lstlisting}
\colorbox{RewriteGreen}{\parbox{0.97\textwidth}{\textbf{Rewrite suggestion.}\par\begingroup\small
I presented a group-parametric \Rocq{} framework for card-based protocols and instantiated it with the action of $\PG$ on eight positions. Theorem A proves executed correctness, recovery parameters $(3,7,8)$, and perfect view and trace privacy for coalitions of at most three players under the uniform group distribution. The tuple $(3,7,8)$ gives the privacy threshold, reconstruction threshold, and total-player count. These security results assume a uniform Boolean prior, the fixed representative dealer, passive adversaries, and the dealer distributions in Section~\ref{sec:model}. The verifier learns the secret. Theorem~\ref{thm:generic-privacy} requires only a three-transitive action, a fixed encoding into decks of distinct cards, and the uniform group distribution. It therefore applies to any instance that supplies these three inputs.
\par\endgroup}}
\Field{Why this rewrite.}{Adapts the baseline's conclusion sequence to the framework, exact guarantees, explicit assumptions, verifier scope, and conditional reuse.}
\Field{Terminology actions.}{Keep: executed correctness, perfect view and trace privacy, passive adversaries, uniform group distribution | Explain: recovery parameters (3,7,8), fixed representative dealer, three-transitive action | Replace: consumes only -> requires only}
\par\smallskip
\needspace{10\baselineskip}\subsubsection*{W-10-P003: partial analogue}
\Field{Location.}{Section 10, prose\_paragraph, lines 1888--1907.}
\Field{Draft purpose.}{Separate settled results from open formalization and security extensions.}
\Field{Qualitative closeness.}{partial analogue. The baseline supplies only the move of ending a conclusion with a concrete research direction.}
\Field{Limits of the analogy.}{The baseline proposes one material-design extension. This unit inventories several formal, cryptographic, and proof-kernel gaps while preserving settled results.}
\Field{Baseline narration used.}{partial.}
\Field{Complexity and jargon.}{Complexity: \Tag{high} $\rightarrow$ \Tag{medium}. Jargon: \Tag{high} $\rightarrow$ \Tag{medium}.}
\textbf{Original WADT unit. Footnotes omitted.}
\begin{lstlisting}
Two extensions of the word-shuffle security results remain open. The
conversion of the $L_1$ bound into a conditional-entropy bound rests on
classical mathematics, since continuity of entropy in total variation is
the Fannes inequality~\cite{Fannes1973,CsiszarKorner2011}. What remains
open here is only its formalization, which the underlying probability
library does not yet provide. The all-decks and shuffle-free dealers need
a word-shuffle mixing statement on their own sample spaces. Another
direction is to connect
bounded-width branching programs to group-based card protocols through
Barrington's theorem~\cite{Barrington1989,DvorakKoucky2021}. This direction
would make the workshop abstract's $\mathrm{NC}^1$ compilation proposal
precise. Further work should add active and compositional security, building
on models for active card protocols~\cite{KochFUN2021}.
The quantitative leakage for reveal sets of four through six cards is
settled in Section~\ref{sec:exact}: the orbit census and the collision
counts are kernel-checked, and the four values are evaluated from the
counts by a committed script. What remains open there is a kernel-level
proof of the mutual-information identity and of the orbit-invariance step
that extends each value from its representative.
\end{lstlisting}
\colorbox{RewriteGreen}{\parbox{0.97\textwidth}{\textbf{Rewrite suggestion.}\par\begingroup\small
The word-shuffle security results leave two direct extensions open. First, converting the $L_1$ bound into a conditional-entropy bound uses classical mathematics. Continuity of entropy in total variation is the Fannes inequality~\cite{Fannes1973,CsiszarKorner2011}. Only its formalization remains open here because the underlying probability library does not yet provide it. Second, the all-decks and shuffle-free dealers need word-shuffle mixing statements on their own sample spaces. Two broader directions would extend the framework. Connecting bounded-width branching programs with group-based card protocols through Barrington's theorem~\cite{Barrington1989,DvorakKoucky2021} would make the workshop abstract's $\mathrm{NC}^1$ compilation proposal precise. Further work should also add active and compositional security by building on models for active card protocols~\cite{KochFUN2021}. A separate proof-completeness issue concerns the quantitative leakage for reveal sets of four through six cards. Section~\ref{sec:exact} settles these values at the current trust level. The orbit census and collision counts are kernel-checked, and a committed script evaluates the four values from these counts. Two kernel-level steps remain open. One is a proof of the mutual-information identity. The other is a proof of the orbit-invariance step that extends each value from its representative.
\par\endgroup}}
\Field{Why this rewrite.}{The rewrite gives the paragraph three explicit layers. It first states the two direct extensions of the word-shuffle results. It then groups the compilation and stronger-security directions as broader framework extensions. Finally, it contrasts the settled leakage computation with the two kernel proofs that remain open. This structure realizes the contract's settled-versus-open function and follows the matched baseline's move from a result to future work grounded in that result. Every original claim, hedge, citation, cross-reference, formula, and trust-scope distinction is preserved.}
\Field{Terminology actions.}{Keep: Fannes inequality, active and compositional security, kernel-checked orbit census and collision counts | Explain: dealer-specific word-shuffle mixing, orbit-invariance step}
\par\smallskip
\clearpage\subsection{Baseline S-06-P009: cyclic-card specification tuple}
\colorbox{BaselineBlue}{\parbox{0.97\textwidth}{\Field{Purpose.}{Assemble the card set, transformation set, symbol set, and vision function into one specification.}\Field{Primary purpose.}{definition}\Field{Narration structure.}{definition\_lead\_in $\rightarrow$ tuple\_assembly $\rightarrow$ display\_handoff}\Field{Reusable move.}{Conclude a componentwise development by packaging the named parts into the formal object.}}}
\needspace{10\baselineskip}\subsubsection*{W-03-TC013: partial analogue}
\Field{Location.}{Section 03, table\_cell, lines 378--379.}
\Field{Draft purpose.}{State that the profile row aggregates the run layout, group action and shuffle law, and reconstruction component.}
\Field{Qualitative closeness.}{partial analogue. The baseline supplies only the move of naming a component by the role it contributes to an assembled object.}
\Field{Limits of the analogy.}{The baseline explicitly names four cyclic-card components. This table cell uses shorthand for three framework rows, so the replacement must spell out run layout, shuffle data, and reconstruction.}
\Field{Baseline narration used.}{partial.}
\Field{Complexity and jargon.}{Complexity: \Tag{low} $\rightarrow$ \Tag{low}. Jargon: \Tag{medium} $\rightarrow$ \Tag{low}.}
\textbf{Original WADT unit. Footnotes omitted.}
\begin{lstlisting}
all of the above
\end{lstlisting}
\colorbox{RewriteGreen}{\parbox{0.97\textwidth}{\textbf{Rewrite suggestion.}\par\begingroup\small
run layout, group action and shuffle distribution, and reconstruction
\par\endgroup}}
\Field{Why this rewrite.}{The rewrite uses the partial analogue's assembly move and replaces the vague reference with the complete component list.}
\Field{Terminology actions.}{Replace: all of the above -> run layout, group action and shuffle distribution, and reconstruction}
\par\smallskip
\needspace{10\baselineskip}\subsubsection*{W-03-TC014: partial analogue}
\Field{Location.}{Section 03, table\_cell, lines 378--379.}
\Field{Draft purpose.}{Name the object produced when the framework components are aggregated.}
\Field{Qualitative closeness.}{partial analogue. The baseline's component inventory provides a structural cue for naming the object formed from the preceding framework rows.}
\Field{Limits of the analogy.}{The baseline assembles a cyclic-card specification and hands it to a display. This cell names a protocol instance and relies on the caption for the four exact record names.}
\Field{Baseline narration used.}{partial.}
\Field{Complexity and jargon.}{Complexity: \Tag{low} $\rightarrow$ \Tag{low}. Jargon: \Tag{medium} $\rightarrow$ \Tag{low}.}
\textbf{Original WADT unit. Footnotes omitted.}
\begin{lstlisting}
one packaged protocol instance
\end{lstlisting}
\colorbox{RewriteGreen}{\parbox{0.97\textwidth}{\textbf{Rewrite suggestion.}\par\begingroup\small
one protocol instance that packages these components
\par\endgroup}}
\Field{Why this rewrite.}{The rewrite uses the partial analogue's completed-object move and explains what packaged means within the concise cell.}
\Field{Terminology actions.}{Keep: protocol instance | Explain: packaged protocol instance}
\par\smallskip
\needspace{10\baselineskip}\subsubsection*{W-03-DIA008: partial analogue}
\Field{Location.}{Section 03, diagram\_text, lines 440--442.}
\Field{Draft purpose.}{Label the executable protocol derived from the completed profile.}
\Field{Qualitative closeness.}{partial analogue. The baseline supplies only the closing assembly move that names the object obtained from formal components. This node applies it to the executable protocol derived from the profile.}
\Field{Limits of the analogy.}{The baseline assembles a card specification directly. This figure node also distinguishes derived output from profile input through color and a dotted edge.}
\Field{Baseline narration used.}{partial.}
\Field{Complexity and jargon.}{Complexity: \Tag{low} $\rightarrow$ \Tag{low}. Jargon: \Tag{medium} $\rightarrow$ \Tag{medium}.}
\textbf{Original WADT unit. Footnotes omitted.}
\begin{lstlisting}
derived protocol\\[-2pt]
      \footnotesize roles, characters, round-trip lemma
\end{lstlisting}
\colorbox{RewriteGreen}{\parbox{0.97\textwidth}{\textbf{Rewrite suggestion.}\par\begingroup\small
derived protocol\\[-2pt] \footnotesize roles, characters as observable profile parameters, round-trip lemma
\par\endgroup}}
\Field{Why this rewrite.}{The rewrite uses the partial analogue's completed-object move. It explains characters as observable profile parameters rather than textual symbols and retains the linking lemma.}
\Field{Terminology actions.}{Keep: derived protocol, round-trip lemma | Explain: characters}
\par\smallskip
\needspace{10\baselineskip}\subsubsection*{W-03-P004: analogous purpose}
\Field{Location.}{Section 03, prose\_paragraph, lines 464--468.}
\Field{Draft purpose.}{Define a completed profile and explain why its components remain tied to one instance.}
\Field{Qualitative closeness.}{analogous purpose. Both units complete a formal package from previously introduced components and then state what the package provides.}
\Field{Limits of the analogy.}{The baseline closes a card specification. The WADT paragraph gives a prose name to a filled record and leads into three proof obligations.}
\Field{Baseline narration used.}{adapted.}
\Field{Complexity and jargon.}{Complexity: \Tag{medium} $\rightarrow$ \Tag{low}. Jargon: \Tag{high} $\rightarrow$ \Tag{medium}.}
\textbf{Original WADT unit. Footnotes omitted.}
\begin{lstlisting}
I call a filled \coqin{MonodromyProfile} a profile. Packing these
components in one profile keeps the participants, verifier,
recovery map, privacy threshold, and shuffle bound tied to the same
instance. Each profile supplies proofs of the following three obligations.
\end{lstlisting}
\colorbox{RewriteGreen}{\parbox{0.97\textwidth}{\textbf{Rewrite suggestion.}\par\begingroup\small
I call a completed profile record a profile. It keeps the participants, verifier, recovery map that decodes the secret, privacy threshold, and shuffle bound tied to one instance. Every profile proves three obligations.
\par\endgroup}}
\Field{Why this rewrite.}{The rewrite adapts the baseline's component-assembly move. It replaces the code-facing record phrase, explains the recovery map, and announces the three obligations after the consistency benefit.}
\Field{Terminology actions.}{Keep: profile, privacy threshold, shuffle bound | Explain: recovery map | Replace: filled MonodromyProfile -> completed profile record}
\par\smallskip
\clearpage\subsection{Baseline S-06-P010: cyclic-card operation analogies}
\colorbox{BaselineBlue}{\parbox{0.97\textwidth}{\Field{Purpose.}{Define cyclic-card operations by separating those inherited from binary cards from those inherited from dihedral cards.}\Field{Primary purpose.}{comparison}\Field{Narration structure.}{comparison\_setup $\rightarrow$ first\_correspondence $\rightarrow$ second\_correspondence}\Field{Reusable move.}{Avoid repeating operation definitions by naming the precise source model for each operation family.}}}
\needspace{10\baselineskip}\subsubsection*{W-09-P004: partial analogue}
\Field{Location.}{Section 09, prose\_paragraph, lines 1818--1824.}
\Field{Draft purpose.}{Position the interpreter relative to established operational semantics.}
\Field{Qualitative closeness.}{partial analogue. The baseline supplies the correspondence move of anchoring a new mechanism to a familiar one and naming the changed behavior.}
\Field{Limits of the analogy.}{The baseline defines inherited operation families and closes an appendix. This unit is literature positioning and continues related work.}
\Field{Baseline narration used.}{partial.}
\Field{Complexity and jargon.}{Complexity: \Tag{medium} $\rightarrow$ \Tag{low}. Jargon: \Tag{medium} $\rightarrow$ \Tag{low}.}
\textbf{Original WADT unit. Footnotes omitted.}
\begin{lstlisting}
The field's standard operational semantics is the abstract-machine model
of Mizuki and Shizuya~\cite{MizukiShizuya2014}, which fixes the sequences
of shuffles, turns, and permutations a protocol may apply. The interpreter
of this paper plays the corresponding role for group-parametric protocols,
with the shuffle drawn from a word distribution over the group's
generators.
\end{lstlisting}
\colorbox{RewriteGreen}{\parbox{0.97\textwidth}{\textbf{Rewrite suggestion.}\par\begingroup\small
Mizuki and Shizuya's abstract-machine model provides the field's standard operational semantics~\cite{MizukiShizuya2014}. It fixes the sequences of shuffles, turns, and permutations that a protocol may apply. This paper's interpreter provides the analogous execution model for group-parametric protocols. It draws the shuffle from a word distribution over the group's generators.
\par\endgroup}}
\Field{Why this rewrite.}{Uses only the baseline's correspondence move to align the earlier operational model with this paper's interpreter and changed shuffle source.}
\Field{Terminology actions.}{Keep: operational semantics, interpreter, word distribution | Explain: group-parametric protocol | Replace: plays the corresponding role -> provides the analogous execution model}
\par\smallskip
\clearpage\subsection{No baseline analogue}
\colorbox{BaselineBlue}{\parbox{0.97\textwidth}{\textbf{Purpose.} These WADT units have no sound primary analogue in the Shinagawa baseline. Their suggestions preserve the local unit form without importing a baseline narration pattern.}}
\needspace{10\baselineskip}\subsubsection*{W-02-DIA001: no baseline analogue}
\Field{Location.}{Section 02, diagram\_text, lines 310--311.}
\Field{Draft purpose.}{Label the ideal distribution at the start of the upper proof path.}
\Field{Qualitative closeness.}{no baseline analogue. No eligible baseline paragraph serves as an atomic diagram-node label for one distribution.}
\Field{Limits of the analogy.}{A prose definition would add development and a bridge that this visual label does not contain.}
\Field{Baseline narration used.}{not applicable.}
\Field{Complexity and jargon.}{Complexity: \Tag{low} $\rightarrow$ \Tag{low}. Jargon: \Tag{low} $\rightarrow$ \Tag{low}.}
\textbf{Original WADT unit. Footnotes omitted.}
\begin{lstlisting}
uniform group\\distribution $U_G$
\end{lstlisting}
\colorbox{RewriteGreen}{\parbox{0.97\textwidth}{\textbf{Rewrite suggestion.}\par\begingroup\small
uniform group\\distribution $U_G$
\par\endgroup}}
\Field{Why this rewrite.}{Keeps the atomic diagram label because its compact mathematical role has no paragraph analogue.}
\Field{Terminology actions.}{Keep: uniform group distribution}
\par\smallskip
\needspace{10\baselineskip}\subsubsection*{W-02-DIA002: no baseline analogue}
\Field{Location.}{Section 02, diagram\_text, lines 311--312.}
\Field{Draft purpose.}{Label the exact privacy result at the end of the upper proof path.}
\Field{Qualitative closeness.}{no baseline analogue. No eligible baseline paragraph has the atomic visual function of naming one theorem endpoint.}
\Field{Limits of the analogy.}{Result prose would obscure the node's role as the terminus of one arrow.}
\Field{Baseline narration used.}{not applicable.}
\Field{Complexity and jargon.}{Complexity: \Tag{low} $\rightarrow$ \Tag{low}. Jargon: \Tag{low} $\rightarrow$ \Tag{low}.}
\textbf{Original WADT unit. Footnotes omitted.}
\begin{lstlisting}
Theorem A: perfect view\\and trace privacy
\end{lstlisting}
\colorbox{RewriteGreen}{\parbox{0.97\textwidth}{\textbf{Rewrite suggestion.}\par\begingroup\small
Theorem A: perfect view\\and trace privacy
\par\endgroup}}
\Field{Why this rewrite.}{Keeps the theorem endpoint as a concise diagram node with no added derivation.}
\Field{Terminology actions.}{Keep: perfect view and trace privacy}
\par\smallskip
\needspace{10\baselineskip}\subsubsection*{W-02-DIA005: no baseline analogue}
\Field{Location.}{Section 02, diagram\_text, lines 316--317.}
\Field{Draft purpose.}{Label the transfer step in the lower proof path.}
\Field{Qualitative closeness.}{no baseline analogue. No eligible baseline paragraph functions as a compact intermediate label in a proof-path diagram.}
\Field{Limits of the analogy.}{A paragraph analogue would require explanatory claims absent from this node.}
\Field{Baseline narration used.}{not applicable.}
\Field{Complexity and jargon.}{Complexity: \Tag{medium} $\rightarrow$ \Tag{low}. Jargon: \Tag{medium} $\rightarrow$ \Tag{low}.}
\textbf{Original WADT unit. Footnotes omitted.}
\begin{lstlisting}
endpoint and\\product-distribution transfers
\end{lstlisting}
\colorbox{RewriteGreen}{\parbox{0.97\textwidth}{\textbf{Rewrite suggestion.}\par\begingroup\small
transfer to card endpoints\\and secret-prior products
\par\endgroup}}
\Field{Why this rewrite.}{Names both transfer targets directly so the intermediate diagram step is readable without paragraph-level detail.}
\Field{Terminology actions.}{Keep: product-distribution transfer | Explain: endpoint transfer -> transfer to card endpoints}
\par\smallskip
\needspace{10\baselineskip}\subsubsection*{W-02-DIA006: no baseline analogue}
\Field{Location.}{Section 02, diagram\_text, lines 318--319.}
\Field{Draft purpose.}{Label the approximate privacy result at the end of the lower proof path.}
\Field{Qualitative closeness.}{no baseline analogue. No eligible baseline paragraph has the same atomic endpoint-label function.}
\Field{Limits of the analogy.}{Result paragraphs contain derivation or interpretation that the diagram node does not express.}
\Field{Baseline narration used.}{not applicable.}
\Field{Complexity and jargon.}{Complexity: \Tag{low} $\rightarrow$ \Tag{low}. Jargon: \Tag{low} $\rightarrow$ \Tag{low}.}
\textbf{Original WADT unit. Footnotes omitted.}
\begin{lstlisting}
Theorem B: view and\\trace bounds at $2^{-39}$
\end{lstlisting}
\colorbox{RewriteGreen}{\parbox{0.97\textwidth}{\textbf{Rewrite suggestion.}\par\begingroup\small
Theorem B: view and\\trace error at most $2^{-39}$
\par\endgroup}}
\Field{Why this rewrite.}{States the two observation bounds and their numerical limit in one compact endpoint label.}
\Field{Terminology actions.}{Keep: view and trace bounds}
\par\smallskip
\needspace{10\baselineskip}\subsubsection*{W-02-CAP001: no baseline analogue}
\Field{Location.}{Section 02, caption, lines 323--326.}
\Field{Draft purpose.}{Explain how to read the figure's two proof paths.}
\Field{Qualitative closeness.}{no baseline analogue. No eligible baseline paragraph performs this figure-specific caption function.}
\Field{Limits of the analogy.}{Roadmap prose can order sections but cannot replace the caption's direct mapping from visual paths to theorem endpoints.}
\Field{Baseline narration used.}{not applicable.}
\Field{Complexity and jargon.}{Complexity: \Tag{low} $\rightarrow$ \Tag{low}. Jargon: \Tag{medium} $\rightarrow$ \Tag{low}.}
\textbf{Original WADT unit. Footnotes omitted.}
\begin{lstlisting}
The two proof paths and the two headline theorems. The upper
  path ends in Theorem A. The lower path transports the mixing certificate
  to Theorem B.
\end{lstlisting}
\colorbox{RewriteGreen}{\parbox{0.97\textwidth}{\textbf{Rewrite suggestion.}\par\begingroup\small
Two proof paths lead to the headline theorems. The upper path ends at Theorem A. The lower path transfers the mixing certificate from the word distribution to the view and trace bounds of Theorem B.
\par\endgroup}}
\Field{Why this rewrite.}{Uses the caption's own contract to map each visual path directly to its theorem endpoint.}
\Field{Terminology actions.}{Keep: mixing certificate | Explain: transports the mixing certificate -> transfers the certificate to view and trace bounds}
\par\smallskip
\needspace{10\baselineskip}\subsubsection*{W-03-CAP002: no baseline analogue}
\Field{Location.}{Section 03, caption, lines 456--461.}
\Field{Draft purpose.}{Make the architecture figure readable by decoding object classes and edge styles.}
\Field{Qualitative closeness.}{no baseline analogue. No eligible baseline paragraph decodes both object colors and typed edges in an architecture figure.}
\Field{Limits of the analogy.}{Roadmap and definition paragraphs can order concepts, but they do not explain colors, line styles, or dependency edges.}
\Field{Baseline narration used.}{not applicable.}
\Field{Complexity and jargon.}{Complexity: \Tag{low} $\rightarrow$ \Tag{low}. Jargon: \Tag{medium} $\rightarrow$ \Tag{low}.}
\textbf{Original WADT unit. Footnotes omitted.}
\begin{lstlisting}
Architecture of the group-parametric card-protocol framework.
  Blue boxes are profile records, green boxes are supporting records, and
  the purple box is the derived protocol. Solid arrows are required
  dependencies, the dashed arrow is optional, and the dotted arrow is
  derivation.
\end{lstlisting}
\colorbox{RewriteGreen}{\parbox{0.97\textwidth}{\textbf{Rewrite suggestion.}\par\begingroup\small
Architecture of the group-parametric card-protocol framework. Blue boxes denote profile records, green boxes denote auxiliary records, and the purple box denotes the derived protocol. Solid arrows mark required dependencies. The dashed arrow marks an optional dependency. The dotted arrow marks derivation.
\par\endgroup}}
\Field{Why this rewrite.}{No baseline analogue applies. The rewrite follows the caption's own contract by mapping each color and line style to exactly one role in short sentences.}
\Field{Terminology actions.}{Keep: profile record, derived protocol, required dependency | Explain: derivation}
\par\smallskip
\needspace{10\baselineskip}\subsubsection*{W-03-TC031: no baseline analogue}
\Field{Location.}{Section 03, table\_cell, lines 511--512.}
\Field{Draft purpose.}{Name the two group carriers that realize the spectral, asymptotic-only witness row.}
\Field{Qualitative closeness.}{no baseline analogue. No eligible baseline paragraph guides a cell that only lists two group carriers without making a rhetorical or inferential move.}
\Field{Limits of the analogy.}{Definition examples elsewhere provide context for groups, but using them as a primary analogue would invent a function absent from this nominative cell.}
\Field{Baseline narration used.}{not applicable.}
\Field{Complexity and jargon.}{Complexity: \Tag{low} $\rightarrow$ \Tag{low}. Jargon: \Tag{medium} $\rightarrow$ \Tag{low}.}
\textbf{Original WADT unit. Footnotes omitted.}
\begin{lstlisting}
$S_5$, $S_5\times S_5$
\end{lstlisting}
\colorbox{RewriteGreen}{\parbox{0.97\textwidth}{\textbf{Rewrite suggestion.}\par\begingroup\small
$S_5$ and the direct product $S_5\times S_5$
\par\endgroup}}
\Field{Why this rewrite.}{No baseline analogue applies. The rewrite follows the cell's own contract and explains the relation between the two mathematical carriers.}
\Field{Terminology actions.}{Keep: S5, S5 times S5 | Explain: direct product}
\par\smallskip
\needspace{10\baselineskip}\subsubsection*{W-03-P010: no baseline analogue}
\Field{Location.}{Section 03, prose\_paragraph, lines 560--584.}
\Field{Draft purpose.}{Explain the record-level privacy theorem, show how the main instance discharges it, and prove that both obligations are necessary.}
\Field{Qualitative closeness.}{no baseline analogue. No eligible baseline paragraph combines record-level theorem packaging, instance discharge, two sharpness counterexamples, and a formal exclusion of another protocol family.}
\Field{Limits of the analogy.}{Security examples and limitation paragraphs each cover only one move. A single primary would hide the paragraph's theorem-audit function.}
\Field{Baseline narration used.}{not applicable.}
\Field{Complexity and jargon.}{Complexity: \Tag{high} $\rightarrow$ \Tag{high}. Jargon: \Tag{high} $\rightarrow$ \Tag{medium}.}
\textbf{Original WADT unit. Footnotes omitted.}
\begin{lstlisting}
Theorem~\ref{thm:generic-privacy} is also stated once at the record level.
One statement quantifies over any profile: if the group of position
permutations realized by the profile's action is $t$-transitive and every
encoded deck has distinct cards, then the view of every coalition of at
most $t$ positions is independent of the secret, under the record's
share-count and deck-length compatibility premise. The $\PG$
instance discharges the two obligations with its three-transitivity lemma
and its distinct-deck lemma. Both obligations are necessary rather than
decorative. Weakening the coalition bound to $t+1$ is refuted at the $\PG$
instance by the explicit four-position coalition of
Section~\ref{sec:pgl}, and dropping the distinct-deck
obligation is refuted by a profile over the same group whose encoder deals
one constant card to every position and leaks the secret to a single
position. The record-level
statement keeps the theorem's two-valued secret, and the constant-deck
counterexample shows the distinct-cards hypothesis cannot be dropped, so
the exclusion of the two-symbol five-card family from this theorem is
forced by the mathematics rather than by a formalization choice.
\end{lstlisting}
\colorbox{RewriteGreen}{\parbox{0.97\textwidth}{\textbf{Rewrite suggestion.}\par\begingroup\small
Theorem~\ref{thm:generic-privacy} is also stated once for every profile. Suppose that the position-permutation group realized by a profile is $t$-transitive. Suppose also that every encoded deck has distinct cards and that the profile's share count matches its deck length. Then every coalition of at most $t$ positions has a view independent of the secret. The $\PG$ instance proves the first two obligations with its three-transitivity and distinct-deck lemmas. Both obligations are necessary for the theorem. An explicit four-position coalition from Section~\ref{sec:pgl} refutes the conclusion when the coalition bound is weakened to $t+1$. A second profile over the same group deals one constant card to every position. This repeated-card encoder leaks the secret to one position, so it refutes the theorem without the distinct-deck obligation. The record-level statement retains a two-valued secret. The constant-deck counterexample shows that distinct cards cannot be omitted. The mathematics therefore excludes the two-symbol five-card family from this theorem.
\par\endgroup}}
\Field{Why this rewrite.}{No baseline analogue applies. The rewrite follows the unit's own theorem-counterexample-scope contract, excludes both footnotes, explains size compatibility and the two-valued scope, and keeps both sharpness arguments.}
\Field{Terminology actions.}{Keep: t-transitive action, distinct-deck obligation, constant-deck counterexample | Explain: share-count and deck-length compatibility, two-valued secret | Replace: necessary rather than decorative -> necessary for the theorem}
\par\smallskip
\needspace{10\baselineskip}\subsubsection*{W-04-DIA003: no baseline analogue}
\Field{Location.}{Section 04, diagram\_text, lines 640--641.}
\Field{Draft purpose.}{State the player count and one-card allocation in the run diagram.}
\Field{Qualitative closeness.}{no baseline analogue. No eligible baseline paragraph performs this compact allocation-label function inside an execution diagram.}
\Field{Limits of the analogy.}{Prose about protocol resources would add explanation that this nominative diagram label does not carry.}
\Field{Baseline narration used.}{not applicable.}
\Field{Complexity and jargon.}{Complexity: \Tag{low} $\rightarrow$ \Tag{low}. Jargon: \Tag{low} $\rightarrow$ \Tag{low}.}
\textbf{Original WADT unit. Footnotes omitted.}
\begin{lstlisting}
five, one card each
\end{lstlisting}
\colorbox{RewriteGreen}{\parbox{0.97\textwidth}{\textbf{Rewrite suggestion.}\par\begingroup\small
five players, one card each
\par\endgroup}}
\Field{Why this rewrite.}{Uses the unit's allocation-label contract and makes the omitted noun explicit without adding a claim.}
\Field{Terminology actions.}{Keep: one card each}
\par\smallskip
\needspace{10\baselineskip}\subsubsection*{W-04-DIA025: no baseline analogue}
\Field{Location.}{Section 04, diagram\_text, lines 764--765.}
\Field{Draft purpose.}{Label one plotted reveal case with its leakage value and status as the maximum.}
\Field{Qualitative closeness.}{no baseline analogue. No eligible baseline paragraph serves as a numeric annotation within a leakage diagram.}
\Field{Limits of the analogy.}{Result prose would overdevelop this atomic visual label.}
\Field{Baseline narration used.}{not applicable.}
\Field{Complexity and jargon.}{Complexity: \Tag{low} $\rightarrow$ \Tag{low}. Jargon: \Tag{low} $\rightarrow$ \Tag{low}.}
\textbf{Original WADT unit. Footnotes omitted.}
\begin{lstlisting}
$0.811$ bits, the cap
\end{lstlisting}
\colorbox{RewriteGreen}{\parbox{0.97\textwidth}{\textbf{Rewrite suggestion.}\par\begingroup\small
$0.811$ bits, maximum leakage
\par\endgroup}}
\Field{Why this rewrite.}{Uses the atomic annotation contract and explains the cap as maximum leakage.}
\Field{Terminology actions.}{Keep: 0.811 bits | Explain: cap}
\par\smallskip
\needspace{10\baselineskip}\subsubsection*{W-04-DIA031: no baseline analogue}
\Field{Location.}{Section 04, diagram\_text, lines 770--771.}
\Field{Draft purpose.}{Label the second maximum-leakage case in the figure.}
\Field{Qualitative closeness.}{no baseline analogue. No eligible baseline paragraph provides a repeated numeric annotation for a separate visual case.}
\Field{Limits of the analogy.}{A prose result would obscure the label's role in showing equal maxima across two rows.}
\Field{Baseline narration used.}{not applicable.}
\Field{Complexity and jargon.}{Complexity: \Tag{low} $\rightarrow$ \Tag{low}. Jargon: \Tag{low} $\rightarrow$ \Tag{low}.}
\textbf{Original WADT unit. Footnotes omitted.}
\begin{lstlisting}
$0.811$ bits, the cap
\end{lstlisting}
\colorbox{RewriteGreen}{\parbox{0.97\textwidth}{\textbf{Rewrite suggestion.}\par\begingroup\small
$0.811$ bits, maximum leakage
\par\endgroup}}
\Field{Why this rewrite.}{Uses the atomic annotation contract for the second visual case and preserves the repeated maximum value.}
\Field{Terminology actions.}{Keep: 0.811 bits | Explain: cap}
\par\smallskip
\needspace{10\baselineskip}\subsubsection*{W-04-CAP002: no baseline analogue}
\Field{Location.}{Section 04, caption, lines 772--776.}
\Field{Draft purpose.}{Explain the leakage figure's cases, card-back notation, and evidential limits.}
\Field{Qualitative closeness.}{no baseline analogue. No eligible baseline paragraph combines a figure legend with a warning about which visual features determine the values.}
\Field{Limits of the analogy.}{General example prose cannot replace the caption's visual-decoding and nonclaim functions.}
\Field{Baseline narration used.}{not applicable.}
\Field{Complexity and jargon.}{Complexity: \Tag{medium} $\rightarrow$ \Tag{low}. Jargon: \Tag{medium} $\rightarrow$ \Tag{low}.}
\textbf{Original WADT unit. Footnotes omitted.}
\begin{lstlisting}
Mutual-information leakage for selected reveal cases. Blue card
  backs mark unrevealed positions. Each value depends only on which
  positions are revealed, never on the faces drawn here, which show one
  possible outcome.
\end{lstlisting}
\colorbox{RewriteGreen}{\parbox{0.97\textwidth}{\textbf{Rewrite suggestion.}\par\begingroup\small
Mutual-information leakage for selected reveal cases. Blue card backs mark unrevealed positions. Each value depends only on the revealed positions. The displayed faces show one possible outcome and do not determine the value.
\par\endgroup}}
\Field{Why this rewrite.}{Uses the caption's own visual-decoding contract. It separates the legend, determinant, and illustrative-face limitation.}
\Field{Terminology actions.}{Keep: mutual-information leakage, revealed positions | Explain: possible outcome}
\par\smallskip
\needspace{10\baselineskip}\subsubsection*{W-05-CAP001: no baseline analogue}
\Field{Location.}{Section 05, caption, lines 1015--1019.}
\Field{Draft purpose.}{Decode the positions, card values, shading, and secret classes in the encoding figure.}
\Field{Qualitative closeness.}{no baseline analogue. No eligible baseline paragraph combines a figure legend with row-by-row secret-class identification.}
\Field{Limits of the analogy.}{Definition or example prose cannot replace the caption's direct mapping from visual marks to mathematical objects.}
\Field{Baseline narration used.}{not applicable.}
\Field{Complexity and jargon.}{Complexity: \Tag{medium} $\rightarrow$ \Tag{low}. Jargon: \Tag{low} $\rightarrow$ \Tag{low}.}
\textbf{Original WADT unit. Footnotes omitted.}
\begin{lstlisting}
The encoded representatives. Gray labels are card positions,
  with position seven representing $\infty$. Boxed numbers are card values,
  and shading marks the four heart values. The row $D_0$ represents the
  harmonic secret $0$, and $D_1$ represents the equianharmonic secret $1$.
\end{lstlisting}
\colorbox{RewriteGreen}{\parbox{0.97\textwidth}{\textbf{Rewrite suggestion.}\par\begingroup\small
The two encoded representatives. Gray labels mark card positions, and position seven represents $\infty$. Boxed numbers are card values. Shading marks the four heart-labeled card values. The row $D_0$ represents the harmonic secret $0$. The row $D_1$ represents the equianharmonic secret $1$.
\par\endgroup}}
\Field{Why this rewrite.}{Uses the caption's own visual contract and assigns one mathematical role to each visual feature.}
\Field{Terminology actions.}{Keep: harmonic, equianharmonic, position seven}
\par\smallskip
\needspace{10\baselineskip}\subsubsection*{W-05-DIA007: no baseline analogue}
\Field{Location.}{Section 05, diagram\_text, lines 1135--1136.}
\Field{Draft purpose.}{Label the player group in the execution diagram.}
\Field{Qualitative closeness.}{no baseline analogue. No eligible baseline paragraph has the atomic diagram-label function of recording a participant count and allocation.}
\Field{Limits of the analogy.}{Protocol prose would add operations and therefore distort this node.}
\Field{Baseline narration used.}{not applicable.}
\Field{Complexity and jargon.}{Complexity: \Tag{low} $\rightarrow$ \Tag{low}. Jargon: \Tag{low} $\rightarrow$ \Tag{low}.}
\textbf{Original WADT unit. Footnotes omitted.}
\begin{lstlisting}
eight, one card each
\end{lstlisting}
\colorbox{RewriteGreen}{\parbox{0.97\textwidth}{\textbf{Rewrite suggestion.}\par\begingroup\small
eight players, one card each
\par\endgroup}}
\Field{Why this rewrite.}{Keeps the atomic participant and allocation label with no added protocol prose.}
\Field{Terminology actions.}{Keep: one card each}
\par\smallskip
\needspace{10\baselineskip}\subsubsection*{W-05-DIA009: no baseline analogue}
\Field{Location.}{Section 05, diagram\_text, lines 1140--1141.}
\Field{Draft purpose.}{Label the dealer's encoding step.}
\Field{Qualitative closeness.}{no baseline analogue. No eligible baseline paragraph is an atomic action label inside a flow diagram.}
\Field{Limits of the analogy.}{A method paragraph would supply sequence and explanation absent from the node.}
\Field{Baseline narration used.}{not applicable.}
\Field{Complexity and jargon.}{Complexity: \Tag{low} $\rightarrow$ \Tag{low}. Jargon: \Tag{low} $\rightarrow$ \Tag{low}.}
\textbf{Original WADT unit. Footnotes omitted.}
\begin{lstlisting}
encode secret $s$\\[-1pt] as the deck $D_s$
\end{lstlisting}
\colorbox{RewriteGreen}{\parbox{0.97\textwidth}{\textbf{Rewrite suggestion.}\par\begingroup\small
encode secret $s$ as deck $D_s$
\par\endgroup}}
\Field{Why this rewrite.}{Keeps the diagram node as one direct encoding action and its output.}
\Field{Terminology actions.}{Keep: encode}
\par\smallskip
\needspace{10\baselineskip}\subsubsection*{W-05-DIA010: no baseline analogue}
\Field{Location.}{Section 05, diagram\_text, lines 1142--1143.}
\Field{Draft purpose.}{Label the dealer's shuffle step.}
\Field{Qualitative closeness.}{no baseline analogue. No eligible baseline paragraph performs this compact flow-node role.}
\Field{Limits of the analogy.}{Shuffle-definition prose includes domains and conditions that the node intentionally omits.}
\Field{Baseline narration used.}{not applicable.}
\Field{Complexity and jargon.}{Complexity: \Tag{low} $\rightarrow$ \Tag{low}. Jargon: \Tag{medium} $\rightarrow$ \Tag{low}.}
\textbf{Original WADT unit. Footnotes omitted.}
\begin{lstlisting}
sample a shuffle $g$,\\[-1pt] rearrange to $g\star D_s$
\end{lstlisting}
\colorbox{RewriteGreen}{\parbox{0.97\textwidth}{\textbf{Rewrite suggestion.}\par\begingroup\small
sample shuffle $g$ and\\rearrange deck to $g\star D_s$
\par\endgroup}}
\Field{Why this rewrite.}{States the sampled group element and the resulting deck in one compact diagram action.}
\Field{Terminology actions.}{Keep: shuffle | Explain: group action on a deck -> rearranges the encoded deck to \$g\textbackslash{}star D\_s\$}
\par\smallskip
\needspace{10\baselineskip}\subsubsection*{W-05-DIA012: no baseline analogue}
\Field{Location.}{Section 05, diagram\_text, lines 1152--1153.}
\Field{Draft purpose.}{Label the verifier's decoding step.}
\Field{Qualitative closeness.}{no baseline analogue. No eligible baseline paragraph functions as the final action label of a diagram.}
\Field{Limits of the analogy.}{An interpretation paragraph would add reasoning beyond the node's single operation.}
\Field{Baseline narration used.}{not applicable.}
\Field{Complexity and jargon.}{Complexity: \Tag{low} $\rightarrow$ \Tag{low}. Jargon: \Tag{medium} $\rightarrow$ \Tag{low}.}
\textbf{Original WADT unit. Footnotes omitted.}
\begin{lstlisting}
decode the orbit label of\\[-1pt] the heart positions $\to s$
\end{lstlisting}
\colorbox{RewriteGreen}{\parbox{0.97\textwidth}{\textbf{Rewrite suggestion.}\par\begingroup\small
decode the orbit label of\\the heart-position set $\to s$
\par\endgroup}}
\Field{Why this rewrite.}{Keeps the node focused on the verifier's invariant, action, and output while explaining which positions form the decoder input.}
\Field{Terminology actions.}{Keep: orbit label | Explain: heart positions -> the heart-position set used by the decoder}
\par\smallskip
\needspace{10\baselineskip}\subsubsection*{W-05-CAP002: no baseline analogue}
\Field{Location.}{Section 05, caption, lines 1154--1160.}
\Field{Draft purpose.}{State the execution participants, decoder input, and shuffle law shown in the figure.}
\Field{Qualitative closeness.}{no baseline analogue. No eligible baseline paragraph combines execution-figure decoding with model-dependent visual interpretation and a cross-figure relation.}
\Field{Limits of the analogy.}{A protocol overview would not preserve the caption's direct tie to the diagram.}
\Field{Baseline narration used.}{not applicable.}
\Field{Complexity and jargon.}{Complexity: \Tag{medium} $\rightarrow$ \Tag{low}. Jargon: \Tag{high} $\rightarrow$ \Tag{medium}.}
\textbf{Original WADT unit. Footnotes omitted.}
\begin{lstlisting}
One $\PG$ run: eight players, one card each, no player inputs,
  and the verifier decodes the orbit label of the heart positions. The
  shuffle $g$ is drawn from the uniform distribution $U_G$ in the
  uniform-shuffle model and from the word distribution $\worddist$ in the
  word-shuffle model. This is the no-input counterpart of
  Figure~\ref{fig:fivecard-run}.
\end{lstlisting}
\colorbox{RewriteGreen}{\parbox{0.97\textwidth}{\textbf{Rewrite suggestion.}\par\begingroup\small
One $\PG$ run with eight players and one card per player. The players supply no private inputs. The verifier decodes the orbit label of the heart-position set. In the uniform-shuffle model, $g$ follows the uniform distribution $U_G$. In the word-shuffle model, $g$ follows the word distribution $\worddist$. This run is the no-input counterpart of Figure~\ref{fig:fivecard-run}.
\par\endgroup}}
\Field{Why this rewrite.}{Uses the caption's own contract to state the actors, no-input condition, decoder input, model-dependent law, and cross-figure relation in order.}
\Field{Terminology actions.}{Keep: uniform distribution on the group, word distribution | Explain: no-input counterpart -> the players supply no private inputs, orbit label -> the label of the heart-position set}
\par\smallskip
\needspace{10\baselineskip}\subsubsection*{W-05-CAP003: no baseline analogue}
\Field{Location.}{Section 05, caption, lines 1208--1211.}
\Field{Draft purpose.}{Explain the generator permutation table.}
\Field{Qualitative closeness.}{no baseline analogue. No eligible baseline paragraph decodes the rows and graphical marks of this generator figure.}
\Field{Limits of the analogy.}{Definition prose would not replace the caption's visual legend.}
\Field{Baseline narration used.}{not applicable.}
\Field{Complexity and jargon.}{Complexity: \Tag{medium} $\rightarrow$ \Tag{low}. Jargon: \Tag{medium} $\rightarrow$ \Tag{low}.}
\textbf{Original WADT unit. Footnotes omitted.}
\begin{lstlisting}
The three generating maps on the projective points. Gray labels
  are positions, and the box at position $i$ shows the image point $g(i)$.
  The symbol $\infty$ is position seven.
\end{lstlisting}
\colorbox{RewriteGreen}{\parbox{0.97\textwidth}{\textbf{Rewrite suggestion.}\par\begingroup\small
The three generating maps on the projective points. Each row shows one elementary group transformation. Gray labels mark positions. The box at position $i$ contains its image point $g(i)$. Position seven represents $\infty$.
\par\endgroup}}
\Field{Why this rewrite.}{Uses the caption's own visual contract to define the rows, labels, boxes, and infinity convention.}
\Field{Terminology actions.}{Keep: image point | Explain: generating map -> one elementary group transformation}
\par\smallskip
\needspace{10\baselineskip}\subsubsection*{W-05-STM009: no baseline analogue}
\Field{Location.}{Section 05, statement, lines 1255--1262.}
\Field{Draft purpose.}{State three-transitivity of the projective linear action.}
\Field{Qualitative closeness.}{no baseline analogue. No eligible baseline paragraph states a universal three-coordinate transport theorem without explanation or calculation.}
\Field{Limits of the analogy.}{Proof and result paragraphs in the baseline add derivations or application conclusions that this terse theorem does not contain.}
\Field{Baseline narration used.}{not applicable.}
\Field{Complexity and jargon.}{Complexity: \Tag{medium} $\rightarrow$ \Tag{low}. Jargon: \Tag{low} $\rightarrow$ \Tag{low}.}
\textbf{Original WADT unit. Footnotes omitted.}
\begin{lstlisting}
For any ordered triples $(x_1,x_2,x_3)$ and $(y_1,y_2,y_3)$ of distinct
projective points, there is a $g\in\PG$ such that
\[
  g(x_1)=y_1, \qquad g(x_2)=y_2, \qquad g(x_3)=y_3.
\]
\end{lstlisting}
\colorbox{RewriteGreen}{\parbox{0.97\textwidth}{\textbf{Rewrite suggestion.}\par\begingroup\small
For any ordered triples $(x_1,x_2,x_3)$ and $(y_1,y_2,y_3)$ of distinct projective points, there exists $g\in\PG$ such that
\[
  g(x_1)=y_1,
  \qquad
  g(x_2)=y_2,
  \qquad
  g(x_3)=y_3.
\]
\par\endgroup}}
\Field{Why this rewrite.}{Uses the WADT statement contract alone and keeps the quantified three-transitivity result terse. The excluded footnote does not enter the suggestion.}
\Field{Terminology actions.}{Keep: ordered triples, projective points}
\par\smallskip
\needspace{10\baselineskip}\subsubsection*{W-05-P025: no baseline analogue}
\Field{Location.}{Section 05, prose\_paragraph, lines 1274--1276.}
\Field{Draft purpose.}{Point readers to the formal source index for this and the next sections.}
\Field{Qualitative closeness.}{no baseline analogue. No eligible baseline paragraph serves only as a cross-reference to a formal-source index.}
\Field{Limits of the analogy.}{Roadmap prose discusses mathematical content rather than provenance lookup.}
\Field{Baseline narration used.}{not applicable.}
\Field{Complexity and jargon.}{Complexity: \Tag{low} $\rightarrow$ \Tag{low}. Jargon: \Tag{low} $\rightarrow$ \Tag{low}.}
\textbf{Original WADT unit. Footnotes omitted.}
\begin{lstlisting}
Table~\ref{tab:source-index} at the end of Section~\ref{sec:mixing}
indexes the formal sources for this section and the two that follow.
\end{lstlisting}
\colorbox{RewriteGreen}{\parbox{0.97\textwidth}{\textbf{Rewrite suggestion.}\par\begingroup\small
Table~\ref{tab:source-index} at the end of Section~\ref{sec:mixing} lists the formal source locations for this section and the next two sections.
\par\endgroup}}
\Field{Why this rewrite.}{Uses the WADT provenance contract alone and names the table, location, and exact coverage directly.}
\Field{Terminology actions.}{Keep: formal sources}
\par\smallskip
\needspace{10\baselineskip}\subsubsection*{W-06-P004: no baseline analogue}
\Field{Location.}{Section 06, prose\_paragraph, lines 1310--1315.}
\Field{Draft purpose.}{Explain the arrangement-level recovery cutoff independently of the implemented decoder.}
\Field{Qualitative closeness.}{no baseline analogue. No eligible baseline paragraph analyzes a reconstruction threshold by combining unique completion at seven positions with ambiguity through six.}
\Field{Limits of the analogy.}{Worked protocol and matching-relation examples do not provide this two-sided threshold argument.}
\Field{Baseline narration used.}{not applicable.}
\Field{Complexity and jargon.}{Complexity: \Tag{medium} $\rightarrow$ \Tag{medium}. Jargon: \Tag{medium} $\rightarrow$ \Tag{low}.}
\textbf{Original WADT unit. Footnotes omitted.}
\begin{lstlisting}
Recovery has a finer threshold than the implemented decoder exposes. Any
seven revealed positions determine the missing card because a valid deck uses
all eight card values once. Hence seven revealed positions determine the
orbit class. By contrast, every set of at most six revealed positions is
consistent with valid arrangements from both secret classes.
\end{lstlisting}
\colorbox{RewriteGreen}{\parbox{0.97\textwidth}{\textbf{Rewrite suggestion.}\par\begingroup\small
The structural recovery threshold is finer than the implemented decoder interface. Any seven revealed positions determine the missing card because every valid deck uses each of the eight card values exactly once. The seven cards therefore determine the orbit class. By contrast, any set of at most six revealed positions is compatible with valid arrangements from both secret classes.
\par\endgroup}}
\Field{Why this rewrite.}{No baseline analogue applies. The rewrite follows the paragraph's own threshold argument and explains compatibility with both classes while preserving the decoder boundary.}
\Field{Terminology actions.}{Keep: recovery threshold, orbit class | Explain: consistent with both secret classes}
\par\smallskip
\needspace{10\baselineskip}\subsubsection*{W-06-STM002: no baseline analogue}
\Field{Location.}{Section 06, statement, lines 1332--1342.}
\Field{Draft purpose.}{State the exact recovery threshold and show that privacy fails at four cards.}
\Field{Qualitative closeness.}{no baseline analogue. No eligible baseline paragraph states both the six-versus-seven recovery threshold and an explicit four-card mutual-information counterexample.}
\Field{Limits of the analogy.}{Security example paragraphs address one worked protocol, not this combined sharp threshold theorem.}
\Field{Baseline narration used.}{not applicable.}
\Field{Complexity and jargon.}{Complexity: \Tag{medium} $\rightarrow$ \Tag{medium}. Jargon: \Tag{high} $\rightarrow$ \Tag{medium}.}
\textbf{Original WADT unit. Footnotes omitted.}
\begin{lstlisting}
Every reveal set of at most six positions is compatible with valid
arrangements from both secret classes, and every set of seven revealed
positions determines the secret class. The view of the coalition of the
four heart positions in the identity arrangement depends on the secret and
has positive mutual information with it.
\end{lstlisting}
\colorbox{RewriteGreen}{\parbox{0.97\textwidth}{\textbf{Rewrite suggestion.}\par\begingroup\small
Every reveal set of at most six positions is compatible with valid arrangements from both secret classes. Every set of seven revealed positions determines the secret class. For the identity arrangement, the coalition consisting of the four heart positions has a view that depends on the secret and has positive mutual information with it.
\par\endgroup}}
\Field{Why this rewrite.}{No baseline analogue applies. The rewrite follows the statement's own threshold-then-witness contract, excludes the footnote, explains compatibility, and separates the three formal claims.}
\Field{Terminology actions.}{Keep: mutual information, four-heart coalition, secret class | Explain: compatible arrangement}
\par\smallskip
\needspace{10\baselineskip}\subsubsection*{W-06-CAP001: no baseline analogue}
\Field{Location.}{Section 06, caption, lines 1365--1369.}
\Field{Draft purpose.}{Explain the recovery-ramp figure and the meaning of each track.}
\Field{Qualitative closeness.}{no baseline analogue. No eligible baseline paragraph performs the combined caption and visual-legend function for two threshold tracks, a leakage witness, and an ambiguity range.}
\Field{Limits of the analogy.}{Comparison paragraphs can state two cutoffs but cannot guide the figure's track-by-track decoding.}
\Field{Baseline narration used.}{not applicable.}
\Field{Complexity and jargon.}{Complexity: \Tag{low} $\rightarrow$ \Tag{low}. Jargon: \Tag{medium} $\rightarrow$ \Tag{low}.}
\textbf{Original WADT unit. Footnotes omitted.}
\begin{lstlisting}
The recovery ramp $(t,r,n)=(3,7,8)$. The view track shows
  perfect privacy for coalitions of at most three positions and the leaking
  four-position witness. The reveal track shows ambiguity through six
  positions and determination from seven.
\end{lstlisting}
\colorbox{RewriteGreen}{\parbox{0.97\textwidth}{\textbf{Rewrite suggestion.}\par\begingroup\small
The recovery ramp $(t,r,n)=(3,7,8)$. The view track shows perfect privacy for coalitions of at most three positions and leakage for the four-position witness. The reveal track shows ambiguity through six positions and determination from seven.
\par\endgroup}}
\Field{Why this rewrite.}{No baseline analogue applies. The rewrite follows the caption's own tuple-view-reveal sequence and explains the witness directly.}
\Field{Terminology actions.}{Keep: recovery ramp, view track, reveal track | Explain: leaking four-position witness}
\par\smallskip
\needspace{10\baselineskip}\subsubsection*{W-06-P012: no baseline analogue}
\Field{Location.}{Section 06, prose\_paragraph, lines 1507--1509.}
\Field{Draft purpose.}{Point the reader to the assumptions supporting the preceding results.}
\Field{Qualitative closeness.}{no baseline analogue. No eligible baseline paragraph consists solely of a formal trust-base pointer for a theorem family.}
\Field{Limits of the analogy.}{Roadmaps and scope paragraphs provide fuller content. Assigning one as primary would overstate this unit's rhetorical development.}
\Field{Baseline narration used.}{not applicable.}
\Field{Complexity and jargon.}{Complexity: \Tag{low} $\rightarrow$ \Tag{low}. Jargon: \Tag{medium} $\rightarrow$ \Tag{low}.}
\textbf{Original WADT unit. Footnotes omitted.}
\begin{lstlisting}
The trust base for these results is stated in
Section~\ref{sec:instances}.
\end{lstlisting}
\colorbox{RewriteGreen}{\parbox{0.97\textwidth}{\textbf{Rewrite suggestion.}\par\begingroup\small
Section~\ref{sec:instances} states the trust base for Theorem A.
\par\endgroup}}
\Field{Why this rewrite.}{No baseline analogue applies. The rewrite follows the unit's own reference contract and replaces the vague result reference with Theorem A.}
\Field{Terminology actions.}{Keep: trust base | Explain: these results}
\par\smallskip
\needspace{10\baselineskip}\subsubsection*{W-07-TC006: no baseline analogue}
\Field{Location.}{Section 07, table\_cell, lines 1671--1672.}
\Field{Draft purpose.}{Label the construction claim that counts the equianharmonic and harmonic orbit classes.}
\Field{Qualitative closeness.}{no baseline analogue. No eligible baseline paragraph serves as a source-index cell for equianharmonic and harmonic orbit counts.}
\Field{Limits of the analogy.}{Eligible paragraphs can explain orbit classes, but they cannot replace this compact label for the two class-count results.}
\Field{Baseline narration used.}{not applicable.}
\Field{Complexity and jargon.}{Complexity: \Tag{medium} $\rightarrow$ \Tag{low}. Jargon: \Tag{medium} $\rightarrow$ \Tag{medium}.}
\textbf{Original WADT unit. Footnotes omitted.}
\begin{lstlisting}
equianharmonic and harmonic counts
\end{lstlisting}
\colorbox{RewriteGreen}{\parbox{0.97\textwidth}{\textbf{Rewrite suggestion.}\par\begingroup\small
equianharmonic and harmonic orbit counts
\par\endgroup}}
\Field{Why this rewrite.}{Uses the table cell's own index-label contract and makes the counted objects explicit.}
\Field{Terminology actions.}{Keep: equianharmonic and harmonic counts | Explain: equianharmonic and harmonic}
\par\smallskip
\needspace{10\baselineskip}\subsubsection*{W-07-TC007: no baseline analogue}
\Field{Location.}{Section 07, table\_cell, lines 1671--1673.}
\Field{Draft purpose.}{Give the two Rocq locators for the equianharmonic and harmonic orbit counts.}
\Field{Qualitative closeness.}{no baseline analogue. No eligible baseline paragraph serves as a source-index cell for the two Rocq results for the orbit split and its complement.}
\Field{Limits of the analogy.}{An eligible paragraph cannot preserve the row's direct pairing between the two orbit-count declarations and the mathematical label in W-07-TC006.}
\Field{Baseline narration used.}{not applicable.}
\Field{Complexity and jargon.}{Complexity: \Tag{medium} $\rightarrow$ \Tag{medium}. Jargon: \Tag{medium} $\rightarrow$ \Tag{medium}.}
\textbf{Original WADT unit. Footnotes omitted.}
\begin{lstlisting}
\coqin{orbit\_class\_split},
      \coqin{orbit\_class\_split\_complement}
\end{lstlisting}
\colorbox{RewriteGreen}{\parbox{0.97\textwidth}{\textbf{Rewrite suggestion.}\par\begingroup\small
\coqin{orbit\_class\_split},       \coqin{orbit\_class\_split\_complement}
\par\endgroup}}
\Field{Why this rewrite.}{Uses the source-index contract and preserves both necessary Rocq locators exactly.}
\Field{Terminology actions.}{Keep: orbit\_class\_split, orbit\_class\_split\_complement | Explain: orbit split and its complement}
\par\smallskip
\needspace{10\baselineskip}\subsubsection*{W-07-TC008: no baseline analogue}
\Field{Location.}{Section 07, table\_cell, lines 1673--1674.}
\Field{Draft purpose.}{Label the construction claim that the encoder returns the orbit class supplied to it.}
\Field{Qualitative closeness.}{no baseline analogue. No eligible baseline paragraph serves as a source-index cell for the claim that the encoder returns the orbit class.}
\Field{Limits of the analogy.}{A paragraph would add exposition to a row whose sole task is to name the encoder-correctness property.}
\Field{Baseline narration used.}{not applicable.}
\Field{Complexity and jargon.}{Complexity: \Tag{low} $\rightarrow$ \Tag{low}. Jargon: \Tag{medium} $\rightarrow$ \Tag{low}.}
\textbf{Original WADT unit. Footnotes omitted.}
\begin{lstlisting}
encoder returns its class
\end{lstlisting}
\colorbox{RewriteGreen}{\parbox{0.97\textwidth}{\textbf{Rewrite suggestion.}\par\begingroup\small
encoder returns its orbit class
\par\endgroup}}
\Field{Why this rewrite.}{Uses the table cell's own claim-label contract and names the class more precisely.}
\Field{Terminology actions.}{Keep: encoder returns its class | Explain: class}
\par\smallskip
\needspace{10\baselineskip}\subsubsection*{W-07-TC013: no baseline analogue}
\Field{Location.}{Section 07, table\_cell, lines 1677--1678.}
\Field{Draft purpose.}{Label the recovery claim for the orbit class after executing the PGL protocol.}
\Field{Qualitative closeness.}{no baseline analogue. No eligible baseline paragraph serves as a source-index cell for executed recovery of the orbit class.}
\Field{Limits of the analogy.}{A paragraph can discuss recovery, but it cannot keep this row's compact focus on executed class recovery.}
\Field{Baseline narration used.}{not applicable.}
\Field{Complexity and jargon.}{Complexity: \Tag{low} $\rightarrow$ \Tag{low}. Jargon: \Tag{medium} $\rightarrow$ \Tag{low}.}
\textbf{Original WADT unit. Footnotes omitted.}
\begin{lstlisting}
executed class recovery
\end{lstlisting}
\colorbox{RewriteGreen}{\parbox{0.97\textwidth}{\textbf{Rewrite suggestion.}\par\begingroup\small
executed recovery of the orbit class
\par\endgroup}}
\Field{Why this rewrite.}{Uses the table cell's own claim-label contract and states the recovered object directly.}
\Field{Terminology actions.}{Keep: executed class recovery | Explain: executed}
\par\smallskip
\needspace{10\baselineskip}\subsubsection*{W-07-TC014: no baseline analogue}
\Field{Location.}{Section 07, table\_cell, lines 1677--1678.}
\Field{Draft purpose.}{Identify the Rocq theorem that proves executed recovery of the orbit class.}
\Field{Qualitative closeness.}{no baseline analogue. No eligible baseline paragraph serves as a source-index cell for the Rocq result for executed class recovery.}
\Field{Limits of the analogy.}{Replacing this locator with prose would remove the direct source link for the executed-recovery claim in W-07-TC013.}
\Field{Baseline narration used.}{not applicable.}
\Field{Complexity and jargon.}{Complexity: \Tag{low} $\rightarrow$ \Tag{low}. Jargon: \Tag{medium} $\rightarrow$ \Tag{medium}.}
\textbf{Original WADT unit. Footnotes omitted.}
\begin{lstlisting}
\coqin{pgl27\_run\_recovers\_class}
\end{lstlisting}
\colorbox{RewriteGreen}{\parbox{0.97\textwidth}{\textbf{Rewrite suggestion.}\par\begingroup\small
\coqin{pgl27\_run\_recovers\_class}
\par\endgroup}}
\Field{Why this rewrite.}{Uses the source-index contract and preserves the necessary Rocq locator exactly.}
\Field{Terminology actions.}{Keep: pgl27\_run\_recovers\_class | Explain: run recovers class}
\par\smallskip
\needspace{10\baselineskip}\subsubsection*{W-07-TC015: no baseline analogue}
\Field{Location.}{Section 07, table\_cell, lines 1678--1679.}
\Field{Draft purpose.}{Label the reconstruction claim that seven revealed cards determine both the full deck and its orbit class.}
\Field{Qualitative closeness.}{no baseline analogue. No eligible baseline paragraph serves as a source-index cell for recovery of the deck and class from seven cards.}
\Field{Limits of the analogy.}{A paragraph would obscure the row's joint seven-card conclusion about both deck and class reconstruction.}
\Field{Baseline narration used.}{not applicable.}
\Field{Complexity and jargon.}{Complexity: \Tag{medium} $\rightarrow$ \Tag{low}. Jargon: \Tag{medium} $\rightarrow$ \Tag{medium}.}
\textbf{Original WADT unit. Footnotes omitted.}
\begin{lstlisting}
seven cards determine deck and class
\end{lstlisting}
\colorbox{RewriteGreen}{\parbox{0.97\textwidth}{\textbf{Rewrite suggestion.}\par\begingroup\small
seven cards determine the deck and orbit class
\par\endgroup}}
\Field{Why this rewrite.}{Uses the table cell's own threshold-result contract and names both recovered objects.}
\Field{Terminology actions.}{Keep: seven cards determine deck and class | Explain: determine deck and class}
\par\smallskip
\needspace{10\baselineskip}\subsubsection*{W-07-TC016: no baseline analogue}
\Field{Location.}{Section 07, table\_cell, lines 1679--1681.}
\Field{Draft purpose.}{Give the two Rocq locators for seven-card deck and class reconstruction.}
\Field{Qualitative closeness.}{no baseline analogue. No eligible baseline paragraph serves as a source-index cell for the two Rocq results for seven-card reconstruction.}
\Field{Limits of the analogy.}{A prose match would not preserve the one-to-one links from the seven-card premise to its separate deck and class declarations.}
\Field{Baseline narration used.}{not applicable.}
\Field{Complexity and jargon.}{Complexity: \Tag{medium} $\rightarrow$ \Tag{medium}. Jargon: \Tag{medium} $\rightarrow$ \Tag{medium}.}
\textbf{Original WADT unit. Footnotes omitted.}
\begin{lstlisting}
\coqin{pgl27\_seven\_reveal\_determines},
      \coqin{pgl27\_seven\_reveal\_class}
\end{lstlisting}
\colorbox{RewriteGreen}{\parbox{0.97\textwidth}{\textbf{Rewrite suggestion.}\par\begingroup\small
\coqin{pgl27\_seven\_reveal\_determines},       \coqin{pgl27\_seven\_reveal\_class}
\par\endgroup}}
\Field{Why this rewrite.}{Uses the source-index contract and preserves both seven-card recovery locators exactly.}
\Field{Terminology actions.}{Keep: pgl27\_seven\_reveal\_determines, pgl27\_seven\_reveal\_class | Explain: seven-reveal determination}
\par\smallskip
\needspace{10\baselineskip}\subsubsection*{W-07-TC017: no baseline analogue}
\Field{Location.}{Section 07, table\_cell, lines 1681--1682.}
\Field{Draft purpose.}{Label the boundary claim that every reveal of at most six cards remains ambiguous.}
\Field{Qualitative closeness.}{no baseline analogue. No eligible baseline paragraph serves as a source-index cell for ambiguity of all reveals with at most six cards.}
\Field{Limits of the analogy.}{A paragraph would expand the quantified six-card ambiguity boundary beyond the table's source-index function.}
\Field{Baseline narration used.}{not applicable.}
\Field{Complexity and jargon.}{Complexity: \Tag{low} $\rightarrow$ \Tag{low}. Jargon: \Tag{medium} $\rightarrow$ \Tag{low}.}
\textbf{Original WADT unit. Footnotes omitted.}
\begin{lstlisting}
at most six cards remain ambiguous
\end{lstlisting}
\colorbox{RewriteGreen}{\parbox{0.97\textwidth}{\textbf{Rewrite suggestion.}\par\begingroup\small
at most six revealed cards remain ambiguous
\par\endgroup}}
\Field{Why this rewrite.}{Uses the table cell's own boundary contract and makes the reveal condition explicit.}
\Field{Terminology actions.}{Keep: at most six cards remain ambiguous | Explain: ambiguous}
\par\smallskip
\needspace{10\baselineskip}\subsubsection*{W-07-TC018: no baseline analogue}
\Field{Location.}{Section 07, table\_cell, lines 1682--1684.}
\Field{Draft purpose.}{Give the specific and general Rocq locators for reveal ambiguity through six cards.}
\Field{Qualitative closeness.}{no baseline analogue. No eligible baseline paragraph serves as a source-index cell for the Rocq results for six-card and general reveal ambiguity.}
\Field{Limits of the analogy.}{Prose cannot replace the paired source links for the six-card boundary theorem and its general ambiguity form.}
\Field{Baseline narration used.}{not applicable.}
\Field{Complexity and jargon.}{Complexity: \Tag{medium} $\rightarrow$ \Tag{medium}. Jargon: \Tag{medium} $\rightarrow$ \Tag{medium}.}
\textbf{Original WADT unit. Footnotes omitted.}
\begin{lstlisting}
\coqin{pgl27\_six\_reveal\_ambiguous},
      \coqin{pgl27\_reveal\_ambiguous}
\end{lstlisting}
\colorbox{RewriteGreen}{\parbox{0.97\textwidth}{\textbf{Rewrite suggestion.}\par\begingroup\small
\coqin{pgl27\_six\_reveal\_ambiguous},       \coqin{pgl27\_reveal\_ambiguous}
\par\endgroup}}
\Field{Why this rewrite.}{Uses the source-index contract and preserves the boundary and general ambiguity locators exactly.}
\Field{Terminology actions.}{Keep: pgl27\_six\_reveal\_ambiguous, pgl27\_reveal\_ambiguous | Explain: six-card and general reveal ambiguity}
\par\smallskip
\needspace{10\baselineskip}\subsubsection*{W-07-TC019: no baseline analogue}
\Field{Location.}{Section 07, table\_cell, lines 1684--1685.}
\Field{Draft purpose.}{Label the sharpness evidence that four-card views depend on the secret and leak information.}
\Field{Qualitative closeness.}{no baseline analogue. No eligible baseline paragraph serves as a source-index cell for four-card dependence and positive leakage.}
\Field{Limits of the analogy.}{A paragraph would separate the dependence fact from the leakage consequence that this table label keeps together.}
\Field{Baseline narration used.}{not applicable.}
\Field{Complexity and jargon.}{Complexity: \Tag{medium} $\rightarrow$ \Tag{low}. Jargon: \Tag{medium} $\rightarrow$ \Tag{medium}.}
\textbf{Original WADT unit. Footnotes omitted.}
\begin{lstlisting}
four-card dependence and leakage
\end{lstlisting}
\colorbox{RewriteGreen}{\parbox{0.97\textwidth}{\textbf{Rewrite suggestion.}\par\begingroup\small
four-card view dependence and leakage
\par\endgroup}}
\Field{Why this rewrite.}{Uses the table cell's own result-label contract and names the observed view.}
\Field{Terminology actions.}{Keep: four-card dependence and leakage | Explain: dependence and leakage}
\par\smallskip
\needspace{10\baselineskip}\subsubsection*{W-07-TC020: no baseline analogue}
\Field{Location.}{Section 07, table\_cell, lines 1685--1686.}
\Field{Draft purpose.}{Give the Rocq locators for four-card dependence and four-card leakage.}
\Field{Qualitative closeness.}{no baseline analogue. No eligible baseline paragraph serves as a source-index cell for the Rocq results for four-card dependence and leakage.}
\Field{Limits of the analogy.}{A prose analogue would hide which declaration proves dependence and which declaration proves leakage at four cards.}
\Field{Baseline narration used.}{not applicable.}
\Field{Complexity and jargon.}{Complexity: \Tag{medium} $\rightarrow$ \Tag{medium}. Jargon: \Tag{medium} $\rightarrow$ \Tag{medium}.}
\textbf{Original WADT unit. Footnotes omitted.}
\begin{lstlisting}
\coqin{pgl27\_view\_dep\_k4}, \coqin{pgl27\_view\_leak\_k4}
\end{lstlisting}
\colorbox{RewriteGreen}{\parbox{0.97\textwidth}{\textbf{Rewrite suggestion.}\par\begingroup\small
\coqin{pgl27\_view\_dep\_k4}, \coqin{pgl27\_view\_leak\_k4}
\par\endgroup}}
\Field{Why this rewrite.}{Uses the source-index contract and preserves the dependence and leakage locators exactly.}
\Field{Terminology actions.}{Keep: pgl27\_view\_dep\_k4, pgl27\_view\_leak\_k4 | Explain: view dependence and view leakage}
\par\smallskip
\needspace{10\baselineskip}\subsubsection*{W-07-TC022: no baseline analogue}
\Field{Location.}{Section 07, table\_cell, lines 1688--1689.}
\Field{Draft purpose.}{Label the view-privacy results for the fixed representative dealer.}
\Field{Qualitative closeness.}{no baseline analogue. No eligible baseline paragraph serves as a source-index cell for view privacy for the fixed representative dealer.}
\Field{Limits of the analogy.}{A paragraph cannot substitute for this compact row label that fixes both the representative dealer and the view observable.}
\Field{Baseline narration used.}{not applicable.}
\Field{Complexity and jargon.}{Complexity: \Tag{low} $\rightarrow$ \Tag{low}. Jargon: \Tag{medium} $\rightarrow$ \Tag{medium}.}
\textbf{Original WADT unit. Footnotes omitted.}
\begin{lstlisting}
fixed representative views
\end{lstlisting}
\colorbox{RewriteGreen}{\parbox{0.97\textwidth}{\textbf{Rewrite suggestion.}\par\begingroup\small
fixed-representative views
\par\endgroup}}
\Field{Why this rewrite.}{Uses the table cell's own scope-label contract and keeps the dealer restriction concise.}
\Field{Terminology actions.}{Keep: fixed representative views | Explain: fixed representative}
\par\smallskip
\needspace{10\baselineskip}\subsubsection*{W-07-TC023: no baseline analogue}
\Field{Location.}{Section 07, table\_cell, lines 1689--1690.}
\Field{Draft purpose.}{Give the Rocq locators for fixed-representative view independence and its leakage bound.}
\Field{Qualitative closeness.}{no baseline analogue. No eligible baseline paragraph serves as a source-index cell for the Rocq results for fixed-representative views.}
\Field{Limits of the analogy.}{Prose would weaken the direct mapping from fixed-representative views to the independence and leakage-bound declarations.}
\Field{Baseline narration used.}{not applicable.}
\Field{Complexity and jargon.}{Complexity: \Tag{medium} $\rightarrow$ \Tag{medium}. Jargon: \Tag{medium} $\rightarrow$ \Tag{medium}.}
\textbf{Original WADT unit. Footnotes omitted.}
\begin{lstlisting}
\coqin{pgl27\_view\_indep}, \coqin{pgl27\_view\_leakage\_le}
\end{lstlisting}
\colorbox{RewriteGreen}{\parbox{0.97\textwidth}{\textbf{Rewrite suggestion.}\par\begingroup\small
\coqin{pgl27\_view\_indep}, \coqin{pgl27\_view\_leakage\_le}
\par\endgroup}}
\Field{Why this rewrite.}{Uses the source-index contract and preserves the view-independence and leakage-bound locators exactly.}
\Field{Terminology actions.}{Keep: pgl27\_view\_indep, pgl27\_view\_leakage\_le | Explain: view independence and leakage bound}
\par\smallskip
\needspace{10\baselineskip}\subsubsection*{W-07-TC024: no baseline analogue}
\Field{Location.}{Section 07, table\_cell, lines 1690--1691.}
\Field{Draft purpose.}{Label the trace-privacy results for the fixed representative dealer.}
\Field{Qualitative closeness.}{no baseline analogue. No eligible baseline paragraph serves as a source-index cell for trace privacy for the fixed representative dealer.}
\Field{Limits of the analogy.}{A paragraph cannot replace this table label's exact combination of fixed-dealer scope and execution traces.}
\Field{Baseline narration used.}{not applicable.}
\Field{Complexity and jargon.}{Complexity: \Tag{low} $\rightarrow$ \Tag{low}. Jargon: \Tag{medium} $\rightarrow$ \Tag{medium}.}
\textbf{Original WADT unit. Footnotes omitted.}
\begin{lstlisting}
fixed representative traces
\end{lstlisting}
\colorbox{RewriteGreen}{\parbox{0.97\textwidth}{\textbf{Rewrite suggestion.}\par\begingroup\small
fixed-representative traces
\par\endgroup}}
\Field{Why this rewrite.}{Uses the table cell's own scope-label contract and keeps the trace observable concise.}
\Field{Terminology actions.}{Keep: fixed representative traces | Explain: trace}
\par\smallskip
\needspace{10\baselineskip}\subsubsection*{W-07-TC025: no baseline analogue}
\Field{Location.}{Section 07, table\_cell, lines 1691--1693.}
\Field{Draft purpose.}{Give the Rocq locators for individual and coalition trace secrecy under the fixed representative dealer.}
\Field{Qualitative closeness.}{no baseline analogue. No eligible baseline paragraph serves as a source-index cell for the Rocq results for fixed-representative trace secrecy.}
\Field{Limits of the analogy.}{A prose analogue would not retain the row's separate source links for individual and coalition trace secrecy.}
\Field{Baseline narration used.}{not applicable.}
\Field{Complexity and jargon.}{Complexity: \Tag{medium} $\rightarrow$ \Tag{medium}. Jargon: \Tag{medium} $\rightarrow$ \Tag{medium}.}
\textbf{Original WADT unit. Footnotes omitted.}
\begin{lstlisting}
\coqin{pgl27\_trace\_secrecy},
      \coqin{pgl27\_coalition\_trace\_secrecy}
\end{lstlisting}
\colorbox{RewriteGreen}{\parbox{0.97\textwidth}{\textbf{Rewrite suggestion.}\par\begingroup\small
\coqin{pgl27\_trace\_secrecy},       \coqin{pgl27\_coalition\_trace\_secrecy}
\par\endgroup}}
\Field{Why this rewrite.}{Uses the source-index contract and preserves both trace-secrecy locators exactly.}
\Field{Terminology actions.}{Keep: pgl27\_trace\_secrecy, pgl27\_coalition\_trace\_secrecy | Explain: individual and coalition traces}
\par\smallskip
\needspace{10\baselineskip}\subsubsection*{W-07-TC026: no baseline analogue}
\Field{Location.}{Section 07, table\_cell, lines 1693--1694.}
\Field{Draft purpose.}{Label the view- and trace-privacy family for the all-decks dealer.}
\Field{Qualitative closeness.}{no baseline analogue. No eligible baseline paragraph serves as a source-index cell for view and trace privacy for the all-decks dealer.}
\Field{Limits of the analogy.}{A paragraph would unpack the all-decks dealer, views, and traces instead of serving as their shared source-index label.}
\Field{Baseline narration used.}{not applicable.}
\Field{Complexity and jargon.}{Complexity: \Tag{medium} $\rightarrow$ \Tag{low}. Jargon: \Tag{medium} $\rightarrow$ \Tag{medium}.}
\textbf{Original WADT unit. Footnotes omitted.}
\begin{lstlisting}
all-decks views and traces
\end{lstlisting}
\colorbox{RewriteGreen}{\parbox{0.97\textwidth}{\textbf{Rewrite suggestion.}\par\begingroup\small
all-decks views and traces
\par\endgroup}}
\Field{Why this rewrite.}{Uses the table cell's own dealer-scope contract and retains both all-decks observable families.}
\Field{Terminology actions.}{Keep: all-decks views and traces | Explain: all-decks dealer}
\par\smallskip
\needspace{10\baselineskip}\subsubsection*{W-07-TC027: no baseline analogue}
\Field{Location.}{Section 07, table\_cell, lines 1694--1697.}
\Field{Draft purpose.}{Give the three Rocq locators for all-decks view, trace, and coalition-trace secrecy.}
\Field{Qualitative closeness.}{no baseline analogue. No eligible baseline paragraph serves as a source-index cell for the Rocq results for all-decks view and trace secrecy.}
\Field{Limits of the analogy.}{Prose cannot preserve the direct three-way mapping to all-decks view, trace, and coalition-secrecy declarations.}
\Field{Baseline narration used.}{not applicable.}
\Field{Complexity and jargon.}{Complexity: \Tag{medium} $\rightarrow$ \Tag{medium}. Jargon: \Tag{medium} $\rightarrow$ \Tag{medium}.}
\textbf{Original WADT unit. Footnotes omitted.}
\begin{lstlisting}
\coqin{pgl27\_view\_indep\_alldecks},
      \coqin{pgl27\_alldecks\_trace\_secrecy},
      \coqin{pgl27\_alldecks\_coalition\_secrecy}
\end{lstlisting}
\colorbox{RewriteGreen}{\parbox{0.97\textwidth}{\textbf{Rewrite suggestion.}\par\begingroup\small
\coqin{pgl27\_view\_indep\_alldecks},       \coqin{pgl27\_alldecks\_trace\_secrecy},       \coqin{pgl27\_alldecks\_coalition\_secrecy}
\par\endgroup}}
\Field{Why this rewrite.}{Uses the source-index contract and preserves all three all-decks locators exactly.}
\Field{Terminology actions.}{Keep: pgl27\_view\_indep\_alldecks, pgl27\_alldecks\_trace\_secrecy, pgl27\_alldecks\_coalition\_secrecy | Explain: all-decks view, trace, and coalition secrecy}
\par\smallskip
\needspace{10\baselineskip}\subsubsection*{W-07-TC028: no baseline analogue}
\Field{Location.}{Section 07, table\_cell, lines 1697--1698.}
\Field{Draft purpose.}{Label the view- and trace-privacy family for the shuffle-free dealer.}
\Field{Qualitative closeness.}{no baseline analogue. No eligible baseline paragraph serves as a source-index cell for view and trace privacy for the shuffle-free dealer.}
\Field{Limits of the analogy.}{A paragraph would explain the shuffle-free dealer model rather than label its view and trace results in the index.}
\Field{Baseline narration used.}{not applicable.}
\Field{Complexity and jargon.}{Complexity: \Tag{medium} $\rightarrow$ \Tag{low}. Jargon: \Tag{medium} $\rightarrow$ \Tag{medium}.}
\textbf{Original WADT unit. Footnotes omitted.}
\begin{lstlisting}
shuffle-free views and traces
\end{lstlisting}
\colorbox{RewriteGreen}{\parbox{0.97\textwidth}{\textbf{Rewrite suggestion.}\par\begingroup\small
shuffle-free views and traces
\par\endgroup}}
\Field{Why this rewrite.}{Uses the table cell's own dealer-scope contract and retains the view and trace claims for the shuffle-free dealer.}
\Field{Terminology actions.}{Keep: shuffle-free views and traces | Explain: shuffle-free dealer}
\par\smallskip
\needspace{10\baselineskip}\subsubsection*{W-07-TC029: no baseline analogue}
\Field{Location.}{Section 07, table\_cell, lines 1698--1701.}
\Field{Draft purpose.}{Give the three Rocq locators for shuffle-free view, trace, and coalition secrecy.}
\Field{Qualitative closeness.}{no baseline analogue. No eligible baseline paragraph serves as a source-index cell for the Rocq results for shuffle-free view and trace secrecy.}
\Field{Limits of the analogy.}{A prose analogue would lose the direct links to the deck-prior view, individual-trace, and coalition-trace declarations.}
\Field{Baseline narration used.}{not applicable.}
\Field{Complexity and jargon.}{Complexity: \Tag{medium} $\rightarrow$ \Tag{medium}. Jargon: \Tag{medium} $\rightarrow$ \Tag{medium}.}
\textbf{Original WADT unit. Footnotes omitted.}
\begin{lstlisting}
\coqin{pgl27\_view\_indep\_deck\_prior},
      \coqin{pgl27\_deck\_trace\_secrecy},
      \coqin{pgl27\_deck\_coalition\_secrecy}
\end{lstlisting}
\colorbox{RewriteGreen}{\parbox{0.97\textwidth}{\textbf{Rewrite suggestion.}\par\begingroup\small
\coqin{pgl27\_view\_indep\_deck\_prior},       \coqin{pgl27\_deck\_trace\_secrecy},       \coqin{pgl27\_deck\_coalition\_secrecy}
\par\endgroup}}
\Field{Why this rewrite.}{Uses the source-index contract and preserves all three shuffle-free locators exactly.}
\Field{Terminology actions.}{Keep: pgl27\_view\_indep\_deck\_prior, pgl27\_deck\_trace\_secrecy, pgl27\_deck\_coalition\_secrecy | Explain: deck-prior view, trace, and coalition secrecy}
\par\smallskip
\needspace{10\baselineskip}\subsubsection*{W-07-TC031: no baseline analogue}
\Field{Location.}{Section 07, table\_cell, lines 1703--1704.}
\Field{Draft purpose.}{Label the checked result that bounds mixing of the finite generator-word distribution.}
\Field{Qualitative closeness.}{no baseline analogue. No eligible baseline paragraph serves as a source-index cell for the finite-word mixing certificate.}
\Field{Limits of the analogy.}{A paragraph can explain finite-word mixing, but it cannot replace this concise label for the checked certificate.}
\Field{Baseline narration used.}{not applicable.}
\Field{Complexity and jargon.}{Complexity: \Tag{low} $\rightarrow$ \Tag{low}. Jargon: \Tag{low} $\rightarrow$ \Tag{low}.}
\textbf{Original WADT unit. Footnotes omitted.}
\begin{lstlisting}
word mixing certificate
\end{lstlisting}
\colorbox{RewriteGreen}{\parbox{0.97\textwidth}{\textbf{Rewrite suggestion.}\par\begingroup\small
word mixing certificate
\par\endgroup}}
\Field{Why this rewrite.}{Uses the table cell's own certificate-label contract and retains the established term.}
\Field{Terminology actions.}{Keep: word mixing certificate | Explain: word mixing}
\par\smallskip
\needspace{10\baselineskip}\subsubsection*{W-07-TC032: no baseline analogue}
\Field{Location.}{Section 07, table\_cell, lines 1703--1704.}
\Field{Draft purpose.}{Identify the Rocq theorem that contains the finite-word mixing certificate.}
\Field{Qualitative closeness.}{no baseline analogue. No eligible baseline paragraph serves as a source-index cell for the Rocq result for word mixing.}
\Field{Limits of the analogy.}{Prose would remove the direct source locator attached to the word-mixing certificate in W-07-TC031.}
\Field{Baseline narration used.}{not applicable.}
\Field{Complexity and jargon.}{Complexity: \Tag{low} $\rightarrow$ \Tag{low}. Jargon: \Tag{medium} $\rightarrow$ \Tag{medium}.}
\textbf{Original WADT unit. Footnotes omitted.}
\begin{lstlisting}
\coqin{pgl27\_word\_mixing}
\end{lstlisting}
\colorbox{RewriteGreen}{\parbox{0.97\textwidth}{\textbf{Rewrite suggestion.}\par\begingroup\small
\coqin{pgl27\_word\_mixing}
\par\endgroup}}
\Field{Why this rewrite.}{Uses the source-index contract and preserves the word-mixing locator exactly.}
\Field{Terminology actions.}{Keep: pgl27\_word\_mixing | Explain: word mixing}
\par\smallskip
\needspace{10\baselineskip}\subsubsection*{W-07-TC033: no baseline analogue}
\Field{Location.}{Section 07, table\_cell, lines 1704--1705.}
\Field{Draft purpose.}{Label the two transfers that carry the mixing bound to card endpoints and secret-product distributions.}
\Field{Qualitative closeness.}{no baseline analogue. No eligible baseline paragraph serves as a source-index cell for the endpoint and secret-product transfer claims.}
\Field{Limits of the analogy.}{A paragraph would develop the transfer argument instead of naming its endpoint and product roles as one indexed pair.}
\Field{Baseline narration used.}{not applicable.}
\Field{Complexity and jargon.}{Complexity: \Tag{medium} $\rightarrow$ \Tag{low}. Jargon: \Tag{low} $\rightarrow$ \Tag{low}.}
\textbf{Original WADT unit. Footnotes omitted.}
\begin{lstlisting}
endpoint and product transfers
\end{lstlisting}
\colorbox{RewriteGreen}{\parbox{0.97\textwidth}{\textbf{Rewrite suggestion.}\par\begingroup\small
endpoint and product transfers
\par\endgroup}}
\Field{Why this rewrite.}{Uses the table cell's own paired-result contract and preserves both transfer roles.}
\Field{Terminology actions.}{Keep: endpoint and product transfers | Explain: endpoint transfer and product transfer}
\par\smallskip
\needspace{10\baselineskip}\subsubsection*{W-07-TC034: no baseline analogue}
\Field{Location.}{Section 07, table\_cell, lines 1705--1706.}
\Field{Draft purpose.}{Give the Rocq locators for the endpoint transfer and the secret-product transfer.}
\Field{Qualitative closeness.}{no baseline analogue. No eligible baseline paragraph serves as a source-index cell for the Rocq results for endpoint and product transfer.}
\Field{Limits of the analogy.}{A prose analogue cannot keep the row's direct association between the endpoint and joint transfer roles and their two declarations.}
\Field{Baseline narration used.}{not applicable.}
\Field{Complexity and jargon.}{Complexity: \Tag{medium} $\rightarrow$ \Tag{medium}. Jargon: \Tag{medium} $\rightarrow$ \Tag{medium}.}
\textbf{Original WADT unit. Footnotes omitted.}
\begin{lstlisting}
\coqin{pgl27\_endpoint\_mixing}, \coqin{pgl27\_joint\_mixing}
\end{lstlisting}
\colorbox{RewriteGreen}{\parbox{0.97\textwidth}{\textbf{Rewrite suggestion.}\par\begingroup\small
\coqin{pgl27\_endpoint\_mixing}, \coqin{pgl27\_joint\_mixing}
\par\endgroup}}
\Field{Why this rewrite.}{Uses the source-index contract and preserves both transfer locators exactly.}
\Field{Terminology actions.}{Keep: pgl27\_endpoint\_mixing, pgl27\_joint\_mixing | Explain: endpoint mixing and joint mixing}
\par\smallskip
\needspace{10\baselineskip}\subsubsection*{W-07-TC035: no baseline analogue}
\Field{Location.}{Section 07, table\_cell, lines 1706--1707.}
\Field{Draft purpose.}{Label the approximate privacy bounds for word-shuffle coalition views and executed traces.}
\Field{Qualitative closeness.}{no baseline analogue. No eligible baseline paragraph serves as a source-index cell for the word-shuffle view and trace privacy bounds.}
\Field{Limits of the analogy.}{A paragraph would expand the approximate privacy argument rather than label its view and trace bounds together.}
\Field{Baseline narration used.}{not applicable.}
\Field{Complexity and jargon.}{Complexity: \Tag{medium} $\rightarrow$ \Tag{low}. Jargon: \Tag{low} $\rightarrow$ \Tag{low}.}
\textbf{Original WADT unit. Footnotes omitted.}
\begin{lstlisting}
word view and trace bounds
\end{lstlisting}
\colorbox{RewriteGreen}{\parbox{0.97\textwidth}{\textbf{Rewrite suggestion.}\par\begingroup\small
word view and trace bounds
\par\endgroup}}
\Field{Why this rewrite.}{Uses the table cell's own paired-bound contract and preserves both observables.}
\Field{Terminology actions.}{Keep: word view and trace bounds | Explain: view and trace bounds}
\par\smallskip
\needspace{10\baselineskip}\subsubsection*{W-07-TC036: no baseline analogue}
\Field{Location.}{Section 07, table\_cell, lines 1707--1709.}
\Field{Draft purpose.}{Give the Rocq locators for word-shuffle view and trace indistinguishability.}
\Field{Qualitative closeness.}{no baseline analogue. No eligible baseline paragraph serves as a source-index cell for the Rocq results for word view and trace indistinguishability.}
\Field{Limits of the analogy.}{Prose would obscure the separate declaration links for word-shuffle view and trace indistinguishability.}
\Field{Baseline narration used.}{not applicable.}
\Field{Complexity and jargon.}{Complexity: \Tag{medium} $\rightarrow$ \Tag{medium}. Jargon: \Tag{medium} $\rightarrow$ \Tag{medium}.}
\textbf{Original WADT unit. Footnotes omitted.}
\begin{lstlisting}
\coqin{pgl27\_word\_view\_indist},
      \coqin{pgl27\_word\_trace\_indist}
\end{lstlisting}
\colorbox{RewriteGreen}{\parbox{0.97\textwidth}{\textbf{Rewrite suggestion.}\par\begingroup\small
\coqin{pgl27\_word\_view\_indist},       \coqin{pgl27\_word\_trace\_indist}
\par\endgroup}}
\Field{Why this rewrite.}{Uses the source-index contract and preserves both indistinguishability locators exactly.}
\Field{Terminology actions.}{Keep: pgl27\_word\_view\_indist, pgl27\_word\_trace\_indist | Explain: view and trace indistinguishability}
\par\smallskip
\needspace{10\baselineskip}\subsubsection*{W-07-TC038: no baseline analogue}
\Field{Location.}{Section 07, table\_cell, lines 1709--1710.}
\Field{Draft purpose.}{Identify the Rocq companion theorem that compares the word-shuffle view law with its ideal product law.}
\Field{Qualitative closeness.}{no baseline analogue. No eligible baseline paragraph serves as a source-index cell for the Rocq result for ideal proximity of the view law.}
\Field{Limits of the analogy.}{A paragraph cannot replace the single source locator for the ideal-proximity companion result.}
\Field{Baseline narration used.}{not applicable.}
\Field{Complexity and jargon.}{Complexity: \Tag{low} $\rightarrow$ \Tag{low}. Jargon: \Tag{medium} $\rightarrow$ \Tag{medium}.}
\textbf{Original WADT unit. Footnotes omitted.}
\begin{lstlisting}
\coqin{pgl27\_view\_mixing}
\end{lstlisting}
\colorbox{RewriteGreen}{\parbox{0.97\textwidth}{\textbf{Rewrite suggestion.}\par\begingroup\small
\coqin{pgl27\_view\_mixing}
\par\endgroup}}
\Field{Why this rewrite.}{Uses the source-index contract and preserves the ideal-proximity locator exactly.}
\Field{Terminology actions.}{Keep: pgl27\_view\_mixing | Explain: view mixing}
\par\smallskip
\needspace{10\baselineskip}\subsubsection*{W-07-TC040: no baseline analogue}
\Field{Location.}{Section 07, table\_cell, lines 1710--1711.}
\Field{Draft purpose.}{Identify the Rocq theorem that proves executed correctness for every sampled generator word.}
\Field{Qualitative closeness.}{no baseline analogue. No eligible baseline paragraph serves as a source-index cell for the Rocq result for wordwise executed correctness.}
\Field{Limits of the analogy.}{Prose would remove the direct source link for correctness of protocol execution under a generator word.}
\Field{Baseline narration used.}{not applicable.}
\Field{Complexity and jargon.}{Complexity: \Tag{low} $\rightarrow$ \Tag{low}. Jargon: \Tag{medium} $\rightarrow$ \Tag{medium}.}
\textbf{Original WADT unit. Footnotes omitted.}
\begin{lstlisting}
\coqin{pgl27\_word\_run\_recovers}
\end{lstlisting}
\colorbox{RewriteGreen}{\parbox{0.97\textwidth}{\textbf{Rewrite suggestion.}\par\begingroup\small
\coqin{pgl27\_word\_run\_recovers}
\par\endgroup}}
\Field{Why this rewrite.}{Uses the source-index contract and preserves the wordwise-correctness locator exactly.}
\Field{Terminology actions.}{Keep: pgl27\_word\_run\_recovers | Explain: word run recovers}
\par\smallskip
\needspace{10\baselineskip}\subsubsection*{W-07-CAP001: no baseline analogue}
\Field{Location.}{Section 07, caption, lines 1713--1714.}
\Field{Draft purpose.}{Identify the table as the Rocq source index for the PGL results.}
\Field{Qualitative closeness.}{no baseline analogue. No eligible baseline paragraph serves as a terse caption for a formal-source index.}
\Field{Limits of the analogy.}{A roadmap paragraph would add an explanatory sequence absent from this caption.}
\Field{Baseline narration used.}{not applicable.}
\Field{Complexity and jargon.}{Complexity: \Tag{low} $\rightarrow$ \Tag{low}. Jargon: \Tag{medium} $\rightarrow$ \Tag{low}.}
\textbf{Original WADT unit. Footnotes omitted.}
\begin{lstlisting}
Rocq sources for the $\PG$ results.
\end{lstlisting}
\colorbox{RewriteGreen}{\parbox{0.97\textwidth}{\textbf{Rewrite suggestion.}\par\begingroup\small
Rocq sources for the $\PG$ results.
\par\endgroup}}
\Field{Why this rewrite.}{Uses the caption's own source-index contract and retains its exact mathematical scope.}
\Field{Terminology actions.}{Keep: Rocq, PGL results}
\par\smallskip
\needspace{10\baselineskip}\subsubsection*{W-08-TC015: no baseline analogue}
\Field{Location.}{Section 08, table\_cell, lines 1748--1749.}
\Field{Draft purpose.}{State the perfect-privacy threshold for the S5 instance.}
\Field{Qualitative closeness.}{no baseline analogue. No eligible baseline paragraph serves as one atomic threshold entry in a comparison table.}
\Field{Limits of the analogy.}{A prose threshold explanation would add quantifiers and context supplied by the row and headers.}
\Field{Baseline narration used.}{not applicable.}
\Field{Complexity and jargon.}{Complexity: \Tag{low} $\rightarrow$ \Tag{low}. Jargon: \Tag{low} $\rightarrow$ \Tag{low}.}
\textbf{Original WADT unit. Footnotes omitted.}
\begin{lstlisting}
fewer than five
\end{lstlisting}
\colorbox{RewriteGreen}{\parbox{0.97\textwidth}{\textbf{Rewrite suggestion.}\par\begingroup\small
fewer than five positions
\par\endgroup}}
\Field{Why this rewrite.}{Keeps the table cell as one strict privacy threshold and makes its counted object explicit.}
\Field{Terminology actions.}{Keep: fewer than five}
\par\smallskip
\needspace{10\baselineskip}\subsubsection*{W-08-TC018: no baseline analogue}
\Field{Location.}{Section 08, table\_cell, lines 1749--1750.}
\Field{Draft purpose.}{State the recovery structure for the product instance.}
\Field{Qualitative closeness.}{no baseline analogue. No eligible baseline paragraph functions as a compact recovery entry whose meaning depends on table headers.}
\Field{Limits of the analogy.}{Definition prose would duplicate context outside the cell.}
\Field{Baseline narration used.}{not applicable.}
\Field{Complexity and jargon.}{Complexity: \Tag{low} $\rightarrow$ \Tag{low}. Jargon: \Tag{medium} $\rightarrow$ \Tag{low}.}
\textbf{Original WADT unit. Footnotes omitted.}
\begin{lstlisting}
two additive 5-of-5 shares
\end{lstlisting}
\colorbox{RewriteGreen}{\parbox{0.97\textwidth}{\textbf{Rewrite suggestion.}\par\begingroup\small
two additive 5-of-5 shares, each reconstructed from all five positions
\par\endgroup}}
\Field{Why this rewrite.}{Keeps the compact table role while stating that each additive share requires all five positions.}
\Field{Terminology actions.}{Keep: five-of-five | Explain: additive share -> each share is reconstructed from all five positions}
\par\smallskip
\needspace{10\baselineskip}\subsubsection*{W-08-TC019: no baseline analogue}
\Field{Location.}{Section 08, table\_cell, lines 1750--1751.}
\Field{Draft purpose.}{State the perfect-privacy cutoff for the product instance.}
\Field{Qualitative closeness.}{no baseline analogue. No eligible baseline paragraph is an atomic table value for a pile-specific privacy threshold.}
\Field{Limits of the analogy.}{A security paragraph would require the dealer law and observation model, which the cell inherits from its row context.}
\Field{Baseline narration used.}{not applicable.}
\Field{Complexity and jargon.}{Complexity: \Tag{medium} $\rightarrow$ \Tag{low}. Jargon: \Tag{medium} $\rightarrow$ \Tag{low}.}
\textbf{Original WADT unit. Footnotes omitted.}
\begin{lstlisting}
fewer than five per pile
\end{lstlisting}
\colorbox{RewriteGreen}{\parbox{0.97\textwidth}{\textbf{Rewrite suggestion.}\par\begingroup\small
fewer than five cards in each five-card pile
\par\endgroup}}
\Field{Why this rewrite.}{Preserves the strict privacy cutoff and explains that the ten-card product splits into two piles.}
\Field{Terminology actions.}{Keep: fewer than five per pile | Explain: pile -> one of the two five-card groups}
\par\smallskip
\needspace{10\baselineskip}\subsubsection*{W-08-TC022: no baseline analogue}
\Field{Location.}{Section 08, table\_cell, lines 1751--1752.}
\Field{Draft purpose.}{State the PGL recovery ramp.}
\Field{Qualitative closeness.}{no baseline analogue. No eligible baseline paragraph serves as a bare parameter tuple in a comparison table.}
\Field{Limits of the analogy.}{Explaining the tuple in prose would change the cell's compact data role.}
\Field{Baseline narration used.}{not applicable.}
\Field{Complexity and jargon.}{Complexity: \Tag{low} $\rightarrow$ \Tag{low}. Jargon: \Tag{medium} $\rightarrow$ \Tag{low}.}
\textbf{Original WADT unit. Footnotes omitted.}
\begin{lstlisting}
$(t,r,n)=(3,7,8)$
\end{lstlisting}
\colorbox{RewriteGreen}{\parbox{0.97\textwidth}{\textbf{Rewrite suggestion.}\par\begingroup\small
$(t,r,n)=(3,7,8)$: privacy threshold, reconstruction threshold, and player count
\par\endgroup}}
\Field{Why this rewrite.}{Keeps the exact recovery-ramp tuple and explains the registered order of its three parameters within the cell.}
\Field{Terminology actions.}{Keep: recovery ramp | Explain: parameters t, r, and n -> privacy threshold, reconstruction threshold, and player count}
\par\smallskip
\needspace{10\baselineskip}\subsubsection*{W-08-TC023: no baseline analogue}
\Field{Location.}{Section 08, table\_cell, lines 1751--1752.}
\Field{Draft purpose.}{State the PGL perfect-privacy scope.}
\Field{Qualitative closeness.}{no baseline analogue. No eligible baseline paragraph is a single table entry for one coalition bound.}
\Field{Limits of the analogy.}{A privacy paragraph would include a distribution and proof condition absent from this cell.}
\Field{Baseline narration used.}{not applicable.}
\Field{Complexity and jargon.}{Complexity: \Tag{low} $\rightarrow$ \Tag{low}. Jargon: \Tag{low} $\rightarrow$ \Tag{low}.}
\textbf{Original WADT unit. Footnotes omitted.}
\begin{lstlisting}
coalitions of at most three
\end{lstlisting}
\colorbox{RewriteGreen}{\parbox{0.97\textwidth}{\textbf{Rewrite suggestion.}\par\begingroup\small
view privacy for coalitions of at most three
\par\endgroup}}
\Field{Why this rewrite.}{Preserves the inclusive coalition cutoff and identifies this cell as the static-view guarantee.}
\Field{Terminology actions.}{Keep: coalitions of at most three}
\par\smallskip
\needspace{10\baselineskip}\subsubsection*{W-08-TC024: no baseline analogue}
\Field{Location.}{Section 08, table\_cell, lines 1752--1753.}
\Field{Draft purpose.}{State the PGL trace-transport scope.}
\Field{Qualitative closeness.}{no baseline analogue. No eligible baseline paragraph functions as one trace-coverage cell tied to a column heading.}
\Field{Limits of the analogy.}{Prose would risk presenting trace privacy as a separate theorem instead of a transport entry.}
\Field{Baseline narration used.}{not applicable.}
\Field{Complexity and jargon.}{Complexity: \Tag{low} $\rightarrow$ \Tag{low}. Jargon: \Tag{low} $\rightarrow$ \Tag{low}.}
\textbf{Original WADT unit. Footnotes omitted.}
\begin{lstlisting}
coalitions of at most three
\end{lstlisting}
\colorbox{RewriteGreen}{\parbox{0.97\textwidth}{\textbf{Rewrite suggestion.}\par\begingroup\small
trace transport for coalitions of at most three
\par\endgroup}}
\Field{Why this rewrite.}{Preserves the same inclusive scope while identifying this cell as transport to executed traces rather than a second security property.}
\Field{Terminology actions.}{Keep: coalitions of at most three}
\par\smallskip
\needspace{10\baselineskip}\subsubsection*{W-08-TC029: no baseline analogue}
\Field{Location.}{Section 08, table\_cell, lines 1762--1763.}
\Field{Draft purpose.}{State the den Boer mixing result.}
\Field{Qualitative closeness.}{no baseline analogue. No eligible baseline paragraph has the atomic comparison-table role of pairing one shuffle count with exactness.}
\Field{Limits of the analogy.}{A result paragraph would require the protocol and distribution context supplied by the row.}
\Field{Baseline narration used.}{not applicable.}
\Field{Complexity and jargon.}{Complexity: \Tag{low} $\rightarrow$ \Tag{low}. Jargon: \Tag{low} $\rightarrow$ \Tag{low}.}
\textbf{Original WADT unit. Footnotes omitted.}
\begin{lstlisting}
exact after one uniform cut
\end{lstlisting}
\colorbox{RewriteGreen}{\parbox{0.97\textwidth}{\textbf{Rewrite suggestion.}\par\begingroup\small
exact mixing after one uniform cut
\par\endgroup}}
\Field{Why this rewrite.}{Keeps both the single physical operation and its zero-error mixing result in one concise cell.}
\Field{Terminology actions.}{Keep: exact after one uniform cut}
\par\smallskip
\needspace{10\baselineskip}\subsubsection*{W-08-TC035: no baseline analogue}
\Field{Location.}{Section 08, table\_cell, lines 1764--1765.}
\Field{Draft purpose.}{State the spectral mixing route for S5.}
\Field{Qualitative closeness.}{no baseline analogue. No eligible baseline paragraph is a compact method entry in a comparison matrix.}
\Field{Limits of the analogy.}{Proof prose would add the inequality and endpoint law, which are outside this cell.}
\Field{Baseline narration used.}{not applicable.}
\Field{Complexity and jargon.}{Complexity: \Tag{low} $\rightarrow$ \Tag{low}. Jargon: \Tag{medium} $\rightarrow$ \Tag{low}.}
\textbf{Original WADT unit. Footnotes omitted.}
\begin{lstlisting}
$L=286$ spectral route
\end{lstlisting}
\colorbox{RewriteGreen}{\parbox{0.97\textwidth}{\textbf{Rewrite suggestion.}\par\begingroup\small
word length $L=286$ via a spectral inequality
\par\endgroup}}
\Field{Why this rewrite.}{Keeps the exact word length and names the inequality-based mixing method directly.}
\Field{Terminology actions.}{Keep: word length 286 | Explain: spectral route -> a mixing estimate obtained from a spectral inequality}
\par\smallskip
\needspace{10\baselineskip}\subsubsection*{W-08-TC036: no baseline analogue}
\Field{Location.}{Section 08, table\_cell, lines 1764--1765.}
\Field{Draft purpose.}{State the extra trust premise for the S5 result.}
\Field{Qualitative closeness.}{no baseline analogue. No eligible baseline paragraph serves as one trust-base table entry.}
\Field{Limits of the analogy.}{An assumption paragraph would explain the premise and affected theorem beyond the cell's role.}
\Field{Baseline narration used.}{not applicable.}
\Field{Complexity and jargon.}{Complexity: \Tag{low} $\rightarrow$ \Tag{low}. Jargon: \Tag{medium} $\rightarrow$ \Tag{low}.}
\textbf{Original WADT unit. Footnotes omitted.}
\begin{lstlisting}
imported Rayleigh bound
\end{lstlisting}
\colorbox{RewriteGreen}{\parbox{0.97\textwidth}{\textbf{Rewrite suggestion.}\par\begingroup\small
imported Rayleigh bound for spectral mixing
\par\endgroup}}
\Field{Why this rewrite.}{Keeps the additional analytic premise and states the exact proof role it supplies.}
\Field{Terminology actions.}{Keep: imported Rayleigh bound | Explain: Rayleigh bound -> the premise for the spectral mixing inequality}
\par\smallskip
\needspace{10\baselineskip}\subsubsection*{W-08-TC038: no baseline analogue}
\Field{Location.}{Section 08, table\_cell, lines 1765--1766.}
\Field{Draft purpose.}{State the product-instance mixing behavior and global limit.}
\Field{Qualitative closeness.}{no baseline analogue. No eligible baseline paragraph has this atomic two-part table role under a mixing column.}
\Field{Limits of the analogy.}{A limitation paragraph would require the invariant pile structure and target law supplied outside the cell.}
\Field{Baseline narration used.}{not applicable.}
\Field{Complexity and jargon.}{Complexity: \Tag{medium} $\rightarrow$ \Tag{low}. Jargon: \Tag{high} $\rightarrow$ \Tag{medium}.}
\textbf{Original WADT unit. Footnotes omitted.}
\begin{lstlisting}
within-pile decay, global floor $1$
\end{lstlisting}
\colorbox{RewriteGreen}{\parbox{0.97\textwidth}{\textbf{Rewrite suggestion.}\par\begingroup\small
within-pile decay, exact global $L_1$ floor $1$ because cards never change piles
\par\endgroup}}
\Field{Why this rewrite.}{Preserves the contrast between local decay and the exact global obstruction while stating the cause of the floor.}
\Field{Terminology actions.}{Keep: global floor one, within-pile decay | Explain: floor -> the distance cannot fall below one because cards never change piles}
\par\smallskip
\needspace{10\baselineskip}\subsubsection*{W-08-TC039: no baseline analogue}
\Field{Location.}{Section 08, table\_cell, lines 1766--1767.}
\Field{Draft purpose.}{State the inherited trust premise for the product instance.}
\Field{Qualitative closeness.}{no baseline analogue. No eligible baseline paragraph is a single inherited-assumption table entry.}
\Field{Limits of the analogy.}{Comparison prose would add the source and consequences that the row context already supplies.}
\Field{Baseline narration used.}{not applicable.}
\Field{Complexity and jargon.}{Complexity: \Tag{low} $\rightarrow$ \Tag{low}. Jargon: \Tag{medium} $\rightarrow$ \Tag{low}.}
\textbf{Original WADT unit. Footnotes omitted.}
\begin{lstlisting}
inherited Rayleigh bound
\end{lstlisting}
\colorbox{RewriteGreen}{\parbox{0.97\textwidth}{\textbf{Rewrite suggestion.}\par\begingroup\small
Rayleigh bound inherited from the $S_5$ instance
\par\endgroup}}
\Field{Why this rewrite.}{Keeps the shared analytic premise and identifies the component instance from which the product case inherits it.}
\Field{Terminology actions.}{Keep: Rayleigh bound | Explain: inherited -> reused from the S5 instance}
\par\smallskip
\needspace{10\baselineskip}\subsubsection*{W-08-CAP001: no baseline analogue}
\Field{Location.}{Section 08, caption, lines 1770--1772.}
\Field{Draft purpose.}{State the table's comparison scope and mixing metric.}
\Field{Qualitative closeness.}{no baseline analogue. No eligible baseline paragraph combines a table-title function with a column-specific metric note.}
\Field{Limits of the analogy.}{A roadmap paragraph would not preserve the caption's direct relation to the table headers.}
\Field{Baseline narration used.}{not applicable.}
\Field{Complexity and jargon.}{Complexity: \Tag{low} $\rightarrow$ \Tag{low}. Jargon: \Tag{low} $\rightarrow$ \Tag{low}.}
\textbf{Original WADT unit. Footnotes omitted.}
\begin{lstlisting}
Coverage of the five protocol instances. The mixing column uses
  the $L_1$ distance in Equation~\ref{eq:l1-definition}.
\end{lstlisting}
\colorbox{RewriteGreen}{\parbox{0.97\textwidth}{\textbf{Rewrite suggestion.}\par\begingroup\small
Proof coverage across the five protocol instances. The mixing column uses the $L_1$ distance defined in Equation~\ref{eq:l1-definition}.
\par\endgroup}}
\Field{Why this rewrite.}{Uses the caption's own contract to state the comparison scope and metric reference directly.}
\Field{Terminology actions.}{Keep: \$L\_1\$ distance}
\par\smallskip
\needspace{10\baselineskip}\subsubsection*{W-08-P003: no baseline analogue}
\Field{Location.}{Section 08, prose\_paragraph, lines 1775--1785.}
\Field{Draft purpose.}{Disclose the checker and assumptions that support each family of formal results.}
\Field{Qualitative closeness.}{no baseline analogue. No eligible baseline paragraph audits assumptions by result family or distinguishes kernel checking from imported premises.}
\Field{Limits of the analogy.}{Background and scope paragraphs do not provide trust accounting. Using them would hide the distinction among shared principles, one added axiom, and checked computation.}
\Field{Baseline narration used.}{not applicable.}
\Field{Complexity and jargon.}{Complexity: \Tag{high} $\rightarrow$ \Tag{medium}. Jargon: \Tag{high} $\rightarrow$ \Tag{medium}.}
\textbf{Original WADT unit. Footnotes omitted.}
\begin{lstlisting}
The checker and the axioms that a result's verification rests on are
called its trust base. The probability, view, and trace results use three
classical principles
from the MathComp probability stack: propositional extensionality,
dependent functional extensionality, and constructive indefinite
description. The $S_5$ and $S_5\times S_5$ mixing bounds additionally
import a Rayleigh-quotient premise as an axiom. The $\PG$ group, orbit,
cardinality, correctness, and recovery results are closed under the global
context. The \Rocq{} kernel checks all accepted proof terms, including
finite computations performed by its virtual machine.
\end{lstlisting}
\colorbox{RewriteGreen}{\parbox{0.97\textwidth}{\textbf{Rewrite suggestion.}\par\begingroup\small
A result's trust base consists of its checker and the axioms on which its verification depends. The probability, view, and trace results use three classical principles from the MathComp probability stack. Propositional extensionality equates logically equivalent propositions. Dependent functional extensionality equates dependent functions that are pointwise equal. Constructive indefinite description selects a witness from a proof that one exists. The $S_5$ and $S_5\times S_5$ mixing bounds also import a Rayleigh-quotient premise as an axiom. This analytic premise is used only by the spectral bounds. The $\PG$ group, orbit, cardinality, correctness, and recovery results are proved without additional local axioms beyond the global assumptions. The \Rocq{} kernel checks every accepted proof term. It also checks the results of finite computations performed by the virtual machine, which evaluates those computations.
\par\endgroup}}
\Field{Why this rewrite.}{Uses the WADT trust-accounting contract alone to define the trust base, explain each principle, isolate the added spectral axiom, and state the kernel's role.}
\Field{Terminology actions.}{Keep: trust base, axiom, Rocq kernel, proof term | Explain: propositional extensionality -> equates logically equivalent propositions, dependent functional extensionality -> equates dependent functions that are pointwise equal, constructive indefinite description -> selects a witness from a proof that one exists, Rayleigh-quotient premise -> an analytic premise used only by the spectral bounds, virtual machine -> evaluates the finite computations whose results the kernel checks | Replace: closed under the global context -> proved without additional local axioms beyond the global assumptions}
\par\smallskip
\needspace{10\baselineskip}\subsubsection*{W-10-P004: no baseline analogue}
\Field{Location.}{Section 10, prose\_paragraph, lines 1910--1914.}
\Field{Draft purpose.}{Disclose language-model assistance and assign scholarly responsibility.}
\Field{Qualitative closeness.}{no baseline analogue. No eligible baseline paragraph performs a publication-ethics disclosure or allocates responsibility between an author and a drafting tool.}
\Field{Limits of the analogy.}{Scientific conclusion or scope prose cannot guide this accountability function.}
\Field{Baseline narration used.}{not applicable.}
\Field{Complexity and jargon.}{Complexity: \Tag{low} $\rightarrow$ \Tag{low}. Jargon: \Tag{low} $\rightarrow$ \Tag{low}.}
\textbf{Original WADT unit. Footnotes omitted.}
\begin{lstlisting}
The author used LLM-generated proof scripts and documentation drafts under
human specification, review, and responsibility. The author retains
responsibility for the mathematical specification, claims, citations, and
final text.
\end{lstlisting}
\colorbox{RewriteGreen}{\parbox{0.97\textwidth}{\textbf{Rewrite suggestion.}\par\begingroup\small
I used LLM-generated proof scripts and documentation drafts under human specification and review. I remain responsible for the mathematical specification, claims, citations, and final text.
\par\endgroup}}
\Field{Why this rewrite.}{Uses the disclosure's own accountability contract and applies first-person singular voice while preserving the tool use, human review, and responsibility statement.}
\Field{Terminology actions.}{Keep: proof scripts, mathematical specification | Explain: LLM}
\par\smallskip
\needspace{10\baselineskip}\subsubsection*{W-10-P005: no baseline analogue}
\Field{Location.}{Section 10, prose\_paragraph, lines 1917--1918.}
\Field{Draft purpose.}{Declare the absence of competing interests.}
\Field{Qualitative closeness.}{no baseline analogue. No eligible baseline paragraph is a competing-interests declaration.}
\Field{Limits of the analogy.}{A scientific conclusion would add content and change this formal disclosure.}
\Field{Baseline narration used.}{not applicable.}
\Field{Complexity and jargon.}{Complexity: \Tag{low} $\rightarrow$ \Tag{low}. Jargon: \Tag{low} $\rightarrow$ \Tag{low}.}
\textbf{Original WADT unit. Footnotes omitted.}
\begin{lstlisting}
The author declares no competing interests.
\end{lstlisting}
\colorbox{RewriteGreen}{\parbox{0.97\textwidth}{\textbf{Rewrite suggestion.}\par\begingroup\small
I declare no competing interests.
\par\endgroup}}
\Field{Why this rewrite.}{Uses the disclosure's own contract and applies first-person singular voice without changing the declaration.}
\Field{Terminology actions.}{Keep: competing interests}
\par\smallskip
\section{Status}
This analysis covers all 299 included Shinagawa units and all 225 included WADT rewrite units. The rewrite suggestions are table-only proposals. The live WADT manuscript and bibliography are outside this artifact and remain unchanged.
\end{document}
